\documentclass[a4paper,11pt]{article}

\usepackage{amsmath}
\usepackage{amssymb}
\usepackage{amsthm}
\usepackage{hyperref}
\usepackage[hmargin=2cm]{geometry}
\usepackage{fvextra}  % for breakable verbatim
\DefineVerbatimEnvironment{Verbatim}{Verbatim}{breaklines,fontsize=\small}
\fvset{breaklines}

\usepackage{microtype}  % character protrusion and font expansion for better line breaks
\emergencystretch=3em
\Urlmuskip=0mu plus 2mu

% Allow hyphenation in tt font
\usepackage{hyphenat}
\usepackage{newunicodechar}
\newunicodechar{β}{\ensuremath{\beta}}
\newunicodechar{ℤ}{\ensuremath{\mathbb{Z}}}
\newunicodechar{ℝ}{\ensuremath{\mathbb{R}}}
\newunicodechar{≤}{\ensuremath{\le}}
\newunicodechar{≥}{\ensuremath{\ge}}
\newunicodechar{ℓ}{\ensuremath{\ell}}
\newunicodechar{₀}{\ensuremath{_0}}
\newunicodechar{₁}{\ensuremath{_1}}
\newunicodechar{₂}{\ensuremath{_2}}
\newunicodechar{₃}{\ensuremath{_3}}
\newunicodechar{∂}{\ensuremath{\partial}}
\newunicodechar{∈}{\ensuremath{\in}}

\newtheorem{theorem}{Theorem}
\newtheorem{lemma}[theorem]{Lemma}
\newtheorem{definition}[theorem]{Definition}
\newtheorem{proposition}[theorem]{Proposition}
\newtheorem{corollary}[theorem]{Corollary}
\newtheorem{axiom}[theorem]{Axiom}
\newtheorem{remark}[theorem]{Remark}

\newcommand{\LeanLambda}{\texorpdfstring{\ensuremath{\Lambda}}{Lambda}}

\title{Ising Model Formalization: A Mathematical Guide to the Lean 4 Code}
\author{ising-model project}
\date{2026-04-17}

\begin{document}
\maketitle

\begin{abstract}
This document is a guide to reading the Lean~4 source code of the
\texttt{ising-model} project. Rather than reproducing full proofs,
it explains \emph{how to navigate the codebase}: which files to read
in which order, what each file's role is, and where the key ideas live.
Mathematical statements are included to orient the reader, but the
primary goal is to make the formal code accessible.

The project covers Glimm--Jaffe, \emph{Quantum Physics: A Functional
Integral Point of View} (2nd ed., 1987), with particular emphasis on
the Ising-model discussions in \S4.1--4.6 (correlation inequalities,
Lee--Yang, free energy), \S5.1--5.4 (phase transitions, Peierls
argument), \S10.2--10.5 (conditioning), and the infinite-volume limit
theorems in \S4.2 (Thm~4.2.3) and \S4.6 (Prop~4.6.1).
\end{abstract}

\tableofcontents

%======================================================================
\section{Formalization regimes}
%======================================================================

We distinguish three formalization regimes in this project.

\subsection{Finite-volume}

The Ising model on a fixed finite interaction graph
$G : \texttt{SimpleGraph}~\iota$ with $[\texttt{Fintype}~\iota]$.
Most of the project lives here: correlation inequalities (GKS, FKG,
GHS), Lee--Yang, the partition-polynomial nonvanishing, the Peierls
bound, and the real-analyticity of $f(h)$ for $h > 0$.

\subsection{Discretized infinite-volume}

A fixed finite ambient type $\iota$ (so $[\texttt{Fintype}~\iota]$)
with a monotone increasing family of subgraphs $G_1 \le G_2 \le \cdots$.
``Growing'' the lattice then corresponds to adding edges.  The
GJ~\S4.2 Thm~4.2.3 and GJ~\S4.6 Prop~4.6.1 are formalized in this
regime: the proof mechanism (GKS + monotonicity + boundedness) is
identical to GJ, but the ambient lattice is finite.  Every
\texttt{*\_monotone\_subgraph} / \texttt{*\_convergent\_subgraph}
result in the codebase lives in this regime.

\subsection{Genuine infinite-volume (\texttt{AmbientLattice.lean})}

An ambient type $V : \texttt{Type}$ \emph{without} a
$[\texttt{Fintype}~V]$ assumption, and finite volumes $\Lambda :
\texttt{Finset}~V$.  Correlations on $\Lambda$ are defined by
instantiating the existing framework on the Fintype $(\uparrow\!\Lambda)$
using the induced subgraph $G.\texttt{induce}~(\uparrow\!\Lambda)$.
An \emph{exhaustion} $\Lambda_n \uparrow V$ is a monotone sequence
of finite volumes eventually containing any prescribed finite set;
this is the setting for the true thermodynamic limit.

At present \texttt{AmbientLattice.lean} establishes the types, the
basic lemmas (positivity of $Z_\Lambda$, $|\langle\sigma^A\rangle_\Lambda|
\le 1$, GKS-I on $\Lambda$), the \texttt{Exhaustion} structure, and
the eventually-defined \texttt{correlationAlongExhaustion}.  The
convergence theorem along an exhaustion is not yet proved in this
regime (it would require a compatibility argument across different
$\Lambda$'s).

%======================================================================
\section{How to read this codebase}
%======================================================================

\subsection{File dependency graph}

Read the files in this order. Each file imports only the files above it.

\begin{Verbatim}
Basic.lean                   Spin, Config, IsingParams
  |
Hamiltonian.lean             H(sigma), edgeSpin, spin flip
  |
GibbsMeasure.lean            Z, Gibbs expectation, correlation
  |
  +-- Inequalities/
  |     NonnegCorrelations.lean   HNC property, exp decomposition
  |     GKS.lean                  GKS-I, GKS-II
  |     FKG.lean                  FKG inequality
  |     GHS.lean                  truncated functions, GHS, Cor 4.3.3-5
  |
  +-- Asano.lean                  MultilinPoly, Asano contraction
  |     LeeYang.lean              Lee-Yang (Harcos/Ruelle, uses MultilinPoly only)
  |
  +-- InfiniteVolume.lean         boundedness, monotonicity, convergence
  |     FreeEnergy.lean           free energy, analyticity, Lee-Yang (Ising)
  |     Conditioning.lean         beta-monotonicity, Z bounds, reflection pos.
  |
  +-- PhaseTransition.lean        phase decomposition, mean field, symmetry
  |
  +-- Peierls.lean                Peierls bound, spontaneous magnetization
  |
  +-- Lattice.lean                lattice graph definitions
  |
  +-- ContinuousSpin/Phi4.lean    phi^4 identities (axioms)
  |
  +-- AmbientLattice.lean         ambient V : Type*, Lambda : Finset V, exhaustion
\end{Verbatim}

\subsection{Recommended reading paths}

\paragraph{Path 1: GKS inequalities (simplest).}
\texttt{Basic.lean} $\to$ \texttt{Hamiltonian.lean} $\to$
\texttt{GibbsMeasure.lean} $\to$ \texttt{NonnegCorrelations.lean} $\to$
\texttt{GKS.lean}.
\emph{Key idea:} the Boltzmann weight factors as a product of
$\cosh + \sinh \cdot \sigma_i\sigma_j$ terms, each with non-negative
coefficients (HNC). Fourier--Walsh expansion then gives GKS-I/II.

\paragraph{Path 2: Lee-Yang theorem.}
\texttt{Asano.lean} (for \texttt{MultilinPoly} type) $\to$
\texttt{LeeYang.lean} $\to$ \texttt{FreeEnergy.lean}.

\emph{Key idea:} the Harcos/Ruelle approach defines a polynomial
from a Hermitian matrix, uses a torus identity to relate complementary
coefficients, and applies the maximum modulus principle.
Induction on $n$ eliminates one variable at a time.
Asano contraction is \emph{not} used.

\paragraph{Path 3: GHS and truncated functions.}
\texttt{GHS.lean} (imports \texttt{InfiniteVolume.lean}).

\emph{Key idea:} the Lebowitz inequality (axiom from $\varphi^4$ theory)
bounds the 3-point function, giving $\langle\sigma_i;\sigma_j;\sigma_k\rangle \le 0$.
Spin-flip symmetry at $h=0$ kills odd correlations.

\paragraph{Path 4: Peierls argument.}
\texttt{Peierls.lean}.

\emph{Key idea:} the Hamiltonian--boundary identity
$H = -J(|E| - 2|\partial\sigma|)$ rewrites the Boltzmann weight.
Contour counting gives the Peierls bound.

\subsection{Conventions used throughout}

\begin{itemize}
  \item \textbf{Type variables.} $\iota$ is the site type
    (\texttt{[Fintype $\iota$] [DecidableEq $\iota$]}), $G$ is a
    \texttt{SimpleGraph $\iota$} with \texttt{[Fintype G.edgeSet]}.
  \item \textbf{Field genericity.} Definitions use a generic ordered
    field $K$; main theorems specialize to $\mathbb{R}$.
  \item \textbf{Sign convention.} $H = -J\sum_e s_i s_j - h\sum_i s_i$
    (negative coupling = ferromagnetic).
  \item \textbf{Boltzmann weight.} $w(\sigma) = \exp(-\beta H(\sigma))$
    (\texttt{boltzmannWeight}).
  \item \textbf{Axioms.} Four axioms from $\varphi^4$ measure theory
    are used: \texttt{phi4\_single\_\allowbreak site\_nonneg},
    \texttt{lebowitz\_\allowbreak third/\allowbreak four/\allowbreak inductive}.
    These have complete mathematical proofs but require measure-theoretic
    infrastructure not yet available in Lean/Mathlib.
  \item \textbf{Edge representation.} Edges are \texttt{Sym2 $\iota$}
    (unordered pairs). The edge spin product
    \texttt{edgeSpin $\sigma$ $e$} $= s(\sigma_i) \cdot s(\sigma_j)$
    is defined via \texttt{Sym2.lift}.
\end{itemize}

\subsection{How the key proofs work in Lean}

Many proofs follow a common pattern:
\begin{enumerate}
  \item \textbf{Unfold definitions} to expose the sum
    $Z^{-1} \sum_\sigma (\cdots)$.
  \item \textbf{Factor the Boltzmann weight} using
    $\exp(\alpha s) = \cosh\alpha + s\sinh\alpha$
    (\texttt{exp\_edgeSpin\_decomp}).
  \item \textbf{Expand products} to get a polynomial in spin products
    $\sigma^S$.
  \item \textbf{Apply non-negativity} of each coefficient
    (HNC, \texttt{hasNonnegCorrelations}).
  \item \textbf{Close with} \texttt{nlinarith}, \texttt{linarith},
    or \texttt{Finset.sum\_nonneg}.
\end{enumerate}

%======================================================================
\section{Basic definitions (\texttt{Basic.lean})}

\par\noindent\textbf{Prerequisites:} None (this is the root file).

\textbf{Role:} Defines the alphabet of the Ising model --- spin values,
configurations, parameters. Everything else builds on these types.

\textbf{What to look for:} the \texttt{Spin} type and its arithmetic
(\texttt{toSign}, \texttt{flip}, \texttt{sign}), and the
\texttt{Ferromagnetic} structure ($J \ge 0$, $h \ge 0$, $\beta > 0$).

\begin{definition}[Spin]
$\mathrm{Spin} = \{\mathrm{up}, \mathrm{down}\}$ with sign function
$s : \mathrm{Spin} \to \mathbb{Z}$, $s(\mathrm{up})=+1$,
$s(\mathrm{down})=-1$.
\end{definition}

\noindent\textbf{Lean:} \texttt{inductive Spin}, \texttt{Spin.toSign}

\begin{definition}[Configuration]
A configuration on a type $\iota$ is a function
$\sigma : \iota \to \mathrm{Spin}$.
\end{definition}

\noindent\textbf{Lean:} \texttt{abbrev Config ($\iota$) := $\iota$ $\to$ Spin}

\begin{definition}[Ising parameters]
Over an ordered field $K$: coupling $J : K$, external field $h : K$,
inverse temperature $\beta : K$.
Ferromagnetic condition: $J \geq 0$, $h \geq 0$, $\beta > 0$.
\end{definition}

\noindent\textbf{Lean:} \texttt{structure IsingParams},
\texttt{structure Ferromagnetic}

\section{Hamiltonian (\texttt{Hamiltonian.lean})}

\par\noindent\textbf{Prerequisites:} \texttt{Basic.lean}.

\textbf{Role:} Defines the energy function $H(\sigma)$ and
the edge spin product. The spin flip symmetry theorem lives here.

\textbf{What to look for:} \texttt{edgeSpin} (the $\pm 1$ product
for an edge, defined via \texttt{Sym2.lift}) and
\texttt{hamiltonian\_flip\_eq} (the $h=0$ symmetry).

\begin{definition}[Ising Hamiltonian]
For a simple graph $G$ on $\iota$,
\[
  H_G(\sigma) = -J \sum_{\{i,j\} \in E(G)} s(\sigma_i)\,s(\sigma_j)
  - h \sum_{i \in \iota} s(\sigma_i).
\]
\end{definition}

\noindent\textbf{Lean:} \texttt{interactionEnergy} $+$
\texttt{externalFieldEnergy} $=$ \texttt{hamiltonian}

Key auxiliary definitions (in \texttt{Basic.lean}):
\begin{itemize}
  \item \texttt{Spin.flip}: $\mathrm{up} \leftrightarrow \mathrm{down}$
  \item \texttt{Spin.sign $K$ $s$}: cast $s(\cdot)$ into the field $K$
  \item \texttt{Config.flip}, \texttt{Config.flipAt}: global and single-site spin flip
\end{itemize}
Key auxiliary definitions (in \texttt{Hamiltonian.lean}):
\begin{itemize}
  \item \texttt{edgeSpin}: $s(\sigma_i) \cdot s(\sigma_j)$ for an edge
    $\{i,j\}$, defined via \texttt{Sym2.lift}
\end{itemize}

\begin{lemma}[Spin flip symmetry --- \texttt{hamiltonian\_flip\_eq}]
If $h = 0$, then $H_G(\bar{\sigma}) = H_G(\sigma)$ where
$\bar{\sigma}_i = \mathrm{flip}(\sigma_i)$.
\end{lemma}

\begin{proof}
Each edge factor satisfies $(-s_i)(-s_j) = s_i s_j$
(\texttt{edgeSpin\_flip}), and the field term vanishes since $h=0$.
\end{proof}

\section{Partition function and Gibbs measure (\texttt{GibbsMeasure.lean})}

\par\noindent\textbf{Prerequisites:} \texttt{Hamiltonian.lean}.

\textbf{Role:} Defines the three central objects: partition function $Z$,
Gibbs expectation $\langle\cdot\rangle$, and correlation function
$\langle\sigma^A\rangle$. Also contains the spin product algebra
($\sigma^A \cdot \sigma^C = \sigma^{A \triangle C}$, $(\sigma^A)^2 = 1$).

\textbf{What to look for:} \texttt{partitionFunction\_pos} ($Z > 0$)
and \texttt{spinProduct\_mul} (the key algebraic identity for GKS proofs).

\begin{definition}[Partition function]
$Z_G = \sum_{\sigma : \iota \to \mathrm{Spin}} \exp(-\beta\, H_G(\sigma))$.
\end{definition}

\begin{theorem}[$Z > 0$ --- \texttt{partitionFunction\_pos}]
$Z_G > 0$.
\end{theorem}

\begin{proof}
Each summand $\exp(\cdot) > 0$ (\texttt{Real.exp\_pos}), and the
index set $\iota \to \mathrm{Spin}$ is nonempty
(\texttt{Finset.univ\_nonempty}).
Apply \texttt{Finset.sum\_pos}.
\end{proof}

\begin{definition}[Gibbs expectation, correlation]
$\langle F \rangle = Z^{-1} \sum_\sigma F(\sigma)\,e^{-\beta H}$, \quad
$\langle \sigma^A \rangle = \langle \prod_{i \in A} s(\sigma_i) \rangle$.
\end{definition}

\noindent\textbf{Lean:} \texttt{gibbsExpectation},
\texttt{spinProduct}, \texttt{correlation}

\subsection{Spin product algebra}

\begin{lemma}[\texttt{spinProduct\_sq}]
$(\sigma^A)^2 = 1$ for any $A$.
\end{lemma}

\begin{lemma}[\texttt{spinProduct\_mul}]
$\sigma^A \cdot \sigma^C = \sigma^{A \bigtriangleup C}$.
\end{lemma}

\begin{proof}
Write $\sigma^A = \prod_{i \in \iota}
(\text{if } i \in A \text{ then } s_i \text{ else } 1)$
(using \texttt{Finset.prod\_filter}).
Multiply pointwise (\texttt{Finset.prod\_mul\_distrib}).
At each site: if $i \in A \cap C$, $s_i \cdot s_i = 1$ (by
\texttt{Spin.toSign\_sq}); otherwise the factor is $s_i$ or $1$,
matching $A \bigtriangleup C$.
\end{proof}

\section{GKS-I proof (\texttt{Inequalities/GKS.lean})}

\par\noindent\textbf{Prerequisites:} \texttt{GibbsMeasure.lean},
\texttt{NonnegCorrelations.lean}.

\textbf{Role:} Proves the first and second Griffiths inequalities.

\textbf{What to look for:} The proof reduces GKS to the HNC property
of the Boltzmann weight (\texttt{boltzmannWeight\_hasNonnegCorrelations}).
Once HNC is established, GKS-I is a one-liner.

\subsection{Reduction to non-negative correlations}

\begin{theorem}[GKS-I --- \texttt{gks\_first}]
For ferromagnetic parameters, $\langle \sigma^A \rangle \geq 0$
for all $A$.
\end{theorem}

The proof factors as:
\begin{enumerate}
  \item Since $Z > 0$, it suffices to show the numerator
    $\sum_\sigma \sigma^A e^{-\beta H} \geq 0$
    (\texttt{gibbsExpectation\_nonneg\_of\_numerator\_nonneg}).
  \item This follows from a property of $e^{-\beta H}$ called
    \emph{non-negative correlations}
    (\texttt{gks\_numerator\_nonneg}).
\end{enumerate}

\subsection{Non-negative correlations (HNC)}

\begin{definition}[\texttt{HasNonnegCorrelations}]
A function $f : (\iota \to \mathrm{Spin}) \to \mathbb{R}$ has
\textbf{non-negative correlations} (HNC) if
$\sum_\sigma \sigma^S \cdot f(\sigma) \geq 0$ for every
$S \subseteq \iota$.
\end{definition}

\begin{lemma}[Base case --- \texttt{hasNonnegCorrelations\_one}]
$f \equiv 1$ has HNC.
\end{lemma}

\begin{proof}
For $S = \emptyset$: $\sum_\sigma 1 = 2^{|\iota|} > 0$.
For $S \neq \emptyset$: $\sum_\sigma \sigma^S = 0$.

The vanishing is proved via the involution
$\mathrm{flipAt}(j)$ for any $j \in S$
(\texttt{sum\_config\_spinProduct\_eq\_zero}).
Each pair $(\sigma, \mathrm{flipAt}(j, \sigma))$ contributes
$\sigma^S + (-\sigma^S) = 0$ because flipping at $j \in S$ negates
the spin product (\texttt{spinProduct\_flipAt\_neg}).
The involution is applied via \texttt{Finset.sum\_ninvolution}.
\end{proof}

\begin{lemma}[Preservation --- \texttt{hasNonnegCorrelations\_mul}]
\label{lem:mul}
If $f$ has HNC and $a, b \geq 0$, then
$\sigma \mapsto f(\sigma)(a + b\,\sigma^C)$ has HNC.
\end{lemma}

\begin{proof}
\[
  \sum_\sigma \sigma^S f(\sigma)(a + b\,\sigma^C)
  = a \underbrace{\sum_\sigma \sigma^S f(\sigma)}_{\geq 0}
  + b \underbrace{\sum_\sigma \sigma^{S \bigtriangleup C} f(\sigma)}_{\geq 0}
\]
using $\sigma^S \sigma^C = \sigma^{S \bigtriangleup C}$
(\texttt{spinProduct\_mul}).
\end{proof}

\begin{lemma}[Iterated product --- \texttt{hasNonnegCorrelations\_finset\_prod},
\texttt{hasNonnegCorrelations\_mul\_prod}]
A product of factors $(a_k + b_k \sigma^{C_k})$ with $a_k, b_k \geq 0$
has HNC. More generally, multiplying an HNC function by such a product
preserves HNC.
\end{lemma}

\begin{proof}
By induction on the index set, applying Lemma~\ref{lem:mul} at each step.
\end{proof}

\subsection{Factoring the Boltzmann weight}

\begin{lemma}[\texttt{exp\_sign\_decomp}, \texttt{exp\_edgeSpin\_decomp}]
For $x \in \{+1, -1\}$ and $\alpha \geq 0$:
$e^{\alpha x} = \cosh\alpha + x\,\sinh\alpha$,
with $\cosh\alpha > 0$ and $\sinh\alpha \geq 0$.
\end{lemma}

\begin{lemma}[\texttt{bwFactored\_eq\_boltzmannWeight}]
\[
  e^{-\beta H_G(\sigma)}
  = \prod_{\{i,j\} \in E(G)} (\cosh(\beta J) + \sinh(\beta J)\,s_i s_j)
  \cdot \prod_{i \in \iota} (\cosh(\beta h) + \sinh(\beta h)\,s_i).
\]
\end{lemma}

\begin{proof}
$-\beta H = \beta J \sum_e s_i s_j + \beta h \sum_i s_i$.
Apply $\exp(\text{sum}) = \prod\exp$ (\texttt{Real.exp\_sum}),
then decompose each factor.
\end{proof}

\subsection{Conclusion}

\begin{theorem}[\texttt{boltzmannWeight\_hasNonnegCorrelations}]
$e^{-\beta H_G}$ has HNC for ferromagnetic parameters.
\end{theorem}

\begin{proof}
The factored form is a product of $(a+b\,\sigma^C)$ terms with
$a,b \geq 0$. Start from $f \equiv 1$ (HNC by base case),
multiply by each edge factor, then by each site factor. HNC is
preserved at each step by Lemma~\ref{lem:mul}.
\end{proof}

\noindent
GKS-I follows: $\langle \sigma^A \rangle
= Z^{-1}\sum_\sigma \sigma^A e^{-\beta H} \geq 0$
since $Z > 0$ and the numerator $\geq 0$ by HNC.

\section{GKS-II proof (\texttt{Inequalities/GKS.lean})}

\noindent\textbf{What to look for:} The proof introduces the
\emph{doubled system} --- a product of two independent copies of the
Ising model --- and a $45^\circ$ rotation $(t, q) = (\sigma + \tau,
\sigma - \tau)/\sqrt{2}$. The key step is showing that the doubled
Boltzmann weight has the HNC property in the $(t, q)$ variables.

The second Griffiths inequality is proved by the duplicate variable trick,
following Friedli--Velenik~\cite{friedli-velenik}, Theorem~3.49~(3.55),
pp.~127--128.

\begin{theorem}[GKS-II --- \texttt{gks\_second}]
For ferromagnetic parameters,
$\langle\sigma^A\sigma^B\rangle \geq \langle\sigma^A\rangle\langle\sigma^B\rangle$
for all $A, B$.
\end{theorem}

\subsection{Proof outline}

\begin{enumerate}
\item Introduce an independent copy $\chi$ of $\sigma$. Then
  \[
    Z^2\bigl(\langle\sigma^{A\triangle B}\rangle
      - \langle\sigma^A\rangle\langle\sigma^B\rangle\bigr)
    = \sum_{\omega,\omega'} \omega^A(\omega^B - \omega'^B)\,w(\omega)\,w(\omega').
  \]
  (\texttt{duplicateSum\_eq})

\item Change variables: $t_i = \omega_i \cdot \omega'_i$
  (\texttt{Spin.mul}). For fixed $\omega$, this is a bijection on
  $\mathrm{Config}\;\iota$. Under this substitution:
  \begin{itemize}
    \item $\omega^A(\omega^B - \omega'^B) = \omega^{A\triangle B}(1 - t^B)$
    \item $w(\omega)\,w(\omega\cdot t) = \mathrm{modifiedWeight}(t, \omega)$
      where each coupling $K_C$ is replaced by $K_C(1 + t^C) \geq 0$.
  \end{itemize}
  After swapping sums:
  \[
    \sum_t (1 - t^B) \sum_\omega \omega^{A\triangle B}\,
      \mathrm{modifiedWeight}(t, \omega).
  \]
  (\texttt{duplicateSum\_eq\_changed})

\item For each fixed $t$:
  \begin{itemize}
    \item $(1 - t^B) \geq 0$ since $t^B \in \{-1,+1\}$
      (\texttt{one\_sub\_spinProduct\_nonneg}).
    \item The inner sum $\geq 0$ by
      \texttt{hasNonnegCorrelations\_edge\_site\_product} applied to
      the modified couplings $K_C(1 + t^C)$
      (\texttt{modifiedWeight\_nonneg\_corr}).
  \end{itemize}

\item Since $Z^2 > 0$ and the duplicate sum $\geq 0$, GKS-II follows.
\end{enumerate}

\subsection{Key new definitions}

\begin{itemize}
  \item \texttt{Spin.mul}: spin multiplication ($\mathrm{up}$ acts as
    identity, $\mathrm{down}$ acts as flip).
    $\mathrm{toSign}(a \cdot b) = \mathrm{toSign}(a) \cdot \mathrm{toSign}(b)$.
  \item \texttt{modifiedWeight}: Boltzmann weight with couplings
    $K_C(1 + t^C)$.
  \item \texttt{hasNonnegCorrelations\_edge\_site\_product}: shared
    lemma for GKS-I and GKS-II (product of $\exp(K \cdot \text{spin})$
    factors with $K \geq 0$ has HNC).
\end{itemize}

\section{Numerical tests (\texttt{test/IsingModel/GKSTest.lean})}

A Float-based reference implementation verifies:
\begin{itemize}
  \item $Z > 0$ for 4 graph types
  \item Spin flip symmetry ($h = 0$)
  \item GKS-I for 6 parameter sets (2--4 sites, all subsets)
  \item GKS-II ($\langle\sigma^A\sigma^B\rangle \geq
    \langle\sigma^A\rangle\langle\sigma^B\rangle$)
    for 6 parameter sets (2--4 sites, all subset pairs)
  \item GKS-I and GKS-II violation for anti-ferromagnetic $J < 0$
\end{itemize}

\section{FKG inequality (\texttt{Inequalities/FKG.lean})}

\par\noindent\textbf{Prerequisites:} \texttt{GKS.lean}.

\textbf{Role:} Proves $\langle fg \rangle \ge \langle f \rangle\langle g \rangle$
for monotone $f, g$.

\textbf{What to look for:} The Holley criterion approach. The proof
shows that the Gibbs measure satisfies the FKG lattice condition
by reducing to GKS-II.

The FKG inequality is proved by applying Mathlib's
\texttt{Combinatorics.\allowbreak SetFamily.\allowbreak FourFunctions.\allowbreak fkg}
to the Ising Boltzmann weight, following
Friedli--Velenik~\cite{friedli-velenik}, Theorem~3.21
(p.~98) and Theorem~3.50 (pp.~128--130).

\begin{theorem}[\texttt{fkg\_ising}]
For ferromagnetic parameters and monotone nondecreasing nonneg
$f, g : \mathrm{Config}\;\iota \to \mathbb{R}$,
$\langle fg \rangle \geq \langle f \rangle \langle g \rangle$.
\end{theorem}

\subsection{Proof outline}
\begin{enumerate}
\item Define a \texttt{LinearOrder} on \texttt{Spin}: $\mathrm{down} < \mathrm{up}$
  ($-1 < +1$). The product order on $\mathrm{Config}\;\iota = \iota \to
  \mathrm{Spin}$ makes it a \texttt{DistribLattice}.
\item Prove log-supermodularity of the Boltzmann weight
  (\texttt{boltzmannWeight\_log\_supermodular}):
  $w(\sigma)w(\sigma') \leq w(\sigma \vee \sigma')w(\sigma \wedge \sigma')$.
  This reduces to the edge inequality
  $\sigma_i\sigma_j + \sigma'_i\sigma'_j \leq
  (\sigma_i \vee \sigma'_i)(\sigma_j \vee \sigma'_j) +
  (\sigma_i \wedge \sigma'_i)(\sigma_j \wedge \sigma'_j)$
  (\texttt{spin\_edge\_supermodular}, verified by exhaustive case analysis).
\item Apply Mathlib's \texttt{fkg} with $\mu = w$, then convert from
  unnormalized sums to \texttt{gibbsExpectation} form.
\end{enumerate}

\section{Asano Contraction (\texttt{Asano.lean})}

\par\noindent\textbf{Prerequisites:} None (self-contained complex analysis).

\textbf{Role:} Defines \texttt{MultilinPoly} (multilinear polynomials
over $\mathbb{C}$) and the Asano contraction operation. This is the
algebraic engine behind the Lee-Yang theorem.

\textbf{What to look for:} \texttt{MultilinPoly.eval} and
\texttt{asanoContract}. The key theorem is that contraction
preserves nonvanishing on the unit polydisk.

The Asano contraction merges two variables in a multilinear polynomial,
keeping only the ``both present'' and ``both absent'' parts.
It is an independent result in algebraic combinatorics; our Lee-Yang proof
(Section~\ref{sec:lee-yang}) uses the Harcos/Ruelle approach instead.

Reference: \cite[Proposition~3.44, pp.~123--124]{friedli-velenik}.

\begin{definition}[Multilinear polynomial]
A multilinear polynomial over $\mathbb{C}$ with variables indexed by a
finite type $\iota$ is a function $p : \mathcal{P}(\iota) \to \mathbb{C}$,
evaluated as $p(z) = \sum_{X \subseteq \iota} p(X) \prod_{i \in X} z_i$.
In Lean: \texttt{Finset $\iota$ $\to$ $\mathbb{C}$}.
\end{definition}

\begin{definition}[Asano contraction]
Given $p$ and distinct $i, j \in \iota$, the contraction of $j$ into $i$
is $q(X) = p(X \cup \{j\})$ if $i \in X$ and $j \notin X$;
$q(X) = p(X)$ if $i \notin X$ and $j \notin X$;
$q(X) = 0$ if $j \in X$.
\end{definition}

\begin{theorem}[Asano contraction preserves non-vanishing]
\label{thm:asano}
If $p$ does not vanish on the open unit polydisk
$\{z : |z_k| < 1\;\forall k\}$, then neither does its Asano contraction.
Reference: \cite[Proposition~3.44, pp.~123--124]{friedli-velenik}.
\end{theorem}

The proof uses a bilinear non-vanishing lemma:

\begin{lemma}[Bilinear non-vanishing]
If $f(z,w) = azw + bw + cz + d \neq 0$ for all $|z|, |w| < 1$,
then $az + d \neq 0$ for all $|z| < 1$.
\end{lemma}

\begin{proof}[Proof sketch]
Show $\|a\| \leq \|d\|$ via the parallelogram law.
For $|z'| < 1$, the affine function $w \mapsto (az'+b)w + (cz'+d)$
has no zero with $|w| < 1$, giving $\|az'+b\| \leq \|cz'+d\|$.
Squaring and averaging at $z' = \pm t$ (real) via the parallelogram law:
$\mathrm{normSq}(a) \cdot t^2 + \mathrm{normSq}(b) \leq
\mathrm{normSq}(c) \cdot t^2 + \mathrm{normSq}(d)$.
Combining the $z$- and $w$-directions and applying a linear function argument
on $[0,1)$ yields $\mathrm{normSq}(a) \leq \mathrm{normSq}(d)$,
hence $\|a\| \leq \|d\|$. Then $az + d = 0$ would require
$|z| = |d/a| \geq 1$, contradicting $|z| < 1$.
\end{proof}

The contraction theorem follows by decomposing the polynomial evaluation
$p(z[i \to u, j \to w])$ as a bilinear form $auw + bw + cu + d$ in $(u, w)$,
where the contraction's evaluation equals $a \cdot z_i + d$.
The decomposition uses a product extraction lemma
(\texttt{prod\_upd}) and a sum reindexing via the involution
$X \mapsto X \triangle \{j\}$ (\texttt{Fintype.sum\_equiv}).

\begin{remark}
When Mathlib provides a native implementation of Asano contraction or
multilinear polynomials, the definitions here should be replaced.
\end{remark}

\section{Lee-Yang circle theorem (\texttt{LeeYang.lean})}

\par\noindent\textbf{Prerequisites:} \texttt{Asano.lean} (for the
\texttt{MultilinPoly} type only; Asano contraction is \emph{not} used).

\textbf{Role:} Proves that the Ising partition polynomial has no
zeros on the open unit polydisk via the Harcos/Ruelle approach.

\textbf{What to look for:} The proof uses a Hermitian matrix formulation,
induction on $n$, and the maximum modulus principle --- not Asano contraction.
Key steps: (1) define \texttt{leeYangPoly} for Hermitian $A$ with
$|a_{ij}| \le 1$; (2) torus identity relating $\alpha$ and
$\overline{\beta}$; (3) induction eliminating the last variable;
(4) maximum modulus to bound $|\beta/\alpha|$ on the polydisk.
\texttt{lee\_yang\_circle} is the final theorem.
\label{sec:lee-yang}

The Lee-Yang circle theorem states that the Ising partition polynomial
has no zeros on the open unit polydisk. Our formalization follows the
Harcos/Ruelle approach~\cite{harcos, ruelle}.

\begin{definition}[Lee-Yang polynomial]
For an $n\times n$ matrix $A$ over $\mathbb{C}$, define
\[
  f_A(z) = \sum_{S\subseteq [n]}
    \Bigl(\prod_{i\in S}\prod_{j\notin S} a_{ij}\Bigr)\prod_{k\in S} z_k.
\]
\end{definition}

\begin{theorem}[Ruelle~{\cite[Theorem~1, p.~589]{ruelle}};
  Harcos~{\cite[Theorem~3, p.~2]{harcos}};
  \texttt{leeYangPoly\_nonvanishing}]
If $A$ is Hermitian with $|a_{ij}|\le 1$, then $f_A(z)\ne 0$
for all $z$ with $|z_k|<1$.
\end{theorem}

The proof is by induction on $n$. The inductive step decomposes
$f_A(z)=\beta + z_n\cdot\alpha$ and shows $\|\alpha\|\le\|\beta\|$
via:
\begin{enumerate}
\item \textbf{Torus identity} (\texttt{torus\_identity}):
  On $|v_k|=1$,
  $\alpha(v) = (\prod v_k)\cdot\overline{\beta(v)}$,
  hence $\|\alpha\|=\|\beta\|$.
\item \textbf{Approximation}: Scale the last column by $t<1$
  so that $\beta_t\ne 0$ on the closed polydisk.
\item \textbf{Iterated max modulus} (\texttt{iterated\_ratio}):
  By single-variable maximum modulus
  (\texttt{one\_var\_max\_ratio}, via \texttt{DiffContOnCl}),
  iterate over each variable to get $\|\alpha_t/\beta_t\|\le 1$.
\item \textbf{Limit $t\to 1$}: Both sides are continuous in $t$;
  $\{t:\|\alpha_t\|\le\|\beta_t\|\}$ is closed and contains
  $[0,1)$, hence contains~$1$.
\end{enumerate}

The Ising application (\texttt{lee\_yang\_circle}) constructs a
Hermitian matrix from the edge list and coupling constants, reducing
the combinatorial Ising model to the matrix formulation.

\section{Infinite volume limit (\texttt{InfiniteVolume.lean})}

\par\noindent\textbf{Prerequisites:} \texttt{GKS.lean}.

\textbf{Role:} Proves boundedness ($|\langle\sigma^A\rangle| \le 1$),
monotonicity in $J$ and $h$, and convergence of correlations.
Also contains the Walsh--Fourier inversion theorem.

\textbf{What to look for:} \texttt{abs\_correlation\_le\_one} (used
everywhere), \texttt{correlation\_monotone\_J/H} (used for free energy
monotonicity), and the covariance non-negativity lemma
\texttt{cov\_hnc\_boltzmann\_nonneg}.

The convergence of correlation functions as the lattice grows to infinity.
Reference: \cite[{\S4.2, pp.~58--59}]{glimm-jaffe}.

\begin{proposition}[\texttt{abs\_correlation\_le\_one}, Proposition 4.2.2, p.~58]
For the Ising model, $|\langle\sigma^A\rangle| \le 1$.
\end{proposition}

\begin{proof}
$|\sigma^A| = 1$ (since $\sigma_i = \pm 1$), so
$|\langle\sigma^A\rangle| \le \langle|\sigma^A|\rangle = 1$.
\end{proof}

\begin{theorem}[\texttt{cov\_hnc\_boltzmann\_nonneg};
  cf.\ \cite{glimm-jaffe}, {\S4.2, p.~58}]
\label{thm:cov-hnc}
For an HNC function $f$ and ferromagnetic Boltzmann weight $w$,
\[
  \Bigl(\sum_\sigma \sigma^B f(\sigma) w(\sigma)\Bigr)
  \Bigl(\sum_\sigma w(\sigma)\Bigr)
  \ge
  \Bigl(\sum_\sigma \sigma^B w(\sigma)\Bigr)
  \Bigl(\sum_\sigma f(\sigma) w(\sigma)\Bigr).
\]
\end{theorem}

\begin{proof}
Fourier expand $f(\sigma) = \sum_S \hat{c}_S \sigma^S$ where
$\hat{c}_S = |\iota|^{-1} \sum_\tau \sigma^S(\tau) f(\tau) \ge 0$
by HNC (\texttt{walsh\_fourier\_inversion}).
Then $\sigma^B f(\sigma) = \sum_S \hat{c}_S \sigma^{B \triangle S}$
using $\sigma^B \sigma^S = \sigma^{B \triangle S}$ (\texttt{spinProduct\_mul}).
The covariance becomes $\sum_S \hat{c}_S \cdot \text{bracket}_S$ where
\[
  \text{bracket}_S =
  \Bigl(\sum_\sigma \sigma^{B\triangle S} w\Bigr)\Bigl(\sum_\sigma w\Bigr)
  - \Bigl(\sum_\sigma \sigma^B w\Bigr)\Bigl(\sum_\sigma \sigma^S w\Bigr).
\]
Since $\hat{c}_S \ge 0$ and $\text{bracket}_S = Z^2(\langle\sigma^{B\triangle S}\rangle
- \langle\sigma^B\rangle\langle\sigma^S\rangle) \ge 0$ by GKS-II
(\texttt{gks\_second}; \cite[Theorem~3.49, p.~127]{friedli-velenik}),
each term is non-negative.
\end{proof}

\begin{proposition}[\texttt{correlation\_monotone\_J}, Proposition 4.2.1, p.~58]
$\langle\sigma^B\rangle$ is monotone increasing in $J$ on $[0,\infty)$.
\end{proposition}

\begin{proof}
For $0 \le J_1 \le J_2$, define the reweighting factor
$R(\sigma) = \prod_e \exp(\beta(J_2 - J_1)\,\text{edgeSpin}(\sigma,e))$.
Then $\exp(-\beta H_{J_2}(\sigma)) = R(\sigma)\exp(-\beta H_{J_1}(\sigma))$
(\texttt{correlation\_reweighting\_nonneg}), and $R$ has HNC by
\texttt{hasNonnegCorrelations\_edge\_site\_product} since all edge
coefficients $\beta(J_2 - J_1) \ge 0$.
The algebraic identity
$\langle\sigma^B\rangle_{J_2} - \langle\sigma^B\rangle_{J_1}
= [\text{num}_2\,\text{den}_1 - \text{num}_1\,\text{den}_2]
  / (\text{den}_1\,\text{den}_2)$
reduces to a covariance of $R$ with $\sigma^B$ under the $J_1$ Boltzmann
weight, which is $\ge 0$ by Theorem~\ref{thm:cov-hnc}.
\end{proof}

\subsection{Walsh orthogonality and Fourier inversion}

The spin products $\{\sigma^S : S \subseteq \iota\}$ form an orthogonal
basis for functions on $\mathrm{Config}\;\iota = \iota \to \{+1,-1\}$.

\begin{lemma}[\texttt{walsh\_orthogonality}]
For $S \ne T$, $\sum_\sigma \sigma^S \sigma^T = 0$.
\end{lemma}

\begin{proof}
$\sigma^S \sigma^T = \sigma^{S \triangle T}$ by \texttt{spinProduct\_mul}.
Since $S \ne T$, $S \triangle T \ne \emptyset$, so
$\sum_\sigma \sigma^{S \triangle T} = 0$ by the flip involution
(\texttt{sum\_config\_spinProduct\_eq\_zero}).
\end{proof}

\begin{lemma}[\texttt{walsh\_normalization}]
$\sum_\sigma (\sigma^S)^2 = 2^{|\iota|}$.
\end{lemma}

\begin{proof}
$(\sigma^S)^2 = \sigma^{S \triangle S} = \sigma^\emptyset = 1$.
So $\sum_\sigma 1 = |\mathrm{Config}\;\iota| = 2^{|\iota|}$.
\end{proof}

\begin{theorem}[\texttt{walsh\_completeness}]
$\sum_S \sigma^S(\tau) \cdot \sigma^S(\sigma) =
\begin{cases} 2^{|\iota|} & \text{if } \sigma = \tau, \\
0 & \text{if } \sigma \ne \tau. \end{cases}$
\end{theorem}

\begin{proof}
Define $\omega_i = \sigma_i \cdot \tau_i$ (componentwise spin
multiplication). Then $\sigma^S(\tau) \cdot \sigma^S(\sigma)
= \omega^S$ by \texttt{Spin.toSign\_mul}. The sum
$\sum_S \omega^S = \prod_i (1 + \omega_i)$ by
\texttt{Finset.prod\_add\_one}. If $\sigma = \tau$ then each
$\omega_i = +1$ so each factor is $2$; if $\sigma \ne \tau$
then $\exists i$ with $\omega_i = -1$ so the product vanishes.
\end{proof}

\begin{theorem}[\texttt{walsh\_fourier\_inversion}]
Any $f : \mathrm{Config}\;\iota \to \mathbb{R}$ has the expansion
$f(\sigma) = \sum_S \hat{c}_S \sigma^S$ where
$\hat{c}_S = 2^{-|\iota|} \sum_\tau \sigma^S(\tau) f(\tau)$.
\end{theorem}

\begin{proof}
Substitute the expansion and apply Walsh completeness:
\[
  \sum_S \hat{c}_S \sigma^S(\sigma)
  = 2^{-|\iota|} \sum_\tau f(\tau)
    \underbrace{\sum_S \sigma^S(\tau) \sigma^S(\sigma)}_{= \, 2^{|\iota|} [\tau = \sigma]}
  = f(\sigma).
\]
This uses \texttt{Finset.sum\_comm} to swap the order of summation
and \texttt{walsh\_completeness} for the Kronecker delta.
\end{proof}

\begin{theorem}[\texttt{correlation\_convergent},
  Theorem 4.2.3, \cite{glimm-jaffe}, p.~59]
For $h \ge 0$ and $\beta > 0$, the sequence
$n \mapsto \langle\sigma^B\rangle_{(G, n, h, \beta)}$ converges.
The limit equals $\sup_n \langle\sigma^B\rangle_{(G, n, h, \beta)}$.
\end{theorem}

\begin{proof}
The sequence $n \mapsto \langle\sigma^B\rangle_{(G, n, h, \beta)}$ is:
\begin{itemize}
  \item monotone increasing by Proposition~4.2.1
    (\texttt{correlation\_monotone\_J});
  \item bounded above by $1$ since
    $\langle\sigma^B\rangle \le |\langle\sigma^B\rangle| \le 1$
    (\texttt{correlationJ\_le\_one} via \texttt{abs\_correlation\_le\_one}).
\end{itemize}
By the monotone convergence theorem (\texttt{tendsto\_atTop\_ciSup}),
the sequence converges to $\sup_n \langle\sigma^B\rangle_{(G, n, h, \beta)}$.
\end{proof}

Auxiliary results:
\begin{itemize}
  \item \texttt{correlationJ\_nonneg}: $\langle\sigma^B\rangle \ge 0$
    for $J \ge 0$ (from GKS-I).
  \item \texttt{correlationJ\_le\_one}: $\langle\sigma^B\rangle \le 1$
    (from $|\langle\sigma^B\rangle| \le 1$).
\end{itemize}

$\mathbb{Z}^d$ direct wrappers
(\texttt{Concrete/LatticeGraphCorrelation.lean}):
\texttt{correlationJ\_\{nonneg,le\_one\}\_latticeGraph} and
\texttt{correlation\_convergent\{,\_h,\_beta\}\_latticeGraph}
are thin pass-throughs of the above at the
\texttt{Ambient.inducedGraph (latticeGraph d)\,$\Lambda$}
subgraph on any finite $\Lambda \subset \mathbb{Z}^d$.
Only \texttt{correlation\_convergent\_latticeGraph} is the direct
GJ~\S4.2 Thm.~4.2.3 specialization ($J \to \infty$ along $\mathbb{N}$);
\texttt{correlation\_convergent\_h\_latticeGraph} is its
Prop.~4.2.4 companion ($h \to \infty$), and
\texttt{correlation\_convergent\_beta\_latticeGraph} is a repo-internal
monotone-bounded extension in $\beta \to \infty$ (not in GJ).

Parallel $\mathbb{Z}^d$ wrappers for the derived \S5 quantities
(\texttt{Concrete/LatticeGraphCorrelation/MagnetizationConvergent.lean},
narrowed out of \texttt{Magnetization.lean} in PR~\#2030):
\texttt{magnetization\_convergent\_\{J,h,beta\}\_latticeGraph},
\texttt{truncated2\_convergent\_\{J,h,beta,subgraph\}\_latticeGraph},
\texttt{susceptibility\_convergent\_subgraph\_latticeGraph}, and
\texttt{magnetization\_total\_convergent\_subgraph\_latticeGraph}
pass through the abstract
\texttt{IsingModel.\{magnetization,truncated2,susceptibility,magnetization\_total\}\_convergent\_*}
(\texttt{IsingModel/PhaseTransition.lean}). The $J,h,\beta$ wrappers
fix the graph to \texttt{Ambient.inducedGraph (latticeGraph d)\,$\Lambda$};
the \texttt{\_subgraph} wrappers only fix $\iota = \operatorname{\uparrow}\Lambda$
(the subgraph parameter $G_n$ on the $\Lambda$-induced type remains
unconstrained and is \emph{not} forced to be $\le$ \texttt{latticeGraph d}).

Susceptibility family (GJ \S5.3, pp.~77–80) at $\mathbb{Z}^d$
(\texttt{Concrete/LatticeGraphCorrelation/MagnetizationSusceptibility.lean},
narrowed out of \texttt{Magnetization.lean} in PR~\#2004):
\texttt{susceptibility\_\{apply,nonneg,J\_zero,beta\_zero\}\_latticeGraph}
and \texttt{susceptibility\_convergent\_\{J,h,beta\}\_latticeGraph}
are thin pass-throughs of the abstract
\texttt{IsingModel.susceptibility\_*} on
\texttt{Ambient.inducedGraph (latticeGraph d)\,$\Lambda$}. The
$J=0$ closed form reads $\chi_i = t(1-t)$ with $t = \tanh(\beta h)$,
inheriting the Finset-based \texttt{truncated2} diagonal caveat
($\{i,i\}$ collapses to $\{i\}$, so the term is $t - t^2$ rather
than the physical $1 - t^2$). In the same batch,
\texttt{eta\_nonneg\_finite\_vol\_latticeGraph} is the
finite-volume slice of GJ Thm~17.7.1 ($\eta \ge 0$) at
$\mathbb{Z}^d$, derived from GKS-II via
\texttt{truncated2\_nonneg}.

{\raggedright
The concrete $\mathbb{Z}^d$ Lambda-layer susceptibility regularity and
parameter-direction convergence wrappers now live in
\path{IsingModel/Concrete/LatticeGraphCorrelation/SusceptibilityLambda.lean}.
\par}

\begin{remark}
The original Glimm--Jaffe formulation uses lattice growth
$\Lambda \uparrow \mathbb{Z}^d$. In the finite-volume setting,
this reduces to increasing coupling constants from $0$ to $J$.
The lattice-growth version ($\Lambda \uparrow$) is formalized
below via the subgraph order on a fixed ambient lattice.
\end{remark}

\subsection{Lattice-growth version of Theorem 4.2.3}

The Glimm--Jaffe proof of Theorem~4.2.3 reads
``increasing $\Lambda$ amounts to increasing certain $J_A$,
and so $\langle\xi^B\rangle$ is monotone increasing in $\Lambda$''.
We formalize this by taking a fixed ambient \texttt{Fintype} $\iota$
and working with subgraphs $G_1 \le G_2$ of the interaction graph:
new edges $e \in E(G_2) \setminus E(G_1)$ correspond exactly to the
coupling $J_A$ going from $0$ to $\beta J$.

\begin{theorem}[\texttt{correlation\_monotone\_subgraph},
  Theorem 4.2.3 (lattice version), \cite{glimm-jaffe}, p.~59]
For a ferromagnetic Ising model on the ambient lattice $\iota$
and subgraphs $G_1 \le G_2$,
$\langle \sigma^A\rangle_{G_1} \le \langle\sigma^A\rangle_{G_2}$.
\end{theorem}

\begin{proof}
Let $R(\sigma) := \prod_{e \in E(G_2)\setminus E(G_1)} \exp(\beta J \sigma_i\sigma_j)$.
\begin{itemize}
  \item \textbf{Factorization.} Since $H_{G_2}(\sigma) - H_{G_1}(\sigma) = -J\sum_{E(G_2)\setminus E(G_1)} \sigma_i\sigma_j$,
    we have $w_{G_2}(\sigma) = R(\sigma)\, w_{G_1}(\sigma)$
    (\texttt{boltzmannWeight\_subgraph\_factor}).
  \item \textbf{HNC.} Each factor decomposes as $\cosh(\beta J) + \sinh(\beta J)\,\sigma_i\sigma_j$
    with both coefficients $\ge 0$; by
    \texttt{hasNonnegCorrelations\_finset\_prod},
    $R$ has non-negative correlations
    (\texttt{hasNonnegCorrelations\_edge\_prod\_of\_finset}).
  \item \textbf{Covariance.} Apply \texttt{cov\_hnc\_boltzmann\_nonneg}
    with base graph $G_1$ and function $R$:
    \[
      \bigl(\textstyle\sum_\sigma \sigma^A R(\sigma) w_1\bigr)Z_{G_1}
      \ge \bigl(\sum_\sigma \sigma^A w_1\bigr)\bigl(\sum_\sigma R(\sigma) w_1\bigr).
    \]
    Substituting $R w_1 = w_2$ yields $Z_{G_1} N_{G_2}(A) \ge Z_{G_2} N_{G_1}(A)$,
    i.e., $\langle\sigma^A\rangle_{G_1} \le \langle\sigma^A\rangle_{G_2}$.
\end{itemize}
\end{proof}

\begin{theorem}[\texttt{correlation\_convergent\_subgraph}]
For any increasing sequence of subgraphs $G_1 \le G_2 \le \cdots$ of the ambient lattice,
the correlation sequence $n \mapsto \langle\sigma^A\rangle_{G_n}$ converges.
\end{theorem}

\begin{proof}
Monotone increasing by \texttt{correlation\_monotone\_subgraph},
bounded above by $1$ via \texttt{correlation\_le\_one}, hence converges by
\texttt{tendsto\_atTop\_ciSup}.
\end{proof}

Named specializations:
\begin{itemize}
  \item \texttt{magnetization\_convergent\_subgraph}: $\langle\sigma_i\rangle_{G_n}$
        converges (case $A = \{i\}$).
  \item \texttt{twoPoint\_convergent\_subgraph}: $\langle\sigma_i\sigma_j\rangle_{G_n}$
        converges (case $A = \{i,j\}$).
\end{itemize}

\subsection{Free energy infinite volume convergence (Proposition 4.6.1)}

GJ Proposition 4.6.1 (p.~68): for the ferromagnetic pair interaction
$H = -J\sum_{\langle i,j\rangle}\sigma_i\sigma_j - h\sum_i\sigma_i$
(strictly ferromagnetic), the free energy per site
$f_\Lambda = |\Lambda|^{-1}\ln Z_\Lambda$ converges as $\Lambda\uparrow\mathbb{Z}^d$.
The discretized (fixed ambient lattice $\iota$, growing subgraph) form:

\begin{theorem}[\texttt{partitionFunction\_monotone\_subgraph}]
For ferromagnetic $p$ and $G_1\le G_2$, $Z_{G_1}\le Z_{G_2}$.
\end{theorem}

\begin{proof}
Factor $Z_{G_2} = \sum_\sigma R(\sigma) w_{G_1}(\sigma)$ with
$R(\sigma)=\prod_{e\in E(G_2)\setminus E(G_1)}\exp(\beta J \sigma_i\sigma_j)$.
Use $\exp(x)\ge 1+x$ to get
$R(\sigma)\ge 1 + \beta J \sum_{e\in E(G_2)\setminus E(G_1)}\sigma_i\sigma_j$.
Then $\sum_\sigma R w_{G_1} \ge Z_{G_1} + \beta J \sum_e \sum_\sigma \sigma_i\sigma_j w_{G_1}$,
and each inner sum is $\ge 0$ by GKS-I (via \texttt{boltzmannWeight\_hasNonnegCorrelations})
applied to the pair $\{i,j\}$. Hence $Z_{G_2}\ge Z_{G_1}$.
\end{proof}

\begin{theorem}[\texttt{freeEnergy\_monotone\_subgraph}]
For ferromagnetic $p$ and $G_1\le G_2$, $f_{G_1}\le f_{G_2}$.
\end{theorem}

\begin{proof}
$|\iota|$ is fixed; apply $\log$ (monotone) to $Z_{G_1}\le Z_{G_2}$.
\end{proof}

\subsection{\texorpdfstring{$\beta$-monotonicity and $\beta \to \infty$ convergence}{beta-monotonicity and beta to infinity convergence}}

Rescale identities:
\[
  \langle\sigma^A\rangle_{(J,h,\beta)} = \langle\sigma^A\rangle_{(\beta J, \beta h, 1)}
  \quad\text{(\texttt{correlation\_rescale\_beta})}
\]
\[
  f(J,h,\beta) = f(\beta J, \beta h, 1)
  \quad\text{(\texttt{freeEnergy\_rescale\_beta})}
\]

\begin{theorem}[\texttt{correlation\_monotone\_beta}]
For $J, h \ge 0$, the correlation function is monotone increasing in
$\beta$ on $(0, \infty)$.
\end{theorem}

\begin{proof}
Apply the rescale identity; then $\beta_1 \le \beta_2$ gives
$\beta_1 J \le \beta_2 J$ and $\beta_1 h \le \beta_2 h$ at both
endpoints.  Apply \texttt{correlation\_monotone\_J} and
\texttt{correlation\_monotone\_h}.
\end{proof}

\begin{theorem}[\texttt{correlation\_convergent\_beta}]
For $J, h \ge 0$ and any finite set $A$, the sequence
$\langle\sigma^A\rangle_{(J,h,n+1)}$ converges as $n \to \infty$.
\end{theorem}

\begin{proof}
Monotone by \texttt{correlation\_monotone\_beta}, bounded above by
$1$ via \texttt{correlation\_le\_one}. Converges by
\texttt{tendsto\_atTop\_ciSup}.
\end{proof}

\begin{theorem}[\texttt{freeEnergy\_monotone\_beta}]
For $J, h \ge 0$, the free energy $f(J,h,\beta)$ is monotone in $\beta$
on $(0,\infty)$.
\end{theorem}

\begin{proof}
Apply the rescale identity, then \texttt{freeEnergy\_monotone\_J} and
\texttt{freeEnergy\_monotone\_h}.
\end{proof}

\begin{remark}
The free energy does not converge as $\beta\to\infty$
(it grows linearly in $\beta$, since $\log Z \sim \beta(J|E| + h|\iota|)$
at low temperature).
\end{remark}

Specializations to magnetization:
\texttt{magnetization\_monotone\_beta},
\texttt{magnetization\_convergent\_beta} (both via $A = \{i\}$).

\subsection{Convergence as \texorpdfstring{$h \to \infty$}{h to infinity}}

\begin{theorem}[\texttt{correlation\_convergent\_h}]
For $J \ge 0$, $\beta > 0$, the sequence
$n \mapsto \langle\sigma^A\rangle_{(J,n,\beta)}$ converges.
\end{theorem}

\begin{proof}
Monotone by \texttt{correlation\_monotone\_h} (Prop.~4.2.4),
bounded above by $1$ via \texttt{correlation\_le\_one}.
Apply \texttt{tendsto\_atTop\_ciSup}.
\end{proof}

Specialization: \texttt{magnetization\_convergent\_h} via $A = \{i\}$.

\begin{theorem}[\texttt{freeEnergy\_convergent\_subgraph}, = Prop 4.6.1]
For any increasing subgraph sequence $G_n\uparrow$ on a fixed finite ambient
lattice $\iota$ and ferromagnetic $p$, $n\mapsto f_{G_n}$ converges.
\end{theorem}

\begin{proof}
Monotone by \texttt{freeEnergy\_monotone\_subgraph}.
Bounded above by $f_{\top}$ (free energy on the complete graph $\top$)
via \texttt{le\_top}. Converges by \texttt{tendsto\_atTop\_ciSup}.
\end{proof}

\section{Free energy (\texttt{FreeEnergy.lean})}

The free energy and its monotonicity properties.
Reference: \cite[{\S4.6, pp.~67--70}]{glimm-jaffe}.

\begin{definition}[\texttt{freeEnergy}, (4.6.1), p.~67]
The free energy per site: $f = |\iota|^{-1} \ln Z$.
Well-defined since $Z > 0$ (\texttt{partitionFunction\_pos}).
\end{definition}

\begin{theorem}[\texttt{partitionFunction\_monotone\_h};
  cf.\ \cite{glimm-jaffe}, {\S4.6, p.~67}]
For $J \ge 0$, $\beta > 0$, $0 \le h_1 \le h_2$:
$Z(J, h_1, \beta) \le Z(J, h_2, \beta)$.
\end{theorem}

\begin{proof}
Rewrite $Z(h_2) = \sum_\sigma R(\sigma) w_1(\sigma)$ where
$R(\sigma) = \prod_i \exp(\beta(h_2-h_1)\,\text{sign}(\sigma_i))$ and
$w_1 = e^{-\beta H_{h_1}}$.
Using $\exp(x) \ge 1 + x$:
\[
  R(\sigma) = \exp\Bigl(\beta(h_2 - h_1)\sum_i \sigma_i\Bigr)
  \ge 1 + \beta(h_2 - h_1)\sum_i \sigma_i.
\]
Multiply by $w_1(\sigma)$ and sum:
$Z(h_2) \ge Z(h_1) + \beta(h_2-h_1)\sum_i Z(h_1)\langle\sigma_i\rangle_{h_1}$.
By GKS-I (\texttt{boltzmannWeight\_hasNonnegCorrelations}):
$\langle\sigma_i\rangle_{h_1} \ge 0$ for each $i$, so
$Z(h_2) \ge Z(h_1)$.
\end{proof}

\begin{theorem}[\texttt{partitionFunction\_monotone\_J};
  cf.\ \cite{glimm-jaffe}, {\S4.6, p.~67}]
For $h \ge 0$, $\beta > 0$, $0 \le J_1 \le J_2$:
$Z(J_1, h, \beta) \le Z(J_2, h, \beta)$.
\end{theorem}

\begin{proof}
Same technique as $h$-monotonicity with
$R(\sigma) = \exp(\beta(J_2-J_1)\sum_e \sigma_i\sigma_j)$.
Using $\exp(x) \ge 1+x$ and GKS-I
($\langle\sigma_i\sigma_j\rangle \ge 0$ for each edge).
\end{proof}

\begin{theorem}[\texttt{freeEnergy\_monotone\_h}, \texttt{freeEnergy\_monotone\_J}]
The free energy $f = |\iota|^{-1}\ln Z$ is monotone increasing
in $h$ on $[0,\infty)$ and in $J$ on $[0,\infty)$.
Immediate from $Z$ monotonicity and $\ln$ monotonicity
(\texttt{Real.log\_le\_log}).
\end{theorem}

\subsection{Free energy analyticity (Theorem 4.6.2)}

\begin{theorem}[\texttt{partitionFunctionH\_analyticAt}]
The partition function $Z(h) = \sum_\sigma \exp(a(\sigma) + b(\sigma) h)$
is real-analytic in $h$ at every point.  Each Boltzmann weight
$\exp(a + bh)$ is real-analytic (exponential of affine),
and a finite sum of real-analytic functions is real-analytic.
\end{theorem}

\begin{theorem}[\texttt{freeEnergyH\_analyticOn}; Thm~4.6.2, finite vol.]
The free energy $f(h) = |\iota|^{-1} \ln Z(h)$ is real-analytic on $(0,\infty)$.
Since $Z(h) > 0$ for all $h$ (\texttt{partitionFunction\_pos}), $\ln Z$
is defined and real-analytic (\texttt{AnalyticAt.log}).
\end{theorem}

\begin{proof}
$Z$ is a finite sum of exponentials of affine functions,
hence real-analytic (\texttt{fun\_prop}).  Since $Z > 0$,
$\ln Z$ is real-analytic.  Multiplication by the constant
$|\iota|^{-1}$ preserves analyticity.
\end{proof}

The same argument gives $Z(J)$ and $f(J)$ real-analytic
(\texttt{partitionFunctionJ\_analyticAt}, \texttt{freeEnergyJ\_analyticOn}).

\subsection{Configuration--subset bijection}

The Lee-Yang polynomial sums over subsets $X \subseteq \iota$
while the partition function sums over configurations
$\sigma : \iota \to \{\pm 1\}$.
The bijection \texttt{configFinsetEquiv} identifies $\sigma$ with
its ``down spin'' set $X = \{i : \sigma_i = -1\}$.

\begin{definition}[\texttt{configToFinset}, \texttt{finsetToConfig}]
$\sigma \mapsto \{i : \sigma_i = \mathrm{down}\}$ and
$X \mapsto (\lambda i.\;\text{if } i \in X \text{ then down else up})$.
These are mutual inverses (\texttt{configFinsetEquiv}).
\end{definition}

The connection identity
(\cite[{(3.63)--(3.65), pp.~122--123}]{friedli-velenik}):
\[
  Z(J, h, \beta) = e^{\beta J|E| + \beta h N} \cdot
  P\bigl(e^{-2\beta h}, \ldots, e^{-2\beta h}\bigr)
\]
where $P = \texttt{isingEdgePoly}$, $z_i = e^{-2\beta h}$,
$t_e = e^{-2\beta J}$.
Combined with Lee-Yang ($P(z) \ne 0$ for $|z_i| < 1$), this gives
$Z(h) \ne 0$ for complex $h$ with $|e^{-2\beta h}| < 1$.

\begin{definition}[\texttt{graphToEdgeList}]
Converts a \texttt{SimpleGraph} with uniform coupling $t$ to the
edge list format used by \texttt{lee\_yang\_circle}:
each edge $e$ is represented as $(\texttt{Quot.out}\;e)$ with coupling $t$.
\end{definition}

\begin{theorem}[\texttt{isingEdgePoly\_nonvanishing\_of\_graph},
  Theorem 4.6.2 {\cite[pp.~64--70]{glimm-jaffe}},
  {\cite[Theorem~3.43, pp.~122--127]{friedli-velenik}}]
For the Ising model on graph $G$ with coupling $t \in [0,1)$
(i.e., $J > 0$), the partition polynomial does not vanish on the
open unit polydisk: $P(z) \ne 0$ for all $|z_k| < 1$.
\end{theorem}

\begin{proof}
Construct the edge list via \texttt{graphToEdgeList}, verify the
conditions (distinct endpoints from \texttt{SimpleGraph.Adj.ne},
coupling bounds from $0 \le t < 1$), and apply
\texttt{lee\_yang\_circle}.
\end{proof}

\section{GHS inequality (\texttt{Inequalities/GHS.lean})}

\par\noindent\textbf{Prerequisites:} \texttt{InfiniteVolume.lean}.

\textbf{Role:} Defines truncated correlation functions ($u_2$, $u_3$,
$u_4$) and proves the GHS inequality $u_3 \le 0$, Cor.~4.3.3
($U_4 \le 0$ for $h=0$), and Cor.~4.3.5 (inductive bound).
Also contains \texttt{correlation\_odd\_vanish} ($h=0$, odd $|A|$).

\textbf{What to look for:} The Lebowitz axioms
(\texttt{lebowitz\_third/four/inductive}) are the only axioms in the
file. Everything else is proved. The spin-flip symmetry proof
(\texttt{correlation\_odd\_vanish}) uses a bijection argument:
pair $\sigma$ with $\mathrm{flip}(\sigma)$ to show $S = -S$.

The Griffiths--Hurst--Sherman inequality for the Ising model.
Reference: Ellis, \emph{Entropy, Large Deviations, and Statistical
Mechanics}, \S V.3, pp.~143--146.

\begin{definition}[\texttt{truncated2}, \texttt{truncated3}]
The truncated 2-point and 3-point functions (Ursell functions):
\begin{align*}
  u_2(i,j) &= \langle\sigma_i\sigma_j\rangle - \langle\sigma_i\rangle\langle\sigma_j\rangle, \\
  u_3(i,j,k) &= \langle\sigma_i\sigma_j\sigma_k\rangle
    - \langle\sigma_i\rangle\langle\sigma_j\sigma_k\rangle
    - \langle\sigma_j\rangle\langle\sigma_i\sigma_k\rangle
    - \langle\sigma_k\rangle\langle\sigma_i\sigma_j\rangle
    + 2\langle\sigma_i\rangle\langle\sigma_j\rangle\langle\sigma_k\rangle.
\end{align*}
\end{definition}

\begin{theorem}[\texttt{truncated2\_nonneg}]
$u_2(i,j) \ge 0$ for ferromagnetic parameters. Proved via GKS-II
for $i \ne j$ and GKS-I $+$ $|\langle\sigma_i\rangle| \le 1$ for $i = j$.
\end{theorem}

\begin{theorem}[\texttt{lebowitz\_third}; axiom]
For ferromagnetic parameters with $h \ge 0$ and distinct $i, j, k$:
\[
  \langle\sigma_i\sigma_j\sigma_k\rangle + \langle\sigma_i\rangle\langle\sigma_j\sigma_k\rangle
  \le \langle\sigma_i\sigma_j\rangle\langle\sigma_k\rangle + \langle\sigma_i\sigma_k\rangle\langle\sigma_j\rangle.
\]
\end{theorem}

The Lebowitz third inequality is proved for continuous $\varphi^4$ spins
via \texttt{phi4\_single\_site\_nonneg} (Theorem~4.3.1 in Glimm--Jaffe),
then transferred to Ising spins by the approximation
$d\mu = e^{-\lambda(\xi^2-1)^2}\,d\xi \to \frac{1}{2}(\delta_{+1}+\delta_{-1})$
as $\lambda \to \infty$.

\textbf{Note:} The per-site factorization in Ellis \S V.3 (Lemma~V.3.2)
does \emph{not} hold for discrete Ising spins---the all-odd parity case
gives $\sum \alpha\beta\gamma\delta\, e^{2h\alpha} = -8\cosh(2h) < 0$.
The continuous-spin route is essential.

\begin{theorem}[\texttt{ghs\_inequality}; proved from \texttt{lebowitz\_third}]
For ferromagnetic parameters with $h \ge 0$ and distinct $i, j, k$:
$u_3(i,j,k) \le 0$.
\end{theorem}

\begin{proof}
Substituting the Lebowitz bound into the definition of $u_3$:
\[
  u_3(i,j,k) \le -2\langle\sigma_i\rangle \cdot u_2(j,k) \le 0,
\]
where the last step uses GKS-I ($\langle\sigma_i\rangle \ge 0$) and
\texttt{truncated2\_nonneg} ($u_2(j,k) \ge 0$).
In Lean, \texttt{nlinarith} closes the goal from
\texttt{lebowitz\_third}, \texttt{gks\_first}, and
\texttt{truncated2\_nonneg}.
\end{proof}

\subsection{Spin-flip symmetry and odd correlations}

\begin{theorem}[\texttt{spinProduct\_flip}]
For any set $A$ and configuration $\sigma$:
$\sigma^A(\mathrm{flip}\;\sigma) = (-1)^{|A|} \cdot \sigma^A(\sigma)$.
\end{theorem}

\begin{theorem}[\texttt{correlation\_odd\_vanish}]
For $h = 0$ and $|A|$ odd: $\langle\sigma^A\rangle = 0$.
\end{theorem}

\begin{proof}
Set $S = \sum_\sigma \sigma^A(\sigma) \cdot w(\sigma)$.
Since $H(\mathrm{flip}\;\sigma) = H(\sigma)$ when $h = 0$
(\texttt{hamiltonian\_flip\_eq}), we have $w(\mathrm{flip}\;\sigma) = w(\sigma)$.
The flip involution $\sigma \mapsto \mathrm{flip}\;\sigma$ is a bijection on
configurations, so
$S = \sum_\sigma \sigma^A(\mathrm{flip}\;\sigma) \cdot w(\mathrm{flip}\;\sigma)
   = (-1)^{|A|} S = -S$,
hence $S = 0$.
\end{proof}

\subsection{Cor.~4.3.3: truncated 4-point function}

\begin{definition}[\texttt{truncated4}]
The truncated (connected) 4-point function:
\[
  U_4(i,j,k,l) = \langle\sigma_i\sigma_j\sigma_k\sigma_l\rangle
  - \langle\sigma_i\sigma_j\rangle\langle\sigma_k\sigma_l\rangle
  - \langle\sigma_i\sigma_k\rangle\langle\sigma_j\sigma_l\rangle
  - \langle\sigma_i\sigma_l\rangle\langle\sigma_j\sigma_k\rangle.
\]
\end{definition}

\begin{theorem}[\texttt{lebowitz\_four}; axiom]
For ferromagnetic parameters with $h \ge 0$ and four distinct sites $i,j,k,l$:
\begin{multline*}
  \langle\sigma_i\sigma_j\sigma_k\sigma_l\rangle
  + \langle\sigma_i\sigma_j\rangle\langle\sigma_k\sigma_l\rangle \\
  \le \langle\sigma_i\sigma_k\rangle\langle\sigma_j\sigma_l\rangle
  + \langle\sigma_i\sigma_l\rangle\langle\sigma_j\sigma_k\rangle
  + \langle\sigma_i\sigma_k\sigma_l\rangle\langle\sigma_j\rangle
  + \langle\sigma_j\sigma_k\sigma_l\rangle\langle\sigma_i\rangle.
\end{multline*}
\end{theorem}

\begin{theorem}[\texttt{cor\_4\_3\_3}; proved]
For $h = 0$ and four distinct sites, $U_4(i,j,k,l) \le 0$.
\end{theorem}

\begin{proof}
From \texttt{lebowitz\_four}, with $h = 0$ the odd-cardinality correlations
vanish: $\langle\sigma_i\rangle = \langle\sigma_j\rangle = 0$ and
$\langle\sigma_{ikl}\rangle = \langle\sigma_{jkl}\rangle = 0$.
The Lebowitz bound reduces to
$\langle\sigma_{ijkl}\rangle + \langle\sigma_{ij}\rangle\langle\sigma_{kl}\rangle
\le \langle\sigma_{ik}\rangle\langle\sigma_{jl}\rangle
  + \langle\sigma_{il}\rangle\langle\sigma_{jk}\rangle$.
Since $\langle\sigma_{ij}\rangle, \langle\sigma_{kl}\rangle \ge 0$ by GKS-I,
this gives $U_4 \le -2\langle\sigma_{ij}\rangle\langle\sigma_{kl}\rangle \le 0$.
\end{proof}

\subsection[Cor.~4.3.5: n-point inductive bound]{Cor.~4.3.5: $n$-point inductive bound}

\begin{theorem}[\texttt{lebowitz\_inductive}; axiom]
For ferromagnetic $h \ge 0$, a set $S$ of sites, and two sites $j, k \notin S$
with $j \ne k$:
\[
  \langle\sigma_{S \cup \{j,k\}}\rangle \le
  \langle\sigma_S\rangle\langle\sigma_j\sigma_k\rangle
  + \sum_{T \subseteq S}
    \langle\sigma_{T \cup \{j\}}\rangle\langle\sigma_{(S\setminus T) \cup \{k\}}\rangle.
\]
\end{theorem}

This follows from the general Lebowitz inequality (Cor.~4.3.2 in
Glimm--Jaffe) applied with $A = S$, $B = \{j,k\}$, and dropping
non-positive terms.  Iterating gives Cor.~4.3.5: the $n$-point function
is bounded by $(n-1)!$ times the sum over partial matchings of products
of 2-point and 1-point correlations.

\section{\texorpdfstring{$\varphi^4$}{phi4} inequalities (\texttt{ContinuousSpin/Phi4.lean})}

The $\varphi^4$ correlation inequalities
\cite[{\S4.3, pp.~59--62}]{glimm-jaffe}
provide Lebowitz-type bounds for continuous spin models with
single-spin distribution $d\mu = e^{-\lambda\xi^4 - \sigma\xi^2}\,d\xi$.

The formalization includes:
\begin{itemize}
\item The quartic identity for the orthogonal transformation
  $(\xi,\chi,\xi',\chi') \mapsto (\alpha,\beta,\gamma,\delta)$
  (\texttt{quartic\_identity}, proved by \texttt{ring}).
\item Sum-of-squares, linear, and inner-product identities
  (orthogonality of the transformation).
\item The positivity lemma $t\cdot\sinh(ct) \ge 0$ for $c \ge 0$
  (\texttt{mul\_sinh\_nonneg}).
\item Odd-function integral vanishing (\texttt{integral\_odd\_eq\_zero}).
\end{itemize}

\paragraph{Axiom.}
One measure-theoretic property is stated as axiom (mathematically
proved but not formalized in Lean due to the heavy measure theory assembly):
\begin{itemize}
\item \texttt{phi4\_single\_\allowbreak site\_nonneg}: non-negativity of the
  symmetrized 4D integral: for $Q$ even and $c \ge 0$,
  \[
    \int \alpha^k\beta^l\gamma^m\delta^n \,
    e^{-Q + c\alpha\beta\gamma\delta} \ge 0.
  \]
\end{itemize}
The proof strategy: write $e^x = \cosh(x) + \sinh(x)$ with $x = c\alpha\beta\gamma\delta$.
The $\cosh$ term is even in each variable, so odd-power monomials integrate to zero
(\texttt{integral\_odd\_eq\_zero}); for all-even $(k,l,m,n)$ the integrand is
$\ge 0$ (even powers $\times$ $\cosh \ge 0$ $\times$ $e^{-Q} \ge 0$).
The $\sinh$ term is odd in each variable, so even-power monomials integrate to zero;
for all-odd $(k,l,m,n)$ the integrand is
$\alpha^{k-1}\beta^{l-1}\gamma^{m-1}\delta^{n-1} \cdot
(\alpha\beta\gamma\delta \cdot \sinh(c\alpha\beta\gamma\delta)) \cdot e^{-Q} \ge 0$
by \texttt{mul\_sinh\_nonneg}.
Formalizing this requires the interchange of finite symmetrization sums with
nested Bochner integrals.

\section{Pure and mixed phases (\texttt{PhaseTransition.lean})}

Formalization of Glimm--Jaffe, \S5.1 (pp.~72--74).

\begin{theorem}[\texttt{truncated2\_le\_one}]
For ferromagnetic parameters: $0 \le \langle\sigma_i;\sigma_j\rangle \le 1$.
The lower bound is GKS-II. The upper bound: $\langle\sigma_i\sigma_j\rangle \le 1$
and $\langle\sigma_i\rangle\langle\sigma_j\rangle \ge 0$ by GKS-I.
\end{theorem}

\begin{theorem}[\texttt{mixed\_phase\_truncated2}; eq.~(5.1.5)]
$M^2 - (M(2\alpha - 1))^2 = 4\alpha(1-\alpha)M^2$.
\end{theorem}

\begin{theorem}[\texttt{mixed\_phase\_pure\_iff}]
For $0 \le \alpha \le 1$ and $M \ne 0$:
$4\alpha(1-\alpha)M^2 = 0$ iff $\alpha = 0$ or $\alpha = 1$.
\end{theorem}

These formalize the algebraic content of \S5.1: the truncated 2-point
function measures the deviation from a pure phase, and
its asymptotic value $4\alpha(1-\alpha)M^2$ vanishes iff the state is pure.

\subsection{Mean field theory (\S5.2)}

The mean field free energy density for the Ising model
(Glimm--Jaffe, \S5.2, eq.~(5.2.3)) is
$\varphi(m) = -\frac{1}{2}Jzm^2 - hm + \beta^{-1}\,\mathrm{entropy}(m)$.

\begin{definition}[\texttt{meanFieldEnergy}]
The interaction part: $-\frac{1}{2}Jzm^2 - hm$.
\end{definition}

\begin{theorem}[\texttt{meanFieldEnergy\_neg}]
At $h = 0$: $\varphi(-m) = \varphi(m)$ (spin-flip symmetry).
\end{theorem}

\begin{theorem}[\texttt{meanField\_zero\_solution}]
$m = 0$ is always a solution of the mean field equation
$m = \tanh(\beta(Jzm + h))$ when $h = 0$, since $\tanh(0) = 0$.
\end{theorem}

\begin{theorem}[\texttt{tanh\_odd}]
$\tanh(-x) = -\tanh(x)$: if $m^*$ solves the mean field equation
at $h = 0$, so does $-m^*$.
\end{theorem}

\subsection{Symmetry breaking (\S5.3)}

\begin{definition}[\texttt{magnetization}]
The order parameter: $M(i) = \langle\sigma_i\rangle$ (eq.~(5.3.5)).
\end{definition}

\begin{definition}[\texttt{susceptibility}]
$\chi(i) = \sum_j \langle\sigma_i;\sigma_j\rangle = \sum_j \texttt{truncated2}(i,j)$.
\end{definition}

\begin{theorem}[\texttt{susceptibility\_nonneg}]
$\chi(i) \ge 0$ for ferromagnetic parameters, by \texttt{truncated2\_nonneg}.
\end{theorem}

\begin{theorem}[\texttt{magnetization\_zero\_at\_h\_zero}]
$M(i) = 0$ when $h = 0$: the $\mathbb{Z}_2$ symmetry is unbroken in finite
volume. Uses \texttt{correlation\_odd\_vanish}.
\end{theorem}

Symmetry breaking occurs only in the infinite volume limit:
$M_+ = \lim_{h \downarrow 0} \langle\sigma_i\rangle_h > 0$ for $\beta > \beta_c$.
The concavity $d^2 M/dh^2 \le 0$ for $h \ge 0$ follows from the GHS inequality
(Cor.~4.3.4), as noted on p.~80 of \cite{glimm-jaffe}.

\subsubsection{\texorpdfstring{$\mathbb{Z}_2$ $h$-symmetry across the tower}{Z2 h-symmetry across the tower}}

The reflection $h \mapsto -h$ combined with the global spin-flip
$\sigma \mapsto -\sigma$ is a symmetry of the Ising Hamiltonian. The
statements in this subsubsection are derived from this symmetry in
combination with the basic \texttt{correlation} / \texttt{magnetization}
/ \texttt{susceptibility} definitions; the citation to
Glimm--Jaffe~\S5.3, pp.~77--80 (2nd ed., Springer, 1987) is at the
section level, as background for the $\mathbb{Z}_2$ $h$-symmetry
framework, rather than as a source of each individual closed form.
On correlations, the framework yields the capstone identity:
\begin{theorem}[\texttt{correlation\_neg\_h}, PR~\#754]
$\langle \sigma^A \rangle_{(J,-h,\beta)}
 = (-1)^{|A|}\,\langle \sigma^A \rangle_{(J,h,\beta)}$.
\end{theorem}
\noindent The proof uses the flip reindexing of numerator and denominator:
\texttt{hamiltonian\_neg\_h} gives $H_{(J,-h,\beta)}(\sigma)
= H_{(J,h,\beta)}(\sigma.\mathrm{flip})$, \texttt{spinProduct\_flip}
gives $\prod_{i\in A}\sigma.\mathrm{flip}(i) = (-1)^{|A|}\prod_{i\in A}\sigma(i)$,
and \texttt{partitionFunction\_neg\_h} shows $Z(-h) = Z(h)$.

Specializations:
\begin{itemize}
\item \texttt{magnetization\_neg\_h}: $M(-h) = -M(h)$ ($|A|=1$).
\item \texttt{truncated\{2,3,4\}\_neg\_h} (PRs~\#755--\#758): the Ursell
      coefficients inherit the parity of their total degree. E.g.\
      \texttt{truncated2\_neg\_h} gives $\langle\sigma_i;\sigma_j\rangle(-h)
      = \langle\sigma_i;\sigma_j\rangle(h)$ for $i\neq j$ (even total).
\item \texttt{correlation\_eq\_abs\_h\_of\_even\_card} (PR~\#760): for
      $|A|$ even, $\langle\sigma^A\rangle_{|h|} = \langle\sigma^A\rangle_{h}$.
\end{itemize}
The Ambient $\Lambda$-layer, along-exhaustion, and infinite-volume
counterparts (even $|A|$ in the $\infty$-volume case, since
$\sup(-f) \ne -\sup(f)$ for odd $|A|$) appear in PRs~\#762--\#766.

\begin{theorem}[\texttt{susceptibility\_neg\_h}, PR~\#767]
$\chi(\langle J,-h,\beta\rangle; i) = \chi(\langle J,h,\beta\rangle; i)
 - 2\,M(\langle J,h,\beta\rangle; i)$.
\end{theorem}
\noindent The off-diagonal sum $\sum_{j\neq i}\texttt{truncated2}(i,j)$ is
invariant by \texttt{truncated2\_neg\_h}; the diagonal term, via the
Finset collapse $\{i,i\}=\{i\}$, contributes $\texttt{truncated2}(-h,i,i)
- \texttt{truncated2}(h,i,i) = -2\,M(h)$ by the $|A|=1$ case of
\texttt{correlation\_neg\_h}.

Under ferromagnetism ($0\le J$, $0<\beta$), the magnetization at $|h|$
coincides with the absolute value:
\begin{theorem}[\texttt{abs\_magnetization\_eq\_magnetization\_abs\_h}, PR~\#769]
$|M(\langle J,h,\beta\rangle;i)| = M(\langle J,|h|,\beta\rangle;i)$
under $0\le J$, $0<\beta$.
\end{theorem}
\noindent The $\Lambda$-layer counterpart is
\begin{theorem}[\texttt{abs\_magnetization\LeanLambda\_eq\_magnetization\LeanLambda\_abs\_h}, PR~\#772]
$|M_\Lambda(\langle J,h,\beta\rangle;i)| = M_\Lambda(\langle J,|h|,\beta\rangle;i)$
under $0\le J$, $0<\beta$, for any finite volume $\Lambda \subseteq V$ and
any site $i \in \Lambda$. Proved by the same
\texttt{abs\_choice}-based case split, using
\texttt{magnetization\LeanLambda\_nonneg} (ferromagnetic, via GKS-I at
$\Lambda$-level) and \texttt{magnetization\LeanLambda\_neg\_h}.
\end{theorem}
\noindent The along-exhaustion counterpart (pointwise per stage) is
likewise an equality, and it is the bridge to the infinite-volume
statement:
\begin{theorem}[\texttt{abs\_magnetizationAlongExhaustion\_eq\_magnetizationAlongExhaustion\_abs\_h}, PR~\#773]
For every exhaustion stage $n$,
$|M_\text{along}(\langle J,h,\beta\rangle;i,n)|
= M_\text{along}(\langle J,|h|,\beta\rangle;i,n)$
under $0 \le J$, $0 < \beta$. Same case split via
\texttt{abs\_choice}, using
\texttt{magnetizationAlongExhaustion\_nonneg} and the new helper
\texttt{magnetizationAlongExhaustion\_neg\_h} (specialization of
\texttt{correlationAlongExhaustion\_neg\_h} at $A = \{i\}$).
\end{theorem}

The naive lift to the thermodynamic-limit $M_\infty$ \emph{fails} as an
equality for this project's sup-based definition
$M_\infty := \sup_n M_\text{along}(n)$. The failure mode is the
odd-$|A|$ obstruction already flagged in
\texttt{correlationInfinite\_neg\_h\_of\_even\_card}: at $h < 0$
ferromagnetic, covered stages give $M_\text{along}(n) \le 0$
(via \texttt{magnetizationAlongExhaustion\_neg\_h} plus the finite/\(\Lambda\)
ferromagnetic positivity at $|h|$), so the supremum is bounded above
by the contributions of any \emph{missed} stages (where
$M_\text{along}(n) = 0$ by the ``if $A \subseteq \Lambda_n$ then
$\ldots$ else $0$'' convention). Concretely, if there exists
$n$ with $i \notin \Lambda_n$, then $M_\infty(h) = 0$ but
$M_\infty(|h|) > 0$, so the equality fails. Since
\texttt{Exhaustion} does not require such a missed stage, this is an
obstruction/example rather than a universal identity; still, the
equality cannot hold in general under this definition.
Glimm--Jaffe's standard thermodynamic limit
(e.g.\ $\pm$BC or $h \to 0^\pm$) restores the equality; upgrading the
definition is out of scope here.

What does survive is the one-sided inequality, which is the natural
$\infty$-volume odd-$|A|$ result compatible with the sup-based setup:
\begin{theorem}[\texttt{abs\_magnetizationInfinite\_le\_magnetizationInfinite\_abs\_h}, PR~\#773]
$|M_\infty(\langle J,h,\beta\rangle;\,i)|
 \le M_\infty(\langle J,|h|,\beta\rangle;\,i)$
under $0 \le J$, $0 < \beta$. Proved via $\mathtt{abs\_le}$ plus two
$\mathtt{ciSup}$-level pointwise comparisons driven by the along-exhaustion
pointwise equality above.
\end{theorem}

Two further $\infty$-volume sign-control statements pin down the
$h \le 0$ behavior of the sup-based $M_\infty$:
\begin{theorem}[\texttt{magnetizationInfinite\_nonpos\_of\_nonpos\_h}, PR~\#774]
Under $0 \le J$, $0 < \beta$, $h \le 0$, and any exhaustion,
$M_\infty(\langle J, h, \beta\rangle; i) \le 0$.
\end{theorem}
\noindent Proof: rewrite $M_\infty = \sup_n M_\text{along}(n)$; each
covered stage is $\le 0$ by \texttt{magnetizationAlongExhaustion\_neg\_h}
plus nonnegativity at $-h$ (ferromagnetic, since $-h \ge 0$), and each
uncovered stage is $0$; close with $\mathtt{ciSup\_le}$.

\begin{theorem}[\texttt{magnetizationInfinite\_eq\_zero\_of\_exists\_stage\_not\_mem}, PR~\#774]
Under $0 \le J$, $0 < \beta$, $h \le 0$, and the extra hypothesis that
some stage misses $i$ (i.e.\ $\exists n_0,\ i \notin \Lambda_{n_0}$),
$M_\infty(\langle J, h, \beta\rangle; i) = 0$.
\end{theorem}
\noindent Proof: the previous theorem gives $M_\infty \le 0$; the
missed stage contributes $M_\text{along}(n_0) = 0 \le M_\infty$ via
\texttt{magnetizationAlongExhaustion\_le\_magnetizationInfinite};
close with \texttt{le\_antisymm}.

Further ``$\Lambda$-layer susceptibility'' and
``$\infty$-volume susceptibility'' refinements are tracked in
issue~\#770.

Finally, a closed form for $\chi$ at $|h|$ that does \emph{not} require
ferromagnetism:
\begin{theorem}[\texttt{susceptibility\_eq\_abs\_h}, PR~\#771]
$\chi(\langle J,|h|,\beta\rangle;i)
 = \chi(\langle J,h,\beta\rangle;i)
   + M(\langle J,|h|,\beta\rangle;i) - M(\langle J,h,\beta\rangle;i)$.
\end{theorem}
\begin{proof}[Proof sketch]
Case $|h| = h$ ($h\ge 0$): $M(|h|) = M(h)$ so the correction vanishes;
both sides equal $\chi(h)$. Case $|h| = -h$ ($h\le 0$): by
\texttt{susceptibility\_neg\_h}, $\chi(-h) = \chi(h) - 2\,M(h)$, and by
\texttt{magnetization\_neg\_h}, $M(-h) = -M(h)$, so the RHS equals
$\chi(h) + (-M(h)) - M(h) = \chi(h) - 2\,M(h)$, matching the LHS.
\end{proof}
\noindent Under ferromagnetism, combining with
\texttt{abs\_magnetization\_eq\_magnetization\_abs\_h} gives the
non-negative closed form
$\chi(|h|) - \chi(h) = |M(h)| - M(h) = 2\max\{0,-M(h)\}$.

\begin{definition}[\texttt{susceptibility\LeanLambda}, PR~\#776]
For $\Lambda \subseteq V$ a finite volume in an infinite ambient $V$ and
$i : \uparrow\Lambda$,
\[
  \chi_\Lambda(G, \Lambda, p; i)
  := \chi\bigl(\mathrm{inducedGraph}\ G\ \Lambda,\ p;\ i\bigr)
  = \sum_{j : \uparrow\Lambda}\mathrm{truncated2}\bigl(\mathrm{inducedGraph}\ G\ \Lambda,\ p;\ i, j\bigr).
\]
This is the $\Lambda$-layer counterpart of \texttt{susceptibility}, matching
the \texttt{magnetization\LeanLambda} / \texttt{correlation\LeanLambda} pattern at the ambient
lattice layer.
\end{definition}

\begin{theorem}[\texttt{susceptibility\LeanLambda\_eq\_abs\_h}, PR~\#776 (A-4, issue~\#770)]
For all $J, h, \beta \in \mathbb R$ and $i : \uparrow\Lambda$,
\[
  \chi_\Lambda(\langle J, |h|, \beta\rangle; i)
  = \chi_\Lambda(\langle J, h, \beta\rangle; i)
    + M_\Lambda(\langle J, |h|, \beta\rangle; i) - M_\Lambda(\langle J, h, \beta\rangle; i),
\]
unconditionally (no ferromagnetic hypothesis required).
\end{theorem}
\begin{proof}[Proof sketch]
Direct lift of \texttt{susceptibility\_eq\_abs\_h} (PR~\#771) along
$\chi_\Lambda := \chi(\mathrm{inducedGraph}\ G\ \Lambda)$ and
$M_\Lambda = M(\mathrm{inducedGraph}\ G\ \Lambda)$. Bundled helpers
\texttt{susceptibility\LeanLambda\_apply}, \texttt{susceptibility\LeanLambda\_nonneg}
(ferromagnetic), \texttt{susceptibility\LeanLambda\_neg\_h}
($\chi_\Lambda(-h) = \chi_\Lambda(h) - 2 M_\Lambda(h)$), and the
$\mathbb{Z}^d$ wrapper
\texttt{susceptibility\LeanLambda\_latticeGraph\_eq\_abs\_h} are included.
The concrete wrapper now lives in
\texttt{CorrelationSymmetry.lean}.
\end{proof}

\begin{definition}[\texttt{susceptibilityAlongExhaustion}, PR~\#777]
For an exhaustion $\Lambda = (\Lambda_n)_{n : \mathbb N}$ of $V$, site $i : V$
and stage $n : \mathbb N$,
\[
  \chi_{\mathrm{along}}(i, n)
  := \begin{cases}
       \chi_{\Lambda_n}(G, \Lambda_n, p;\ \langle i, h_{in}\rangle) & \text{if } i \in \Lambda_n \\
       0 & \text{otherwise.}
     \end{cases}
\]
Along-exhaustion companion of $\chi_\Lambda$, mirroring
\texttt{magnetizationAlongExhaustion}. Prerequisite for a future
\texttt{susceptibilityInfinite}.
\end{definition}

\begin{theorem}[\texttt{susceptibilityAlongExhaustion\_eq\_abs\_h},
PR~\#777 (A-4b, issue~\#770)]
For all $J, h, \beta \in \mathbb R$, $i : V$ and $n : \mathbb N$,
\[
  \chi_{\mathrm{along}}(\langle J, |h|, \beta\rangle; i, n)
  = \chi_{\mathrm{along}}(\langle J, h, \beta\rangle; i, n)
    + M_{\mathrm{along}}(\langle J, |h|, \beta\rangle; i, n) - M_{\mathrm{along}}(\langle J, h, \beta\rangle; i, n),
\]
unconditionally (no ferromagnetic hypothesis required).
\end{theorem}
\begin{proof}[Proof sketch]
Case split on $i \in \Lambda_n$:
\begin{itemize}
  \item \emph{Covered stage}: reduce to
    \texttt{susceptibility\LeanLambda\_eq\_abs\_h} (PR~\#776) at the lifted
    subtype site $\langle i, h_i\rangle$. The magnetization-side
    terms are converted via the helper
    \texttt{magnetizationAlongExhaustion\_of\_mem\_eq\_magnetization\LeanLambda},
    which upgrades the \texttt{liftFinset $\{i\}$}-form returned by
    \texttt{magnetizationAlongExhaustion\_of\_mem} (itself a singleton
    specialization of \texttt{correlationAlongExhaustion\_of\_subset})
    to the \texttt{magnetization\LeanLambda} form via the auxiliary
    \texttt{liftFinset\_singleton}.
  \item \emph{Uncovered stage}: all four terms are $0$ by the
    \texttt{\_of\_not\_mem} unfoldings, and the identity holds.
\end{itemize}
Bundled helpers: \texttt{susceptibilityAlongExhaustion\_apply},
\texttt{\_of\_mem}, \texttt{\_of\_not\_mem}, \texttt{\_nonneg}
(ferromagnetic), \texttt{\_neg\_h}
($\chi_{\mathrm{along}}(-h) = \chi_{\mathrm{along}}(h) - 2 M_{\mathrm{along}}(h)$),
and the $\mathbb Z^d$ wrapper
\texttt{susceptibilityAlongExhaustion\_latticeGraph\_eq\_abs\_h}. The concrete
wrapper now lives in
\texttt{CorrelationSymmetry.lean}.
\end{proof}

\begin{theorem}[\texttt{susceptibilityAlongExhaustion\_le\_abs\_h},
PR~\#778 (A-4c, issue~\#770)]
For $0 \le J$, $0 < \beta$, any $h \in \mathbb R$, $i : V$ and $n : \mathbb N$,
\[
  \chi_{\mathrm{along}}(\langle J, h, \beta\rangle; i, n)
  \le \chi_{\mathrm{along}}(\langle J, |h|, \beta\rangle; i, n).
\]
\end{theorem}
\begin{proof}[Proof sketch]
Case $|h| = h$ ($h \ge 0$): equality by $h = |h|$. Case $|h| = -h$
($h \le 0$): by PR~\#777's \texttt{susceptibilityAlongExhaustion\_eq\_abs\_h},
$\chi_{\mathrm{along}}(|h|) = \chi_{\mathrm{along}}(h) + M_{\mathrm{along}}(|h|) - M_{\mathrm{along}}(h)$.
$\mathrm{Ferromagnetic}\ \langle J, |h|, \beta\rangle$ holds by
$0 \le J$, $0 \le |h|$, $0 < \beta$, giving
$M_{\mathrm{along}}(|h|) \ge 0$. Via
\texttt{magnetizationAlongExhaustion\_neg\_h}, $M_{\mathrm{along}}(h) = -M_{\mathrm{along}}(|h|) \le 0$.
Thus the correction $M_{\mathrm{along}}(|h|) - M_{\mathrm{along}}(h) \ge 0$, and
$\chi_{\mathrm{along}}(|h|) \ge \chi_{\mathrm{along}}(h)$.
The concrete $\mathbb Z^d$ wrapper now lives in
\texttt{CorrelationSymmetry.lean}.
\end{proof}

\begin{definition}[\texttt{susceptibilityInfinite}, PR~\#778]
For an exhaustion $\Lambda = (\Lambda_n)_{n : \mathbb N}$ of $V$ and a
site $i : V$,
\[
  \chi_\infty(G, \Lambda, p; i)
  := \sup_{n : \mathbb N} \chi_{\mathrm{along}}(G, \Lambda, p; i, n)
\]
(\texttt{ciSup} on $\mathbb R$: defaults to $0$ when the sequence is
unbounded above). Unlike $M$ and correlations which are bounded by $1$,
susceptibility is not automatically bounded; at the ferromagnetic
critical point, $\chi_{\mathrm{along}}$ diverges and the \texttt{ciSup}
default kicks in. Theorems comparing $\chi_\infty$ values typically
assume an explicit $\texttt{BddAbove}$ hypothesis in the unbounded
regime. Bundled: \texttt{\_eq\_ciSup}, \texttt{\_apply},
\texttt{\_nonneg} (ferromagnetic).
\end{definition}

\begin{theorem}[\texttt{susceptibilityInfinite\_le\_abs\_h},
PR~\#778 (A-5$'$, issue~\#770)]
For $0 \le J$, $0 < \beta$, any $h \in \mathbb R$, $i : V$, and under
the hypothesis that
$\{\chi_{\mathrm{along}}(\langle J, |h|, \beta\rangle; i, n) : n \in \mathbb N\}$
is bounded above,
\[
  \chi_\infty(\langle J, h, \beta\rangle; i)
  \le \chi_\infty(\langle J, |h|, \beta\rangle; i).
\]
\end{theorem}
\begin{proof}[Proof sketch]
A-4c gives pointwise $\chi_{\mathrm{along}}(h) \le \chi_{\mathrm{along}}(|h|)$
at every $n$. The \texttt{BddAbove} hypothesis on the $|h|$-side and
pointwise comparison gives the same for the $h$-side. Apply
\texttt{ciSup\_mono} (mathlib).
\end{proof}

\begin{theorem}[\texttt{truncated2Infinite\_tendsto\_cofinite\_zero\_of\_summable},
PR~\#779]
For any ambient type $V$, graph $G$, exhaustion $\Lambda$, parameters
$p$, and fixed site $i : V$, if the function
$j \mapsto \mathrm{truncated2Infinite}\ G\ \Lambda\ p\ i\ j$ is summable
over $V$, then
\[
  \mathrm{Tendsto}\bigl(\lambda j,\ \mathrm{truncated2Infinite}\ G\ \Lambda\ p\ i\ j\bigr)
    \ \mathrm{Filter.cofinite}\ (\mathrm{nhds}\ 0).
\]
\end{theorem}
\begin{proof}[Proof sketch]
Direct application of mathlib's \texttt{Summable.tendsto\_cofinite\_zero}.
\end{proof}
\noindent
\emph{Interpretation.} The summability hypothesis is a finiteness
condition on the two-point function summed over the free argument $j$;
in translation-invariant / connected-correlation settings (a pure phase
of a $\mathbb Z^d$ Ising model) it matches the physical ``finite
susceptibility'' condition expected away from the critical line, but in
the general ambient setup of the Lean statement it is simply the
real-analysis condition $\mathrm{Summable}$. $\mathrm{Filter.cofinite}$
on $V$ is the filter of cofinite subsets; on $\mathrm{Fin}\ d \to
\mathbb Z$ (with $d \ge 1$) bounded subsets are finite, so cofinite
aligns with the informal ``$|r| \to \infty$'' filter. This is therefore
a \emph{conditional} cluster decay statement (Glimm--Jaffe \S5.1).
Unconditional exponential decay (Simon--Lieb inequality) is not yet
formalized; this lemma is the elementary real-analysis step ready to be
composed with a future proof of summability. $\mathbb Z^d$ wrapper:
\texttt{truncated2Infinite\_latticeGraph\_tendsto\_cofinite\_zero\_of\_summable}.

\paragraph{$\ell^1$ lattice distance on $\mathbb Z^d$.}
As infrastructure for turning the above cofinite statement into a
distance-based statement (Epic~\#780), PR~\#781 introduces an
$\mathbb N$-valued $\ell^1$ distance on $\mathrm{Fin}\,d \to \mathbb Z$
defined by
\[
  d_{\mathrm{latt}}(x, y) \;=\; \sum_{i \in \mathrm{Fin}\,d}
    \lvert x_i - y_i \rvert_{\mathbb N}, \qquad
  \lvert \cdot \rvert_{\mathbb N} = \texttt{Int.natAbs}.
\]
The basic lemmas \texttt{latticeDistance\_self},
\texttt{latticeDistance\_comm},
\texttt{latticeDistance\_eq\_zero\_iff},
\texttt{latticeDistance\_triangle} give a full pseudo-metric structure,
and \texttt{latticeGraph\_adj\_iff\_latticeDistance\_eq\_one} ties the
new metric to the pre-existing adjacency relation
$\sum_i \lvert x_i - y_i\rvert = 1$ via the identity
$\lvert n \rvert = (\lvert n \rvert_{\mathbb N} : \mathbb Z)$. The
range is kept in $\mathbb N$ rather than in $\mathbb R$ to avoid
$\mathbb N$/$\mathbb Z$/$\mathbb R$ cast bookkeeping in the downstream
summability arguments (Simon--Lieb iteration) that this infrastructure
is meant to support. No mathlib
\texttt{PseudoMetricSpace} instance is registered at this stage.

\paragraph{Finite $\ell^1$ balls and proper-map property.}
PR~\#782 complements the metric definition with the proper-map
companion to PR~\#779's cofinite cluster decay. First,
\texttt{latticeDistance\_le\_finite (d i N)} shows that the
$\ell^1$ ball
\[
  B_N^{\ell^1}(i) \;=\; \{\, j \in \mathrm{Fin}\,d \to \mathbb Z
  \mid d_{\mathrm{latt}}(i, j) \le N \,\}
\]
is finite for every dimension, basepoint, and radius. The proof
uses the coordinatewise bound
$(i_k - j_k)_{\mathbb N} \le d_{\mathrm{latt}}(i, j)$, which places
each $j_k$ in the finite integer interval
$[i_k - N,\, i_k + N]$; the ball then embeds into a finite product
of finite intervals via \texttt{Set.Finite.pi} and
\texttt{Set.finite\_Icc}. Second,
\texttt{tendsto\_latticeDistance\_atTop\_cofinite (d i)} packages
this as a proper-map statement
$\mathrm{Tendsto}\,\bigl(j \mapsto d_{\mathrm{latt}}(i, j)\bigr)\,
\mathrm{Filter.cofinite}\; \mathrm{Filter.atTop}$
by reducing \texttt{Filter.eventually\_cofinite} to finiteness of
$\{\,j \mid d_{\mathrm{latt}}(i, j) < N\,\}$ for every $N$, a
subset of the ball at radius $N$. No hypothesis $d \ge 1$ is
required: at $d = 0$ the domain is a singleton, so
$\mathrm{Filter.cofinite}$ on the source collapses to $\bot$ (in
the mathlib convention, the improper filter containing every
subset) and the Tendsto holds vacuously. Combined with
PR~\#779 this upgrades the cofinite cluster-decay statement into
a distance-based one conditional on summability of
$j \mapsto U_2(i, j)$.

\paragraph{Distance-based cluster decay capstone (PR~\#783).}
PR~\#783 closes the §5.1 cluster-decay infrastructure stack by
combining the previous two ingredients into a single distance-based
statement. The bridge is the filter equality
\[
  \mathrm{Filter.comap}\,\bigl(d_{\mathrm{latt}}(i, \cdot)\bigr)\,
    \mathrm{Filter.atTop} \;=\; \mathrm{Filter.cofinite}
\]
(\texttt{comap\_latticeDistance\_atTop\_eq\_cofinite}). The $\ge$
direction is exactly $\mathrm{Filter.Tendsto.le\_comap}$ applied to
PR~\#782's \texttt{tendsto\_latticeDistance\_atTop\_cofinite}; the
$\le$ direction uses $\mathrm{Set.Finite.bddAbove}$ on $\mathbb N$
to bound $d_{\mathrm{latt}}(i, \cdot)$ on the finite complement
$S^{\mathrm c}$ of any cofinite $S$, then takes
$[M+1, \infty) \in \mathrm{Filter.atTop}$ to witness $S \in
\mathrm{comap}\,d_{\mathrm{latt}}(i, \cdot)\,\mathrm{atTop}$. Then
\[
  \mathrm{Summable}\bigl(j \mapsto U_2(i, j)\bigr)
  \;\Longrightarrow\;
  U_2(i, j) \xrightarrow[d_{\mathrm{latt}}(i, j) \to \infty]{} 0
\]
formally as
\texttt{truncated2Infinite\_latticeGraph\_tendsto\_atTop\_zero\_of\_summable}:
$\mathrm{Tendsto}\bigl(j \mapsto U_2(i, j)\bigr)\,
\mathrm{Filter.comap}\,d_{\mathrm{latt}}(i, \cdot)\,
\mathrm{Filter.atTop}\,(\mathrm{nhds}\,0)$, with the equivalent
$\varepsilon$-$N$ reading: for every $\varepsilon > 0$ there exists
$N \in \mathbb N$ such that $d_{\mathrm{latt}}(i, j) \ge N$
implies $|U_2(i, j)| < \varepsilon$. The Lean proof rewrites the
source filter via the equality lemma and applies PR~\#779's
cofinite version. The statement is a $\mathrm{Summable}$-conditioned
corollary, not a stand-alone Glimm--Jaffe result: it presents the
\S5.1 cluster picture in its distance-based form, with the
$\mathrm{Summable}$ hypothesis a placeholder for the unconditional
summability the Simon--Lieb inequality (Friedli--Velenik
Prop~9.31) is expected to provide in subsequent PRs. This is the
natural infrastructure capstone before tackling that
research-level multi-PR step.

\paragraph{\texttt{LocallyFinite (latticeGraph d)} + ambient
degree bound (PR~\#784).}
On the way to formalising the Simon--Lieb inequality, the ambient
graph $\mathrm{latticeGraph}(d)$ on the infinite vertex set
$\mathrm{Fin}\,d \to \mathbb Z$ is endowed with the typeclass
instance
\[
  \mathrm{Fintype}\,\bigl(\mathrm{latticeGraph}(d).\mathrm{neighborSet}\,v\bigr)
  \quad \text{for every } v
\]
(\texttt{latticeGraph\_neighborSet\_fintype}), which is mathlib's
\texttt{LocallyFinite} property. The witness is
$\mathrm{Fintype.ofFinset}$ on the filter
$(\mathrm{latticeNeighborEnum}\,d\,v).\mathrm{filter}\,
(\mathrm{latticeGraph}(d).\mathrm{Adj}\,v)$, where
$\mathrm{latticeNeighborEnum}$ and the membership lemma
$\mathrm{latticeGraph\_adj\_mem\_neighborEnum}$ already supply the
$2d$-element candidate enumeration. As a corollary,
$\mathrm{latticeGraph}(d).\mathrm{degree}\,v \le 2d$
(\texttt{latticeGraph\_degree\_le}), via
$\mathrm{Finset.card\_le\_card}$ on the inclusion
$\mathrm{neighborFinset}\,v \subseteq \mathrm{latticeNeighborEnum}\,d\,v$
and the existing $\mathrm{latticeNeighborEnum\_card\_le}$. Until
now the project's \texttt{Fintype} reasoning happened only on the
induced subgraph $(\mathrm{inducedGraph}\,(\mathrm{latticeGraph}\,d)\,\Lambda)$
for finite $\Lambda$; with the ambient instance in place, the
next infrastructure layer toward FV Prop~9.31 (outer vertex
boundary of a finite Finset of $\mathbb Z^d$, edge boundary,
$\sup_v \mathrm{degree}$ uniform bounds, etc.) becomes statable
in mathlib's standard idiom.

\paragraph{Outer vertex boundary
\texttt{SimpleGraph.outerVertexBoundary G S} (PR~\#785).}
The Simon--Lieb inequality (FV Prop~9.31)
\[
  \langle \sigma_0 \sigma_x \rangle_+^{\Lambda}
  \;\le\; \sum_{y \in \partial^v_o B}
    \langle \sigma_0 \sigma_y \rangle_+^{B}
    \cdot \langle \sigma_y \sigma_x \rangle_+^{\Lambda}
  \qquad (B \subset \Lambda,\ x \notin B,\ 0 \in B)
\]
indexes its boundary sum by the \emph{outer vertex boundary} of
the inner box $B$. PR~\#785 introduces this notion as a generic
graph-theoretic primitive in mathlib's \texttt{SimpleGraph}
namespace: for any locally finite simple graph $G$ and any Finset
$S$ of vertices,
\[
  \mathrm{SimpleGraph.outerVertexBoundary}\,G\,S
  \;:=\; (S.\mathrm{biUnion}\,G.\mathrm{neighborFinset}) \setminus S
  \;\equiv\; \{\,y \notin S \mid \exists x \in S,\ G.\mathrm{Adj}\,x\,y\,\}.
\]
The four basic lemmas are the membership characterisation
$\mathrm{mem\_outerVertexBoundary\_iff}$, the disjointness
$\mathrm{outerVertexBoundary\_disjoint}$, the empty-set value
$\mathrm{outerVertexBoundary\_empty}$, and the cardinality bound
\[
  |\partial^v_o S| \;\le\; \sum_{x \in S} \deg_G(x)
\]
(\texttt{outerVertexBoundary\_card\_le\_sum\_degrees}). On
$\mathbb Z^d$, the wrapper
$|\partial^v_o S| \le 2 d \cdot |S|$
(\texttt{IsingModel.Ambient.latticeGraph\_outerVertexBoundary\_card\_le})
follows by chaining the generic bound with PR~\#784's per-vertex
$\mathrm{latticeGraph\_degree\_le}$. This is the preparatory
indexing object for the Simon--Lieb statement to plug into
(Epic~\#780, B-iii); the present PR does not yet establish
Simon--Lieb itself.

\paragraph{Inner vertex boundary
\texttt{SimpleGraph.innerVertexBoundary G S} and outer-by-inner
degree bound (PR~\#786).}
The companion of the outer vertex boundary is the inner vertex
boundary, in mathlib's \texttt{SimpleGraph} namespace:
\[
  \mathrm{SimpleGraph.innerVertexBoundary}\,G\,S
  \;:=\; S.\mathrm{filter}\bigl(\lambda x \mapsto \neg (G.\mathrm{neighborFinset}\,x \subseteq S)\bigr)
  \;\equiv\; \{\,x \in S \mid \exists y,\ G.\mathrm{Adj}\,x\,y \land y \notin S\,\}.
\]
Three basic lemmas are
$\mathrm{mem\_innerVertexBoundary\_iff}$,
$\mathrm{innerVertexBoundary\_subset\_self}$ (via
$\mathrm{Finset.filter\_subset}$), and
$\mathrm{innerVertexBoundary\_empty}$. The substantive new
content is the Cheeger-type bound
\[
  |\partial^v_o S| \;\le\; \sum_{x \in \partial^v_i S} \deg_G(x)
\]
(\texttt{outerVertexBoundary\_card\_le\_sum\_degrees\_innerVertexBoundary}).
The proof exhibits the inclusion
$\partial^v_o S \subseteq (\partial^v_i S).\mathrm{biUnion}\,
G.\mathrm{neighborFinset}$ — every $y \in \partial^v_o S$ has
some $x \in S$ adjacent to $y \notin S$, which already places
$x$ in $\partial^v_i S$ — and then chains
$\mathrm{Finset.card\_le\_card}$,
$\mathrm{Finset.card\_biUnion\_le}$, and
$\mathrm{card\_neighborFinset\_eq\_degree}$. On $\mathbb Z^d$
this becomes the elementary max-degree-based linear bound
$|\partial^v_o S| \le 2 d \cdot |\partial^v_i S|$
(\texttt{latticeGraph\_outerVertexBoundary\_card\_le\_two\_mul\_d\_mul\_innerVertexBoundary\_card}),
combining the generic bound with PR~\#784's
$\mathrm{latticeGraph\_degree\_le}$. (This is not the optimal
vertex-isoperimetric inequality on $\mathbb Z^d$, just the
elementary corollary of the per-vertex degree bound.) As with
PR~\#785, this PR is preparatory infrastructure: no Simon--Lieb
result is established here.

\paragraph{Oriented edge boundary
\texttt{SimpleGraph.edgeBoundary G S} (PR~\#787).}
Closes the elementary boundary-primitive infrastructure with the
oriented edge boundary
\[
  \mathrm{SimpleGraph.edgeBoundary}\,G\,S
  \;:=\; (\partial^v_i S).\mathrm{biUnion}\,\bigl(\lambda x.\,
    ((G.\mathrm{neighborFinset}\,x).\mathrm{filter}\,(\cdot \notin S)).\mathrm{image}\,(\lambda y.\,(x, y))\bigr),
\]
of type $\mathrm{Finset}\,(V \times V)$. Each crossing edge
$\{x, y\}$ with $x \in S, y \notin S$ has a unique oriented
representative $(x, y)$, sidestepping the $\mathrm{Sym2}\,V$
quotient handling. The membership characterisation
\[
  (x, y) \in \partial^e S \iff x \in S \land y \notin S \land G.\mathrm{Adj}\,x\,y
\]
($\mathrm{mem\_edgeBoundary\_iff}$) follows by unfolding
$\mathrm{Finset.mem\_biUnion} / \mathrm{mem\_image} /
\mathrm{mem\_filter}$ and the inner-boundary characterisation,
with $\mathrm{Prod.mk.injEq}$ resolving the pair equality. The
empty case $\partial^e \emptyset = \emptyset$ comes from
PR~\#786's $\mathrm{innerVertexBoundary\_empty}$. The
cardinality bound
\[
  |\partial^e S| \;\le\; \sum_{x \in \partial^v_i S} \deg_G(x)
\]
($\mathrm{edgeBoundary\_card\_le\_sum\_degrees\_innerVertexBoundary}$)
chains $\mathrm{Finset.card\_biUnion\_le}$,
$\mathrm{card\_image\_le}$,
$\mathrm{card\_filter\_le}$, and
$\mathrm{card\_neighborFinset\_eq\_degree}$. On $\mathbb Z^d$
this gives
$|\partial^e S| \le 2 d \cdot |\partial^v_i S|$
($\mathrm{latticeGraph\_edgeBoundary\_card\_le\_two\_mul\_d\_mul\_innerVertexBoundary\_card}$),
the linear edge-boundary bound from the per-vertex degree bound.
Together with PRs~\#785 / \#786 this closes the elementary
boundary-primitive infrastructure for any locally finite simple
graph; subsequent Simon--Lieb work (Epic~\#780, B-iii) is
research-level.

\paragraph{Edge boundary card disjoint
$\mathrm{biUnion}$ equality (PR~\#788).}
The previous bound is tightened to an equality
\[
  |\partial^e S| \;=\;
    \sum_{x \in \partial^v_i S} \bigl|\{\,y \in N_G(x) \mid y \notin S\,\}\bigr|
\]
($\mathrm{edgeBoundary\_card\_eq\_sum\_inner\_filter}$). The
$\mathrm{biUnion}$ defining $\partial^e S$ is disjoint across
$x \in \partial^v_i S$ (different first coordinates), and the
embedding $y \mapsto (x, y)$ is injective for fixed $x$, so
$\mathrm{Finset.card\_biUnion}$ applies and
$\mathrm{Finset.card\_image\_of\_injective}$ replaces each
image cardinality with its filter cardinality. Bounding the
filter cardinality by
$|G.\mathrm{neighborFinset}\,x| = \deg_G(x)$ recovers PR~\#787's
inequality.

\paragraph{Outer-side companion equality (PR~\#789).}
The symmetric equality counts crossing edges by their outside
endpoint:
\[
  |\partial^e S| \;=\;
    \sum_{y \in \partial^v_o S} \bigl|\{\,x \in N_G(y) \mid x \in S\,\}\bigr|
\]
($\mathrm{edgeBoundary\_card\_eq\_sum\_outer\_filter}$). The proof
rewrites $\partial^e S$ as a disjoint $\mathrm{biUnion}$ indexed by
$y \in \partial^v_o S$ (via $\mathrm{Finset.ext}$, with forward /
backward using $G.\mathrm{symm}$ to swap the adjacency direction),
then applies the same $\mathrm{Finset.card\_biUnion}$ +
$\mathrm{Finset.card\_image\_of\_injective}$ pair. Combined with
PR~\#788 this yields the double-counting identity
$\sum_{x \in \partial^v_i S} |N_G(x) \setminus S| =
\sum_{y \in \partial^v_o S} |N_G(y) \cap S|$.

\paragraph{\S17.1 / \S17.5 lattice mass foundation (PR~\#793).}
GJ \S17.1 introduces the mass $m(\sigma)$ as the exponential
decay rate (17.1.5) of $\langle\phi(x)\phi(y)\rangle - M(\sigma)^2$.
On the lattice we formalise the foundation of this notion as a
predicate:
\[
  \mathrm{HasExponentialDecay}(d, \Lambda, p, \alpha) \;:=\;
    \exists\, C \ge 0,\ \forall\, i \ne j,\
      |U_2(i, j)| \le C \cdot e^{-\alpha\, d_{\mathrm{latt}}(i, j)}.
\]
($\mathrm{IsingModel.Ambient.HasExponentialDecay}$, in
\texttt{Concrete/LatticeGraphCorrelation/LatticeMassFoundation.lean}.)
Bundled
companions: $\mathrm{HasExponentialDecay\_beta\_zero}$ (at
$\beta = 0$, witness $C = 0$ since $U_2 \equiv 0$),
$\mathrm{HasExponentialDecay\_J\_zero}$ (at $J = 0$ ferromagnetic,
witness $C = 0$ since $U_2 = 0$ off-diagonal), and
$\mathrm{HasExponentialDecay\_mono}$ ($\alpha$-monotonicity:
decreasing $\alpha$ weakens the bound while the constant $C$ is
preserved). The general non-trivial-slice exponential decay
($m(\beta) > 0$ for $\beta < \beta_c$) requires Simon--Lieb
(Issue~\#780, research-level).

\paragraph{$\mathrm{clusterProperty}$ linkage (PR~\#794).}
PR~\#794 closes the link: for $\alpha > 0$,
$\mathrm{HasExponentialDecay}(d, \Lambda, p, \alpha)$ implies
PR~\#792's $\mathrm{clusterProperty}(\mathrm{latticeGraph}\,d, \Lambda, p)$
($\mathrm{clusterProperty\_latticeGraph\_of\_HasExponentialDecay}$).
The proof composes PR~\#782's
$\mathrm{tendsto\_latticeDistance\_atTop\_cofinite}$ with the
mathlib chain
$\mathrm{tendsto\_id.const\_mul\_atTop}$ +
$\mathrm{Real.tendsto\_exp\_neg\_atTop\_nhds\_zero}$ to obtain
$(j \mapsto C \cdot e^{-\alpha\, d_{\mathrm{latt}}(i, j)}) \to 0$
along $\mathrm{cofinite}$, then applies the squeeze
$\mathrm{tendsto\_of\_tendsto\_of\_tendsto\_of\_le\_of\_le'}$
between $-g$ and $g$ via $|U_2| \le g$. Together PR~\#793 and
PR~\#794 give the full §17.1 / §17.5 lattice mass foundation
including the precise relationship between exponential decay and
the cluster-property predicate.

\paragraph{$\mathrm{latticeMass}$ formal definition (PR~\#795).}
The mass / inverse correlation length is finally defined as a
supremum:
\[
  \mathrm{latticeMass}(d, \Lambda, p) \;:=\;
    \sup\bigl(\{\, (\alpha : \mathrm{ENNReal}) \mid \alpha : \mathrm{NNReal},\
       \mathrm{HasExponentialDecay}(d, \Lambda, p, \alpha) \,\}\bigr)
\]
as $\mathrm{IsingModel.Ambient.latticeMass}\,d\,\Lambda\,p :
\mathrm{ENNReal}$. Trivial-slice values
$\mathrm{latticeMass\_top\_of\_beta\_zero}$ and
$\mathrm{latticeMass\_top\_of\_J\_zero}$ both give
$\mathrm{latticeMass} = \top$, matching the physical picture:
no correlation $\Rightarrow$ infinite mass $\Rightarrow$ zero
correlation length. The proof: any candidate upper bound
$b \ne \top$ is exceeded by $\alpha := b.\mathrm{toNNReal} + 1$,
which sits in the supremand by the corresponding
$\mathrm{HasExponentialDecay}$ trivial-slice theorem; this
contradicts $b$ being an upper bound. Sanity:
$\mathrm{latticeMass\_nonneg}$ is $\mathrm{bot\_le}$.
The bridge lemma
$\mathrm{latticeMass\_ge\_of\_HasExponentialDecay}$ records the
direct lower-bound direction of this definition: if a non-negative
real rate $\alpha$ validates
$\mathrm{HasExponentialDecay}(d,\Lambda,p,\alpha)$, then
$\mathrm{ENNReal.ofReal}\,\alpha \le \mathrm{latticeMass}(d,\Lambda,p)$.
The companion $\mathrm{latticeMass\_pos\_of\_HasExponentialDecay}$
specialises this to a strictly positive rate and concludes
$0 < \mathrm{latticeMass}(d,\Lambda,p)$. These package the
$\mathrm{le\_sSup}$ argument needed in the later GJ \S17.5
Lemma~17.5.2 lower-bound bridge.
These foundation declarations were split from the legacy
\texttt{Concrete/LatticeGraphCorrelation/Inequalities.lean} path into
\texttt{Concrete/LatticeGraphCorrelation/LatticeMassFoundation.lean}
in PR~\#1925; the legacy module still imports the child module for
compatibility.
The transfer theorem
$\mathrm{HasExponentialDecay\_transfer\_exhaustion}$ records that,
under ferromagnetic parameters, a validating rate for one exhaustion
also validates the same rate for any other exhaustion. Its proof is
just $\mathrm{truncated2Infinite\_indep\_exhaustion}$ followed by
reuse of the same constant and rate.
The conditional Step~117l bridge
$\mathrm{latticeMass\_ge\_pseudoMassFromParamsAtPair\_of\_decay}$
packages the final supremum step on the lower-bound side of
Glimm--Jaffe \S17.5 Lemma~17.5.2 (2nd ed., pp.~311--312): once the
concrete pair pseudo-mass at $h=0$ is supplied as a validating
$\mathrm{HasExponentialDecay}$ rate, its
$\mathrm{ENNReal.ofReal}$ value is bounded by
$\mathrm{latticeMass}$. The companion
$\mathrm{latticeMass\_pos\_of\_pseudoMassFromParamsAtPair\_decay}$
turns a positive validating pseudo-mass into $0 < \mathrm{latticeMass}$.
The high-temperature Simon--Lieb estimate now has an exhaustion-uniform
form.  The transfer, pseudo-mass comparison, critical-temperature, and
product-summability bridge declarations in this paragraph were moved from
\texttt{Concrete/LatticeGraphCorrelation/Inequalities.lean} to
\texttt{Concrete/LatticeGraphCorrelation/LatticeMassPseudoMassTransfer.lean}
in PR~\#1927 while preserving the legacy import path:
$\mathrm{HasExponentialDecay\_transfer\_high\_temp}$ transfers the
rate $-\log(\beta J\,2d)$ from the cubic exhaustion to any exhaustion,
using the exhaustion-independence of the infinite-volume truncated
two-point function. Consequently
$\mathrm{latticeMass\_ge\_neg\_log\_of\_high\_temp\_exhaustion}$
gives the same explicit lower bound for every exhaustion, and
$\mathrm{latticeMass\_pos\_of\_high\_temp\_exhaustion}$ gives positive
mass for $d \ge 1$ in the strict high-temperature regime. This discharges the
exhaustion-uniform high-temperature decay transfer used by the
Step~117l bridge.
The finite-susceptibility and Lebowitz derivative bound declarations
used by the subsequent Lipschitz layer now live in
\texttt{Concrete/LatticeGraphCorrelation/LatticeMassLebowitzDerivative.lean}
(moved from \texttt{Concrete/LatticeGraphCorrelation/Inequalities.lean}
in PR~\#1928), while the legacy import path remains available.
The high-temperature Lipschitz, compact/open continuity, uniform and locally
uniform convergence, and almost-everywhere differentiability wrappers now live
in
\texttt{Concrete/LatticeGraphCorrelation/LatticeMassHighTempLipschitz.lean}
(moved from \texttt{Concrete/LatticeGraphCorrelation/Inequalities.lean}
in PR~\#1929), again preserving the legacy import path.
The zero-boundary linear bounds, closed-interval and half-open continuity,
zero-included Lipschitz, half-line differentiability, and half-open locally
uniform convergence wrappers now live in
\texttt{Concrete/LatticeGraphCorrelation/LatticeMassHighTempZeroBoundary.lean}
(moved from \texttt{Concrete/LatticeGraphCorrelation/Inequalities.lean}
in PR~\#1930), with the legacy import path preserved.
The high-temperature regularity wrappers for the infinite-volume Ursell
two-point function $\mathrm{truncated2Infinite}$ at $h = 0$ -- continuity on
$\mathrm{Ioo}\,0\,\beta_c$, $\mathrm{Icc}\,0\,b$, and $\mathrm{Ico}\,0\,\beta_c$
in $\beta$ (Step~185) and in $J$ (Step~239), Lipschitz / a.e.\
differentiable / monotone wrappers in $\beta$ (Steps~186--187) and in $J$
(Step~240), and full-neighborhood $\mathrm{ContinuousAt}$ wrappers at every
interior point of $\mathrm{Ioo}\,0\,\beta_c$ and $\mathrm{Ioo}\,0\,J_c$
(Step~241) -- now live in
\texttt{Concrete/LatticeGraphCorrelation/LatticeMassTruncated2HighTemp.lean}
(moved from \texttt{Concrete/LatticeGraphCorrelation/Inequalities.lean}
in PR~\#1931), with the legacy import path preserved.
The comparison form
$\mathrm{HasExponentialDecay\_pseudoMassFromParamsAtPair\_of\_le\_high\_temp\_rate}$
then records the remaining order-theoretic step: if the concrete pair
pseudo-mass is at most $-\log(\beta J\,2d)$, monotonicity of
$\mathrm{HasExponentialDecay}$ turns the transferred Simon--Lieb rate into a
validating pseudo-mass rate. Consequently
$\mathrm{latticeMass\_ge\_pseudoMassFromParamsAtPair\_of\_le\_high\_temp\_rate}$
gives $\mathrm{ENNReal.ofReal}(m^-_{x,z}) \le \mathrm{latticeMass}$, and
$\mathrm{latticeMass\_pos\_of\_pseudoMassFromParamsAtPair\_le\_high\_temp\_rate}$
gives positive lattice mass when the pseudo-mass is positive. The remaining
substantive input is the quantitative comparison of the concrete pseudo-mass
with the transferred Simon--Lieb rate.
The exhaustion-transfer refinements
$\mathrm{HasExponentialDecay\_pseudoMassFromParamsAtPair\_of\_exhaustion\_le\_high\_temp\_rate}$,
$\mathrm{latticeMass\_ge\_pseudoMassFromParamsAtPair\_of\_exhaustion\_le\_high\_temp\_rate}$,
and
$\mathrm{latticeMass\_pos\_of\_pseudoMassFromParamsAtPair\_exhaustion\_le\_high\_temp\_rate}$
allow that comparison to be checked on a reference exhaustion
$\Lambda_0$ while producing the decay and mass conclusions on a target
exhaustion $\Lambda$.  Their cubic-reference specialisations replace
$\Lambda_0$ by $\mathrm{cubicExhaustion}\,d$, using
$\mathrm{pseudoMassFromParamsAtPair\_indep\_exhaustion}$ under ferromagnetic
parameters.  This packages the canonical-exhaustion comparison form needed for
later Step~117l work, but it still does not prove the quantitative comparison
itself.
The reference-rate variants
$\mathrm{HasExponentialDecay\_reference\_pseudoMassFromParamsAtPair\_of\_exhaustion\_le\_high\_temp\_rate}$,
$\mathrm{latticeMass\_ge\_reference\_pseudoMassFromParamsAtPair\_of\_exhaustion\_le\_high\_temp\_rate}$,
and
$\mathrm{latticeMass\_pos\_of\_reference\_pseudoMassFromParamsAtPair\_exhaustion\_le\_high\_temp\_rate}$
state the same target-exhaustion conclusions using the reference pseudo-mass
itself.  The corresponding cubic-reference forms make a later comparison
proved on $\mathrm{cubicExhaustion}\,d$ immediately usable as a target
decay rate and target $\mathrm{latticeMass}$ lower bound for arbitrary
exhaustions.

The profile-comparison refinement
\path{pseudoMassFromParamsAtPair_le_high_temp_rate_of_pseudoMassG_le_corr}
reduces this remaining quantitative comparison to a direct pair-correlation
lower bound.  If
$C_\Lambda(x,z)=\mathrm{correlationInfinite}(\mathrm{latticeGraph}\,d,\Lambda,
\langle J,0,\beta\rangle,\{x,z\})$ lies in $(0,2)$ and
\[
  \mathrm{pseudoMassG}\,\alpha\,r\,(-\log(\beta J\,2d))
    \le C_\Lambda(x,z),
\]
then $m^-_{\Lambda,x,z}\le -\log(\beta J\,2d)$.  The proof uses
$0\le -\log(\beta J\,2d)$ from $0\le J$, $0<\beta$, and
$\beta J\,2d<1$, rewrites $\mathrm{pseudoMassExt}$ to
$\mathrm{pseudoMass}$ on the active interval $(0,2)$, and applies the
implicit characterization
$\mathrm{pseudoMass}(c)\le t \iff \mathrm{pseudoMassG}(\alpha,r,t)\le c$.

The companion profile-bound theorems compose this comparison with the
transferred high-temperature rate.  Their reference-exhaustion and
cubic-reference variants provide the same target decay and target
$\mathrm{latticeMass}$ conclusions when the profile lower bound is checked on
$\Lambda_0$ or on $\mathrm{cubicExhaustion}\,d$.

For the canonical anchored cubic pair, the lower-bound input can be reduced
one step further to a tanh-power estimate.  The theorem
\path{pseudoMassG_le_cubic_correlation_of_le_tanh_pow_dist} uses the
Glimm--Jaffe \S17.5 path lower bound, pp.~304--306,
\[
  \tanh(\beta J)^{\mathrm{latticeDistance}\,d\,0\,z}
    \le \mathrm{twoPointFunction}\,d\,\langle J,0,\beta\rangle\,z,
\]
and the definition of $\mathrm{twoPointFunction}$ as the cubic
$\mathrm{correlationInfinite}$ at $\{0,z\}$.  Hence the concrete numerical
condition
\[
  \mathrm{pseudoMassG}\,\alpha\,r\,(-\log(\beta J\,2d))
    \le \tanh(\beta J)^{\mathrm{latticeDistance}\,d\,0\,z}
\]
implies the pair-correlation lower bound required by the profile-comparison
bridge above.  This is still a conditional reduction toward Lemma~17.5.2
rather than the full lower-bound side of the lemma: it replaces only the direct
infinite-volume correlation lower-bound hypothesis by an explicit tanh-power
comparison, while the active-range condition
$C_{\mathrm{cubic}}(0,z)\in(0,2)$ remains a separate input to the
profile-comparison bridge.

The positivity theorem
\path{correlationInfinite_cubic_pair_pos_of_pseudoMassG_le_tanh_pow_dist}
closes the lower endpoint of that active interval under the same tanh-power
profile condition.  For Lean's total real logarithm, $J\ge 0$, $\beta>0$, and
$\beta J\,2d<1$ give $0\le -\log(\beta J\,2d)$.  Since $r>0$,
$\mathrm{pseudoMassG}$ is strictly positive at this nonnegative rate; the
tanh-power reduction then transfers this strict positivity to
$C_{\mathrm{cubic}}(0,z)$.  The upper endpoint $C_{\mathrm{cubic}}(0,z)<2$ is
provided by the general $\mathrm{correlationInfinite}<2$ API and is still kept
as a separate composition step.

\paragraph{Cubic active range from tanh profile bounds (PR~\#1819).}
The cubic active-range package makes that composition explicit without
unfolding the concrete pseudo-mass wrapper.  The theorem
\path{correlationInfinite_cubic_pair_mem_Ioo_zero_two_of_pseudoMassG_le_tanh_pow_dist}
combines
\path{correlationInfinite_cubic_pair_pos_of_pseudoMassG_le_tanh_pow_dist}
with the universal estimate
\[
  C_{\mathrm{cubic}}(0,z) \le 1 < 2
\]
to prove
\[
  C_{\mathrm{cubic}}(0,z)\in(0,2).
\]
The companion forms
\path{correlationInfinite_cubic_pair_ne_zero_of_pseudoMassG_le_tanh_pow_dist},
\path{correlationInfinite_cubic_pair_mem_Ioc_zero_one_of_pseudoMassG_le_tanh_pow_dist},
and
\path{correlationInfinite_cubic_pair_lt_two_of_pseudoMassG_le_tanh_pow_dist}
record non-vanishing, the sharper interval $(0,1]$, and the isolated upper
endpoint.  These statements provide the active-range input required by the
profile-comparison bridge while avoiding the heavy direct definitional
identification of \(\mathrm{pseudoMassFromParamsAtPair}\) with the anchored
two-point function.

\paragraph{Profile-compatible high-temperature lattice-mass corollaries
(PR~\#1820, moved to a narrower child module in PR~\#1921).}
The small module
\path{Concrete/LatticeGraphCorrelation/CubicPseudoMassTanhProfile.lean}
records target-level corollaries under the same anchored cubic tanh-profile
hypotheses.  The theorem
\path{latticeMass_ge_neg_log_of_high_temp_exhaustion_of_cubic_tanh_profile}
checks the active-range theorem internally, then returns
the transferred Simon--Lieb lower bound
\[
  \mathrm{ofReal}(-\log(\beta J\,2d))
    \le \mathrm{latticeMass}(d,\Lambda,\langle J,0,\beta\rangle).
\]
With \(d\ge 1\) and \(0<\beta J\), the companion
\path{latticeMass_pos_of_high_temp_exhaustion_of_cubic_tanh_profile}
returns \(0<\mathrm{latticeMass}\).  These are deliberately target-level
profile-compatible restatements of the high-temperature transfer theorems:
direct new statements exposing the concrete
\(\mathrm{pseudoMassFromParamsAtPair}\) value still trigger heavy elaboration,
so the concrete pseudo-mass comparison itself remains a separate follow-up.

\paragraph{Anchored cubic pseudo-mass named-rate bridge
(PR~\#1821, moved to a narrower child module in PR~\#1922).}
The basic cubic pseudo-mass module
\path{Concrete/LatticeGraphCorrelation/CubicPseudoMassBasic.lean}
names the concrete anchored cubic pseudo-mass at the pair \((0,z)\) as
\path{cubicOriginPseudoMassFromParamsAtPair}, with arguments
\(h_\alpha\), \(h_r\), \(\beta\), \(J\), and \(z\).
The definition is accompanied by
\path{cubicOriginPseudoMassFromParamsAtPair_eq} and
\path{cubicOriginPseudoMassFromParamsAtPair_nonneg}.  The cubic named-rate
module
\path{Concrete/LatticeGraphCorrelation/CubicPseudoMassNamedRate.lean}
contains the conditional bridge
\path{HasExponentialDecay_cubicOriginPseudoMassFromParamsAtPair_of_le_high_temp_rate}
which states that, once this named rate is bounded by
\(-\log(\beta J\,2d)\), it validates target-exhaustion exponential decay.
Consequently
\path{latticeMass_ge_cubicOriginPseudoMassFromParamsAtPair_of_le_high_temp_rate}
gives the target \(\mathrm{latticeMass}\) lower bound, and
\path{latticeMass_pos_of_cubicOriginPseudoMassFromParamsAtPair_le_high_temp_rate}
gives positive target \(\mathrm{latticeMass}\) from a positive named rate.
This isolates the stable high-temperature transfer step without restating the
full high-arity concrete pseudo-mass expression in the theorem conclusions.
PR~\#1822 adds the transport lemmas
\path{cubicOriginPseudoMassFromParamsAtPair_le_iff},
\path{cubicOriginPseudoMassFromParamsAtPair_lt_iff}, and
\path{cubicOriginPseudoMassFromParamsAtPair_eq_iff}.  PR~\#1825 completes
the same transport layer with
\path{cubicOriginPseudoMassFromParamsAtPair_ge_iff},
\path{cubicOriginPseudoMassFromParamsAtPair_gt_iff}, and
\path{cubicOriginPseudoMassFromParamsAtPair_ne_iff}.  These rewrite
order, equality, and non-equality comparisons between the named rate and the
underlying concrete
\(\mathrm{pseudoMassFromParamsAtPair}\) expression, keeping downstream theorem
statements on the named-rate API while still permitting explicit conversion
when needed.

PR~\#1823 connects that named-rate API to the cubic active-range and profile
lower-bound inputs.  The theorem
\path{cubicOriginPseudoMassFromParamsAtPair_le_high_temp_rate_of_cubic_pseudoMassG_le_corr}
proves that the named rate is at most \(-\log(\beta J\,2d)\) whenever the
anchored cubic pair correlation lies in \((0,2)\) and dominates
\(\mathrm{pseudoMassG}\) at that rate.  The companion
\path{cubicOriginPseudoMassFromParamsAtPair_pos_of_cubic_corr_mem} records
positivity of the named pseudo-mass from the same active-range membership.
These feed
\path{HasExponentialDecay_cubicOriginPseudoMassFromParamsAtPair_of_cubic_pseudoMassG_le_corr},
\path{latticeMass_ge_cubicOriginPseudoMassFromParamsAtPair_of_cubic_pseudoMassG_le_corr},
and
\path{latticeMass_pos_of_cubicOriginPseudoMassFromParamsAtPair_cubic_pseudoMassG_le_corr},
which give the target-exhaustion decay, lower-bound, and positivity
consequences directly at the named anchored cubic pseudo-mass rate without
restating the heavy tanh-profile reduction in these theorem statements.
PR~\#1824 records the reusable active-range positivity layer:
\path{cubicOriginPseudoMassFromParamsAtPair_ne_zero_of_cubic_corr_mem},
\path{ENNReal_ofReal_cubicOriginPseudoMassFromParamsAtPair_pos_of_cubic_corr_mem},
and
\path{latticeMass_pos_of_cubicOriginPseudoMassFromParamsAtPair_cubic_corr_mem_le_high_temp_rate}.
PR~\#1825 adds the matching non-vanishing conveniences
\path{ENNReal_ofReal_cubicOriginPseudoMassFromParamsAtPair_ne_zero_of_cubic_corr_mem},
\path{latticeMass_ne_zero_of_cubicOriginPseudoMassFromParamsAtPair_le_high_temp_rate},
\path{latticeMass_ne_zero_of_cubicOriginPseudoMassFromParamsAtPair_cubic_pseudoMassG_le_corr},
and
\path{latticeMass_ne_zero_of_cubic_corr_mem_le_high_temp_rate}.
Thus active-range membership alone gives a nonzero named pseudo-mass and
positive/nonzero \(\mathrm{ofReal}\) coercions, while the final positive and
nonzero \(\mathrm{latticeMass}\) bridges only ask for the named-rate comparison
with \(-\log(\beta J\,2d)\) or the existing cubic profile lower-bound package.
PR~\#1826 packages the same lower-bound and positivity information as interval
membership.  The closed-interval theorems
\begin{Verbatim}[breaklines=true,breakanywhere=true]
cubicNamedRate_ofReal_mem_Icc_latticeMass_of_le_high_temp_rate
cubicNamedRate_ofReal_mem_Icc_latticeMass_of_cubic_pseudoMassG_le_corr
\end{Verbatim}
cover the direct high-temperature comparison and cubic-profile comparison cases:
\[
  \operatorname{ofReal}(m^-_{\mathrm{cub}}) \in
  [0,\mathrm{latticeMass}],
\]
while the half-open theorems
\begin{Verbatim}[breaklines=true,breakanywhere=true]
cubicNamedRate_ofReal_mem_Ioc_latticeMass_of_pos_le_high_temp_rate
cubicNamedRate_ofReal_mem_Ioc_latticeMass_of_cubic_pseudoMassG_le_corr
cubicNamedRate_ofReal_mem_Ioc_latticeMass_of_corr_mem_le_high_temp_rate
\end{Verbatim}
cover the positive comparison, cubic-profile, and active-range plus
high-temperature comparison cases:
\(\operatorname{ofReal}(m^-_{\mathrm{cub}}) \in
(0,\mathrm{latticeMass}]\) form.  These statements are pure compositions of the
named-rate comparison and active-range positivity bridges, so they avoid the
known heavy tanh-profile elaboration path.
PR~\#1827 adds conjunction forms that package these interval statements with
the validating decay proof itself:
\begin{Verbatim}[breaklines=true,breakanywhere=true]
cubicNamedRate_decay_mem_Icc_of_le_high_temp_rate
cubicNamedRate_decay_mem_Ioc_of_pos_le_high_temp_rate
cubicNamedRate_decay_mem_Icc_of_cubic_pseudoMassG_le_corr
cubicNamedRate_decay_mem_Ioc_of_cubic_pseudoMassG_le_corr
cubicNamedRate_decay_mem_Ioc_of_corr_mem_le_high_temp_rate
\end{Verbatim}
Each theorem returns
\[
  \mathrm{HasExponentialDecay}\bigl(d,\Lambda,\langle J,0,\beta\rangle,
    m^-_{\mathrm{cub}}\bigr)
  \quad\text{and}\quad
  \operatorname{ofReal}(m^-_{\mathrm{cub}})
    \in [0,\mathrm{latticeMass}]
\]
or the corresponding half-open interval
\((0,\mathrm{latticeMass}]\).  The proofs only pair the named-rate decay bridge
with the PR~\#1826 interval package, leaving the tanh-profile reduction outside
the new theorem statements. PR~\#2002 narrows the
\path{*_of_cubic_pseudoMassG_le_corr},
\path{*_of_cubic_corr_mem_le_high_temp_rate},
\path{cubicNamedRate_decay_mem_*_of_cubic_pseudoMassG_le_corr},
\path{cubicNamedRate_decay_mem_*_of_le_high_temp_rate},
\path{cubicNamedRate_decay_mem_Ioc_of_pos_le_high_temp_rate}, and
\path{cubicNamedRate_decay_mem_Ioc_of_corr_mem_le_high_temp_rate}
families (17 wrappers from the PR~\#1820/\#1821/\#1824/\#1825/\#1826/\#1827
phases) into the narrower
\path{Concrete/LatticeGraphCorrelation/CubicPseudoMassNamedRateCorr.lean}
child module.  The legacy import path is preserved.
PR~\#1828 packages the downstream cluster and summability consequences, moved
to the narrower module
\path{Concrete/LatticeGraphCorrelation/CubicPseudoMassClusterSummability.lean}
in PR~\#1923.  The named summand
\path{cubicTruncated2Product} abbreviates
\[
  w \mapsto U_2(x,w)\,U_2(y,w)
\]
on the cubic exhaustion at zero external field, keeping the summability
statements lightweight.  The cluster-property consequences are
\begin{Verbatim}[breaklines=true,breakanywhere=true]
clusterProperty_of_cubicNamedRate_pos_le_high_temp_rate
clusterProperty_of_cubicNamedRate_cubic_pseudoMassG_le_corr
clusterProperty_of_cubicNamedRate_corr_mem_le_high_temp_rate
\end{Verbatim}
and the cubic product-summability consequences are
\begin{Verbatim}[breaklines=true,breakanywhere=true]
summable_truncated2Infinite_prod_of_cubicNamedRate_pos_le_high_temp_rate
summable_truncated2Infinite_prod_of_cubicNamedRate_cubic_pseudoMassG_le_corr
summable_truncated2Infinite_prod_of_cubicNamedRate_corr_mem_le_high_temp_rate
\end{Verbatim}
The proofs compose the named-rate decay bridges with
\path{clusterProperty_latticeGraph_of_HasExponentialDecay} and the Step~127
theorem \path{summable_truncated2Infinite_prod_of_hasExponentialDecay}.
PR~\#1829 packages the corresponding explicit product-sum upper bounds, moved
to the narrower module
\path{Concrete/LatticeGraphCorrelation/CubicPseudoMassProductSum.lean} in
PR~\#1924.  Once the named rate \(m\) supplies a positive exponential-decay
witness, Step~127
\path{tsum_truncated2Infinite_prod_le} yields a constant \(C \ge 0\) such that
\[
  \sum_w \operatorname{cubicTruncated2Product}(d,\beta,J,x,y,w)
  \le (C+1)^2
    \left(2\sum_w e^{-(m/2)d(0,w)}\right)
    e^{-(m/2)d(x,y)/2}.
\]
The three packaged forms are
\begin{Verbatim}[breaklines=true,breakanywhere=true]
exists_tsum_cubicTruncated2Product_le_of_cubicNamedRate_pos_le_high_temp_rate
exists_tsum_cubicTruncated2Product_le_of_cubicNamedRate_cubic_pseudoMassG_le_corr
exists_tsum_cubicTruncated2Product_le_of_cubicNamedRate_corr_mem_le_high_temp_rate
\end{Verbatim}
covering the conditional named-rate comparison, the cubic active-range/profile
input, and active range plus direct named-rate comparison.
PR~\#1830 names the anchored tanh-power profile condition as
\path{cubicTanhProfileBound}.  The unfolding theorem
\path{cubicTanhProfileBound_iff} records that it is exactly
\[
  \mathrm{pseudoMassG}(\alpha,r,-\log(\beta J\,2d))
  \le \tanh(\beta J)^{d(0,z)}.
\]
The bridge theorems
\begin{Verbatim}[breaklines=true,breakanywhere=true]
cubicTanhProfileBound_le_cubic_correlation
correlationInfinite_cubic_pair_pos_of_cubicTanhProfileBound
\end{Verbatim}
feed this named condition into the existing tanh-profile lower-bound and
positivity API.  This keeps future statements from restating the heavy
tanh-profile expression while preserving the same GJ~§17.5 Lemma~17.5.2 input.
PR~\#1831 extends this named predicate to the cubic active-range package:
\begin{Verbatim}[breaklines=true,breakanywhere=true]
correlationInfinite_cubic_pair_mem_Ioo_zero_two_of_cubicTanhProfileBound
correlationInfinite_cubic_pair_ne_zero_of_cubicTanhProfileBound
correlationInfinite_cubic_pair_mem_Ioc_zero_one_of_cubicTanhProfileBound
correlationInfinite_cubic_pair_lt_two_of_cubicTanhProfileBound
\end{Verbatim}
These wrappers expose the active interval \((0,2)\), nonzero correlation, the
sharper interval \((0,1]\), and the isolated upper endpoint \(<2\) from
\path{cubicTanhProfileBound}.  They are direct compositions of the existing
tanh-profile active-range API, so later named-rate statements can use the named
condition without repeating the full tanh-power formula.
PR~\#1832 packages the named tanh-profile predicate as the two bridge inputs
consumed by the anchored cubic profile-to-named-rate comparison:
\begin{Verbatim}[breaklines=true,breakanywhere=true]
cubicTanhProfileBound_cubic_corr_mem_Ioo_and_profile
\end{Verbatim}
It returns active-range membership of the anchored cubic pair correlation
together with the cubic correlation lower bound supplied by
\path{cubicTanhProfileBound_le_cubic_correlation}.  Directly concluding the
named pseudo-mass comparison from \path{cubicTanhProfileBound} still triggers
deterministic elaboration timeouts, so this theorem keeps the public API on the
stable bridge-input layer while avoiding repetition of the tanh-power formula.
PR~\#1833 adds the corresponding eliminator
\begin{Verbatim}[breaklines=true,breakanywhere=true]
cubicTanhProfileBound_elim_cubic_corr_mem_Ioo_and_profile
\end{Verbatim}
which proves an arbitrary proposition once it can be proved from those two
bridge inputs.  This is a proof-engineering wrapper around the PR~\#1832
bundle: downstream statements can receive the active-range membership and
profile lower bound without restating the tanh-power condition or placing the
heavy named-rate comparison directly in the theorem statement.
PR~\#1834 adds named-predicate high-temperature lattice-mass consequences:
\begin{Verbatim}[breaklines=true,breakanywhere=true]
latticeMass_ge_neg_log_of_high_temp_exhaustion_of_cubicTanhProfileBound
latticeMass_pos_of_high_temp_exhaustion_of_cubicTanhProfileBound
latticeMass_ge_neg_log_and_pos_of_high_temp_exhaustion_of_cubicTanhProfileBound
\end{Verbatim}
These are the \path{cubicTanhProfileBound} input forms of the existing final
high-temperature lower-bound and positivity theorems for \path{latticeMass}.
They deliberately keep the conclusion at the stable high-temperature
\(-\log(\beta J\,2d)\) layer: exposing the heavier
\path{cubicOriginPseudoMassFromParamsAtPair} named-rate conclusion directly from
\path{cubicTanhProfileBound} still triggers deterministic elaboration timeouts.
PR~\#1835 names the heavy comparison itself as an irreducible proposition:
\begin{Verbatim}[breaklines=true,breakanywhere=true]
cubicOriginNamedRateLeHighTemp
cubicOriginNamedRateLeHighTemp_of_cubic_pseudoMassG_le_corr
\end{Verbatim}
The proposition abbreviates the comparison
\[
  \mathrm{cubicOriginPseudoMassFromParamsAtPair}(\alpha,r,\beta,J,z)
  \le -\log(\beta J\,2d)
\]
without requiring downstream theorem conclusions to expose that concrete
expression directly.  The active-range/profile-input bridge proves the named
proposition from the same two inputs consumed by the existing profile-to-rate
theorem.  A direct \path{cubicTanhProfileBound} wrapper for this named
proposition still triggers deterministic \path{whnf} timeouts, so it is not part
of the public API.
PR~\#1836 adds downstream consumers for the irreducible proposition:
\begin{Verbatim}[breaklines=true,breakanywhere=true]
cubicOriginPseudoMassFromParamsAtPair_le_high_temp_rate_of_cubicOriginNamedRateLeHighTemp
HasExponentialDecay_cubicOriginPseudoMassFromParamsAtPair_of_cubicOriginNamedRateLeHighTemp
latticeMass_ge_cubicOriginPseudoMassFromParamsAtPair_of_cubicOriginNamedRateLeHighTemp
cubicNamedRate_ofReal_mem_Icc_latticeMass_of_cubicOriginNamedRateLeHighTemp
latticeMass_pos_of_cubicOriginNamedRateLeHighTemp_cubic_corr_mem
cubicNamedRate_ofReal_mem_Ioc_latticeMass_of_cubicOriginNamedRateLeHighTemp
latticeMass_ne_zero_of_cubicOriginNamedRateLeHighTemp_cubic_corr_mem
cubicNamedRate_decay_mem_Icc_of_cubicOriginNamedRateLeHighTemp
cubicNamedRate_decay_mem_Ioc_of_cubicOriginNamedRateLeHighTemp
\end{Verbatim}
These wrappers let downstream proofs use the named proposition directly as the
high-temperature comparison input for decay, \path{latticeMass}, interval, and
nonzero consequences.  The closed interval forms require only the named
proposition; the positive half-open and nonzero forms also take the active-range
membership needed for strict positivity of the anchored cubic pseudo-mass.
PR~\#1838 adds the corresponding direct-positive consumers:
\begin{Verbatim}[breaklines=true,breakanywhere=true]
latticeMass_pos_of_cubicOriginNamedRateLeHighTemp_pos
cubicNamedRate_ofReal_mem_Ioc_latticeMass_of_cubicOriginNamedRateLeHighTemp_pos
latticeMass_ne_zero_of_cubicOriginNamedRateLeHighTemp_pos
cubicNamedRate_decay_mem_Ioc_of_cubicOriginNamedRateLeHighTemp_pos
\end{Verbatim}
These forms use strict positivity of
\path{cubicOriginPseudoMassFromParamsAtPair} as a direct hypothesis, rather
than deriving it from the anchored cubic active range.  They are useful when a
later argument has positivity from another source but still wants to stay on
the irreducible named comparison proposition for the high-temperature bound.
PR~\#1837 extends the same named-proposition input to the cluster and
summability layer:
\begin{Verbatim}[breaklines=true,breakanywhere=true]
clusterProperty_of_cubicOriginNamedRateLeHighTemp_pos
clusterProperty_of_cubicOriginNamedRateLeHighTemp_cubic_corr_mem
summable_truncated2Infinite_prod_of_cubicOriginNamedRateLeHighTemp_pos
summable_truncated2Infinite_prod_of_cubicOriginNamedRateLeHighTemp_cubic_corr_mem
exists_tsum_cubicTruncated2Product_le_of_cubicOriginNamedRateLeHighTemp_pos
exists_tsum_cubicTruncated2Product_le_of_cubicOriginNamedRateLeHighTemp_cubic_corr_mem
\end{Verbatim}
The positive forms accept strict positivity of the named pseudo-mass directly.
The active-range forms derive that positivity from the anchored cubic pair
correlation and use \path{cubicOriginNamedRateLeHighTemp} for the rate
comparison.  This keeps the cluster-property bridge, cubic product summability,
and explicit Step~127 product-sum bound on the same lightweight named
comparison API.
PR~\#1839 adds capstone consequence bundles for the same irreducible named
proposition:
\begin{Verbatim}[breaklines=true,breakanywhere=true]
cubicNamedRate_capstone_bundle_of_cubicOriginNamedRateLeHighTemp_pos
cubicNamedRate_capstone_bundle_of_cubicOriginNamedRateLeHighTemp_cubic_corr_mem
\end{Verbatim}
Each bundle returns the target-exhaustion decay witness, half-open target
interval membership, positive and nonzero target \path{latticeMass}, the target
cluster property, cubic product summability, and the Step~127 product-sum
witness.  The first form takes strict positivity of the named cubic pseudo-mass
directly; the second derives it from anchored cubic active-range membership.
Both forms keep later arguments on \path{cubicOriginNamedRateLeHighTemp} and do
not unfold or restate the heavy concrete comparison.
PR~\#1840 adds the corresponding capstone bundles one layer before the named
proposition:
\begin{Verbatim}[breaklines=true,breakanywhere=true]
cubicNamedRate_capstone_bundle_of_cubic_pseudoMassG_le_corr
cubicNamedRate_capstone_bundle_of_corr_mem_le_high_temp_rate
\end{Verbatim}
The first form consumes anchored cubic active-range membership together with the
cubic profile lower bound
\(\mathrm{pseudoMassG}(\alpha,r,-\log(\beta J\,2d))\le
\mathrm{correlationInfinite}(\{0,z\})\).  The second consumes active-range
membership together with the raw named-rate comparison.  Both forms delegate to
the irreducible named-proposition bundle, so downstream proofs can obtain the
same target decay, interval, \path{latticeMass}, cluster, summability, and
Step~127 witness package without exposing a direct \path{cubicTanhProfileBound}
heavy-comparison wrapper.
PR~\#1841 adds the bundled-input form:
\begin{Verbatim}[breaklines=true,breakanywhere=true]
cubicNamedRate_capstone_bundle_of_cubic_corr_mem_Ioo_and_profile
\end{Verbatim}
The hypothesis is the conjunction produced by
\path{cubicTanhProfileBound_cubic_corr_mem_Ioo_and_profile}: anchored cubic
active-range membership and the cubic profile lower bound.  This gives
downstream proofs a stable way to pass that pair into the capstone bundle
without placing \path{cubicTanhProfileBound} itself in the heavy comparison
statement.

The totalized comparison
\path{pseudoMassExt_twoPointFunction_le_high_temp_rate_of_pseudoMassG_le_tanh_pow_dist}
packages the same tanh-power input without forcing the upper endpoint into the
proof.  It proves
\[
  \mathrm{pseudoMassExt}\bigl(
    \mathrm{twoPointFunction}\,d\,\langle J,0,\beta\rangle\,z\bigr)
    \le -\log(\beta J\,2d).
\]
If the two-point function lies in $(0,2)$, this is the ordinary pseudo-mass
comparison via the implicit characterization of $\mathrm{pseudoMass}$.  If it
does not lie in $(0,2)$, the totalized function is zero, and the comparison
follows from $0\le -\log(\beta J\,2d)$ for Lean's total real logarithm.  The
concrete $\mathrm{pseudoMassFromParamsAtPair}$ wrapper remains a separate
lightweight composition step.

\paragraph{Active two-point pseudo-mass comparison from tanh profile bounds
(PR~\#1817).}
The follow-up
\path{twoPointFunction_mem_Ioo_zero_two_of_pseudoMassG_le_tanh_pow_dist}
records that the same tanh-power profile hypothesis places the anchored
two-point function itself in the active interval $(0,2)$.  The lower endpoint
is obtained by the positivity of $\mathrm{pseudoMassG}$ at the nonnegative
Lean real-log rate and the tanh-distance lower bound, while the upper endpoint
uses the universal estimate
\(\mathrm{twoPointFunction}\le 1 < 2\).
Consequently
\path{pseudoMass_twoPointFunction_le_high_temp_rate_of_pseudoMassG_le_tanh_pow_dist}
removes the totalisation and proves the ordinary pseudo-mass comparison
\[
  \mathrm{pseudoMass}\bigl(
    \mathrm{twoPointFunction}\,d\,\langle J,0,\beta\rangle\,z\bigr)
    \le -\log(\beta J\,2d),
\]
with the required active-interval proof supplied as the implicit argument to
\(\mathrm{pseudoMass}\).  The equality
\path{pseudoMassExt_twoPointFunction_eq_pseudoMass_of_pseudoMassG_le_tanh_pow_dist}
then states that the totalised comparison of PR~\#1816 agrees with the ordinary
pseudo-mass comparison under this profile hypothesis.  This stays on the
anchored two-point API and deliberately does not claim the concrete
\(\mathrm{pseudoMassFromParamsAtPair}\) wrapper.

\paragraph{Positive active two-point pseudo-mass consequences (PR~\#1818).}
The positivity follow-up keeps the same anchored two-point scope.  The theorem
\path{pseudoMass_twoPointFunction_pos_of_pseudoMassG_le_tanh_pow_dist}
applies \(\mathrm{pseudoMass\_pos}\) to the active-interval proof from
PR~\#1817.  The totalized form
\path{pseudoMassExt_twoPointFunction_pos_of_pseudoMassG_le_tanh_pow_dist}
uses \(\mathrm{pseudoMassExt\_pos\_of\_mem}\) with the same active proof, and
\path{pseudoMassExt_twoPointFunction_ne_zero_of_pseudoMassG_le_tanh_pow_dist}
records the corresponding non-vanishing corollary.  These are bookkeeping
consequences for later Step~117l compositions and still do not assert the
concrete \(\mathrm{pseudoMassFromParamsAtPair}\) comparison.

\paragraph{\S5.1 cluster property bundled formalisation (PR~\#792).}
The Glimm--Jaffe \S5.1 cluster property is formalised as a single
predicate
\[
  \mathrm{clusterProperty}\,G\,\Lambda\,p \;:=\;
    \forall\, i \in V,\
    \mathrm{Tendsto}\bigl(j \mapsto U_2(i, j)\bigr)\,
    \mathrm{Filter.cofinite}\,(\mathrm{nhds}\,0)
\]
($\mathrm{IsingModel.Ambient.clusterProperty}$, in
\texttt{AmbientLattice.lean}). The PR bundles the
sufficient condition $\mathrm{clusterProperty\_of\_summable}$
(per-site application of
$\mathrm{truncated2Infinite\_tendsto\_cofinite\_zero\_of\_summable}$),
together with the two trivial-slice proofs
$\mathrm{clusterProperty\_J\_zero}$ (for ferromagnetic
$\langle 0, h, \beta \rangle$, using $U_2 = 0$ off-diagonally and
$\mathrm{Filter.eventuallyEq}$ on $\{i\}^{\mathrm c}$) and
$\mathrm{clusterProperty\_beta\_zero}$ (for any
$\langle J, h, 0 \rangle$, using $U_2 \equiv 0$ at $\beta = 0$).
$\mathbb Z^d$ wrappers
$\mathrm{clusterProperty\_latticeGraph\_\{of\_summable, J\_zero, beta\_zero\}}$
specialise these to the integer lattice graph. The general
(non-trivial) case requires the Simon--Lieb inequality (FV Prop~9.31)
or random-current representation; both are research-level and
remain open (Issue~\#780).

\paragraph{Outer-side companion of PR~\#787's edge bound (PR~\#791).}
The symmetric outer-side bound
\[
  |\partial^e S| \;\le\; \sum_{y \in \partial^v_o S} \deg_G(y)
\]
($\mathrm{edgeBoundary\_card\_le\_sum\_degrees\_outerVertexBoundary}$)
follows from PR~\#789's outer-side equality plus the
$|\mathrm{filter}| \le |\mathrm{neighborFinset}| = \deg$
estimate, mirroring PR~\#787 with $\partial^v_o$ in place of
$\partial^v_i$. On $\mathbb Z^d$ this gives
$|\partial^e S| \le 2 d \cdot |\partial^v_o S|$, and combining
with PR~\#787's inner-side wrapper yields the strictly stronger
combined bound
$|\partial^e S| \le 2 d \cdot \min(|\partial^v_i S|, |\partial^v_o S|)$.

\paragraph{Vertex-by-edge bounds (PR~\#790).}
Two simple-but-useful bounds tying the vertex boundaries to the
edge boundary:
\[
  |\partial^v_i S| \;\le\; |\partial^e S|, \quad
  |\partial^v_o S| \;\le\; |\partial^e S|
\]
($\mathrm{innerVertexBoundary\_card\_le\_edgeBoundary\_card}$,
$\mathrm{outerVertexBoundary\_card\_le\_edgeBoundary\_card}$).
Each follow from the equalities of PRs~\#788 / \#789 by observing
that the very witness making a vertex a boundary vertex is an
element of the corresponding filter, so each summand is $\ge 1$;
sum monotonicity finishes. Together with the existing
edge-boundary bounds these give the linear chain
\[
  |\partial^v_i S|,\, |\partial^v_o S|
  \;\le\; |\partial^e S|
  \;\le\; \deg_{\max} \cdot |\partial^v_i S|,
\]
which on $\mathrm{latticeGraph}(d)$ becomes the chain with
$\deg_{\max} = 2d$.

\subsection{Lattice-growth convergence of \S5 quantities}

Named corollaries of \texttt{correlation\_convergent\_subgraph} (PR~\#64):

\begin{theorem}[\texttt{truncated2\_convergent\_subgraph}]
$\langle\sigma_i;\sigma_j\rangle_{G_n}$ converges for any increasing subgraph
sequence $G_n$ with ferromagnetic parameters.
\end{theorem}

\begin{proof}
Each of $\langle\sigma_i\sigma_j\rangle_{G_n}$, $\langle\sigma_i\rangle_{G_n}$,
$\langle\sigma_j\rangle_{G_n}$ converges; apply \texttt{Tendsto.sub} and
\texttt{Tendsto.mul}.
\end{proof}

\begin{theorem}[\texttt{susceptibility\_convergent\_subgraph}]
$\chi_i(G_n) = \sum_j \langle\sigma_i;\sigma_j\rangle_{G_n}$ converges.
\end{theorem}

\begin{theorem}[\texttt{magnetization\_total\_convergent\_subgraph}]
$\sum_i \langle\sigma_i\rangle_{G_n}$ converges.
\end{theorem}

Both follow from \texttt{tendsto\_finset\_sum} applied to the componentwise
convergence.

\section{Conditioning inequalities (\texttt{Conditioning.lean})}

Lattice analogues of Glimm--Jaffe, \S10.1--10.2 (pp.~193--194).

\begin{theorem}[\texttt{partitionFunction\_beta\_rescale}]
$Z(J, h, \beta) = Z(\beta J, \beta h, 1)$: the partition function
depends on $(J, h, \beta)$ only through $(\beta J, \beta h)$.
\end{theorem}

\begin{theorem}[\texttt{partitionFunction\_monotone\_beta}; Cor.~10.2.3]
For $J \ge 0$, $h \ge 0$, and $0 < \beta_1 \le \beta_2$:
$Z(J, h, \beta_1) \le Z(J, h, \beta_2)$.
\end{theorem}

\begin{proof}
By the rescaling identity, $Z(J,h,\beta) = Z(\beta J, \beta h, 1)$.
Since $\beta J$ and $\beta h$ are monotone in $\beta$ for $J, h \ge 0$,
the result follows from \texttt{partitionFunction\_monotone\_J} and
\texttt{partitionFunction\_monotone\_h}.
\end{proof}

\subsection{Partition function bounds (Cor.~10.3.2)}

\begin{theorem}[\texttt{hamiltonian\_abs\_le}]
$|H(\sigma)| \le |J| \cdot |E| + |h| \cdot |\iota|$, since
$|\text{edgeSpin}| \le 1$ and $|\text{sign}(\sigma_i)| \le 1$.
\end{theorem}

\begin{theorem}[\texttt{partitionFunction\_upper}]
$Z \le 2^{|\iota|} \cdot \exp(|\beta|(|J| \cdot |E| + |h| \cdot |\iota|))$.
\end{theorem}

\begin{theorem}[\texttt{partitionFunction\_lower}]
$\exp(-|\beta|(|J| \cdot |E| + |h| \cdot |\iota|)) \le Z$.
\end{theorem}

\begin{theorem}[\texttt{freeEnergy\_upper\_bound}]
For $|\iota| > 0$,
\[
\texttt{freeEnergy}\,G\,p \;\le\; \log 2
  + \frac{|\beta|(|J| \cdot |E| + |h| \cdot |\iota|)}{|\iota|}.
\]
Take $\log$ of \texttt{partitionFunction\_upper} and use
$|\text{Config}| = 2^{|\iota|}$, then divide by $|\iota|$.
\end{theorem}

\begin{theorem}[\texttt{freeEnergyAlongExhaustion\_upper\_bound}]
For nonempty $\Lambda.\text{volume}\,n$,
\[
\texttt{freeEnergyAlongExhaustion}\,G\,\Lambda\,p\,n \;\le\; \log 2
  + \frac{|\beta|(|J| \cdot |E_n| + |h| \cdot |\Lambda_n|)}{|\Lambda_n|},
\]
specialization of \texttt{freeEnergy\_upper\_bound} at each stage.
\end{theorem}

\subsection{Free-spin / one-body limit \texorpdfstring{$(G = \bot)$}{(G = bottom)}}

Reference: Glimm--Jaffe, opening of \S 4 (free spin).
Regardless of $J$, if the ambient graph is $\bot$ (no edges), the
interaction term disappears.

\begin{theorem}[\texttt{hamiltonian\_bot}]
$H_{\bot, p}(\sigma) = -h \sum_{i \in \iota} \operatorname{sign}(\sigma_i)$.
\texttt{edgeFinset\_bot}: $(\bot).\text{edgeFinset} = \emptyset$, so the
$J$-term is an empty sum and is zero.
\end{theorem}

\begin{theorem}[\texttt{sum\_exp\_spin\_sign}]
$\sum_{s \in \{\pm 1\}} \exp(\beta h \cdot s) = 2 \cosh(\beta h)$.
\end{theorem}

\begin{theorem}[\texttt{partitionFunction\_bot}]
$Z_{\bot, p} = (2 \cosh(\beta h))^{|\iota|}$.
\end{theorem}

\begin{proof}
\texttt{hamiltonian\_bot} + $\exp(-\beta H) = \prod_i \exp(\beta h \cdot \operatorname{sign}(\sigma_i))$.
The finite product-sum rule \texttt{Finset.sum\_prod\_piFinset} gives
$\sum_\sigma \prod_i f(\sigma_i) = \prod_i \sum_s f(s)$.
The one-site sum is closed by \texttt{sum\_exp\_spin\_sign}.
\end{proof}

\begin{theorem}[\texttt{freeEnergy\_bot}]
For $|\iota| > 0$,
\[
  f(\bot, p) = \log(2 \cosh(\beta h)).
\]
Apply \texttt{Real.log\_pow} to \texttt{partitionFunction\_bot} and divide
by $|\iota|$.
\end{theorem}

\begin{theorem}[\texttt{freeEnergy\_bot\_h\_zero}]
For $|\iota| > 0$ and arbitrary $J, \beta \in \mathbb R$,
\[
  f(\bot, \langle J, 0, \beta \rangle) = \log 2.
\]
This follows immediately from \texttt{freeEnergy\_bot} and $\cosh 0 = 1$.
The parameter $J$ is arbitrary because the $J$-term is empty on $\bot$.
It extends PR \#120 \texttt{freeEnergy\_zero\_params} ($J = h = 0$) to
arbitrary $J$.
\end{theorem}

\begin{theorem}[\texttt{partitionFunction\_beta\_zero}]
$Z_G(\langle J, h, 0 \rangle) = |\text{Config}\,\iota|$.
When $-\beta = 0$, the Boltzmann weight is $\exp 0 = 1$.
\end{theorem}

\begin{theorem}[\texttt{freeEnergy\_beta\_zero}]
For $|\iota| > 0$ and arbitrary $J, h, G$,
\[
  f(G, \langle J, h, 0 \rangle) = \log 2.
\]
This is the $\beta$-direction analogue of PR \#120
\texttt{freeEnergy\_zero\_params}.
\texttt{partitionFunction\_beta\_zero} + \texttt{Real.log\_pow} + $|\iota|^{-1}$ cancel.
\end{theorem}

\begin{theorem}[\texttt{freeEnergyAlongExhaustion\_beta\_zero}; PR \#132]
For nonempty $\Lambda.\text{volume}\,n$,
$\texttt{freeEnergyAlongExhaustion}\,G\,\Lambda\,\langle J, h, 0 \rangle\,n = \log 2$.
Specialize by \texttt{change} and definitional unfolding.
\end{theorem}

\begin{theorem}[\texttt{freeEnergyInfinite\_beta\_zero}; PR \#133]
Assume $\Lambda.\text{volume}\,n$ is nonempty for every $n$:
$\texttt{freeEnergyInfinite}\,G\,\Lambda\,\langle J, h, 0 \rangle = \log 2$.
Rewrite to a constant sequence and close with \texttt{Filter.limsup\_const}.
\end{theorem}

\begin{theorem}[\texttt{freeEnergyAlongExhaustion\_zero\_params}; PR \#162]
For nonempty $\Lambda.\text{volume}\,n$,
$\texttt{freeEnergyAlongExhaustion}\,G\,\Lambda\,\langle 0, 0, \beta \rangle\,n = \log 2$.
Use \texttt{change} and definitional unfolding to specialize
\texttt{IsingModel.freeEnergy\_zero\_params}.
\end{theorem}

\begin{theorem}[\texttt{freeEnergyInfinite\_zero\_params}; PR \#162]
Assume $\Lambda.\text{volume}\,n$ is nonempty for every $n$:
$\texttt{freeEnergyInfinite}\,G\,\Lambda\,\langle 0, 0, \beta \rangle = \log 2$.
Reuse the structure of \texttt{freeEnergyInfinite\_beta\_zero} through
\texttt{freeEnergyAlongExhaustion\_zero\_params}. This gives the maximal
entropy value from the $J = h = 0$ pair, where the Hamiltonian is
identically zero.
\end{theorem}

\begin{theorem}[\texttt{partitionFunctionAlongExhaustion\_zero\_params}; PR \#163]
For every $n$,
\[
  \texttt{partitionFunctionAlongExhaustion}\,G\,\Lambda\,\langle 0, 0, \beta \rangle\,n
    = 2^{|\Lambda_n|}.
\]
\texttt{IsingModel.partitionFunction\_zero\_params} is specialized directly,
then combined with
\texttt{card\_config\_eq\_two\_pow} + \texttt{Fintype.card\_coe}.
\end{theorem}

\begin{theorem}[\texttt{log\_partitionFunctionAlongExhaustion\_zero\_params}; PR \#163]
For every $n$,
\[
  \log \texttt{partitionFunctionAlongExhaustion}\,G\,\Lambda\,\langle 0, 0, \beta \rangle\,n
    = |\Lambda_n| \cdot \log 2.
\]
This is the composition of the preceding theorem with \texttt{Real.log\_pow}.
\end{theorem}

\begin{theorem}[\texttt{partitionFunctionAlongExhaustion\_beta\_zero}; PR \#164]
For every $n$ and arbitrary $J, h \in \mathbb R$,
\[
  \texttt{partitionFunctionAlongExhaustion}\,G\,\Lambda\,\langle J, h, 0 \rangle\,n
    = 2^{|\Lambda_n|}.
\]
Specialize \texttt{IsingModel.partitionFunction\_beta\_zero}. At
$\beta = 0$, the Boltzmann weight satisfies $\exp(-\beta H) = 1$.
\end{theorem}

\begin{theorem}[\texttt{log\_partitionFunctionAlongExhaustion\_beta\_zero}; PR \#164]
Under the same assumptions,
\[
  \log \texttt{partitionFunctionAlongExhaustion}\,G\,\Lambda\,\langle J, h, 0 \rangle\,n
    = |\Lambda_n| \cdot \log 2.
\]
\end{theorem}

\begin{theorem}[\texttt{freeEnergy\_ge\_log\_two\_of\_ferromagnetic}; PR \#168]
For $0 < |\iota|$ and ferromagnetic $p$,
\[
  \log 2 \;\le\; f_G(p).
\]
\texttt{freeEnergy\_ge\_log\_two\_cosh}, \texttt{Real.one\_le\_cosh}, and
monotonicity of logarithm. The existing $\Lambda$ version
\texttt{freeEnergy\LeanLambda\_ge\_log\_two} (PR \#125) is refactored here into a
thin wrapper around the base identity.
\end{theorem}

\begin{theorem}[\texttt{freeEnergyInfinite\_eq\_of\_tendsto}; PR \#170]
Assuming
$\texttt{Tendsto}\,(\texttt{freeEnergyAlongExhaustion}\,G\,\Lambda\,p)\,
\texttt{atTop}\,(\mathcal N (L))$,
\[
  \texttt{freeEnergyInfinite}\,G\,\Lambda\,p \;=\; L.
\]
Combine the definition \texttt{freeEnergyInfinite := Filter.limsup} with
\texttt{Filter.Tendsto.limsup\_eq}. This infrastructure identifies the value
of \texttt{freeEnergyInfinite} once Fekete convergence is available.
\end{theorem}

\begin{theorem}[Glimm--Jaffe \S 4.6 Prop 4.6.1 (Fekete convergence);
PR \#192]
Under the following bundled assumptions, expressing a disjoint-tower
exhaustion condition, the free-energy density converges:
\begin{itemize}
  \item $\forall m, n, |\Lambda_{m+n}| = |\Lambda_m| + |\Lambda_n|$
        (cardinality additivity)
  \item $\forall m, n, \log Z_{\Lambda_m} + \log Z_{\Lambda_n}
          \le \log Z_{\Lambda_{m+n}}$ (super-additivity)
  \item $\texttt{freeEnergyAlongExhaustion}\,G\,\Lambda\,p$ is bounded above
  \item $|\Lambda_1| \neq 0$
\end{itemize}
Under these assumptions,
\[
  \texttt{Tendsto}\,(\texttt{freeEnergyAlongExhaustion}\,G\,\Lambda\,p)\,
    \texttt{atTop}\,
    (\mathcal N (\texttt{freeEnergyInfinite}\,G\,\Lambda\,p)).
\]
Proof: $u_n := -\log Z_{\Lambda_n}$ is \texttt{Subadditive} by sign reversal
of super-additivity. Cardinality additivity gives
$|\Lambda_n| = n \cdot |\Lambda_1|$ by induction. Apply mathlib
\texttt{Subadditive.tendsto\_lim} (Fekete) to obtain $u_n / n \to \ell$.
The ratio is
\[
  \texttt{freeEnergyAlongExhaustion}\,G\,\Lambda\,p\,n
    = -(u_n / n) / |\Lambda_1|
\]
for $n \ge 1$, so the \texttt{Tendsto} result transfers.
\texttt{freeEnergyInfinite\_eq\_of\_tendsto} identifies the limit value.

This completes the Fekete stage of Proposition 4.6.1; together with the
existing disjoint-union super-additivity, upper bound, and divergence
infrastructure, the targeted exhaustion class is covered.
\end{theorem}

\begin{theorem}[Support toward GJ \S 4.6 Thm 4.6.2: finite-volume complex $Z$ analyticity; PR \#193]
Preliminary support for GJ Theorem 4.6.2: after extending the finite-volume
partition function $Z_G(J, h, \beta)$ to $(J, h, \beta) \in \mathbb C^3$, it
is entire in each parameter separately:
\[
  Z_G^{\mathbb C}(J, h, \beta) = \sum_\sigma \exp(-\beta \cdot H^{\mathbb C}_G(J, h; \sigma)),
\]
$H^{\mathbb C}_G$ is affine in $(J, h)$. Hence each Boltzmann weight is
$\exp \circ$(affine) and entire, and the finite sum $Z$ is entire.
\texttt{partitionFunctionComplex\_analyticAt\_\{h,J,beta\}}.
This PR only covers separate-parameter analyticity of $Z$. The body of GJ
Theorem 4.6.2, namely complex analyticity of free energy on
$|\mathrm{Im}\,h| < \mathrm{Re}\,h$, joint analyticity, $\log Z$ via the
slit plane, and the infinite-volume Vitali step, is left to later PRs. Note
that in mathlib, \texttt{Complex.log} is analytic only on \texttt{slitPlane}.
\end{theorem}

\begin{theorem}[GJ \S 4.6 Thm 4.6.2 continuation: $Z \in$ slitPlane on real slice; PR \#198]
\texttt{partitionFunctionComplex\_mem\_slitPlane\_of\_real}:
For real $p$, $Z^{\mathbb C}_G(p.J, p.h, p.\beta) \in \texttt{Complex.slitPlane}$.
Proof: PR \#196 real-complex compatibility and \texttt{partitionFunction\_pos}
($Z > 0$) give positive real part and zero imaginary part. This is the
\texttt{Or.inl} branch of the slit-plane condition.
\end{theorem}

\begin{theorem}[GJ \S 4.6 Thm 4.6.2 continuation: \texttt{freeEnergy} real-complex compat; PR \#197]
\texttt{freeEnergy\_ofReal\_eq\_freeEnergyComplex}:
For real parameters $p$,
$((\texttt{freeEnergy}\,G\,p : \mathbb R) : \mathbb C) = \texttt{freeEnergyComplex}\,G\,(J:\mathbb C)\,(h:\mathbb C)\,(\beta:\mathbb C)$.
Proof: use the $Z$ compatibility from PR \#196 and
\texttt{partitionFunction\_pos} ($Z > 0$), then apply
\texttt{Complex.ofReal\_log}.
\end{theorem}

\begin{theorem}[GJ \S 4.6 Thm 4.6.2 continuation: real-complex compatibility; PR \#196]
\texttt{partitionFunction\_ofReal\_eq\_partitionFunctionComplex}:
For real parameters $p = \langle J, h, \beta\rangle$, coercing
$(\texttt{partitionFunction}\,G\,p : \mathbb R)$ to $\mathbb C$ gives
$\texttt{partitionFunctionComplex}\,G\,(J:\mathbb C)\,(h:\mathbb C)\,(\beta:\mathbb C)$.
Proof: \texttt{Complex.ofReal} is a ring homomorphism, so it commutes with
finite sums, exponentials, and each site's spin sign.
\end{theorem}

\begin{theorem}[GJ \S 4.6 Thm 4.6.2 continuation: joint analyticity in $(J,h,\beta)$; PR \#195]
\texttt{partitionFunctionComplex\_analyticAt\_joint}:
The map $(J, h, \beta) \in \mathbb C^3 \mapsto Z_G^{\mathbb C}$ is entire.
\texttt{freeEnergyComplex\_analyticAt\_joint}: $Z \in$ \texttt{slitPlane}
at a point where $f_G^{\mathbb C}$ is also jointly analytic.
Proof: by the definitions $H_J = -J \cdot \text{edgeSum}$ and
$H_h = -h \cdot \text{siteSum}$, the exponent
$-\beta(H_J + H_h) = \beta \cdot (J \cdot \text{edgeSum} + h \cdot \text{siteSum})$
is polynomial in $(J, h, \beta)$. Use
\texttt{analyticAt\_fst}/\texttt{analyticAt\_snd} for projections and compose
with \texttt{add}/\texttt{mul}/\texttt{neg}. Then
\texttt{AnalyticAt.cexp'} plus \texttt{Finset.analyticAt\_fun\_sum} gives
entireness of $Z$, and \texttt{AnalyticAt.clog} plus
\texttt{analyticAt\_const.mul} gives analyticity of $f$.
\end{theorem}

\begin{theorem}[GJ \S 4.6 Thm 4.6.2 continuation: \texttt{freeEnergyComplex} analyticity via slitPlane; PR \#194]
Following PR \#193, where $Z$ is entire in each parameter separately, at a
point with $Z(J, h_0, \beta) \in \texttt{Complex.slitPlane}$,
\[
  \texttt{AnalyticAt}\,\mathbb C\,
    (h \mapsto \texttt{freeEnergyComplex}\,G\,J\,h\,\beta)\,h_0
\]
with analogous statements for the other two parameters. Proof: compose
mathlib \texttt{AnalyticAt.clog}
(\texttt{Mathlib.Analysis.SpecialFunctions.Complex.Analytic}) with
multiplication by a constant.
\end{theorem}

\begin{theorem}[Lee-Yang domain infrastructure (support toward GJ \S 4.6 Thm 4.6.2); PR \#199]
Support lemmas toward the eventual completion of GJ \S 4.6 Thm 4.6.2
(pp.~67--70). The following are independent auxiliary results, not
statements of GJ itself:
\begin{itemize}
\item \texttt{leeYangDomain} $= \{h \in \mathbb C : |\mathrm{Im}\,h| < \mathrm{Re}\,h\}$
 is open, is a subset of \texttt{slitPlane}, and contains the positive
 real axis.
\item \texttt{leeYangFugacity}$(\beta, h) := e^{-2\beta h}$ is entire
 on $\beta, h \in \mathbb C$, and on the Lee-Yang domain with
 $\beta > 0$ it satisfies $\|\cdot\| < 1$ (i.e. maps into the open
 unit disk).
\item \texttt{leeYangNormalization}$(\beta, J, h, |E|, |\iota|)
 := \exp(\beta J |E| + \beta h |\iota|)$ is entire, non-vanishing, and
 has positive real part at real parameters.
\end{itemize}
\end{theorem}

\begin{theorem}[Friedli--Velenik factorisation; PR \#199]
\label{thm:friedli-velenik}
For real $\beta, J$ and complex uniform field $h \in \mathbb C$,
with finite graph $G$, $|\iota| = \texttt{Fintype.card}\,\iota$,
$|E| = G.\texttt{edgeFinset.card}$:
\[
  Z(J, h, \beta)
    = \exp(\beta J |E| + \beta h |\iota|) \cdot P_E(z),
\]
where $z_k = e^{-2\beta h}$ (uniform fugacity vector) and $P_E$ is the
Ising partition polynomial \texttt{isingEdgePoly}
(from \texttt{LeeYang.lean}, with
\texttt{graphToEdgeList}\,$G$\,$(e^{-2\beta J})$ as coefficients).
Proof: transport through the bijection
$\sigma \mapsto X := \{i : \sigma_i = \mathrm{down}\}$
(\texttt{configFinsetEquiv}) and, for each $\sigma$, expand
$\exp(-\beta H(\sigma))$ per-site and per-edge and identify
coefficients.
Lean name: \texttt{partitionFunctionComplex\_eq\_normalization\_mul\_isingEdgePoly}.
Reference: Friedli--Velenik (3.63)--(3.65), pp.~122--123.
\end{theorem}

\begin{theorem}[GJ Thm 4.6.2 non-vanishing half on Lee-Yang domain; PR \#199]
For real ferromagnetic $J > 0$, $\beta > 0$, and
$h \in \texttt{leeYangDomain}$,
$\texttt{partitionFunctionComplex}\,G\,(J:\mathbb C)\,h\,(\beta:\mathbb C)
\ne 0$.
The Friedli--Velenik identity (\ref{thm:friedli-velenik}) expresses
$Z$ as a non-vanishing normalisation times $P_E(z)$; the Lee-Yang
circle theorem
(\texttt{lee\_yang\_circle} / \texttt{isingEdgePoly\_nonvanishing\_of\_graph})
gives non-vanishing of the $P_E$ factor.
Lean name: \texttt{partitionFunctionComplex\_ne\_zero\_on\_leeYangDomain}.

Note: non-vanishing does not directly give $Z \in \texttt{slitPlane}$.
Establishing slit-plane membership via continuous branch selection, and
hence completing the analyticity of \texttt{freeEnergyComplex}, is
handled in a later session.
\end{theorem}

\begin{theorem}[Real-slice slitPlane corollary (preliminary to GJ Thm 4.6.2); PR \#199]
For arbitrary real $J, h_0, \beta$,
$\texttt{AnalyticAt}\,\mathbb C\,
  (h \mapsto \texttt{freeEnergyComplex}\,G\,(J:\mathbb C)\,h\,(\beta:\mathbb C))\,
  (h_0:\mathbb C)$.
This is not the statement of GJ Thm 4.6.2 proper, but a simple
composition of
\texttt{partitionFunctionComplex\_mem\_slitPlane\_of\_real}
(at real parameters $Z > 0$ is real, hence in slit-plane) with
\texttt{freeEnergyComplex\_analyticAt\_h}. No assumption $J > 0$,
$\beta > 0$, or Lee-Yang-domain is required.
Lean name: \texttt{freeEnergyComplex\_analyticAt\_h\_ofReal}.
\end{theorem}

\begin{theorem}[GJ Thm 4.6.2 finite-volume local branch on Lee-Yang domain; PR \#200]
For real ferromagnetic $\beta > 0$, $J > 0$, with
\texttt{[Nonempty }$\iota$\texttt{]}, at any
$h_0 \in \texttt{leeYangDomain}$ there exists an analytic function
$f : \mathbb C \to \mathbb C$ such that
$\texttt{AnalyticAt}\,\mathbb C\,f\,h_0$,
$\exp(|\iota| f(h_0)) = Z(h_0)$, and
$f(h_0) = \texttt{freeEnergyComplex}\,G\,(J:\mathbb C)\,h_0\,(\beta:\mathbb C)$.
Proof chain:
\begin{enumerate}
\item $Z'/Z$ analytic on Lee-Yang (\texttt{logDeriv\_...\_analyticOnNhd\_leeYangDomain}).
\item Morera gives a primitive on each ball contained in the domain
  (\texttt{DifferentiableOn.isExactOn\_ball}, mathlib), producing $g$
  with $g' = Z'/Z$ (\texttt{exists\_logZ\_branch\_on\_ball\_of\_leeYangDomain}).
\item Normalise at the basepoint: $g(h_0) = \log Z(h_0)$
  (\texttt{\_normalised\_}).
\item Show $F(z) := \exp(g(z))/Z(z)$ has zero derivative on the ball
  (chain + quotient rules); the ball is connected, hence $F$ is
  constant. At the centre, $F(h_0) = \exp(\log Z(h_0))/Z(h_0) = 1$,
  so $\exp g = Z$ pointwise on the ball.
\item $g$ is \texttt{DifferentiableOn} an open set $\Rightarrow$
  \texttt{AnalyticOnNhd}.
\item Take a ball around each $h_0$ to obtain a pointwise statement.
\item Set $f = g / |\iota|$ for the local branch of
  \texttt{freeEnergyComplex}.
\end{enumerate}
Lean names: \texttt{exists\_freeEnergyComplex\_analyticAt\_branch\_of\_leeYangDomain},
\texttt{analyticBranch\_freeEnergyComplex\_leeYangDomain}.

Note: the principal branch \texttt{Complex.log}($Z$) is discontinuous
where $Z$ crosses the negative real axis; the local branch $f$ is
continuous across such crossings. The statement
``log $Z$ is analytic'' intended by GJ Thm 4.6.2 is interpreted in the
local-branch sense.
\end{theorem}

\begin{theorem}[$\infty$-vol Vitali lift preparation --- ingredients; PR \#200]
\begin{itemize}
\item \texttt{vitali\_bridge}: a direct wrapper of mathlib's
  \texttt{TendstoLocallyUniformlyOn.differentiableOn}. A locally
  uniform limit of holomorphic functions on an open set is holomorphic.
\item \texttt{norm\_partitionFunctionComplex\_le\_partitionFunction}:
  $|Z(J, h, \beta)| \le Z(J, \mathrm{Re}\,h, \beta)$ (real Ising
  partition function).
\item \texttt{norm\_partitionFunctionComplex\_le\_trivial\_bound}:
  combining with the existing \texttt{partitionFunction\_upper},
  $|Z| \le 2^{|\iota|} \exp(|\beta|(|J||E| + |\mathrm{Re}\,h| \cdot |\iota|))$.
\item \texttt{norm\_complex\_log\_le}:
  $\|{\rm log}\,z\| \le |\log\|z\|| + \pi$.
\item \texttt{norm\_freeEnergyComplex\_le\_trivial\_bound}:
  $\|f_{\rm complex}\| \le (|\log\|Z\|| + \pi)/|\iota|$.
\end{itemize}
These are the ingredients for a Montel-style argument (locally uniform
boundedness + real-axis Fekete convergence) yielding locally uniform
convergence. Next step: formalise locally uniform convergence and
apply \texttt{vitali\_bridge\_leeYangDomain} to obtain analyticity of
the infinite-volume limit $f_\infty$.
\end{theorem}

\begin{theorem}[Conditional $\infty$-volume Vitali assembly for GJ \S 4.6 Thm 4.6.2; PR \#2673]
Let $G$ be an ambient graph with exhaustion $\Lambda_n$, and let
\[
  F_n(h)=\texttt{freeEnergyComplexAlongExhaustion}\,G\,\Lambda\,J\,h\,\beta\,n .
\]
The Lean theorem
\texttt{freeEnergyComplexAlongExhaustion\_vitali\_bridge} states that,
on any open set $U\subseteq\mathbb C$, if every $F_n$ is holomorphic on
$U$ and $F_n\to f$ locally uniformly on $U$, then $f$ is holomorphic on
$U$. The specialization
\texttt{freeEnergyComplexAlongExhaustion\_vitali\_bridge\_leeYangDomain}
sets $U=\texttt{leeYangDomain}$.

The companion theorem
\texttt{freeEnergyComplexAlongExhaustion\_limit\_eq\_freeEnergyInfinite\_at\_real}
adds the real-axis identification. If
\[
  F_n(h)\to f(h)
\]
locally uniformly on the Lee-Yang domain, and the real parameters satisfy
the existing \texttt{DisjointTowerHypotheses} and
\texttt{BoundedEdgeDensity} assumptions, then at a real point
$p.h\in\texttt{leeYangDomain}$,
\[
  f(p.h)=
  \bigl(\texttt{freeEnergyInfinite}\,G\,\Lambda\,p:\mathbb R\bigr):\mathbb C .
\]
The proof uses mathlib's pointwise consequence
\texttt{TendstoLocallyUniformlyOn.tendsto\_at}, the previously
formalised real-axis convergence of
\texttt{freeEnergyComplexAlongExhaustion}, and uniqueness of limits in
the Hausdorff space $\mathbb C$.
\end{theorem}

\begin{theorem}[Per-stage Lee-Yang local branches for the Vitali input; PR \#2674]
Let $G$ be an ambient graph with exhaustion $\Lambda_n$, and assume
that the induced edge set at each finite stage is finite.  For real
ferromagnetic parameters $\beta>0$ and $J>0$, the Lean theorem
\texttt{freeEnergyComplexAlongExhaustion\_exists\_analyticAt\_branch\_leeYangDomain\_stage}
states that every nonempty stage $\Lambda_n$ inherits the finite-volume
Lee-Yang local branch theorem from PR \#200.  More explicitly, for each
$h_0\in\texttt{leeYangDomain}$ there is a function
$f:\mathbb C\to\mathbb C$ such that
\[
  f \text{ is analytic at } h_0,\qquad
  \exp(|\Lambda_n|\,f(h_0)) = Z_{\Lambda_n}(h_0),
\]
and
\[
  f(h_0)=
  \texttt{freeEnergyComplexAlongExhaustion}\,G\,\Lambda\,J\,h_0\,\beta\,n .
\]
The wrapper
\texttt{freeEnergyComplexAlongExhaustion\_analyticBranch\_leeYangDomain\_stage}
packages this in pointwise
$\forall h_0\in\texttt{leeYangDomain}$ form, and
\texttt{freeEnergyComplexAlongExhaustion\_analyticBranch\_leeYangDomain\_all\_stages}
packages the same statement for all stages under
$\forall n,\ \Lambda_n\ne\emptyset$.

The stronger ball-level wrapper
\texttt{freeEnergyComplexAlongExhaustion\_exists\_analyticOnNhd\_branch\_ball\_stage}
is the usable local-branch input for Vitali.  If
$B(h_0,r)\subseteq\texttt{leeYangDomain}$ and $r>0$, it gives
$f:\mathbb C\to\mathbb C$ such that $f$ is analytic on a neighbourhood
of the ball and
\[
  \forall z\in B(h_0,r),\qquad
  \exp(|\Lambda_n|\,f(z))=Z_{\Lambda_n}(z).
\]
The all-stages form
\texttt{freeEnergyComplexAlongExhaustion\_analyticOnNhd\_branch\_ball\_all\_stages}
packages this for every nonempty stage.  The concrete $\mathbb Z^d$
specializations are recorded in
\texttt{Concrete/LatticeGraphCorrelation/PerStageComplex.lean}.

These are local branch inputs for the normal-family/Vitali part of
GJ \S 4.6 Thm 4.6.2; they do not by themselves choose a coherent
global branch family or prove local uniform convergence of the chosen
branches.  The pointwise branch agrees with the principal
\texttt{freeEnergyComplexAlongExhaustion} value at the basepoint $h_0$,
while the local analytic representative may differ from the principal
\texttt{Complex.log} branch away from $h_0$.
Reference: Glimm--Jaffe \cite{glimm-jaffe}, \S 4.6, pp.~67--70.
\end{theorem}

\begin{theorem}[Normalised Lee-Yang branch-ball witnesses for the Vitali input; PR \#2675]
The finite-volume theorem
\texttt{exists\_freeEnergyComplex\_analyticOnNhd\_branch\_ball\_strong}
strengthens the ball-level branch statement by keeping the centre
normalisation in the same witness.  For real ferromagnetic parameters
$\beta>0$, $J>0$, and nonempty finite volume, if
$B(h_0,r)\subseteq\texttt{leeYangDomain}$ and $r>0$, then there is
$f:\mathbb C\to\mathbb C$ such that
\[
  f \text{ is analytic on a neighbourhood of } B(h_0,r),
\]
\[
  \forall z\in B(h_0,r),\qquad
  \exp(|\iota|\,f(z))=Z(z),
\]
and
\[
  f(h_0)=\texttt{freeEnergyComplex}\,G\,J\,h_0\,\beta .
\]
The proof applies the Morera-built log branch
\texttt{exists\_logZ\_analytic\_branch\_on\_ball} and sets
$f(z)=g(z)/|\iota|$; the nonempty-volume hypothesis gives
$|\iota|\ne 0$, while the basepoint equality follows from the stored
normalisation $g(h_0)=\log Z(h_0)$.

The along-exhaustion wrappers
\texttt{freeEnergyComplexAlongExhaustion\_exists\_analyticOnNhd\_branch\_ball\_stage\_strong}
and
\texttt{freeEnergyComplexAlongExhaustion\_analyticOnNhd\_branch\_ball\_all\_stages\_strong}
apply this theorem to the induced graph on $\Lambda_n$ for the same real
ferromagnetic parameter range.  Thus, at every nonempty stage, the same
local witness carries analyticity on the ball,
the ball-wide exponential identity for $Z_{\Lambda_n}$, and agreement
with the principal \texttt{freeEnergyComplexAlongExhaustion} value at
the centre.  This still does not prove coherent branch selection or
locally uniform convergence; it fixes the normalised local input that a
future normal-family/Vitali argument can consume.

Reference: Glimm--Jaffe \cite{glimm-jaffe}, \S 4.6, pp.~67--70.
\end{theorem}

\begin{theorem}[Local branch-family Vitali assembly on Lee--Yang balls; PR \#2676]
The theorem
\texttt{freeEnergyComplexAlongExhaustion\_branchFamily\_vitali\_bridge\_ball}
packages the local Vitali handoff for a chosen branch family.  Let
$F_n:\mathbb C\to\mathbb C$ be functions on a ball $B(h_0,r)$ such that each
$F_n$ is analytic on a neighbourhood of the ball, satisfies the ball-wide
branch identity
\[
  \forall z\in B(h_0,r),\qquad
  \exp(|\Lambda_n|\,F_n(z))=Z_{\Lambda_n}(z),
\]
and is normalised at the centre by
\[
  F_n(h_0)=
  \texttt{freeEnergyComplexAlongExhaustion}\,G\,\Lambda\,J\,h_0\,\beta\,n .
\]
If $F_n$ converges locally uniformly on $B(h_0,r)$ to $f$, then $f$ is
holomorphic on $B(h_0,r)$.  The proof is the existing
\texttt{vitali\_bridge} applied to the analytic part of the branch witness.

The centre-identification theorem
\texttt{freeEnergyComplexAlongExhaustion\_branchFamily\_vitali\_ball\_identified\_at\_center}
specialises the centre to a real parameter $h_0=p.h$.  Local uniform
convergence gives $F_n(p.h)\to f(p.h)$, while the normalisation above rewrites
this sequence to the principal stage free energies.  The real-axis Fekete
theorem
\texttt{freeEnergyComplexAlongExhaustion\_tendsto\_at\_real\_of\_disjointTowerHypotheses}
then gives convergence to
\[
  \texttt{freeEnergyInfinite}\,G\,\Lambda\,p .
\]
Uniqueness of limits identifies
\[
  f(p.h)=\texttt{freeEnergyInfinite}\,G\,\Lambda\,p .
\]
The concrete $\mathbb Z^d$ wrappers live in
\texttt{Concrete/LatticeGraphCorrelation/PerStageComplex.lean}.

This result does not construct the coherent branch family or prove its local
uniform convergence.  It records the exact local theorem that can consume that
future normal-family/Montel output.
Reference: Glimm--Jaffe \cite{glimm-jaffe}, \S 4.6, pp.~67--70.
\end{theorem}

\begin{theorem}[Local-cover branch-family Vitali handoff on Lee--Yang; PR \#2677]
The theorem
\texttt{freeEnergyComplexAlongExhaustion\_branchFamily\_vitali\_localCover}
packages the passage from local balls to the full Lee--Yang domain.  Suppose
that for every $h_0\in\texttt{leeYangDomain}$ there is a radius $r>0$ with
$B(h_0,r)\subseteq\texttt{leeYangDomain}$ and a branch family $F_n$ on that
ball satisfying the PR \#2676 hypotheses:
\[
  F_n\text{ analytic on }B(h_0,r),\qquad
  \exp(|\Lambda_n|F_n(z))=Z_{\Lambda_n}(z),
\]
\[
  F_n(h_0)=
  \texttt{freeEnergyComplexAlongExhaustion}\,G\,\Lambda\,J\,h_0\,\beta\,n,
\]
and $F_n\to f$ locally uniformly on $B(h_0,r)$, with the same limiting
function $f$ for all centres.  Then $f$ is holomorphic on
\texttt{leeYangDomain}.  The proof applies the ball-level theorem at each
point and converts holomorphicity on the containing ball into
within-domain differentiability.

The real-centre theorem
\texttt{freeEnergyComplexAlongExhaustion\_branchFamily\_vitali\_localCover\_real}
adds the Fekete identification at a real Lee--Yang point $p.h$, assuming
\texttt{BoundedEdgeDensity} for $(G,\Lambda)$,
\texttt{DisjointTowerHypotheses} for $(G,\Lambda,p)$, and
$(p.h:\mathbb C)\in\texttt{leeYangDomain}$:
\[
  f(p.h)=\texttt{freeEnergyInfinite}\,G\,\Lambda\,p .
\]
It selects the local ball supplied by the cover at $p.h$ and reuses the
PR \#2676 centre-identification theorem.  The concrete $\mathbb Z^d$ wrappers
live in \texttt{Concrete/LatticeGraphCorrelation/PerStageComplex.lean}.

This result still does not construct the coherent local cover or prove local
uniform convergence; it records the global conditional handoff that such a
future normal-family/Montel argument must satisfy.
Reference: Glimm--Jaffe \cite{glimm-jaffe}, \S 4.6, pp.~67--70.
\end{theorem}

\begin{theorem}[Compact Lee--Yang fugacity gap; PR \#2684]
The theorems
\texttt{exists\_leeYangFugacity\_norm\_le\_lt\_one\_on\_isCompact}
and
\texttt{exists\_leeYangFugacityVec\_norm\_le\_lt\_one\_on\_isCompact}
record the compactness gap for the Lee--Yang fugacity.  If $\beta>0$,
$K$ is compact, and $K\subseteq\texttt{leeYangDomain}$, then there is
$r<1$ such that
\[
  |\exp(-2\beta h)|\le r \qquad(h\in K).
\]
The vector form gives the same bound for every coordinate of
\texttt{leeYangFugacityVec}.

The proof uses continuity of $h\mapsto |\exp(-2\beta h)|$ and the extreme
value theorem on the compact set $K$.  At a maximiser $h_0$, the pointwise
Lee--Yang bound gives $|\exp(-2\beta h_0)|<1$; if $K$ is empty one takes
$r=0$.

This is preparatory for quantitative Lee--Yang polynomial lower bounds.
It does not prove lower normalised-log control or Montel convergence.
Reference: Glimm--Jaffe \cite{glimm-jaffe}, \S 4.6, pp.~67--70, and
Friedli--Velenik \cite{friedli-velenik}, Theorem 3.43, pp.~122--127.
\end{theorem}

\begin{theorem}[Compact-uniform complex partition-function bound; PR \#2678]
The theorem
\texttt{exists\_abs\_re\_le\_on\_isCompact} records the elementary
compactness input needed for the Montel side of the Lee--Yang argument: for
every compact set $K\subseteq\mathbb C$ there is a real number $R\ge 0$ such
that
\[
  |\mathrm{Re}\,h|\le R\qquad (h\in K).
\]
Combining this with the stage-wise estimate from PR \#630 gives
\[
  \|Z_{\Lambda_n}(J,h,\beta)\|
  \le
  2^{|\Lambda_n|}
  \exp\!\left(|\beta|\left(|J|\,|E_{\Lambda_n}|+R|\Lambda_n|\right)\right)
\]
for all stages $n$ and all $h\in K$.  In Lean this packaged form is
\[
  \texttt{norm\_partitionFunctionComplexAlongExhaustion\_le\_on\_isCompact\_stage}.
\]
The concrete $\mathbb Z^d$ wrapper
\texttt{norm\_partitionFunctionComplexAlongExhaustion\_le\_on\_isCompact\_stage\_latticeGraph}
lives in \texttt{Concrete/LatticeGraphCorrelation/PerStageComplex.lean}.

This is a compact-field envelope for the partition functions themselves: it is
uniform in $h\in K$, but the right-hand side still depends on $|\Lambda_n|$ and
$|E_{\Lambda_n}|$.  It is therefore an input to later normalised logarithmic
or free-energy estimates, not a standalone stage-uniform Montel bound and not
yet a construction of coherent locally uniformly convergent branches.
Reference: Glimm--Jaffe \cite{glimm-jaffe}, \S 4.6, pp.~67--70.
\end{theorem}

\begin{theorem}[Upper normalised log bound from compact envelope; PR \#2681]
The theorem
\texttt{exists\_real\_log\_norm\_partitionFunctionComplexAlongExhaustion\_div\_card\_le\_on\_isCompact}
records the upper half of the normalised-log input.  Assume that $K$ is compact,
all stages are nonempty, the exhaustion has bounded edge density
$|E_{\Lambda_n}|\le c|\Lambda_n|$, and
$Z_{\Lambda_n}(J,h,\beta)\ne 0$ for $h\in K$.  Then there is a constant $C$
such that
\[
  \frac{\log\|Z_{\Lambda_n}(J,h,\beta)\|}{|\Lambda_n|}\le C
  \qquad(n\in\mathbb N,\ h\in K).
\]
The proof first uses the compact real-part bound $|\mathrm{Re}\,h|\le R$ and
the envelope of PR \#2678 to get the per-stage estimate
\[
  \frac{\log\|Z_{\Lambda_n}(J,h,\beta)\|}{|\Lambda_n|}
  \le
  \log 2+
  |\beta|\left(|J|\,\frac{|E_{\Lambda_n}|}{|\Lambda_n|}+R\right).
\]
Bounded edge density replaces the edge ratio by the single constant $c$.
The concrete $\mathbb Z^d$ wrappers live in
\texttt{Concrete/LatticeGraphCorrelation/PerStageComplex.lean}.

This is not yet the absolute normalised-log bound used by PR \#2679 and
PR \#2680: it supplies only an upper bound for
$\log\|Z_{\Lambda_n}\|$.  A lower bound for $\log\|Z_{\Lambda_n}\|$ remains a
separate normal-family input.  Reference: Glimm--Jaffe \cite{glimm-jaffe},
\S 4.6, pp.~67--70.
\end{theorem}

\begin{theorem}[Lee--Yang compact upper normalised log bound; PR \#2682]
The theorem
\texttt{exists\_log\_norm\_partitionFunctionComplexAlongExhaustion\_div\_card\_le\_leeYangDomain}
specialises PR \#2681 to compact subsets of the Lee--Yang domain.  Assume
$\beta>0$, $J>0$, all stages are nonempty, the exhaustion has bounded edge
density, and $K\subseteq\texttt{leeYangDomain}$ is compact.  Then there is a
constant $C$ such that
\[
  \frac{\log\|Z_{\Lambda_n}(J,h,\beta)\|}{|\Lambda_n|}\le C
  \qquad(n\in\mathbb N,\ h\in K).
\]
The only new step beyond PR \#2681 is to discharge the nonvanishing hypothesis:
Lee--Yang gives
\[
  Z_{\Lambda_n}(J,h,\beta)\ne 0
  \qquad(h\in K\subseteq\texttt{leeYangDomain})
\]
at every finite stage.  The concrete $\mathbb Z^d$ wrapper lives in
\texttt{Concrete/LatticeGraphCorrelation/PerStageComplex.lean}.

This is still an upper normalised-log estimate only.  It does not prove the
lower control on $\log\|Z_{\Lambda_n}\|$ needed to turn the result into the
absolute normalised-log hypothesis used by PR \#2679, PR \#2680, and
PR \#2683.
Reference: Glimm--Jaffe \cite{glimm-jaffe}, \S 4.6, pp.~67--70.
\end{theorem}

\begin{theorem}[Lee--Yang lower log to locally bounded handoff; PR \#2683]
The theorem
\texttt{exists\_abs\_log\_norm\_partitionFunctionComplexAlongExhaustion\_div\_card\_le\_lower\_leeYang}
combines the compact Lee--Yang upper normalised-log bound of PR \#2682 with a
remaining lower normalised-log hypothesis.  Assume $\beta>0$, $J>0$, all
stages are nonempty, the exhaustion has bounded edge density, and
$K\subseteq\texttt{leeYangDomain}$ is compact.  If there is a constant $L$
such that
\[
  -L \le
  \frac{\log\|Z_{\Lambda_n}(J,h,\beta)\|}{|\Lambda_n|}
  \qquad(n\in\mathbb N,\ h\in K),
\]
then there is a constant $C$ such that
\[
  \frac{\left|\log\|Z_{\Lambda_n}(J,h,\beta)\|\right|}{|\Lambda_n|}
  \le C
  \qquad(n\in\mathbb N,\ h\in K).
\]
The proof takes the maximum of the lower-control constant and the automatic
Lee--Yang upper-control constant, then applies the elementary two-sided
absolute-value bridge
\texttt{abs\_log\_norm\_partitionFunctionComplexAlongExhaustion\_div\_card\_le\_of\_two\_sided\_on\_set}.

The companion theorem
\texttt{exists\_norm\_freeEnergyComplexAlongExhaustion\_le\_lower\_log\_leeYang}
immediately feeds this absolute normalised-log bound into PR \#2680, giving a
constant $C$ with
\[
  \|f_n(h)\|\le C+\pi
  \qquad(n\in\mathbb N,\ h\in K).
\]
The concrete $\mathbb Z^d$ wrappers are
\texttt{exists\_abs\_log\_norm\_div\_card\_le\_lower\_ly\_latticeGraph}
and
\texttt{exists\_norm\_freeEnergyComplexAlongExhaustion\_le\_lower\_log\_ly\_latticeGraph}.

This is still conditional: the lower normalised-log control is isolated as an
explicit hypothesis rather than proved.  Reference: Glimm--Jaffe
\cite{glimm-jaffe}, \S 4.6, pp.~67--70.
\end{theorem}

\begin{theorem}[Normalised log bound handoff; PR \#2679]
The theorem
\texttt{norm\_freeEnergyComplexAlongExhaustion\_le\_of\_abs\_log\_norm\_bound\_stage}
records the exact elementary step from normalised logarithmic control to a
finite-stage complex free-energy bound.  If the stage $\Lambda_n$ is nonempty
and
\[
  \frac{\left|\log\|Z_{\Lambda_n}(J,h,\beta)\|\right|}{|\Lambda_n|}
  \le C,
\]
then
\[
  \|f_n(h)\|
  \le
  C+\frac{\pi}{|\Lambda_n|},
\]
where $f_n$ is the principal complex free energy
\texttt{freeEnergyComplexAlongExhaustion}.  The proof is a direct
specialisation of
\texttt{norm\_freeEnergyComplex\_le\_trivial\_bound} to the induced graph on
$\Lambda_n$, using the principal-log estimate
\[
  \|\log z\|\le |\log\|z\||+\pi .
\]
The setwise companion
\texttt{norm\_freeEnergyComplexAlongExhaustion\_le\_of\_abs\_log\_norm\_bound\_on\_set}
packages one constant $C$ for all stages and all fields in a set $K$.  The
concrete $\mathbb Z^d$ wrappers live in
\texttt{Concrete/LatticeGraphCorrelation/PerStageComplex.lean}.

This is intentionally a conditional handoff.  It identifies the missing
normal-family input after the compact $Z$ envelope of PR \#2678: one still
needs control of the absolute logarithm, including lower control on
\(\|Z_{\Lambda_n}\|\), not merely an upper envelope for
\(\|Z_{\Lambda_n}\|\).  Reference: Glimm--Jaffe \cite{glimm-jaffe},
\S 4.6, pp.~67--70.
\end{theorem}

\begin{theorem}[Locally bounded family from normalised log control; PR \#2680]
The theorem
\texttt{norm\_freeEnergyComplexAlongExhaustion\_le\_of\_abs\_log\_norm\_bound\_on\_set\_uniform}
turns the PR \#2679 stagewise estimate into a stage-independent local bound.
Assume that all stages are nonempty and that a set $K\subseteq\mathbb C$ has a
common normalised-log bound
\[
  \frac{\left|\log\|Z_{\Lambda_n}(J,h,\beta)\|\right|}{|\Lambda_n|}
  \le C
  \qquad (n\in\mathbb N,\ h\in K).
\]
Then, for every stage and every $h\in K$,
\[
  \|f_n(h)\|\le C+\pi .
\]
The proof composes PR \#2679 with the elementary consequence of nonemptiness
\[
  \frac{\pi}{|\Lambda_n|}\le \pi ,
\]
using \texttt{Fintype.card\_pos} in Lean.  The concrete $\mathbb Z^d$ wrapper
lives in \texttt{Concrete/LatticeGraphCorrelation/PerStageComplex.lean}.

This theorem packages the exact locally bounded family shape that a later
Montel/normal-family argument can consume.  It remains conditional on the
normalised-log estimate; no lower control on $Z_{\Lambda_n}$ is proved here.
Reference: Glimm--Jaffe \cite{glimm-jaffe}, \S 4.6, pp.~67--70.
\end{theorem}

\begin{remark}[Status of PR \#200 — merged 2026-04-19, \texttt{52ea2f1}]
PR \#200 was squash-merged into \texttt{main} on 2026-04-19 as
commit \texttt{52ea2f1}. Items now present in \texttt{main}:
Morera-based pointwise local analytic log-$Z$ branch on the Lee-Yang
domain (headline
\texttt{exists\_freeEnergyComplex\_analyticAt\_branch\_of\_leeYangDomain}),
the subdomain result in \texttt{leeYangSubdomain}, the Vitali bridge
wrapping mathlib's
\texttt{TendstoLocallyUniformlyOn.differentiableOn}, and modulus bounds
$|Z(J,h,\beta)|\le Z(J,\mathrm{Re}\,h,\beta)$, trivial upper bound,
$\|\log z\|\le|\log\|z\||+\pi$, and the derived
$\|f_{\text{complex}}\|\le(|\log\|Z\||+\pi)/|\iota|$.
Still not formalised (future PRs): Montel-style locally uniform
convergence of $f_\Lambda$ on Lee-Yang (mathlib lacks Montel;
self-implementation required). Note that
\texttt{vitali\_bridge\_leeYangDomain} itself \emph{is} already in
\texttt{main}; what remains is the locally uniform convergence
\emph{input} to the bridge, after which the bridge immediately
yields analyticity of the infinite-volume limit $f_\infty$ on
Lee-Yang. Separately, a global branch patched from the local
pieces via simple connectivity is also not yet formalised.
The principal-branch \texttt{Complex.log} of $Z$ may itself be
discontinuous where $Z$ crosses the negative real axis; the
local-branch statement is the mathematically correct form of GJ
Thm 4.6.2.
\end{remark}

\begin{theorem}[Subdomain of Lee-Yang on which \texttt{freeEnergyComplex} is analytic; PR \#200]
For real $\beta > 0$, real $J$, and complex $h$ with
$\beta \cdot |\mathrm{Im}\,h| \cdot |\iota| < \pi / 2$ (a Lee-Yang
subdomain),
\[
  \mathrm{Re}\,(\texttt{partitionFunctionComplex}\,G\,(J:\mathbb C)\,h\,(\beta:\mathbb C)) > 0,
\]
so the value lies in \texttt{slitPlane}, and
$\texttt{freeEnergyComplex}$ is analytic in $h$.
Proof: for each $\sigma$, the real part of $\exp(-\beta H(\sigma))$
equals $\exp(a) \cos(b)$
(with $a$ the real part of the exponent and
$b = \beta \cdot \mathrm{Im}\,h \cdot \sum_i \sigma_i$,
$|b| \le \beta |\mathrm{Im}\,h| |\iota| < \pi/2$), hence $\cos > 0$ and
the term is positive. Summing over $\sigma$ gives positivity of
$\mathrm{Re}\,Z$. Lean name:
\texttt{freeEnergyComplex\_analyticAt\_h\_of\_leeYangSubdomain}.

Note: the subdomain shrinks as $\beta |\iota|$ grows, so it is not
directly usable at the infinite volume limit. Extension to the full
Lee-Yang domain requires branch selection and is a TODO of this PR.
\end{theorem}

\begin{theorem}[\texttt{freeEnergyInfinite\_of\_eventually\_const}; PR \#171]
If $\forall^\infty n, f_n = c$, then
$\texttt{freeEnergyInfinite}\,G\,\Lambda\,p = c$.
\texttt{tendsto\_const\_nhds.congr'} turns eventual constancy into a
\texttt{Tendsto} statement, and PR \#170 turns the limsup into the limit.
This generalizes the existing \texttt{\_beta\_zero} and
\texttt{\_zero\_params} results, which assumed every $n$ was nonempty.
\end{theorem}

\begin{theorem}[\texttt{freeEnergy\_J\_zero} family; PR \#174]
At $J = 0$, the interaction term is zero and $Z_G = Z_\bot$, independently
of the graph. For $|\iota| > 0$ and arbitrary $h, \beta$,
\[
  Z_G(\langle 0, h, \beta\rangle) = (2 \cosh(\beta h))^{|\iota|},
  \quad
  f_G(\langle 0, h, \beta\rangle) = \log(2 \cosh(\beta h)).
\]
The same PR also provides the along-exhaustion and infinite-volume lifts
under eventual nonemptiness. This is the fourth zero-parameter closed form
alongside the existing $\beta = 0$, $J = h = 0$, and $\bot$-graph families.
\end{theorem}

\begin{theorem}[\texttt{\_eq\_bot\_at\_J\_zero} along-exhaustion / $\infty$-vol lift; PR \#176]
The base-layer $\bot$-equivalence identities from PR \#175 are extended to
the along-exhaustion and infinite-volume layers:
\[
  \texttt{freeEnergyAlongExhaustion}\,G\,\Lambda\,\langle 0,h,\beta\rangle\,n
    = \texttt{freeEnergyAlongExhaustion}\,\bot\,\Lambda\,\langle 0,h,\beta\rangle\,n,
\]
\[
  \texttt{freeEnergyInfinite}\,G\,\Lambda\,\langle 0,h,\beta\rangle
    = \texttt{freeEnergyInfinite}\,\bot\,\Lambda\,\langle 0,h,\beta\rangle.
\]
The along-exhaustion version uses \texttt{change} and
\texttt{freeEnergy\_eq\_bot\_at\_J\_zero}. The infinite-volume version uses
functoriality of \texttt{Filter.limsup}.
\end{theorem}

\begin{theorem}[\texttt{correlation\_eq\_bot\_at\_J\_zero}; PR \#190]
The $\bot$-equivalence identity chain from PR \#175 is extended to the
correlation layer. For arbitrary $G, A, h, \beta$,
\[
  \texttt{correlation}\,G\,\langle 0, h, \beta\rangle\,A
    = \texttt{correlation}\,\bot\,\langle 0, h, \beta\rangle\,A.
\]
The denominator $Z$ is graph-independent
(\texttt{partitionFunction\_eq\_bot\_at\_J\_zero}),
and the Boltzmann weight in the numerator integral is also graph-independent
(\texttt{hamiltonian\_eq\_bot\_at\_J\_zero}).
Taking the ratio gives graph-independence.
\end{theorem}

\begin{theorem}[\texttt{correlation\_bot\_closed} / \texttt{correlation\_J\_zero}; PR \#191]
This records the closed form for correlations on the $\bot$-graph and lifts
it to arbitrary graphs at $J=0$. For arbitrary $p, A$,
\[
  \texttt{correlation}\,\bot\,p\,A = \tanh(p.\beta \cdot p.h)^{|A|},
\]
Combining with PR \#190 gives, for arbitrary $G, A, h, \beta$,
\[
  \texttt{correlation}\,G\,\langle 0, h, \beta\rangle\,A
    = \tanh(\beta h)^{|A|}.
\]
Proof: on $\bot$, the Boltzmann weight factors site-wise as
$\prod_i \exp(\beta h\,\mathrm{sign}(\sigma_i))$. On the
\texttt{spinProduct} side, \texttt{Finset.prod\_filter} rewrites it as
$\prod_{i \in \iota}\, (\mathrm{if}\ i \in A\ \mathrm{then}\ \mathrm{sign}(\sigma_i)\ \mathrm{else}\ 1)$
Then \texttt{Fintype.sum\_prod\_piFinset} swaps
$\sum_\sigma \prod_i$ to $\prod_i \sum_s$. The site sum is $2\sinh$ when
$i \in A$ (\texttt{sum\_spin\_sign\_exp\_sign}) and $2\cosh$ when
$i \notin A$ (\texttt{sum\_exp\_spin\_sign}). Dividing by
$Z = (2\cosh)^{|\iota|}$ yields
$\prod_i (\mathrm{if}\ i \in A\ \mathrm{then}\ \tanh\ \mathrm{else}\ 1) = \tanh^{|A|}$.
This is the fifth zero-parameter closed form paired with the $Z$/free-energy
results from PR \#175.
\end{theorem}

\begin{theorem}[$\bot$-equivalence core identities at $J=0$; PR \#175]
Refactor the separate ``reduction to $\bot$'' arguments from each layer of
PR \#174 into a single base identity chain:
\begin{align*}
  H_G(\langle 0, h, \beta\rangle, \sigma) &= H_\bot(\langle 0, h, \beta\rangle, \sigma), \\
  Z_G(\langle 0, h, \beta\rangle) &= Z_\bot(\langle 0, h, \beta\rangle), \\
  f_G(\langle 0, h, \beta\rangle) &= f_\bot(\langle 0, h, \beta\rangle).
\end{align*}
(\texttt{hamiltonian\_eq\_bot\_at\_J\_zero} /
\texttt{partitionFunction\_eq\_bot\_at\_J\_zero} /
\texttt{freeEnergy\_eq\_bot\_at\_J\_zero}, with names unified under the
\texttt{\_eq\_bot\_at\_J\_zero} suffix). The main \texttt{\_J\_zero}
statements reduce to these identities composed by \texttt{.trans} with
\texttt{partitionFunction\_bot} and \texttt{freeEnergy\_bot}.
\end{theorem}

\begin{theorem}[\texttt{freeEnergyInfinite\_\{beta\_zero,zero\_params\}\_of\_eventually\_nonempty}; PR \#172]
Under the assumption $\forall^\infty n, \Lambda_n$ is nonempty,
\[
  \texttt{freeEnergyInfinite}\,G\,\Lambda\,\langle J, h, 0\rangle
    \;=\; \texttt{freeEnergyInfinite}\,G\,\Lambda\,\langle 0, 0, \beta\rangle
    \;=\; \log 2.
\]
Feed the finite versions
\texttt{freeEnergyAlongExhaustion\_\{beta\_zero,zero\_params\}}
into PR \#171 \texttt{\_of\_eventually\_const} using
\texttt{filter\_upwards}. Under \texttt{[Nonempty V]}, this holds
automatically by \texttt{Exhaustion.eventually\_volume\_nonempty}.
\end{theorem}

\begin{theorem}[\texttt{freeEnergy\_nonneg\_of\_ferromagnetic}; PR \#169]
For $0 < |\iota|$ and ferromagnetic $p$,
\[
  0 \;\le\; f_G(p).
\]
This follows by transitivity from PR \#168 and \texttt{Real.log\_pos}
($0 < \log 2$). The existing $\Lambda$ version
\texttt{freeEnergy\LeanLambda\_nonneg\_of\_ferromagnetic} (PR \#125) is refactored
here into a thin wrapper around the base result.
\end{theorem}

\begin{theorem}[\texttt{card\_mul\_freeEnergy\_eq\_log\_partitionFunction}; PR \#167]
For $|\iota| > 0$,
\[
  |\iota| \cdot f_G(p) = \log Z_G(p).
\]
Unfold the definition $f_G(p) = |\iota|^{-1} \cdot \log Z_G(p)$
and use \texttt{field\_simp} to cancel $|\iota|$. This is the base layer for
the existing $\Lambda$ version
\texttt{card\_mul\_freeEnergy\LeanLambda\_eq\_log\_partitionFunction\LeanLambda\_of\_nonempty}
(PR \#119); that $\Lambda$ version is refactored here into a thin wrapper.
\end{theorem}

\begin{theorem}[\texttt{log\_partitionFunction\_ge\_card\_mul\_log\_two\_cosh\_of\_ferromagnetic} family; PR \#173]
For ferromagnetic $p$,
\[
  |\iota| \cdot \log(2 \cosh(\beta h)) \;\le\; \log Z_G(p).
\]
This sharpens PR \#165 \texttt{\_ge\_card\_mul\_log\_two}. It provides
three-layer lifts: base, $\Lambda$, and along-exhaustion. Proof:
$Z_G \ge Z_\bot = (2 \cosh(\beta h))^{|\iota|}$, followed by $\log$ and
\texttt{Real.log\_pow}.
\end{theorem}

\begin{theorem}[\texttt{magnetization\_beta\_zero}; PR \#187]
Finite-volume magnetization vanishes at $\beta = 0$:
$\texttt{magnetization}\,G\,\langle J, h, 0\rangle\,i = 0$.
Specialize PR \#182 to a singleton. This is the finite-volume companion to
PR \#186 (\texttt{magnetizationInfinite\_beta\_zero}).
\end{theorem}

\begin{theorem}[\texttt{magnetizationInfinite\_beta\_zero}; PR \#186]
Infinite-volume magnetization vanishes at $\beta = 0$:
$\texttt{magnetizationInfinite}\,G\,\Lambda\,\langle J, h, 0\rangle\,i = 0$.
Specialize PR \#183 \texttt{correlationInfinite\_beta\_zero\_vanish} to the
singleton $A = \{i\}$. In addition to the existing
\texttt{magnetizationInfinite\_zero\_at\_h\_zero} from $h=0$ symmetry, this
formalizes vanishing at $\beta = 0$ (infinite temperature).
\end{theorem}

\begin{theorem}[\texttt{correlation\_empty}; PR \#185]
Normalization of the Gibbs measure:
$\texttt{correlation}\,G\,p\,\emptyset = 1$
holds in all four layers: base, $\Lambda$, along-exhaustion, and
infinite-volume. \texttt{spinProduct\_empty} gives $\sigma^\emptyset = 1$;
unfolding the correlation definition gives
$Z^{-1} \sum_\sigma 1 \cdot w = Z^{-1} \cdot Z = 1$.
\end{theorem}

\begin{theorem}[\texttt{correlation\LeanLambda / correlationAlongExhaustion / correlationInfinite\_beta\_zero\_vanish}; PR \#183]
PR \#182 is lifted through three layers:
\begin{itemize}
\item $\Lambda$ layer: for $A : \texttt{Finset}\,\uparrow\Lambda$ with
  $A.\texttt{Nonempty}$, $\texttt{correlation\LeanLambda} = 0$.
\item Along-exhaustion: both branches, $A \subseteq \Lambda_n$ and its
  negation, give 0.
\item Infinite-volume: \texttt{ciSup\_const} gives the supremum of the
  zero sequence as 0.
\end{itemize}
\end{theorem}

\begin{theorem}[\texttt{correlation\LeanLambda / correlationAlongExhaustion / correlationInfinite\_zero\_params\_vanish}; PR \#189]
PR \#188 is lifted through three layers, as the $J=h=0$ counterpart of
PR \#183.
\end{theorem}

\begin{theorem}[\texttt{correlation\_zero\_params\_vanish\_of\_nonempty\_A}; PR \#188]
Correlation vanishes on the $J = h = 0$ slice: if $A.\texttt{Nonempty}$, then
$\texttt{correlation}\,G\,\langle 0, 0, \beta\rangle\,A = 0$.
\texttt{hamiltonian\_zero\_params} gives Boltzmann weight $= 1$, and
PR \#148 \texttt{sum\_config\_spinProduct\_eq\_zero} makes the numerator
vanish.
This is the companion to PR \#182 ($\beta = 0$).
\end{theorem}

\begin{theorem}[\texttt{correlation\_beta\_zero\_vanish\_of\_nonempty\_A}; PR \#182]
At $\beta = 0$, correlations vanish for nonempty $A$:
\[
  A.\texttt{Nonempty} \Rightarrow
    \texttt{correlation}\,G\,\langle J, h, 0\rangle\,A = 0.
\]
This is the infinite-temperature ($\beta = 0$) slice of the Glimm--Jaffe
\S 4.1 correlation definition. Since the Boltzmann weight is 1, the
numerator is $\sum_\sigma \prod_{i \in A} \sigma_i$, which is 0 by the
existing \texttt{sum\_config\_spinProduct\_eq\_zero}
(NonnegCorrelations.lean, via the \texttt{flipAt} involution). Alongside the
existing $\beta = 0$ closed forms
(\texttt{freeEnergy\_beta\_zero} = $\log 2$, PR \#131;
\texttt{partitionFunction\_beta\_zero}), this completes the correlation
slice.
\end{theorem}

\begin{theorem}[\texttt{freeEnergyInfinite\_\{J\_zero,beta\_zero,zero\_params\}\_of\_nonempty}; PR \#181]
Under \texttt{[Nonempty V]}, these are assumption-free versions. The values
are $\log(2 \cosh(\beta h))$ for $J=0$ and $\log 2$ for $\beta=0$ or
$J=h=0$. They directly combine the existing
\texttt{\_of\_eventually\_nonempty} results with
\texttt{\LeanLambda.eventually\_volume\_nonempty}.
\end{theorem}

\begin{theorem}[\texttt{Ambient.Finset.Nonempty.fintype\_card\_coe\_pos}; PR \#180]
Helper extraction refactor: the repeated pattern in 21+ locations,
$\Lambda.\texttt{Nonempty} \Rightarrow 0 < \texttt{Fintype.card}\,\uparrow\Lambda$
is made into a theorem. The chain
\texttt{rw [Fintype.card\_coe]; exact Finset.card\_pos.mpr hne} is replaced
by one call. The mathematical content is unchanged, and readability improves
in AmbientLattice and AmbientLatticeSum.
\end{theorem}

\begin{theorem}[\texttt{freeEnergyAlongExhaustion\_\{beta\_zero,zero\_params\}\_tendsto\_of\_eventually\_nonempty}; PR \#179]
This is the $\beta=0$ and $J=h=0$ parallel version of the $J=0$ Tendsto
statement from PR \#178:
\[
  \texttt{Tendsto}\,(\cdots \langle J, h, 0\rangle)\,\texttt{atTop}\,(\mathcal N(\log 2))
\]
\[
  \texttt{Tendsto}\,(\cdots \langle 0, 0, \beta\rangle)\,\texttt{atTop}\,(\mathcal N(\log 2))
\]
It uses the same eventually-constant plus
\texttt{tendsto\_const\_nhds.congr'} pattern. This completes the
infinite-volume Tendsto results for all zero-parameter slices.
\end{theorem}

\begin{theorem}[\texttt{freeEnergyAlongExhaustion\_J\_zero\_tendsto\_of\_eventually\_nonempty}; PR \#178]
This is the first nontrivial infinite-volume \texttt{Tendsto} convergence
after the CLAUDE.local.md update bringing infinite systems into scope. Under
$\forall^\infty n, \Lambda_n$ nonempty,
\[
  \texttt{Tendsto}\,(\texttt{freeEnergyAlongExhaustion}\,G\,\Lambda\,\langle 0, h, \beta\rangle)\,
    \texttt{atTop}\,(\mathcal N(\log(2 \cosh(\beta h)))).
\]
The $J = 0$ slice is stagewise constant with value
$\log(2 \cosh(\beta h))$, so it avoids the Fekete step and does not require
translation invariance. Transport by \texttt{tendsto\_const\_nhds.congr'}.
\end{theorem}

\begin{theorem}[\texttt{partitionFunction\_ge\_two\_cosh\_pow\_card\_of\_ferromagnetic} family; PR \#177]
For ferromagnetic $p$, the non-logarithmic form is
\[
  (2 \cosh(\beta h))^{|\iota|} \;\le\; Z_G(p).
\]
This is the exponential image of the log form from PR \#173. Proof:
\texttt{partitionFunction\_bot} plus
\texttt{partitionFunction\_monotone\_subgraph}. It provides three-layer
lifts: base, $\Lambda$, and along-exhaustion.
\end{theorem}

\begin{theorem}[\texttt{partitionFunction\_ge\_two\_pow\_card\_of\_ferromagnetic} family; PR \#165]
For ferromagnetic $p$,
\[
  2^{|\iota|} \;\le\; Z_G(p).
\]
\texttt{partitionFunction\_bot\_ge\_two\_pow\_card}
This combines the $\bot$-graph bound with $\cosh \ge 1$ and
\texttt{partitionFunction\_monotone\_subgraph}. The log form
\texttt{log\_partitionFunction\_ge\_card\_mul\_log\_two\_of\_ferromagnetic}
states $|\iota| \cdot \log 2 \le \log Z_G$, and the same PR also provides
the $\Lambda$ and along-exhaustion lifts.
\end{theorem}

\begin{theorem}[\texttt{inducedGraph\_bot}; PR \#129]
$\texttt{inducedGraph}\,(\bot : \texttt{SimpleGraph}\,V)\,\Lambda
  = (\bot : \texttt{SimpleGraph}\,\uparrow\Lambda)$.
\texttt{induce = comap} + \texttt{SimpleGraph.comap\_bot}.
It is marked \texttt{@[simp]}.
\end{theorem}

\begin{theorem}[\texttt{freeEnergy\_ge\_log\_two\_cosh}]
For a ferromagnetic system with $J \ge 0, h \ge 0, \beta > 0$ and
$|\iota| > 0$,
\[
  \log(2 \cosh(\beta h)) \le \texttt{freeEnergy}\,G\,p.
\]
\texttt{freeEnergy\_bot} identifies the left-hand side with $f(\bot, p)$,
then \texttt{freeEnergy\_monotone\_subgraph} is applied to \texttt{bot\_le}.
This sharpens the uniform $\log 2$ lower bound from PR \#121 using
$\cosh \ge 1$.
\end{theorem}

\begin{theorem}[\texttt{freeEnergyAlongExhaustion\_ge\_log\_two\_cosh}]
The same lower bound holds at each nonempty stage.
Specialization via \texttt{change} + definitional unfolding.
\end{theorem}

\subsection{Spin flip symmetry / \texorpdfstring{$h$}{h}-even (PR \#126)}

\begin{theorem}[\texttt{hamiltonian\_neg\_h}]
\[
  H_{G, (J, -h, \beta)}(\sigma)
    = H_{G, (J, h, \beta)}(\sigma.\text{flip}).
\]
The $J$-term is invariant by \texttt{interactionEnergy\_flip}; the external
field satisfies
$-(-h) \sum \text{sign}(\sigma_i) = h \sum \text{sign}(\sigma_i)
  = -h \sum \text{sign}(\sigma.\text{flip}_i)$.
\end{theorem}

\begin{theorem}[\texttt{partitionFunction\_neg\_h}]
$Z(J, -h, \beta) = Z(J, h, \beta)$.
\texttt{hamiltonian\_neg\_h} plus reindexing along the \texttt{flipEquiv}
bijection.
\end{theorem}

\begin{theorem}[\texttt{freeEnergy\_neg\_h}]
$f(J, -h, \beta) = f(J, h, \beta)$.
The free-energy density is an even function of $h$.
\end{theorem}

\subsection{\texorpdfstring{$|h|$}{|h|}-monotonicity (PR \#127)}

\begin{theorem}[\texttt{freeEnergy\_eq\_abs\_h}]
$f(J, h, \beta) = f(J, |h|, \beta)$.
Split into the cases $|h| = h$ and $|h| = -h$, then use
\texttt{freeEnergy\_neg\_h}.
\end{theorem}

\begin{theorem}[\texttt{freeEnergy\_monotone\_abs\_h}]
For ferromagnetic systems,
$|h_1| \le |h_2| \Rightarrow f(J, h_1, \beta) \le f(J, h_2, \beta)$.
Rewrite both sides to $|h_i|$ using \texttt{freeEnergy\_eq\_abs\_h}, then
apply \texttt{freeEnergy\_monotone\_h} (monotone on \texttt{Ici 0}) from
$|h_i| \ge 0$.
\end{theorem}

\subsubsection{Along-exhaustion specializations (PR \#128)}

\begin{theorem}[\texttt{freeEnergyAlongExhaustion\_neg\_h} /
\texttt{\_eq\_abs\_h} / \texttt{\_monotone\_abs\_h}]
For each $n$, specialize the finite-volume versions directly:
\begin{itemize}
  \item $f_n(J, -h, \beta) = f_n(J, h, \beta)$.
  \item $f_n(J, h, \beta) = f_n(J, |h|, \beta)$.
  \item For ferromagnetic systems with $J \ge 0, \beta > 0$,
        $|h_1| \le |h_2| \Rightarrow f_n(h_1) \le f_n(h_2)$.
\end{itemize}
Each proof uses \texttt{change} and definitional unfolding to apply
\texttt{IsingModel.freeEnergy\_*}.
\end{theorem}

\subsection{Reflection positivity (\S10.4)}

\begin{definition}[\texttt{ReflectionPositive}]
A bilinear form $b$ is reflection-positive if $b(x, x) \ge 0$ for all $x$.
\end{definition}

\begin{theorem}[\texttt{discriminant\_nonneg}]
If $a t^2 + 2bt + c \ge 0$ for all $t \in \mathbb{R}$, then $b^2 \le ac$.
This is the algebraic core of the Schwarz inequality (10.4.2).
Proved using \texttt{discrim\_le\_zero} from Mathlib.
\end{theorem}

\subsection{Multiple reflections (\S10.5--10.6)}

\begin{theorem}[\texttt{iterated\_schwarz\_sq}; Prop.~10.5.2 core]
If $0 \le x$, $0 \le a$, and $x^2 \le ax$, then $x \le a$.
This is the algebraic core of the iterated Schwarz inequality
for multiple reflection bounds.
\end{theorem}

\S10.6 extends the Schwarz inequality to nonsymmetric reflections,
needed only for regularity of $P(\varphi)_2$ fields (not existence).

\subsubsection{Nonsymmetric reflections (\S10.6 algebraic core)}

For a possibly non-symmetric bilinear form $b$ with
$b(x, y) \ne b(y, x)$, expanding $b(x + ty, x + ty)$ gives
$b(y, y)\,t^2 + \bigl(b(x, y) + b(y, x)\bigr)\,t + b(x, x) \ge 0$
for all $t \in \mathbb{R}$, i.e.\ a quadratic with asymmetric linear
coefficient $b_1 + b_2$.

\begin{theorem}[\texttt{nonsymmetric\_discriminant\_mean}]
If $a t^2 + (b_1 + b_2)\,t + c \ge 0$ for all $t$, then
$\bigl((b_1 + b_2) / 2\bigr)^2 \le a \cdot c$.
Proved by reducing to \texttt{discriminant\_nonneg} at
$b := (b_1 + b_2) / 2$.
\end{theorem}

\begin{theorem}[\texttt{nonsymmetric\_mean\_le\_geom\_mean};
\texttt{nonsymmetric\_sum\_abs\_bound}]
$|(b_1 + b_2) / 2| \le \sqrt{a \cdot c}$, equivalently
$|b_1 + b_2| \le 2\sqrt{a \cdot c}$ (under $a, c \ge 0$).
\end{theorem}

\begin{theorem}[\texttt{schwarz\_abs\_bound}]
Symmetric counterpart: $|b| \le \sqrt{a \cdot c}$ from the standard
quadratic $a t^2 + 2bt + c \ge 0$ with $a, c \ge 0$.
\end{theorem}

\begin{theorem}[\texttt{discriminant\_nonneg\_converse};
\texttt{discriminant\_nonneg\_iff}]
Converse of \texttt{discriminant\_nonneg} for $a > 0$:
$b^2 \le a\,c$ implies $a t^2 + 2bt + c \ge 0$ for all $t$.
Proved by completing the square
$a t^2 + 2bt + c = a(t + b/a)^2 + (c - b^2/a)$.
Combined biconditional as \texttt{discriminant\_nonneg\_iff}.
\end{theorem}

Two-variable iteration of the non-symmetric Schwarz gives a
Cauchy-Schwarz-in-product form and cross-variable cube bounds:

\begin{theorem}[\texttt{nonsymmetric\_iterated\_schwarz};
\texttt{nonsymmetric\_two\_le\_sum}]
From $x^2 \le a \cdot b$ with $x, a, b \ge 0$:
$x \le \sqrt{a \cdot b}$ (GM bound) and $2x \le a + b$ (AM bound,
via $(a - b)^2 \ge 0$).
\end{theorem}

\begin{theorem}[\texttt{nonsymmetric\_product\_bound}]
Under $x^2 \le a$, $y^2 \le b$ with $xy \ge 0$:
$xy \le \sqrt{a \cdot b}$.
\end{theorem}

\begin{theorem}[\texttt{nonsymmetric\_cross\_iteration\_x},
\texttt{\_y}; \texttt{nonsymmetric\_cube\_bound\_x}, \texttt{\_y}]
From the cross-bounds $x^2 \le a \cdot y$ and $y^2 \le b \cdot x$:
$x^4 \le a^2 \cdot b \cdot x$ and $y^4 \le a \cdot b^2 \cdot y$.
When $x > 0$ (resp.\ $y > 0$): $x^3 \le a^2 \cdot b$ (resp.\
$y^3 \le a \cdot b^2$).
\end{theorem}

These algebraic pieces are the foundation for non-symmetric reflection
positivity; combined with a bilinear-form expansion hypothesis, they
yield Schwarz-type bounds without requiring $b(x, y) = b(y, x)$.

\paragraph{Capstone: reflection-positive Schwarz.}

\begin{theorem}[\texttt{polarization\_identity}]
For a bilinear form $b : \alpha \to \alpha \to \mathbb{R}$ on an
additive commutative group $\alpha$ (with bilinearity given as four
elementary hypotheses),
\[
  b(x + y, x + y) - b(x - y, x - y) = 2 \cdot (b(x, y) + b(y, x)).
\]
Exposes the symmetrized off-diagonal $b(x, y) + b(y, x)$ as a difference
of diagonal entries, useful even when $b$ is non-symmetric.
\end{theorem}

\begin{theorem}[\texttt{schwarz\_of\_reflection\_positive}; \S10.6 main]
For a bilinear $b : \alpha \to \alpha \to \mathbb{R}$ on an $\mathbb{R}$-module $\alpha$
with \texttt{ReflectionPositive b} ($b(x, x) \ge 0$),
\[
  \left(\frac{b(x, y) + b(y, x)}{2}\right)^2 \le b(x, x) \cdot b(y, y).
\]
Proof: bilinearity expands
$b(x + t\cdot y, x + t\cdot y) = b(y, y)\,t^2 + (b(x, y) + b(y, x))\,t
+ b(x, x)$;
reflection positivity gives this quadratic $\ge 0$ for all $t$;
\texttt{nonsymmetric\_discriminant\_mean} yields the Schwarz bound.
\end{theorem}

\begin{theorem}[Corollaries]
AM-GM form \texttt{reflection\_positive\_mean\_le\_geom\_mean}:
$|(b(x, y) + b(y, x)) / 2| \le \sqrt{b(x, x) \cdot b(y, y)}$,
equivalently \texttt{\_sum\_abs\_bound}:
$|b(x, y) + b(y, x)| \le 2 \sqrt{b(x, x) \cdot b(y, y)}$.

Degenerate case \texttt{\_off\_diag\_zero\_of\_diag\_zero}:
$b(x, x) = 0 \Rightarrow b(x, y) + b(y, x) = 0$
(and the symmetric $\texttt{\_right}$ variant for $b(y, y) = 0$).
\end{theorem}

\begin{theorem}[\texttt{classical\_schwarz\_of\_symmetric\_reflection\_positive}]
Under the additional symmetry $b(x, y) = b(y, x)$, the non-symmetric
bound reduces to the classical Cauchy-Schwarz inequality:
$(b(x, y))^2 \le b(x, x) \cdot b(y, y)$, with abs form
$|b(x, y)| \le \sqrt{b(x, x) \cdot b(y, y)}$, and the sharper
degenerate case $b(x, x) = 0 \Rightarrow b(x, y) = 0$ (for the
individual entry, not just the sum).
\end{theorem}

This completes the §10.6 algebraic core required for non-symmetric
reflection-positivity Schwarz bounds.

\paragraph{Closure operations on ReflectionPositive.}

The class of reflection-positive bilinear forms is closed under
standard cone operations:

\begin{itemize}
\item \texttt{ReflectionPositive.zero}: the zero form is RP.
\item \texttt{ReflectionPositive.const}: the constant form $b(x, y) = c$ with $c \ge 0$ is RP.
\item \texttt{ReflectionPositive.of\_diag\_nonneg}: a form depending only on the first argument, with nonneg values, is RP.
\item \texttt{ReflectionPositive.add}: the sum of two RP forms is RP.
\item \texttt{ReflectionPositive.smul\_nonneg}: a nonneg scalar multiple of an RP form is RP.
\item \texttt{ReflectionPositive.comp}: reparametrization $b'(x, y) := b(g(x), g(y))$ under $g : \beta \to \alpha$ inherits RP.
\item \texttt{ReflectionPositive.sum}: a finite sum $\sum_i b_i$ of RP forms is RP.
\item \texttt{ReflectionPositive.weighted\_sum}: a nonneg-weighted sum $\sum_i c_i b_i$ is RP.
\item \texttt{ReflectionPositive.of\_diag\_eq}: RP depends only on diagonal values (transfer lemma).
\end{itemize}

Together these show that RP bilinear forms form a convex cone closed
under reparametrization and finite sums, supplying building blocks
for constructing concrete RP witnesses in applications.

\section{Phase transitions (\S16.1)}

\begin{theorem}[\texttt{magnetization\_monotone\_h}; \S16.1, p.~283]
The magnetization $M(i) = \langle\sigma_i\rangle$ is monotone increasing
in $h$ on $[0, \infty)$ for ferromagnetic parameters.
Follows from \texttt{correlation\_monotone\_h}.
\end{theorem}

In finite volume, $f(h)$ is real-analytic (\texttt{freeEnergyH\_analyticOn}),
so $M(h)$ is smooth and no phase transitions occur.
Phase transitions appear only in the infinite volume limit:
a first-order transition at $h_0$ corresponds to a jump discontinuity
of $M(h)$ at $h_0$ (\S16.1, p.~283).

\section{Critical exponents (\S17.5, \S17.7)}

\begin{theorem}[\texttt{eta\_nonneg\_finite\_vol}; Thm~17.7.1]
The critical exponent $\eta \ge 0$: in finite volume,
$\langle\sigma_i;\sigma_j\rangle \ge 0$ by GKS-II
(\texttt{truncated2\_nonneg}).
\end{theorem}

The bound $\zeta \ge 0$ follows from $U_4 \le 0$ (\texttt{cor\_4\_3\_3}).
The correlation length definition (inverse mass, \S17.5)
and the continuity theorem (Thm~17.5.1) require the infinite
volume limit and transfer matrix spectral theory.

\section{Cluster expansion (\S18.1--18.3)}

The lattice Ising analogue of the cluster expansion is the
\textbf{high-temperature expansion}, decomposing each Boltzmann factor:
$\exp(\beta J \sigma_i\sigma_j) = \cosh(\beta J)(1 + \tanh(\beta J)\sigma_i\sigma_j)$.
This is already proved as \texttt{exp\_edgeSpin\_decomp}.

\begin{definition}[\texttt{highTempParam}]
$t = \tanh(\beta J)$, the high-temperature expansion parameter.
\end{definition}

\begin{theorem}[\texttt{abs\_highTempParam\_lt\_one}]
$|t| = |\tanh(\beta J)| < 1$ for all finite $\beta J$.
\end{theorem}

\begin{theorem}[\texttt{highTempParam\_lt\_one}]
$t < 1$.
\end{theorem}

\begin{theorem}[\texttt{partitionFunction\_high\_temp\_expansion}; GJ~\S18.3,
cf.\ FV~\S3.7.3 eq.\,(3.42)]
For any Ising parameter $p = (J, h, \beta)$ on a finite simple graph
$G = (\iota, E)$,
\[
Z(G; p) = (\cosh(\beta J))^{|E|}
\sum_{\sigma}
\Bigl(\prod_{\{i,j\} \in E} (1 + \tanh(\beta J)\,\sigma_i\sigma_j)\Bigr)
\exp\!\Bigl(\beta h \sum_i \sigma_i\Bigr).
\]
\end{theorem}

\textit{Proof.} The Hamiltonian factorises as
$-\beta H = \beta J \sum_e \sigma_i\sigma_j + \beta h \sum_i \sigma_i$,
so $e^{-\beta H} = \prod_e e^{\beta J \sigma_i\sigma_j} \cdot
e^{\beta h \sum_i \sigma_i}$.
Since $\sigma_i\sigma_j \in \{\pm 1\}$,
$e^{\beta J \sigma_i\sigma_j} = \cosh(\beta J) + \sinh(\beta J)\sigma_i\sigma_j
= \cosh(\beta J)(1 + \tanh(\beta J)\sigma_i\sigma_j)$
by \texttt{exp\_edgeSpin\_decomp} and $\cosh(\beta J) > 0$.
Multiplying over edges and pulling out $\cosh(\beta J)^{|E|}$ gives the claim.
\hfill$\square$

\begin{theorem}[\texttt{partitionFunction\_high\_temp\_expansion\_h\_zero}; GJ~\S18.3 / FV~\S3.7.3 eq.\,(3.42)]
At zero external field, the high-temperature expansion specialises to
\[
Z(G; J, 0, \beta) = (\cosh(\beta J))^{|E|}
\sum_{\sigma}
\prod_{\{i,j\} \in E} (1 + \tanh(\beta J)\,\sigma_i\sigma_j).
\]
\end{theorem}

\textit{Proof.} Direct specialisation of the previous theorem at $h = 0$,
where the field factor $\exp(\beta h \sum_i \sigma_i) = 1$.
\hfill$\square$

\begin{theorem}[\texttt{partitionFunction\_high\_temp\_expansion\_h\_zero\_closed}; FV~\S3.7.3 eq.\,(3.45)]
For any Ising parameter $(J, 0, \beta)$ on a finite simple graph
$G = (\iota, E)$,
\[
Z(G; J, 0, \beta) = 2^{|\iota|} \cdot (\cosh(\beta J))^{|E|}
\sum_{\substack{X \subseteq E \\ \text{every $v$ has even $X$-degree}}}
  \tanh(\beta J)^{|X|}.
\]
\end{theorem}

\textit{Proof.} Starting from the zero-field expansion
$Z = (\cosh \beta J)^{|E|} \sum_\sigma \prod_e (1 + \tanh(\beta J)\sigma_i\sigma_j)$,
expand each edge product via
$\prod_{e \in E}(1 + f_e) = \sum_{X \subseteq E} \prod_{e \in X} f_e$
(\texttt{Finset.prod\_one\_add}); pull out $\tanh(\beta J)^{|X|}$ and
swap the $\sigma$- and $X$-sums. The inner $\sigma$-sum
$\sum_\sigma \prod_{e \in X} \sigma_i\sigma_j = \prod_v \sum_{s_v}(s_v)^{d_X(v)}$
factorises by Fubini and equals $2^{|\iota|}$ when every
$d_X(v) := |\{e \in X : v \in e\}|$ is even, else $0$. The terms with
odd-degree vertices vanish, leaving the sum over even-degree
subgraphs.
\hfill$\square$

\begin{theorem}[\texttt{correlation\_high\_temp\_expansion\_h\_zero\_closed}; FV~\S3.7.3 eq.\,(3.46)]
For any $A \subseteq \iota$ on a finite simple graph $G = (\iota, E)$,
\[
\langle \sigma_A \rangle_{\beta, 0} =
  \frac{\sum_{X \subseteq E,\, \partial X = A} \tanh(\beta J)^{|X|}}
       {\sum_{X \subseteq E,\, \partial X = \emptyset} \tanh(\beta J)^{|X|}},
\]
where $\partial X := \{ v : \iota \mid d_X(v) \text{ is odd}\}$ is the
boundary of the edge set $X$.
\end{theorem}

\textit{Proof.} The correlation function is
$\langle \sigma_A \rangle = N / Z$ with
$N := \sum_\sigma \sigma_A \cdot e^{-\beta H(\sigma)}$. Repeating the
high-temperature expansion of the previous theorem with the extra
factor $\sigma_A = \prod_{v \in A} (\sigma_v)^{1}$ gives
\[
N = 2^{|\iota|} \cdot (\cosh \beta J)^{|E|}
   \sum_{X \subseteq E,\, \partial X = A} \tanh(\beta J)^{|X|};
\]
the parity step now requires
$d_X(v) + \mathbf{1}_{v \in A}$ to be even at every $v$, i.e.\
$d_X(v) \equiv \mathbf{1}_{v \in A} \pmod{2}$, equivalently
$\partial X = A$. Dividing by $Z$ from the previous theorem
cancels the prefactor $2^{|\iota|} (\cosh \beta J)^{|E|}$.
\hfill$\square$

\paragraph{Wrappers.} The two closed forms are lifted to the
\(\Lambda\)-layer (\texttt{partitionFunction\LeanLambda\_high\_temp\_expansion\_h\_zero\_closed},
\texttt{correlation\LeanLambda\_high\_temp\_expansion\_h\_zero\_closed} in
\texttt{AmbientLattice/Defs.lean}) and to the concrete
\(\mathbb{Z}^d\) lattice
(\texttt{partitionFunction\LeanLambda\_latticeGraph\_high\_temp\_expansion\_h\_zero\_closed},
\texttt{correlation\LeanLambda\_latticeGraph\_high\_temp\_expansion\_h\_zero\_closed}
in \texttt{Concrete/LatticeGraphCorrelation/HighTemperatureBounds.lean}).

\begin{theorem}[\texttt{partitionFunction\_high\_temp\_expansion\_h\_zero\_lower\_bound}]
For \(0 \le \beta J\),
\(Z(G; J, 0, \beta) \ge 2^{|\iota|} \cdot (\cosh(\beta J))^{|E|}.\)
\end{theorem}

\textit{Proof.} In FV (3.45) the empty edge set \(X = \emptyset\) is
trivially even-degree (every vertex has degree \(0\)) and contributes
\(\tanh(\beta J)^0 = 1\) to the sum; every other even-degree subset
contributes a nonneg amount once \(\tanh(\beta J) \ge 0\) is known
(equivalent to \(0 \le \beta J\)). Multiplying the inequality
\(\sum_X \tanh(\beta J)^{|X|} \ge 1\) by the nonneg prefactor
\(2^{|\iota|} (\cosh \beta J)^{|E|}\) gives the claim.
\hfill$\square$

\section{Cluster (polymer / Mayer) expansion (\S18.4--\S18.7)}

\paragraph{Polymer-family identity (\S18.4).} The
even-subgraph sum in FV~(3.45) is reorganised into a sum over
\textbf{vertex-disjoint compatible polymer families}. A
\emph{polymer} is a non-empty connected even-degree edge subset
(\texttt{IsPolymer}); compatibility is vertex-disjointness
(\texttt{IsPolymerVertexDisjoint}). Each even-degree edge set $X$
decomposes uniquely into its connected components
(\texttt{polymerDecomposition}), giving a bijection between
$\{X : \partial X = \emptyset\}$ and
$\mathrm{vdCompatiblePolymerFamilies}\,G$. The capstone identity is
\[
Z(G; J, 0, \beta) = 2^{|\iota|} \cdot (\cosh \beta J)^{|E|}
  \sum_{\Gamma \in \mathrm{vdCompatiblePolymerFamilies}(G)}
    \prod_{P \in \Gamma} \tanh(\beta J)^{|P|}
\]
(\texttt{partitionFunction\_high\_temp\_expansion\_h\_zero\_polymer\_family}).

\paragraph{Mayer / Ursell infrastructure (Steps~576--615).}
The exponential identity $\log \Xi = \sum_{n \ge 1} \frac{1}{n!}
\sum_{\omega} \phi^T(\omega)\, z(t,\omega)$ is formalised via
\texttt{ursellCoefficient} (\(\phi^T(\omega)\),
Step~583), \texttt{clusterSeqActivity} (\(z(t,\omega)\),
Step~581), \texttt{mayerExpansionTerm} (the $n$-th term,
Step~587), and \texttt{mayerPartialSum} (truncation through
cluster size $N$, Step~591). Pair-Ursell formula
$\phi^T(\omega) = -\tfrac{1}{2}$ for incompatible pairs,
$0$ otherwise (Steps~585--586); uniform absolute bound
$|\phi^T(\omega)| \le 2^{\binom{n}{2}}/n!$ (Step~615) via
$|\mathrm{connectedSpanningEdgeSubsets}\,G| \le 2^{|G.\mathrm{edgeFinset}|}$
(Step~602) plus
\texttt{SimpleGraph.card\_edgeFinset\_le\_card\_choose\_two}.
Polymer free energy
$\mathrm{polymerFreeEnergy}\,G\,t := \log(\sum_\Gamma \prod_P t^{|P|})$
(Step~610) is real-analytic on $[0,\infty)$ and connects to the
Ising free energy via
$f(J,0,\beta) = \log 2 + (|E|/|\iota|) \log\cosh(\beta J)
  + \mathrm{polymerFreeEnergy}\,G\,(\tanh \beta J)/|\iota|$
under $0 < |\iota|$, $0 \le \beta J$ (Step~612).

\paragraph{Phase~B base cases (PRs~\#1518--\#1546, \#1553).}
Explicit values of the alternating connected-subgraph sum
$\mathrm{alternatingConnectedSubgraphSum}\,G
:= \sum_{S \in \mathrm{connectedSpanningEdgeSubsets}\,G} (-1)^{|S|}$
were computed by \texttt{decide} (with
\texttt{set\_option maxRecDepth} / \texttt{maxHeartbeats} raised) on:
\begin{itemize}
  \item Complete graphs \(K_3 = 2\), \(K_4 = -6\)
    (PRs~\#1518, \#1519);
  \item Path graphs \(P_n\) for \(n = 3,\dots,8\) with values
    \(1, -1, 1, -1, 1, -1\) (PRs~\#1542--\#1544; \(P_8\) requires
    \texttt{maxRecDepth 8000} + \texttt{maxHeartbeats 1000000});
  \item Cycle graphs \(C_n\) for \(n = 3,\dots,7\) with values
    \(2, -3, 4, -5, 6\) (PRs~\#1545, \#1546, \#1553).
\end{itemize}
The general formula
$\mathrm{alternatingConnectedSubgraphSum}(K_n) = (-1)^{n-1}(n-1)!$
and any case with $n \ge 9$ require chromatic-polynomial /
matrix-tree-theorem machinery not yet in Mathlib, and remain
deferred (Phase~B blocker).

\paragraph{Polymer free energy bounds and characterisations
(PRs~\#1517--\#1566).} A coordinated set of sharpenings was added
in this batch:
\begin{itemize}
  \item \emph{Convergence radius via $\log$ \texttt{HasSum}}
    (\#1517): \texttt{polymerFreeEnergy\_hasSum\_via\_log\_of\_pow\_lt\_two}
    expresses $\mathrm{polymerFreeEnergy}\,G\,t$ as a `HasSum` of
    the natural log expansion under the explicit radius
    $(1+t)^{|E|} < 2$.
  \item \emph{Filter forms} (\#1521, \#1522): the
    \texttt{mayerExpansionTerm} / \texttt{mayerPartialSum} sums
    can be restricted to connected polymer-sequence supports
    (\texttt{\_filter\_connected}) without changing the value.
  \item \emph{Upper bounds} (\#1523--\#1527):
    $\mathrm{polymerFreeEnergy}\,G\,t \le \varepsilon(G,t)$
    (where $\varepsilon := \sum_\Gamma \prod_P t^{|P|} - 1$),
    $\le (1+t)^{|E|} - 1$, $< \log 2$ under
    $(1+t)^{|E|} < 2$, the
    \texttt{polymerFreeEnergy\_high\_temp\_sandwich} 5-tuple, and
    the explicit free-energy correction
    \texttt{freeEnergy\_lt\_log\_two\_plus\_high\_temp\_correction}.
  \item \emph{Strict / iff characterisations}
    (\#1547--\#1550, \#1554--\#1560): under $\varepsilon > 0$,
    $\mathrm{polymerFreeEnergy} < \varepsilon$ (strict; \#1547),
    $= 0 \iff \varepsilon = 0$, $> 0 \iff \varepsilon > 0$
    (\#1548), strict $< (1+t)^{|E|} - 1$ (\#1554, \#1555 tanh
    form), \texttt{vdPolymerFamilies\_sum\_eq\_one\_iff\_eps\_eq\_zero}
    and \texttt{\_gt\_one\_iff\_eps\_pos} (\#1557), bundles of
    strict-mono / strict-pos (\#1559, \#1560)
    including \texttt{StrictMonoOn (Set.Ioi 0)}.
  \item \emph{Pair connectivity} (\#1551, \#1552):
    \texttt{polymerSeqIncompatibilityGraph\_two\_connected\_iff\_incompatible}
    and the connected-restricted form of the $n=2$ Mayer term
    equate to incompatibility.
  \item \emph{$\Lambda$-layer wrappers} (\#1563, \#1564):
    $10 + 8$ theorems lifting the abstract API to the
    $\Lambda$ layer in
    \texttt{IsingModel/AmbientLattice/Analyticity.lean}.
  \item \emph{Ferromagnetic / tanh-monotonicity} (\#1565, \#1566):
    $0 \le J,\, 0 < \beta$ versions of the strict / iff bundle, and
    $\mathrm{polymerFreeEnergy}(\tanh(\beta J))$ monotone (resp.\
    strictly monotone) in $\beta$ and $J$ under the appropriate
    sign assumptions; \#1566 introduces the local helper
    \texttt{real\_tanh\_strictMono}.
\end{itemize}

\paragraph{General-$h$ analyticity capstone (\S18.6,
PRs~\#1528--\#1541).} Per-direction and joint analyticity of the
free energy at \emph{general} external field $h$:
\begin{itemize}
  \item \texttt{freeEnergy\_analyticAt\_\{beta,J,h\}\_general\_h}
    (\#1528--\#1530) and matching
    \texttt{\_analyticOnNhd\_\{beta,J,h\}\_general\_h}: the partial
    real-analyticity of $f(J,h,\beta)$ is now available without
    the $h = 0$ restriction.
  \item \texttt{freeEnergy\_analyticAt\_joint} (\#1531) and
    \texttt{\_analyticOnNhd\_joint} (\#1532) over $\mathrm{Set.univ}$:
    joint analyticity in $(\beta, J, h)$.
  \item \texttt{freeEnergy\_continuous\_joint} +
    \texttt{\_differentiable\_joint} (\#1533).
  \item Lambda-layer joint analyticity in
    \texttt{IsingModel/AmbientLattice/Analyticity.lean} (\#1535).
  \item \texttt{correlation\_analyticAt\_joint} +
    full hierarchy at the Lambda layer (\#1536, \#1537).
  \item \texttt{gibbsExpectation\_analyticAt\_joint} for an arbitrary
    observable $F: \mathrm{Configuration} \to \mathbb{R}$ (\#1539).
  \item \texttt{partitionFunction} continuity / differentiability per
    direction at general $h$ (\#1540) plus Lambda-layer wrappers
    (\#1541).
\end{itemize}

\paragraph{Ferromagnetic forms of \texttt{polymerFreeEnergy}
sandwich + hasSum, all volume layers (\S18.5, PR~\#1573).} Under
$0 \le J,\, 0 < \beta$ each of the §18.5 PR~\#1569 statements
(\texttt{polymerFreeEnergy\_tanh\_high\_temp\_sandwich} and
\texttt{polymerFreeEnergy\_tanh\_hasSum\_via\_log\_of\_pow\_lt\_two})
admits a ferromagnetic version, derived through
\texttt{mul\_nonneg h$\beta$.le hJ} and matched at each of the five
layers (abstract / $\Lambda$-direct / along-exhaustion /
$\mathbb{Z}^d$ $\Lambda$-direct / $\mathbb{Z}^d$ along-exhaustion).
The four $\mathbb{Z}^d$ wrappers use the abbreviated suffix
\texttt{\_ferro} (instead of \texttt{\_ferromagnetic}) so each
declaration fits the linter's 100-character budget. The along-exhaustion
and $\mathbb{Z}^d$ high-temperature wrappers from this family now live in
\texttt{AmbientLattice/SpecialCases/HighTemperature.lean} and
\texttt{Concrete/LatticeGraphCorrelation/HighTemperature.lean};
PR~\#1864 moved only those non-$\Lambda$ wrappers out of the legacy
modules. The $\Lambda$-layer wrappers were already in
\texttt{IsingModel/AmbientLattice/Analyticity.lean} and remain there.

\paragraph{\texttt{vdPolymerFamilies\_sum} bound family completion,
all volume layers (\S18.5, PR~\#1595).} Adds Lambda, along-ex, and
$\mathbb{Z}^d$ wraps for the three abstract bound theorems
\texttt{\_le\_two\_pow},
\texttt{\_le\_one\_plus\_tanh\_pow}, and
\texttt{one\_le\_vdPolymerFamilies\_sum} (which had no wrappers
before). The Lambda wrappers for the four sandwich abstracts already
existed; this PR fills the bare-bound and `\geq 1` gaps. 12 new
thin direct-instantiation wrappers (3 abstracts $\times$ 4
layers). The along-exhaustion and concrete $\mathbb{Z}^d$ wrappers
now live in the matching \texttt{MayerVdBounds.lean} child modules;
the $\Lambda$-layer wrappers live in
\texttt{IsingModel/AmbientLattice/AnalyticityLambdaBasicIdentities.lean}
(narrowed out of \texttt{Analyticity.lean} in PR~\#1986).

\paragraph{\S17.5 Lemma 17.5.2 stub honest naming + docs cleanup
(PR~\#1644).} The previously-named
\texttt{latticeMass\_ge\_pseudoMass} and
\texttt{latticeMass\_le\_constant\_mul\_pseudoMass} in
\texttt{PseudoMass.lean} were misleading aliases of
\texttt{pseudoMass\_pos} and \texttt{discrete\_hls\_constant}
respectively (their conclusions did not mention
\texttt{latticeMass}). They have been renamed to
\texttt{lemma\_17\_5\_2\_pseudoMass\_pos} and
\texttt{lemma\_17\_5\_2\_constant\_exists} to honestly reflect
their content as \emph{helpers} toward the full Lemma 17.5.2,
not the lemma itself. \texttt{docs/index.md} updated: §17.5
Lemma 17.5.2 (full) is now marked as \textbf{Partial}, with the
gap explicitly identified as the missing concrete-to-abstract map
$(d, \Lambda, p) \to (\alpha, r, c)$ and the exponential decay
bounds on $\mathbb{Z}^d$ (Step 117h+).

\paragraph{\S17.5 Lemma 17.5.2 conditional sandwich capstone
(Step 117l, PRs~\#2663--\#2672).} The new module
\texttt{Concrete/LatticeGraphCorrelation/Lemma\_17\_5\_2.lean}
introduces the named capstone interface for
Glimm--Jaffe \S17.5 Lemma~17.5.2 (2nd ed., pp.~311--312):
\[
m^-(\beta) \le m(\beta) \le C \cdot m^-(\beta),
\]
where $m^- := \mathrm{pseudoMassFromParamsAtPair}$ and
$m := \mathrm{latticeMass}$.  The capstone module exposes the following named
declarations:
\begin{itemize}
\item \texttt{Lemma\_17\_5\_2\_UpperBound} -- the \texttt{Prop}-valued
  predicate naming the upper-bound side
  $m \le C \cdot \mathrm{ofReal}\,m^-$ as a hypothesis.  Future
  Lipschitz + HLS work (cf.\ \texttt{pseudoMass\_pow\_succ\_lipschitz}
  and \texttt{discrete\_hls\_constant}) is intended to discharge it
  unconditionally.
\item \texttt{lemma\_17\_5\_2\_lower\_bound\_of\_decay} -- the
  Lemma~17.5.2 lower-bound side, a renaming of the existing decay-based
  bridge \texttt{latticeMass\_ge\_pseudoMassFromParamsAtPair\_of\_decay}
  under the GJ-named identifier.
\item \texttt{lemma\_17\_5\_2\_sandwich\_of\_decay\_and\_upper} -- the
  conditional sandwich capstone, returning both inequalities once a
  validating exponential decay rate and an upper-bound hypothesis are
  supplied.
\item \texttt{lemma\_17\_5\_2\_upper\_bound\_of\_all\_decay\_rates\_le} --
  the final order-theoretic upper-bound assembly: since
  \texttt{latticeMass} is an \texttt{sSup} over validating decay rates, a
  uniform bound on every admissible rate by $C\,\mathrm{ofReal}(m^-)$ gives
  the named \texttt{Lemma\_17\_5\_2\_UpperBound}.
\item \texttt{lemma\_17\_5\_2\_sandwich\_of\_decay\_and\_all\_decay\_rates\_le}
  -- the corresponding sandwich wrapper combining the lower validating
  pseudo-mass decay input with the all-decay-rate upper assembly.
\item \texttt{lemma\_17\_5\_2\_cubic\_high\_temp\_lower\_capstone} and
  \texttt{lemma\_17\_5\_2\_cubic\_high\_temp\_latticeMass\_pos} -- the
  Lemma~17.5.2 lower-bound side specialised to the cubic exhaustion at
  high temperature $\beta J (2d) < 1$, derived as direct projections
  of the existing
  \texttt{cubicNamedRate\_capstone\_bundle\_of\_cubic\_corr\_mem\_Ioo\_and\_profile}
  bundle.  These do not require any upper-bound hypothesis, but they
  remain conditional on active-range membership of the anchored cubic
  pair correlation $\langle\sigma_0\sigma_z\rangle_\infty$ in $(0,2)$
  together with the tanh-power profile lower bound
  $g_\alpha(-\log(\beta J\,2d)) \le \langle\sigma_0\sigma_z\rangle_\infty$.
\item \texttt{lemma\_17\_5\_2\_constant\_existence} -- a strict-positivity
  alias of \texttt{discrete\_hls\_constant} ($2\alpha > d$).
\item \texttt{lemma\_17\_5\_2\_hls\_convolution\_constant} -- the
  positive constant form of the polynomial HLS convolution estimate.
\item \texttt{Lemma\_17\_5\_2\_HLSDenominatorComparison} -- the named
  finite-stage scalar comparison between the high-temperature
  Lebowitz/susceptibility derivative bound and the HLS denominator
  $Kc(\beta)/(h(\beta))^{2\alpha}$.
\item \texttt{Lemma\_17\_5\_2\_InfiniteHLSDenominatorComparison} -- the
  infinite-volume derivative comparison
  $|c'(\beta)|\le Kc(\beta)/(h(\beta))^{2\alpha}$ for
  $c(\beta)=\langle\sigma_x\sigma_z\rangle_\infty$.
\item \texttt{lemma\_17\_5\_2\_beta\_deriv\_abs\_le\_high\_temp} -- the
  finite-stage high-temperature beta-derivative absolute bound
  $|c'|\le J M^2 + 4dJ$, exposing the concrete Lebowitz/susceptibility
  input that must be compared with the HLS pseudo-mass denominator.
\item \texttt{lemma\_17\_5\_2\_hls\_derivative\_hypothesis\_of\_abs\_bound}
  -- the scalar comparison bridge from a concrete absolute derivative bound
  to the HLS denominator hypothesis
  $|c'|\le Kc(\beta)/(h(\beta))^{2\alpha}$.
\item \texttt{lemma\_17\_5\_2\_beta\_hls\_derivative\_hypothesis\_of\_high\_temp}
  -- the finite-stage high-temperature package combining the preceding
  beta-derivative bound with that HLS denominator comparison.
\item \texttt{lemma\_17\_5\_2\_beta\_pseudoMass\_power\_deriv\_le\_of\_high\_temp\_bound}
  -- the finite-stage concrete handoff from the HLS derivative hypothesis to
  \texttt{pseudoMass\_power\_deriv\_le}, giving
  $(h(\beta))^{2\alpha}|h'(\beta)|\le K/\rho$.
\item \texttt{lemma\_17\_5\_2\_beta\_pseudoMass\_power\_deriv\_le\_of\_hls\_constant}
  -- the existential HLS-constant version of the preceding power-derivative
  estimate.
\item \texttt{lemma\_17\_5\_2\_beta\_pseudoMass\_pow\_succ\_deriv\_bound\_of\_high\_temp\_bound}
  -- the corresponding finite-stage derivative bound for
  $\beta\mapsto(h(\beta))^{2\alpha+1}$.
\item \texttt{lemma\_17\_5\_2\_beta\_pseudoMass\_pow\_succ\_deriv\_bound\_of\_hls\_constant}
  -- the HLS-constant existential version of the pointwise
  $(h(\beta))^{2\alpha+1}$ derivative bound.
\item \texttt{lemma\_17\_5\_2\_beta\_pseudoMass\_pow\_succ\_lipschitz\_of\_high\_temp\_bound}
  -- the finite-stage interval Lipschitz estimate for
  $\beta\mapsto(h(\beta))^{2\alpha+1}$, obtained by applying the MVT to the
  pointwise concrete pseudo-mass power derivative bound.
\item \texttt{lemma\_17\_5\_2\_beta\_pseudoMass\_pow\_succ\_lipschitz\_of\_hls\_constant}
  -- the corresponding HLS-constant existential interval estimate, still
  conditional on the pointwise denominator comparison.
\item \texttt{lemma\_17\_5\_2\_infinite\_pseudoMass\_pow\_succ\_deriv\_bound\_of\_hls\_constant}
  and
  \texttt{lemma\_17\_5\_2\_infinite\_pseudoMass\_pow\_succ\_lipschitz\_of\_hls\_constant}
  -- the corresponding infinite-volume pointwise derivative and interval
  Lipschitz handoffs, retaining the HLS convolution inequality for the chosen
  constant and using the exact infinite-volume denominator comparison.
\end{itemize}
PR~\#2664 adds a finite high-temperature bridge discharging the upper-bound
predicate in the cubic active-range setting without adding an axiom.
The theorem
\texttt{lemma\_17\_5\_2\_cubic\_high\_temp\_upper\_bound\_of\_active\_range}
uses the already-proved Step~115 bound
\[
  \mathrm{latticeMass}(d,\Lambda_{\mathrm{cubic}},\langle J,0,\beta\rangle)
  \leq
  \mathrm{ofReal}\bigl(-\log(\tanh(\beta J))\bigr)
\]
together with exhaustion-independence of \texttt{latticeMass} and the
active-range positivity of
\[
  m^- :=
  \texttt{cubicOriginPseudoMassFromParamsAtPair}\,\alpha\,r\,\beta\,J\,z.
\]
Since $m^- > 0$, the finite constant
\[
  C_{\mathrm{fin}} =
  \frac{-\log(\tanh(\beta J))}{m^-}
\]
satisfies
\[
  \mathrm{latticeMass}(d,\Lambda,\langle J,0,\beta\rangle)
  \leq C_{\mathrm{fin}}\,\mathrm{ofReal}(m^-).
\]
In Lean this last algebraic step is
\texttt{ENNReal.ofReal\_div\_of\_pos} followed by
\texttt{ENNReal.div\_mul\_cancel}.  The theorem
\texttt{lemma\_17\_5\_2\_cubic\_high\_temp\_sandwich\_capstone}
then combines this finite upper bridge with
\texttt{lemma\_17\_5\_2\_cubic\_high\_temp\_lower\_capstone}, returning
\texttt{HasExponentialDecay}, the lower inequality
\[
  \mathrm{ofReal}(m^-) \leq
  \mathrm{latticeMass}(d,\Lambda,\langle J,0,\beta\rangle),
\]
and the finite upper inequality above.

This finite constant is not the HLS-uniform constant from the proof of
GJ~Theorem~17.5.1 (p.~312).  It closes the existing Lean upper-bound
predicate in the high-temperature active-range regime using Step~115; the
sharper HLS-uniform upper-bound proof remains a future refinement.

\paragraph{\S5.4 Prop 5.4.2 infinite-volume Peierls bound (liminf form)
(PR~\#1643).} Adds the genuine infinite-volume lift of GJ Prop 5.4.2:
under the per-stage hypotheses of \texttt{prop\_5\_4\_2\_along\_exhaustion},
the canonical $\infty$-vol $+$-BC quantity
\texttt{plusGibbsExpectationLiminf} $G \Lambda \langle J, 0, \beta\rangle B
(\lambda n\,\sigma.\,\mathrm{sign}(\sigma\,i_n))$ satisfies the
sandwich
$0 \le 1 - \mathrm{plusGibbsExpectationLiminf} \le e^{-c\beta}$.
Theorems
\texttt{prop\_5\_4\_2\_plusGibbsExpectationLiminf\_lower\_bound}
(via \texttt{Filter.liminf\_le\_limsup} chain through
\texttt{Filter.limsup\_le\_of\_le} on the per-stage
\texttt{plusGibbsExp}$_n \le 1$) and the full sandwich
\texttt{prop\_5\_4\_2\_plusGibbsExpectationLiminf} combining the
existing upper bound (PR~\#1099) with the new lower bound.

\paragraph{\texttt{partitionFunctionAlongExhaustion} /
\texttt{freeEnergyAlongExhaustion} joint AnalyticAt + AnalyticOnNhd
wrappers (along-ex + $\mathbb{Z}^d$)
(\S17.5--\S17.6, PR~\#1642).} 8 new pointwise / global analytic
wrappers, each a thin instantiation of the existing $\Lambda$-layer
\texttt{partitionFunction\LeanLambda\_*\_joint} /
\texttt{freeEnergy\LeanLambda\_*\_joint} at
$\Lambda := \Lambda.\mathrm{volume}\,n$. The along-exhaustion wrappers
now live in
\texttt{AmbientLattice/SpecialCases/JointAnalyticity.lean}, and the
$\mathbb{Z}^d$ wrappers live in
\texttt{Concrete/LatticeGraphCorrelation/JointAnalyticity.lean}
(moved from the legacy modules in PR~\#1860).

\paragraph{Joint AnalyticAt + AnalyticOnNhd wrappers across
all 4 layers (\S17.5--\S17.6, PR~\#1641).} Adds joint (in
$(\beta, J, h)$) AnalyticAt and AnalyticOnNhd wrappers for
\texttt{magnetization}, \texttt{susceptibility}, and along-ex
\texttt{correlation} across abstract, $\Lambda$, along-exhaustion,
and $\mathbb{Z}^d$ concrete layers. 24 new theorems. The
susceptibility joint AnalyticAt case is proved as
\texttt{Finset.analyticAt\_fun\_sum} over the
\texttt{truncated2} summands. The along-exhaustion wrappers now live
in \texttt{AmbientLattice/SpecialCases/JointAnalyticity.lean}, and the
$\mathbb{Z}^d$ wrappers live in
\texttt{Concrete/LatticeGraphCorrelation/JointAnalyticity.lean}
(moved from the legacy modules in PR~\#1860).

\paragraph{Joint ContinuousAt + DifferentiableAt wrappers across
all 4 layers (\S17.5--\S17.6, PR~\#1640).} Adds joint pointwise
(in $(\beta, J, h)$) ContinuousAt and DifferentiableAt wrappers
for \texttt{correlation}, \texttt{magnetization}, and
\texttt{susceptibility} across abstract, $\Lambda$,
along-exhaustion, and $\mathbb{Z}^d$ concrete layers. 30 new
theorems, each derived from PR~\#1639 via
\texttt{.continuousAt} / \texttt{.differentiableAt}. The
along-exhaustion and $\mathbb{Z}^d$ wrapper layers now live in
\texttt{AmbientLattice/SpecialCases/JointRegularity.lean} and
\texttt{Concrete/LatticeGraphCorrelation/JointRegularity.lean}
(moved from the legacy modules in PR~\#1859). PR~\#2097 further carves
the 12 \texttt{*\_continuousAt\_joint} /
\texttt{*\_differentiableAt\_joint} concrete wrappers out into
\texttt{Concrete/LatticeGraphCorrelation/JointRegularityPointwise.lean}.
PR~\#2098 further carves the 6 along-exhaustion
\texttt{*\_continuous\_joint} / \texttt{*\_differentiable\_joint}
concrete wrappers out into
\texttt{Concrete/LatticeGraphCorrelation/JointRegularityAlongEx.lean}.

\paragraph{Joint Continuous + Differentiable wrappers across
all 4 layers
(\S17.5--\S17.6, PR~\#1639).} Adds joint (in $(\beta, J, h)$)
Continuous and Differentiable wrappers for
\texttt{magnetization}, \texttt{susceptibility}, and
\texttt{correlation} across abstract, $\Lambda$, along-exhaustion,
and $\mathbb{Z}^d$ concrete layers. 26 new theorems total:
abstract \texttt{magnetization\_/susceptibility\_\{continuous,
differentiable\}\_joint} (4); $\Lambda$ counterparts (4); along-ex
general-$G$ \texttt{correlation/magnetization/
susceptibility\_AlongExhaustion\_\{continuous,
differentiable\}\_joint(\_gen)} (6); $\mathbb{Z}^d$ concrete
$\Lambda$- and along-ex-layer
\texttt{*\_latticeGraph\_\{continuous, differentiable\}\_joint}
(12). The susceptibility joint case is proved as a finite sum of
\texttt{truncated2 = corr(\{i,j\}) - corr(\{i\})·corr(\{j\})},
each Continuous / Differentiable joint via
\texttt{correlation\_\{continuous, differentiable\}\_joint}. The
along-exhaustion and $\mathbb{Z}^d$ wrapper layers now live in
\texttt{AmbientLattice/SpecialCases/JointRegularity.lean} and
\texttt{Concrete/LatticeGraphCorrelation/JointRegularity.lean}
(moved from the legacy modules in PR~\#1859).

\paragraph{$\mathbb{Z}^d$ concrete
\texttt{partitionFunctionAlongExhaustion} /
\texttt{freeEnergyAlongExhaustion} ContinuousAt +
DifferentiableAt wrappers
(\S17.5--\S17.6, PR~\#1638).}
\texttt{Concrete/LatticeGraphCorrelation.lean}:
\texttt{partitionFunctionAlongExhaustion\_latticeGraph\_\{continuousAt,
differentiableAt\}\_\{beta\_h\_zero, J\_h\_zero,
beta\_general\_h, J\_general\_h, h, joint\}} (12 thms) and
\texttt{freeEnergyAlongExhaustion\_latticeGraph\_\{continuousAt,
differentiableAt\}\_\{beta, field, J, joint\}} (8 thms). 20 new
pointwise wrappers at $G := \texttt{IsingModel.latticeGraph}\,d$,
lifting the ambient versions from PR~\#1636.

\paragraph{$\mathbb{Z}^d$ concrete along-exhaustion
ContinuousAt + DifferentiableAt wrappers
(\S17.5--\S17.6, PR~\#1637).}
\texttt{Concrete/LatticeGraphCorrelation.lean}:
\texttt{correlationAlongExhaustion\_latticeGraph\_\{continuousAt,
differentiableAt\}\_J} (2),
\texttt{magnetizationAlongExhaustion\_latticeGraph\_\{continuousAt,
differentiableAt\}\_\{beta, field, J\}} (6),
\texttt{susceptibilityAlongExhaustion\_latticeGraph\_\{continuousAt,
differentiableAt\}\_\{beta\_general\_h, field, J\}} (6). 14 new
wrappers at $G := \texttt{IsingModel.latticeGraph}\,d$, lifting the
ambient general-$G$ versions from PR~\#1635.

\paragraph{\texttt{partitionFunctionAlongExhaustion} /
\texttt{freeEnergyAlongExhaustion} ContinuousAt + DifferentiableAt
wrappers (ambient)
(\S17.5--\S17.6, PR~\#1636).}
\texttt{SpecialCases.lean}:
\texttt{partitionFunctionAlongExhaustion\_\{continuousAt,
differentiableAt\}\_\{beta\_h\_zero, J\_h\_zero,
beta\_general\_h, J\_general\_h, h, joint\}} (12 thms) and
\texttt{freeEnergyAlongExhaustion\_\{continuousAt,
differentiableAt\}\_\{beta, field, J, joint\}} (8 thms). 20 new
pointwise wrappers, each derived from the existing whole-$\mathbb{R}$
wrapper via \texttt{.continuousAt} / \texttt{.differentiableAt}.

\paragraph{Along-exhaustion ContinuousAt + DifferentiableAt
wrappers (ambient general $G$)
(\S17.5--\S17.6, PR~\#1635).}
\texttt{correlationAlongExhaustion\_\{continuousAt,
differentiableAt\}\_J\_gen} (2 thms in
\texttt{JDerivative.lean}); \texttt{magnetizationAlongExhaustion\_
\{continuousAt, differentiableAt\}\_\{beta, field, J\}} (6 thms;
$\beta$ at general $h$) and
\texttt{susceptibilityAlongExhaustion\_\{continuousAt,
differentiableAt\}\_\{beta, field, J\}\_gen} (6 thms; in
\texttt{SpecialCases.lean}). 14 new pointwise wrappers, each
derived from the existing whole-$\mathbb{R}$ Continuous /
Differentiable via \texttt{(continuous \_).continuousAt} or
\texttt{(differentiable \_).differentiableAt}.

\paragraph{\texttt{susceptibilityAlongExhaustion} \texttt{HasDerivAt}
wrappers (ambient + $\mathbb{Z}^d$ concrete)
(\S17.5--\S17.6, PR~\#1634).} Adds \texttt{HasDerivAt} along-ex
wrappers for \texttt{susceptibilityAlongExhaustion} in all four
parameter directions. Ambient (general $G$, split across
\texttt{BetaDerivative.lean} / \texttt{JDerivative.lean} /
\texttt{FieldDerivative.lean}):
\texttt{susceptibilityAlongExhaustion\_hasDerivAt\_\{beta\_gen,
beta\_general\_h\_gen, J\_gen, field\_gen\}} (4 thms).
$\mathbb{Z}^d$ concrete:
\texttt{susceptibilityAlongExhaustion\_latticeGraph\_hasDerivAt\_*}
(4 thms). 8 new wrappers in existence form
$\exists d \in \mathbb{R},\,\mathrm{HasDerivAt}\,f\,d\,x$. The
ambient layer uses \texttt{\_gen} suffix to avoid potential clash
with $\mathbb{Z}^d$-specialized variants.

\paragraph{\texttt{magnetizationAlongExhaustion} $\beta$-direction
+ general-$G$ \texttt{susceptibilityAlongExhaustion}
Continuous/Differentiable wrappers
(\S17.5--\S17.6, PR~\#1633).} Ambient
(\texttt{SpecialCases.lean}):
\texttt{magnetizationAlongExhaustion\_\{continuous,
differentiable\}\_beta} (general $h$, 2) and
\texttt{susceptibilityAlongExhaustion\_\{continuous,
differentiable\}\_\{beta, field, J\}\_gen} (general $G$, general
$h$, 8). The \texttt{\_gen} suffix on susceptibility distinguishes
from the $\mathbb{Z}^d$-specialized variants in
\texttt{Concrete/LatticeGraphCorrelation/Inequalities.lean} which
share the unsuffixed names with \texttt{(IsingModel.latticeGraph
d)} baked in. $\mathbb{Z}^d$ concrete
(\texttt{LatticeGraphCorrelation.lean}):
\texttt{magnetizationAlongExhaustion\_latticeGraph\_\{continuous,
differentiable\}\_beta} (2). 12 new wrappers.

\paragraph{Per-parameter
\texttt{freeEnergyAlongExhaustion} Continuous/Differentiable
wrappers (along-ex + $\mathbb{Z}^d$ concrete)
(\S17.5--\S17.6, PR~\#1632).} Previously only joint Continuous /
Differentiable wrappers existed. This PR adds per-parameter
versions at both ambient and $\mathbb{Z}^d$ concrete layers,
mirroring the existing per-parameter coverage of
\texttt{partitionFunctionAlongExhaustion}. Ambient
(\texttt{SpecialCases.lean}):
\texttt{freeEnergyAlongExhaustion\_\{continuous, differentiable\}\_\{beta, field, J\}}
(6, with $\beta$ at general $h$).
$\mathbb{Z}^d$ concrete (\texttt{LatticeGraphCorrelation.lean}):
\texttt{freeEnergyAlongExhaustion\_latticeGraph\_\{continuous,
differentiable\}\_\{beta, field, J\}} (6). 12 new thin
instantiations of the corresponding
\texttt{freeEnergy\LeanLambda\_\{continuous, differentiable\}\_*} at
$\Lambda := \Lambda.\mathrm{volume}\,n$.

\paragraph{Along-exhaustion + $\mathbb{Z}^d$ concrete
\texttt{partitionFunctionAlongExhaustion} /
\texttt{freeEnergyAlongExhaustion} \texttt{HasDerivAt} wrappers
(\S17.5--\S17.6, PR~\#1631).} Specializes the $\Lambda$-layer
\texttt{hasDerivAt\_\{partitionFunction\LeanLambda, freeEnergy\LeanLambda\}\_*} family
(PRs~\#1623, \#1624) at $\Lambda := \Lambda.\mathrm{volume}\,n$.
Along-ex layer:
\texttt{partitionFunctionAlongExhaustion\_hasDerivAt\_\{beta, J,
field\}} (3),
\texttt{freeEnergyAlongExhaustion\_hasDerivAt\_\{beta\_general\_h,
J, field\}} (3) split across
\texttt{AmbientLattice/\{BetaDerivative, JDerivative,
FieldDerivative\}.lean}.
$\mathbb{Z}^d$ concrete:
\texttt{partitionFunctionAlongExhaustion\_latticeGraph\_hasDerivAt\_*}
(3) and
\texttt{freeEnergyAlongExhaustion\_latticeGraph\_hasDerivAt\_*}
(3) in
\texttt{Concrete/LatticeGraphCorrelation.lean}.
12 new theorems in existence form
$\exists c \in \mathbb{R},\,\mathrm{HasDerivAt}\,f\,c\,x$.

\paragraph{$\mathbb{Z}^d$ concrete (\texttt{latticeGraph} $d$)
along-exhaustion \texttt{HasDerivAt} wrappers
(\S17.5--\S17.6, PR~\#1630).} Direct instantiations at
$G := \texttt{IsingModel.latticeGraph}\,d$ of the
along-exhaustion \texttt{hasDerivAt} family
(\texttt{AmbientLattice/\{BetaDerivative, JDerivative,
FieldDerivative\}.lean}; PR~\#1628 + earlier):
\texttt{correlationAlongExhaustion\_latticeGraph\_hasDerivAt\_*}
(4) and
\texttt{magnetizationAlongExhaustion\_latticeGraph\_hasDerivAt\_*}
(4) for $\beta$ (h $=$ 0 and general h), $J$, and $h$ directions.
8 new theorems in existence form
$\exists c \in \mathbb{R},\,\mathrm{HasDerivAt}\,f\,c\,x$.

\paragraph{$\mathbb{Z}^d$ concrete (\texttt{latticeGraph} $d$)
$\Lambda$-layer \texttt{HasDerivAt} wrappers
(\S17.5--\S17.6, PR~\#1629).} Direct instantiations at
$G := \texttt{IsingModel.latticeGraph}\,d$ of the
$\Lambda$-layer \texttt{hasDerivAt} family
(PRs~\#1619, \#1623--\#1628):
\texttt{hasDerivAt\_correlation\LeanLambda\_latticeGraph\_*} (4),
\texttt{hasDerivAt\_freeEnergy\LeanLambda\_latticeGraph\_*} (3),
\texttt{hasDerivAt\_partitionFunction\LeanLambda\_latticeGraph\_*} (3),
\texttt{hasDerivAt\_boltzmannWeight\LeanLambda\_latticeGraph\_*} (3),
\texttt{magnetization\LeanLambda\_latticeGraph\_hasDerivAt\_*} (4),
\texttt{susceptibility\LeanLambda\_latticeGraph\_hasDerivAt\_*} (4).
21 new theorems in existence form
$\exists c \in \mathbb{R},\,\mathrm{HasDerivAt}\,f\,c\,x$, since
the explicit derivative formulas at the $\mathbb{Z}^d$ concrete
layer would be long and unwieldy; consumers needing the explicit
formula call \texttt{Ambient.hasDerivAt\_*\LeanLambda\_*} directly.

\paragraph{\texttt{correlationAlongExhaustion} /
\texttt{magnetizationAlongExhaustion} HasDerivAt in $J$ / $h$
/ $\beta$ along-exhaustion layer
(\S17.5--\S17.6, PR~\#1628).} Adds the missing
\texttt{hasDerivAt} directions for the along-exhaustion layer.
New files
\texttt{IsingModel/AmbientLattice/JDerivative.lean} and
\texttt{IsingModel/AmbientLattice/FieldDerivative.lean}
(mirroring \texttt{BetaDerivative.lean}):
\texttt{correlationAlongExhaustion\_hasDerivAt\_\{J, field\}},
\texttt{magnetizationAlongExhaustion\_hasDerivAt\_\{J, field\}}.
$\beta$-direction wrappers in
\texttt{BetaDerivative.lean}:
\texttt{magnetizationAlongExhaustion\_hasDerivAt\_beta} (at $h = 0$),
\texttt{magnetizationAlongExhaustion\_hasDerivAt\_beta\_general\_h\_gen}.
6 new theorems with existence-form derivative
$\exists d,\,\mathrm{HasDerivAt}\,f\,d\,x$; subset case lifts to
finite-volume \texttt{hasDerivAt\_correlation\_*}, non-subset is
constant zero. The corresponding \texttt{Continuous*} /
\texttt{Differentiable*} along-ex wrappers already exist in
\texttt{BetaDerivative.lean} under the \texttt{\_gen} suffix.

\paragraph{\texttt{magnetization} / \texttt{susceptibility}
HasDerivAt in $h$ + \texttt{magnetization} HasDerivAt in $\beta$
at $h = 0$, abstract + $\Lambda$ layers
(\S17.5--\S17.6, PR~\#1627).} Completes the §17.5--§17.6
single-site \texttt{hasDerivAt} suite at the $\Lambda$ layer
(the \texttt{\_J} and \texttt{\_beta\_general\_h} directions
were added in PR~\#1619; this PR adds the \texttt{\_field} and
\texttt{\_beta} ($h = 0$) directions). Abstract layer
(\texttt{PhaseTransition.lean}):
\texttt{magnetization\_hasDerivAt\_field}
($\partial_h \langle\sigma_i\rangle = \beta \cdot
(\langle\sigma_i \cdot M\rangle - \langle\sigma_i\rangle
\langle M\rangle)$),
\texttt{magnetization\_hasDerivAt\_beta} (at $h = 0$, edge-sum
form), \texttt{susceptibility\_hasDerivAt\_field} (sum-of-
\texttt{truncated2} h-derivatives via
\texttt{truncated2\_hasDerivAt\_field}). $\Lambda$-layer
(\texttt{AmbientLattice/Defs.lean}):
\texttt{magnetization\LeanLambda\_hasDerivAt\_field},
\texttt{magnetization\LeanLambda\_hasDerivAt\_beta} (at $h = 0$),
\texttt{susceptibility\LeanLambda\_hasDerivAt\_field}. 6 new theorems
(3 abstract + 3 thin $\Lambda$ wrappers); standard
thermodynamic-derivative identities (Glimm--Jaffe §17.5 / §17.6).

\paragraph{\texttt{hasDerivAt\_correlation\LeanLambda} in $\beta / J / h$
at the $\Lambda$-layer (\S17.5--\S17.6, \S4.6, PR~\#1626).}
Lifts four explicit-derivative correlation abstracts to the
$\Lambda$ layer with covariance-form derivatives:
\texttt{hasDerivAt\_correlation\LeanLambda\_beta} (at $h = 0$),
\texttt{\_beta\_general\_h},
\texttt{\_J}, \texttt{\_field}. 4 new thin wrappers.

\paragraph{\texttt{hasDerivAt\_boltzmannWeight\LeanLambda} in $\beta / J / h$
at the $\Lambda$-layer (\S17.5--\S17.6, \S4.6, PR~\#1625).}
Lifts the three per-configuration Boltzmann-weight derivative
abstracts to the $\Lambda$ layer:
\texttt{hasDerivAt\_boltzmannWeight\LeanLambda\_beta},
\texttt{\_J}, \texttt{\_field}. Building blocks for the
partition-function / free-energy variants (PR~\#1623, \#1624).
3 new thin wrappers.

\paragraph{\texttt{hasDerivAt\_partitionFunction\LeanLambda} in $\beta / J
/ h$ at the $\Lambda$-layer (\S17.5--\S17.6, \S4.6, PR~\#1624).}
Lifts three explicit-derivative abstracts to the $\Lambda$
layer: \texttt{hasDerivAt\_partitionFunction\LeanLambda\_beta}
($\partial_\beta Z = \sum_\sigma -H \cdot w$),
\texttt{\_J} ($\partial_J Z = \sum_\sigma \beta \cdot \sum_e
\mathrm{edgeSpin} \cdot w$), and \texttt{\_field}
($\partial_h Z = \sum_\sigma \beta \cdot M \cdot w$). 3 new thin
wrappers.

\paragraph{\texttt{hasDerivAt\_freeEnergy\LeanLambda} in $\beta / J / h$ at
the $\Lambda$-layer (\S17.5--\S17.6, \S4.6, PR~\#1623).} Lifts
three explicit-derivative thermodynamic-identity abstracts to
the $\Lambda$ layer:
\texttt{hasDerivAt\_freeEnergy\LeanLambda\_beta\_general\_h}
($\partial_\beta f = (|\!\uparrow\!\Lambda|)^{-1} \cdot \langle
-H\rangle$),
\texttt{hasDerivAt\_freeEnergy\LeanLambda\_J}
($\partial_J f = (|\!\uparrow\!\Lambda|)^{-1} \cdot \langle
\beta \cdot \sum_e \mathrm{edgeSpin}\rangle$), and
\texttt{hasDerivAt\_freeEnergy\LeanLambda\_field}
($\partial_h f = (|\!\uparrow\!\Lambda|)^{-1} \cdot \langle
\beta \cdot M\rangle$). 3 new thin wrappers.

\paragraph{\texttt{correlation\LeanLambda} / \texttt{susceptibility\LeanLambda}
ContinuousAt + DifferentiableAt at the $\Lambda$-layer
(\S17.5--\S17.6, PR~\#1622).} Point-wise versions
complementing the whole-$\mathbb{R}$ \texttt{Continuous} /
\texttt{Differentiable} wrappers (PR~\#1620, \#1616):
\texttt{correlation\LeanLambda\_continuousAt\_beta},
\texttt{\_continuousAt\_field},
\texttt{\_differentiableAt\_field},
\texttt{susceptibility\LeanLambda\_continuousAt\_beta},
\texttt{\_differentiableAt\_beta},
\texttt{\_continuousAt\_field},
\texttt{\_differentiableAt\_field}. 7 new thin wrappers.

\paragraph{\texttt{correlation\LeanLambda} Continuous + Differentiable in
$\beta / h / J$, $\Lambda$-layer (\S17.5--\S17.6, PR~\#1620).}
Lifts eight correlation regularity abstracts to the $\Lambda$
layer:
\texttt{correlation\LeanLambda\_\{continuous, differentiable\}\_beta} (at
$h = 0$),
\texttt{\_\{continuous, differentiable\}\_beta\_general\_h},
\texttt{\_\{continuous, differentiable\}\_field}, and
\texttt{\_\{continuous, differentiable\}\_J}. 8 new thin
direct-instantiation wrappers; along-exhaustion versions
already exist in
\texttt{Concrete/LatticeGraphCorrelation/Inequalities.lean} and
\texttt{AmbientLattice/BetaDerivative.lean}.

\paragraph{\texttt{susceptibility} / \texttt{magnetization}
HasDerivAt in $\beta/J$, $\Lambda$-layer (\S17.5--\S17.6,
PR~\#1619).} Lifts five HasDerivAt abstracts to the $\Lambda$
layer with explicit derivative formulas:
\texttt{susceptibility\LeanLambda\_hasDerivAt\_beta} (at $h = 0$),
\texttt{\_hasDerivAt\_beta\_general\_h},
\texttt{\_hasDerivAt\_J} (sum over $j : \uparrow\!\Lambda$ of
\texttt{deriv (truncated2 \ldots) \ldots}),
\texttt{magnetization\LeanLambda\_hasDerivAt\_J} (sum over edges of
\texttt{(inducedGraph G $\Lambda$).edgeFinset}), and
\texttt{magnetization\LeanLambda\_hasDerivAt\_beta\_general\_h}. 5 new
thin direct-instantiation wrappers; explicit derivative
formulas are inherently tied to the ambient site set, so the
along-exhaustion and $\mathbb{Z}^d$ versions are deferred until
a use case requires them.

\paragraph{\texttt{susceptibility} parameter-direction convergent
($\beta/h/J \to \infty$), all volume layers (\S5.3, PR~\#1618).}
Lifts three convergence abstracts to all 4 volume layers:
\texttt{susceptibility\_convergent\_beta} ($\beta \to \infty$ at
$J, h \ge 0$),
\texttt{\_convergent\_h} ($h \to \infty$ at $J \ge 0$,
$\beta > 0$), \texttt{\_convergent\_J} ($J \to \infty$ at
$h \ge 0$, $\beta > 0$). 12 new thin direct-instantiation
wrappers (3 abstracts $\times$ 4 layers). The along-ex
wrappers use `\_param` suffix to distinguish from the existing
volume-direction
\texttt{susceptibilityAlongExhaustion\_convergent}, and case-split
on $i \in \Lambda.\mathrm{volume}\ n$ (witness $L = 0$ when not
in volume).

\paragraph{\texttt{magnetization} parameter-direction convergent
($\beta/h/J \to \infty$), $\Lambda$ + along-ex + $\mathbb{Z}^d$
along-ex (\S5.3, PR~\#1617).} Lifts three convergence abstracts
to the missing volume layers ($\mathbb{Z}^d$ $\Lambda$-direct
versions already exist in
\texttt{correlation}$\Lambda \ldots \{i\}$ form):
\texttt{magnetization\_convergent\_beta} ($\beta \to \infty$ at
$J, h \ge 0$),
\texttt{\_convergent\_h} ($h \to \infty$ at $J \ge 0$, $\beta >
0$), \texttt{\_convergent\_J} ($J \to \infty$ at $h \ge 0$,
$\beta > 0$). 9 new thin direct-instantiation wrappers (3
abstracts $\times$ 3 missing layers). All return
$\exists L,\ \mathrm{Tendsto}\ (\ldots\ \mathrm{atTop})\ (\mathrm{nhds}\ L)$.
The along-ex wrappers case-split on
$i \in \Lambda.\mathrm{volume}\ n$: when $i \in$, the value
reduces to the $\Lambda$ wrap; otherwise it is constantly $0$
(witness $L = 0$).

\paragraph{\texttt{susceptibility} Continuous + Differentiable in
$h$ / $J$, $\Lambda$ and $\mathbb{Z}^d$ $\Lambda$ layers (\S17.5--\S17.6,
PR~\#1616).} Lifts four \texttt{susceptibility} regularity
abstracts to the $\Lambda$ and $\mathbb{Z}^d$ $\Lambda$ layers
(along-ex $\mathbb{Z}^d$ wraps already exist in
\texttt{Inequalities.lean} under the unprefixed name):
\texttt{\_continuous\_field}, \texttt{\_differentiable\_field}
(in $h$), \texttt{\_continuous\_J},
\texttt{\_differentiable\_J} (in $J$). 8 new thin
direct-instantiation wrappers (4 abstracts $\times$ 2 missing
layers).

\paragraph{\texttt{magnetization} Continuous + Differentiable in
$h$ / $J$, all volume layers (\S17.5--\S17.6, PR~\#1615).} Lifts
four \texttt{magnetization} regularity abstracts to the four
volume layers (\textit{Lambda}, along-ex, $\mathbb{Z}^d$ \textit{Lambda},
$\mathbb{Z}^d$ along-ex):
\texttt{\_continuous\_field}, \texttt{\_differentiable\_field}
(in $h$),
\texttt{\_continuous\_J}, \texttt{\_differentiable\_J} (in $J$).
16 new thin direct-instantiation wrappers (4 abstracts $\times$
4 layers). The along-ex wrappers case-split on
$i \in \Lambda.\mathrm{volume}\ n$: when $i \in$, the value
reduces to the \textit{Lambda} wrap; otherwise it is constantly $0$.

\paragraph{Mayer filter-connected $n=0/1/2$ + $\varepsilon^n$
expansion + \texttt{mayerPartialSum\_analyticOnNhd}, all volume
layers (\S18.5, PR~\#1613).} Lifts five Mayer / $\varepsilon(t)$
lemmas to the four volume layers (\textit{Lambda}, along-ex,
$\mathbb{Z}^d$ \textit{Lambda}, $\mathbb{Z}^d$ along-ex):
\texttt{mayerPartialSum\_analyticOnNhd}
(\texttt{AnalyticOnNhd ℝ \_ Set.univ}),
\texttt{vdPolymerFamilies\_sum\_minus\_one\_pow}
($\varepsilon(t)^k$ as multi-$\Gamma$ piFinset sum),
\texttt{mayerExpansionTerm\_filter\_connected\_zero}
(filter at $n=0$ is $\emptyset$),
\texttt{\_filter\_connected\_one} (filter at $n=1$ equals full
piFinset), and \texttt{\_two\_filter\_connected\_eq\_incompat}
(filter-connected $\Leftrightarrow$ pair-incompatible at $n=2$).
20 new thin direct-instantiation wrappers (5 abstracts $\times$
4 layers). The longest $\mathbb{Z}^d$ along-ex name
abbreviates \texttt{\_filter\_connected\_eq\_incompat} to
\texttt{\_filter\_conn\_eq\_incompat} to fit under the
100-character linter budget.

\paragraph{\S18.4--\S18.6 capstones ($Z$ formula +
\texttt{freeEnergy} decomposition + log $2$ / $N=1$ corollaries),
all volume layers (PR~\#1612).} Lifts six capstone formulas to the
four volume layers (\textit{Lambda}, along-ex, $\mathbb{Z}^d$
\textit{Lambda}, $\mathbb{Z}^d$ along-ex):
\texttt{partitionFunction\_high\_temp\_expansion\_h\_zero\_polymer\_family}
(\S18.4 main: $Z = 2^{|\iota|} \cdot \cosh^{|E|} \cdot
\sum_\Gamma \prod \tanh^{|P|}$),
\texttt{\_closed\_evenSubgraphs} (FV (3.45) form),
\texttt{freeEnergy\_eq\_polymerFreeEnergy} and
\texttt{\_ferromagnetic} (\S18.6 decomposition $f = \log 2 +
(|E|/|\iota|) \cdot \log \cosh + \mathrm{pFE} / |\iota|$),
\texttt{freeEnergy\_eq\_log\_two\_at\_betaJ\_zero} ($f = \log 2$
at $\beta J = 0$), and \texttt{mayerPartialSum\_one\_at\_one}
($\mathrm{mayerPartialSum}\,G\,1\,1 = |\mathrm{allPolymers}\,G|$).
24 new thin direct-instantiation wrappers (6 abstracts $\times$
4 layers). The hypothesis $0 < |\iota|$ at the Lambda / along-ex
layers is discharged via \texttt{Finset.Nonempty.fintype\_card\_coe\_pos}.
The $\mathbb{Z}^d$ along-ex ferromagnetic wrap uses the
abbreviated \texttt{\_ferro} suffix in
\texttt{freeEnergyAlongExhaustion\_latticeGraph\_eq\_polymerFreeEnergy\_ferro}
to fit under the 100-character linter budget.

\paragraph{\texttt{partitionFunction}
Continuous/Differentiable per-direction at general $h$ +
\texttt{freeEnergy} joint Continuous/Differentiable, completion
to all volume layers (\S18.6, PR~\#1611).} Lifts eight abstract
regularity theorems to the missing along-exhaustion +
$\mathbb{Z}^d$ layers (Lambda-layer wraps already exist via PRs
\#1541, \#1533):
\texttt{partitionFunction\_continuous\_beta\_general\_h},
\texttt{\_continuous\_J\_general\_h},
\texttt{\_differentiable\_beta\_general\_h},
\texttt{\_differentiable\_J\_general\_h},
\texttt{\_continuous\_h},
\texttt{\_differentiable\_h}, plus
\texttt{freeEnergy\_continuous\_joint} and
\texttt{\_differentiable\_joint}. 24 new thin
direct-instantiation wrappers (8 abstracts $\times$ 3 missing
layers).

\paragraph{\texttt{partitionFunction} joint
Continuous/Differentiable + general-$h$ analyticAt, all volume
layers (\S18.6, PR~\#1610).} Lifts five abstract
\texttt{partitionFunction} regularity theorems to the four
volume layers (\textit{Lambda}, along-ex, $\mathbb{Z}^d$ \textit{Lambda},
$\mathbb{Z}^d$ along-ex):
\texttt{\_continuous\_joint},
\texttt{\_differentiable\_joint},
\texttt{\_analyticAt\_beta\_general\_h},
\texttt{\_analyticAt\_J\_general\_h}, and
\texttt{\_analyticAt\_h}. 20 new thin direct-instantiation
wrappers (5 abstracts $\times$ 4 layers). Complements the
existing joint analyticity (PR~\#1535, joint analyticAt) and
$h$-zero regularity (PR~\#1608) wrappers.

\paragraph{\texttt{freeEnergy} per-direction analyticAt +
analyticOnNhd, all volume layers (\S18.6, PR~\#1609).} Lifts
ten abstract \texttt{freeEnergy} per-direction analyticity
theorems to the four volume layers (\textit{Lambda}, along-ex,
$\mathbb{Z}^d$ \textit{Lambda}, $\mathbb{Z}^d$ along-ex):
\texttt{\_analyticAt\_beta\_h\_zero},
\texttt{\_analyticAt\_J\_h\_zero},
\texttt{\_analyticOnNhd\_beta\_h\_zero},
\texttt{\_analyticOnNhd\_J\_h\_zero},
\texttt{\_analyticAt\_beta\_general\_h},
\texttt{\_analyticAt\_J\_general\_h},
\texttt{\_analyticAt\_h},
\texttt{\_analyticOnNhd\_beta\_general\_h},
\texttt{\_analyticOnNhd\_J\_general\_h}, and
\texttt{\_analyticOnNhd\_h}. 40 new thin
direct-instantiation wrappers (10 abstracts $\times$ 4 layers).
Complements the existing joint-direction
\texttt{freeEnergy\LeanLambda\_analyticAt\_joint} (PR~\#1535) at the same
volume layers.

\paragraph{\texttt{partitionFunction} \texttt{h\_zero} regularity,
all volume layers (\S18.6, PR~\#1608).} Lifts eight abstract
\texttt{partitionFunction\_*\_h\_zero} regularity theorems to the
four volume layers (\textit{Lambda}, along-ex, $\mathbb{Z}^d$
\textit{Lambda}, $\mathbb{Z}^d$ along-ex):
\texttt{\_continuous\_beta\_h\_zero},
\texttt{\_continuous\_J\_h\_zero},
\texttt{\_differentiable\_beta\_h\_zero},
\texttt{\_differentiable\_J\_h\_zero},
\texttt{\_analyticAt\_beta\_h\_zero},
\texttt{\_analyticAt\_J\_h\_zero},
\texttt{\_analyticOnNhd\_beta\_h\_zero}, and
\texttt{\_analyticOnNhd\_J\_h\_zero}. 32 new thin
direct-instantiation wrappers (8 abstracts $\times$ 4 layers).
These provide the polymer-expansion-derived h-zero regularity of
$Z(J, 0, \beta)$ in $\beta / J$ at every volume layer, alongside
the existing general-$h$ wrappers. The along-exhaustion and concrete
$\mathbb{Z}^d$ wrappers from this family now live in the matching
\texttt{PartitionFunctionRegularity.lean} child modules under
\path{AmbientLattice/SpecialCases} and
\path{Concrete/LatticeGraphCorrelation}; the $\Lambda$-layer wrappers
live in
\path{IsingModel/AmbientLattice/AnalyticityLambdaSection186.lean}
(narrowed out of \texttt{Analyticity.lean} in PR~\#1992).

The along-exhaustion and concrete $\mathbb{Z}^d$ joint/general-$h$
\texttt{partitionFunction} wrappers now live in the matching
\texttt{PartitionFunctionGeneralAnalyticity.lean} child modules; the
$\Lambda$-layer wrappers live in
\path{IsingModel/AmbientLattice/AnalyticityLambdaSection186.lean}
(narrowed out of \texttt{Analyticity.lean} in PR~\#1992).

The along-exhaustion and concrete $\mathbb{Z}^d$
\texttt{partitionFunction}/\texttt{freeEnergy} Continuous and Differentiable
wrappers now live in matching \texttt{PartitionFreeEnergyRegularity.lean}
child modules; the $\Lambda$-layer wrappers remain in
\path{IsingModel/AmbientLattice/Analyticity.lean}.

The along-exhaustion and concrete $\mathbb{Z}^d$ high-temperature
\texttt{partitionFunction}/\texttt{freeEnergy} capstone wrappers now live in
matching \texttt{HighTemperatureCapstones.lean} child modules; the
$\Lambda$-layer wrappers live in
\path{IsingModel/AmbientLattice/AnalyticityLambdaCapstones.lean}
(narrowed out of \texttt{Analyticity.lean} in PR~\#1993).

The along-exhaustion and concrete $\mathbb{Z}^d$ Mayer filter-connected and
epsilon-power wrappers now live in matching
\texttt{MayerFilterConnected.lean} child modules; the $\Lambda$-layer wrappers
live in
\path{IsingModel/AmbientLattice/AnalyticityLambdaCapstones.lean}
(narrowed out of \texttt{Analyticity.lean} in PR~\#1993).

The along-exhaustion and concrete $\mathbb{Z}^d$ magnetization regularity
wrappers now live in matching \texttt{MagnetizationRegularity.lean} child
modules; the $\Lambda$-layer wrappers remain in
\path{IsingModel/AmbientLattice/Defs.lean}.

The along-exhaustion and concrete $\mathbb{Z}^d$ magnetization convergence
wrappers now live in matching \texttt{MagnetizationConvergence.lean} child
modules; the general $\Lambda$-layer wrappers remain in
\path{IsingModel/AmbientLattice/Defs.lean}, and the concrete $\Lambda$-direct
wrappers remain in the legacy concrete module.

The along-exhaustion and concrete $\mathbb{Z}^d$ susceptibility convergence
wrappers now live in matching \texttt{SusceptibilityConvergence.lean} child
modules; the general $\Lambda$-layer wrappers remain in
\path{IsingModel/AmbientLattice/Defs.lean}, and the concrete $\Lambda$-direct
wrappers remain in the legacy concrete module.

{\raggedright
The along-exhaustion partition-function trivial-slice closed forms now live in
\path{IsingModel/AmbientLattice/SpecialCases/PartitionFunctionClosedForms.lean},
with the concrete $\mathbb{Z}^d$ $\Lambda$- and along-exhaustion
wrappers in
\path{IsingModel/Concrete/LatticeGraphCorrelation/PartitionFunctionClosedForms.lean};
PR~\#2095 carves the six cubicExhaustion-$\Lambda$ closed-form wrappers
out into
\path{IsingModel/Concrete/LatticeGraphCorrelation/PartitionFunctionClosedFormsCubicLambda.lean}.
PR~\#2096 carves the six cubicExhaustion-alongExhaustion closed-form
wrappers out into
\path{IsingModel/Concrete/LatticeGraphCorrelation/PartitionFunctionClosedFormsCubicAlongEx.lean}.
The lower $\Lambda$-level abstract closed forms remain in
\path{IsingModel/AmbientLattice/Defs.lean} and the corresponding log
$\Lambda$ forms remain in \path{IsingModel/AmbientLatticeSum.lean}.
\par}

{\raggedright
The along-exhaustion partition-function h-symmetry and absolute-field wrappers
now live in
\path{IsingModel/AmbientLattice/SpecialCases/PartitionFunctionSymmetry.lean}.
The concrete $\mathbb{Z}^d$ $\Lambda$-, cubic-, and along-exhaustion wrappers,
including the corresponding log forms, now live in
\path{IsingModel/Concrete/LatticeGraphCorrelation/PartitionFunctionSymmetry.lean}.
\par}

{\raggedright
The concrete $\mathbb{Z}^d$ free-energy closed-form, lower-bound,
monotonicity, h-symmetry, and absolute-field wrappers
now live in
\path{IsingModel/Concrete/LatticeGraphCorrelation/FreeEnergySpecialCases.lean}.
The bottom-graph ($\bot \le \mathrm{latticeGraph}\,d$) comparison
wrappers and the matching
\texttt{*\_latticeGraph\_monotone\_ambient\_subgraph} family now live in
\path{IsingModel/Concrete/LatticeGraphCorrelation/TwoPointAmbientBot.lean}
(narrowed out of \path{TwoPoint.lean} in PR~\#2009); PR~\#2092 further
carves the seven $\Lambda$ + along-exhaustion
\texttt{\{freeEnergy,partitionFunction\}} / \texttt{correlation} /
\texttt{log\_partitionFunction}$\Lambda$ \texttt{*\_bot\_le\_latticeGraph}
wrappers out into
\path{IsingModel/Concrete/LatticeGraphCorrelation/TwoPointAmbientBotLambda.lean}.
PR~\#2093 further
carves the seven $\Lambda$ + along-exhaustion + correlationInfinite
\texttt{\{freeEnergy,partitionFunction\}} / \texttt{correlation}
\texttt{*\_latticeGraph\_monotone\_ambient\_subgraph}
wrappers out into
\path{IsingModel/Concrete/LatticeGraphCorrelation/TwoPointAmbientBotMonotone.lean}.
PR~\#2094 further
carves the five \texttt{magnetization*} / \texttt{spontaneous*}
\texttt{*\_bot\_le\_latticeGraph} wrappers out into
\path{IsingModel/Concrete/LatticeGraphCorrelation/TwoPointAmbientBotMagnetization.lean}. The 7
$\mathbb{Z}^d$ $\mathrm{correlationInfinite\_latticeGraph\_*}$
wrappers (\texttt{\_le\_one}, \texttt{\_nonneg},
\texttt{\_indep\_exhaustion},
\texttt{\_cubicExhaustion\_monotone\_h}, \texttt{\_beta}, \texttt{\_J},
\texttt{\_gks\_second}) now live in
\path{IsingModel/Concrete/LatticeGraphCorrelation/TwoPointCorrelationInfinite.lean}
(narrowed out of \path{TwoPoint.lean} in PR~\#2025). The 5
$\mathbb{Z}^d$ magnetization monotonicity wrappers
(\texttt{spontaneousMagnetization\_latticeGraph\_cubicExhaustion\_monotone\_\{J,beta\}},
\texttt{magnetizationInfinite\_latticeGraph\_cubicExhaustion\_monotone\_\{J,h,beta\}})
now live in
\path{IsingModel/Concrete/LatticeGraphCorrelation/TwoPointMagnetizationMonotone.lean}
(narrowed out of \path{TwoPoint.lean} in PR~\#2026). The 6
$\mathbb{Z}^d$ \texttt{truncated3Infinite\_latticeGraph\_swap\_\{ij,jk,ik\}}
and \texttt{truncated4Infinite\_latticeGraph\_swap\_\{ij,jk,kl\}}
symmetry wrappers now live in
\path{IsingModel/Concrete/LatticeGraphCorrelation/TwoPointTruncatedSwaps.lean}
(narrowed out of \path{TwoPoint.lean} in PR~\#2027). The 5
$\mathbb{Z}^d$ \texttt{truncated3/4TwoPoint} trivial-slice wrappers
(\texttt{truncated3TwoPoint\_h\_zero\_of\_distinct},
\texttt{\_J\_zero\_of\_distinct}, \texttt{\_beta\_zero},
\texttt{truncated4TwoPoint\_J\_zero\_of\_distinct}, \texttt{\_beta\_zero})
now live in
\path{IsingModel/Concrete/LatticeGraphCorrelation/TwoPointTruncatedTrivialSlices.lean}
(narrowed out of \path{TwoPoint.lean} in PR~\#2028). The 3
$\mathbb{Z}^d$ \texttt{twoPointFunction} trivial-slice wrappers
(\texttt{twoPointFunction\_zero\_params}, \texttt{\_beta\_zero},
\texttt{\_J\_zero\_of\_ne\_zero}) now live in
\path{IsingModel/Concrete/LatticeGraphCorrelation/TwoPointFunctionTrivialSlices.lean}
(narrowed out of \path{TwoPoint.lean} in PR~\#2029).
The corresponding ambient along-exhaustion closed forms and h-symmetry wrappers
remain in \path{IsingModel/AmbientLattice/SpecialCases/FreeEnergy.lean}.
\par}

{\raggedright
The concrete $\mathbb{Z}^d$ finite-volume graph, spin-algebra, bottom-graph,
and Hamiltonian symmetry wrappers now live in
\path{IsingModel/Concrete/LatticeGraphCorrelation/FiniteVolumeBasics.lean};
PR~\#2103 carves the five Walsh-basis / spin-config wrappers out into
\path{IsingModel/Concrete/LatticeGraphCorrelation/FiniteVolumeBasicsWalsh.lean};
PR~\#2104 carves the seven spinProduct / edgeSpin algebra wrappers
out into
\path{IsingModel/Concrete/LatticeGraphCorrelation/FiniteVolumeBasicsSpin.lean};
PR~\#2105 carves the five Hamiltonian flip / negative-field
wrappers out into
\path{IsingModel/Concrete/LatticeGraphCorrelation/FiniteVolumeBasicsHamiltonian.lean}.
\par}

{\raggedright
The concrete $\mathbb{Z}^d$ finite-volume HNC, Gibbs-expectation, and
correlation monotonicity and convergence wrappers now live in
\path{IsingModel/Concrete/LatticeGraphCorrelation/FiniteVolumeCorrelationMonotonicity.lean}.
\par}

{\raggedright
The concrete $\mathbb{Z}^d$ finite-volume extension graph comparison and
correlation equality wrappers now live in
\path{IsingModel/Concrete/LatticeGraphCorrelation/FiniteVolumeExtensions.lean}.
\par}

{\raggedright
The concrete $\mathbb{Z}^d$ Lambda-layer correlation and magnetization
convergence and monotonicity wrappers now live in
\path{IsingModel/Concrete/LatticeGraphCorrelation/LambdaCorrelationMonotonicity.lean}.
\par}

{\raggedright
The concrete $\mathbb{Z}^d$ along-exhaustion correlation boundedness,
eventuality, and infinite-volume limit wrappers now live in
\path{IsingModel/Concrete/LatticeGraphCorrelation/CorrelationExhaustionLimits.lean}.
\par}

{\raggedright
The concrete $\mathbb{Z}^d$ finite-volume, along-exhaustion, and
infinite-volume translation-invariance wrappers now live in
\path{IsingModel/Concrete/LatticeGraphCorrelation/Translation.lean}.
\par}

{\raggedright
The concrete $\mathbb{Z}^d$ partition/free-energy disjoint-union monotonicity
and superadditivity wrappers now live in
\path{IsingModel/Concrete/LatticeGraphCorrelation/PartitionFreeEnergySuperadditivity.lean}.
\par}

{\raggedright
The concrete $\mathbb{Z}^d$ partition-function along-exhaustion volume /
parameter monotonicity, positivity, divergence, and infinite-volume free-energy
positivity wrappers now live in
\path{IsingModel/Concrete/LatticeGraphCorrelation/PartitionExhaustionBounds.lean}.
\par}

{\raggedright
The concrete $\mathbb{Z}^d$ free-energy along-exhaustion,
log partition-function along-exhaustion, and Lambda-layer partition-function
parameter monotonicity wrappers now live in
\path{IsingModel/Concrete/LatticeGraphCorrelation/PartitionFreeEnergyMonotonicity.lean}.
\par}

{\raggedright
The concrete $\mathbb{Z}^d$ partition/free-energy lower-bound,
nonnegativity, infinite-volume bridge, and uniform-bound wrappers now live in
\path{IsingModel/Concrete/LatticeGraphCorrelation/PartitionFreeEnergyBounds.lean}.
\par}

{\raggedright
The concrete $\mathbb{Z}^d$ finite-volume correlation and
truncated-correlation inequality, odd-vanishing, and trivial-slice wrappers now
live in
\path{IsingModel/Concrete/LatticeGraphCorrelation/FiniteVolumeCorrelationInequalities.lean}.
\par}

{\raggedright
The concrete $\mathbb{Z}^d$ finite-volume Boltzmann-weight, Hamiltonian,
partition-function, and free-energy bound wrappers now live in
\path{IsingModel/Concrete/LatticeGraphCorrelation/FiniteVolumeEnergyBounds.lean}.
\par}

{\raggedright
The concrete $\mathbb{Z}^d$ finite-volume Hamiltonian closed forms, direct
energy and partition-function bounds, and spinProduct helper wrappers now live
in
\path{IsingModel/Concrete/LatticeGraphCorrelation/EnergyClosedForms.lean}.
\par}

{\raggedright
The concrete $\mathbb{Z}^d$ direct finite-volume partition-function
monotonicity, trivial-slice, and negative-field symmetry wrappers now live in
\path{IsingModel/Concrete/LatticeGraphCorrelation/FiniteVolumePartition.lean}.
\par}

{\raggedright
The concrete $\mathbb{Z}^d$ direct finite-volume partition-function
absolute-field, positivity, and ferromagnetic lower-bound wrappers now live in
\path{IsingModel/Concrete/LatticeGraphCorrelation/FiniteVolumePartitionBounds.lean}.
\par}

{\raggedright
The concrete $\mathbb{Z}^d$ correlation, magnetization, and
truncated-correlation h-symmetry and absolute-field wrappers now live in
\path{IsingModel/Concrete/LatticeGraphCorrelation/CorrelationSymmetry.lean}.
\par}

{\raggedright
The concrete $\mathbb{Z}^d$ Peierls along-exhaustion wrappers
\texttt{prop\_5\_4\_2\_along\_exhaustion\_latticeGraph} and
\texttt{prop\_5\_4\_2\_limsup\_le\_latticeGraph} now live in
\path{IsingModel/Concrete/LatticeGraphCorrelation/Peierls.lean}, while the
abstract exhaustion statements remain in \path{IsingModel/PeierlsInfinite.lean}.
\par}

The analogous \texttt{freeEnergy} per-direction analyticity wrappers
now live in matching \texttt{FreeEnergyAnalyticity.lean} child modules;
their $\Lambda$-layer wrappers live in
\path{IsingModel/AmbientLattice/AnalyticityLambdaSection186.lean}
(narrowed out of \texttt{Analyticity.lean} in PR~\#1992).

\paragraph{vdSum sandwich/monotone + $\varepsilon$ bound +
\texttt{polymerFreeEnergy(tanh)} bound + $\log 2$, all volume
layers (\S18.5, PR~\#1607).} Lifts six abstract bound/regularity
lemmas to the four volume layers (\textit{Lambda}, along-ex,
$\mathbb{Z}^d$ \textit{Lambda}, $\mathbb{Z}^d$ along-ex):
\texttt{vdPolymerFamilies\_sum\_sandwich\_of\_nonneg}
($1 \le \mathrm{vdSum} \le (1+t)^{|E|}$),
\texttt{\_monotoneOn\_Ici\_zero},
\texttt{\_minus\_one\_le\_of\_nonneg}
($\varepsilon(t) \le (1+t)^{|E|} - 1$),
\texttt{polymerFreeEnergy\_tanh\_le\_eps\_of\_betaJ\_nonneg},
\texttt{\_tanh\_le\_pow\_sub\_one\_of\_betaJ\_nonneg}, and
\texttt{\_tanh\_lt\_log\_two\_of\_pow\_lt\_two}. 24 new thin
direct-instantiation wrappers. The longest ℤ$^d$ along-ex name
abbreviates the suffix \texttt{\_of\_betaJ\_nonneg} to
\texttt{\_betaJ} to fit under the 100-character linter budget.
The along-exhaustion and concrete $\mathbb{Z}^d$ wrappers now live in the
matching \texttt{PolymerFreeEnergyHighTemperatureBounds.lean} child modules
under \texttt{AmbientLattice/SpecialCases} and
\texttt{Concrete/LatticeGraphCorrelation}; the $\Lambda$-layer wrappers live
in \texttt{IsingModel/AmbientLattice/AnalyticityLambdaPfeSharpening.lean}
(narrowed out of \texttt{Analyticity.lean} in PR~\#1991).

\paragraph{$\varepsilon(t)$ nonneg + non-tanh
\texttt{polymerFreeEnergy} sharpening, all volume layers (\S18.5,
PR~\#1606).} Lifts six abstract sharpening lemmas to the four
volume layers (\textit{Lambda}, along-ex, $\mathbb{Z}^d$ \textit{Lambda},
$\mathbb{Z}^d$ along-ex):
\texttt{vdPolymerFamilies\_sum\_minus\_one\_nonneg\_of\_nonneg}
($0 \le \varepsilon(t)$ for $t \ge 0$),
\texttt{\_minus\_one\_pow\_at\_zero}
($\varepsilon(0)^n = 0$ for $n \ge 1$),
\texttt{polymerFreeEnergy\_eq\_zero\_iff\_eps\_eq\_zero},
\texttt{\_pos\_iff\_eps\_pos},
\texttt{\_lt\_eps\_iff\_eps\_pos}, and
\texttt{\_lt\_pow\_sub\_one\_of\_eps\_pos}. 24 new thin
direct-instantiation wrappers.
The along-exhaustion and concrete $\mathbb{Z}^d$ wrappers now live in the
matching \texttt{PolymerFreeEnergyEpsilonSharpening.lean} child modules under
\texttt{AmbientLattice/SpecialCases} and
\texttt{Concrete/LatticeGraphCorrelation} (moved from the legacy modules in
PR~\#1882); the $\Lambda$-layer wrappers live in
\texttt{IsingModel/AmbientLattice/AnalyticityLambdaPfeSharpening.lean}
(narrowed out of \texttt{Analyticity.lean} in PR~\#1991).

\paragraph{\texttt{polymerFreeEnergy} tanh sharpening + $\beta/J$
strict-mono under polymers $\ne \emptyset$, all volume layers
(\S18.5, PR~\#1605).} Lifts nine abstract \texttt{polymerFreeEnergy}
theorems to the four volume layers (\textit{Lambda}, along-ex,
$\mathbb{Z}^d$ \textit{Lambda}, $\mathbb{Z}^d$ along-ex):
five tanh sharpening lemmas under $0 \le \beta J$
(\texttt{\_tanh\_lt\_eps\_iff\_eps\_pos},
\texttt{\_tanh\_eq\_zero\_iff\_eps\_eq\_zero},
\texttt{\_tanh\_pos\_iff\_eps\_pos},
\texttt{\_tanh\_lt\_eps\_of\_eps\_pos},
\texttt{\_tanh\_lt\_pow\_sub\_one\_of\_eps\_pos}) plus four
$\beta/J$-strict-mono lemmas under polymers nonempty
(\texttt{\_tanh\_lt\_of\_lt\_in\_beta},
\texttt{\_tanh\_lt\_of\_lt\_in\_J},
\texttt{\_tanh\_strictMonoOn\_beta},
\texttt{\_tanh\_strictMonoOn\_J}). 36 new thin
direct-instantiation wrappers. Several long
$\mathbb{Z}^d$ along-ex names abbreviate
\texttt{\_of\_polymers\_nonempty} to \texttt{\_polymers\_nonempty}
(or further to \texttt{\_polymers}) to fit under the
100-character linter budget.
The along-exhaustion and concrete $\mathbb{Z}^d$ wrappers now live in the
matching \texttt{PolymerFreeEnergyTanhSharpening.lean} child modules under
\texttt{AmbientLattice/SpecialCases} and
\texttt{Concrete/LatticeGraphCorrelation} (moved from the legacy modules in
PR~\#1881); the $\Lambda$-layer wrappers live in
\texttt{IsingModel/AmbientLattice/AnalyticityLambdaPfeSharpening.lean}
(narrowed out of \texttt{Analyticity.lean} in PR~\#1991).

\paragraph{\texttt{polymerFreeEnergy} / \texttt{vdSum} tanh
ferromagnetic iff family, all volume layers (\S18.5, PR~\#1604).}
Lifts nine ferromagnetic abstract iff / inequality lemmas to the
four volume layers (\textit{Lambda}, along-ex, $\mathbb{Z}^d$
\textit{Lambda}, $\mathbb{Z}^d$ along-ex). All require
$0 \le \beta$ and $0 \le J$:
\texttt{\_tanh\_lt\_eps\_iff\_eps\_pos\_ferromagnetic},
\texttt{\_tanh\_eq\_zero\_iff\_eps\_eq\_zero\_ferromagnetic},
\texttt{\_tanh\_pos\_iff\_eps\_pos\_ferromagnetic},
\texttt{\_tanh\_pos\_iff\_ferromagnetic},
\texttt{\_tanh\_eq\_zero\_iff\_ferromagnetic},
\texttt{vdPolymerFamilies\_sum\_tanh\_gt\_one\_iff\_ferromagnetic},
\texttt{\_eq\_one\_iff\_ferromagnetic},
\texttt{\_tanh\_lt\_pow\_sub\_one\_of\_eps\_pos\_ferromagnetic},
\texttt{\_tanh\_lt\_eps\_of\_eps\_pos\_ferromagnetic}.
36 new thin direct-instantiation wrappers; all use the
abbreviated \texttt{\_ferro} suffix in place of
\texttt{\_ferromagnetic} to fit under the 100-character linter
budget. The along-exhaustion and concrete $\mathbb{Z}^d$ wrappers
now live in
\texttt{AmbientLattice/SpecialCases/MayerTanhFerromagneticIff.lean}
and
\texttt{Concrete/LatticeGraphCorrelation/MayerTanhFerromagneticIff.lean}
(moved from the legacy modules in PR~\#1880). The Lambda-layer
wrappers live in
\texttt{IsingModel/AmbientLattice/AnalyticityLambdaTanhFerroIff.lean}
(narrowed out of \texttt{Analyticity.lean} in PR~\#1990).

\paragraph{Strict-mono / strict-pos under polymers $\ne
\emptyset$, all volume layers (\S18.5, PR~\#1603).} Lifts ten
abstract strict-monotone / strict-positive lemmas to the four
volume layers (\textit{Lambda}, along-ex, $\mathbb{Z}^d$
\textit{Lambda}, $\mathbb{Z}^d$ along-ex). All require the
hypothesis $\mathrm{allPolymers}\ G \ne \emptyset$:
\texttt{vdPolymerFamilies\_sum\_lt\_of\_lt\_of\_polymers\_nonempty},
\texttt{\_strictMonoOn\_of\_polymers\_nonempty} (over
$[0,\infty)$),
\texttt{polymerFreeEnergy\_pos\_of\_t\_pos\_of\_polymers\_nonempty},
\texttt{vdPolymerFamilies\_sum\_gt\_one\_of\_t\_pos\_of\_polymers\_nonempty},
\texttt{\_minus\_one\_pos\_of\_t\_pos\_of\_polymers\_nonempty},
\texttt{polymerFreeEnergy\_tanh\_pos\_of\_tanh\_pos\_of\_polymers\_nonempty},
\texttt{vdPolymerFamilies\_sum\_tanh\_gt\_one\_of\_tanh\_pos\_of\_polymers\_nonempty},
\texttt{\_minus\_one\_tanh\_pos\_of\_tanh\_pos\_of\_polymers\_nonempty},
\texttt{polymerFreeEnergy\_strictMonoOn\_Ioi\_zero\_of\_polymers\_nonempty},
\texttt{vdPolymerFamilies\_sum\_strictMonoOn\_Ioi\_zero\_of\_polymers\_nonempty}.
40 new thin direct-instantiation wrappers. The longest name
(ℤ$^d$ along-ex \texttt{\_minus\_one\_tanh\_pos}) drops the
second \texttt{\_of\_tanh\_pos\_} from its name to fit under the
100-character linter budget; the hypothesis is preserved in the
\texttt{h\_tanh\_pos} parameter. The along-exhaustion and concrete
$\mathbb{Z}^d$ wrappers now live in
\texttt{AmbientLattice/SpecialCases/MayerStrictPositivity.lean}
and
\texttt{Concrete/LatticeGraphCorrelation/MayerStrictPositivity.lean}
(moved from the legacy modules in PR~\#1879). The Lambda-layer
wrappers live in
\texttt{IsingModel/AmbientLattice/AnalyticityLambdaEpsilonIff.lean}
(narrowed out of \texttt{Analyticity.lean} in PR~\#1989).

\paragraph{$\varepsilon(t)$ / \texttt{polymerFreeEnergy}
positivity / zero iff family, all volume layers (\S18.5,
PR~\#1602).} Lifts six abstract iff characterisations to the
four volume layers (\textit{Lambda}, along-ex, $\mathbb{Z}^d$
\textit{Lambda}, $\mathbb{Z}^d$ along-ex):
\texttt{vdPolymerFamilies\_sum\_minus\_one\_pos\_iff}
($0 < \varepsilon(t) \Leftrightarrow 0 < t \wedge
\mathrm{allPolymers}\ G \ne \emptyset$),
\texttt{\_eq\_zero\_iff}, \texttt{\_tanh\_pos\_iff},
\texttt{\_tanh\_eq\_zero\_iff},
\texttt{polymerFreeEnergy\_tanh\_pos\_iff}, and
\texttt{polymerFreeEnergy\_tanh\_eq\_zero\_iff}. 24 new thin
direct-instantiation wrappers. The along-exhaustion and concrete
$\mathbb{Z}^d$ wrappers now live in
\texttt{AmbientLattice/SpecialCases/MayerEpsilonPositivity.lean}
and
\texttt{Concrete/LatticeGraphCorrelation/MayerEpsilonPositivity.lean}.
The Lambda-layer wrappers live in
\texttt{IsingModel/AmbientLattice/AnalyticityLambdaEpsilonIff.lean}
(narrowed out of \texttt{Analyticity.lean} in PR~\#1989).

\paragraph{$\varepsilon(t)$ infrastructure + Mayer-term sign at
$n = 1, 2$ + \texttt{allPolymers} empty on edgeless, all volume
layers (\S18.5, PR~\#1601).} Lifts seven abstract Mayer-expansion
theorems to the four volume layers (\textit{Lambda}, along-ex,
$\mathbb{Z}^d$ \textit{Lambda}, $\mathbb{Z}^d$ along-ex):
\texttt{mayerExpansionTerm\_one\_nonneg\_of\_nonneg},
\texttt{\_two\_nonpos\_of\_nonneg},
\texttt{vdPolymerFamilies\_sum\_minus\_one\_at\_zero}
($\varepsilon(0) = 0$),
\texttt{\_continuous}, \texttt{\_analyticAt},
\texttt{\_lt\_one\_eventually}, and
\texttt{allPolymers\_eq\_empty\_of\_edgeFinset\_empty}.
28 new thin direct-instantiation wrappers. The along-exhaustion and
concrete \texttt{\_analyticAt} wrappers for
\texttt{vdPolymerFamilies\_sum\_minus\_one} now live in
\texttt{AmbientLattice/SpecialCases/VdPolymerFamiliesAnalyticity.lean}
and
\texttt{Concrete/LatticeGraphCorrelation/VdPolymerFamiliesAnalyticity.lean};
the along-exhaustion and concrete non-analytic wrappers from this family
now live in
\texttt{AmbientLattice/SpecialCases/MayerEpsilonInfrastructure.lean}
and
\texttt{Concrete/LatticeGraphCorrelation/MayerEpsilonInfrastructure.lean}
(moved from the legacy modules in PR~\#1877);
the $\Lambda$-layer wrappers live in
\texttt{IsingModel/AmbientLattice/AnalyticityLambdaMayerRecurrenceEpsilon.lean}
(narrowed out of \texttt{Analyticity.lean} in PR~\#1988); the
\texttt{vdPolymerFamilies\_sum\_minus\_one\_analyticAt} Lambda-layer
wrapper specifically was moved from the legacy modules in PR~\#1863
before its current narrow location.

\paragraph{Mayer recurrence + \texttt{polymerFreeEnergy} hasSum +
$\varepsilon(t) \to 0$ tendsto, all volume layers (\S18.5,
PR~\#1600).} Lifts five abstract Mayer-expansion theorems to the
four volume layers (\textit{Lambda}, along-ex, $\mathbb{Z}^d$
\textit{Lambda}, $\mathbb{Z}^d$ along-ex):
\texttt{mayerPartialSum\_succ} (recurrence in $N$),
\texttt{mayerExpansionTerm\_eq\_mayerPartialSum\_diff},
\texttt{polymerFreeEnergy\_hasSum\_via\_log} (under $|\varepsilon|
< 1$),
\texttt{polymerFreeEnergy\_hasSum\_via\_log\_eventually} (true
in some neighbourhood of $t = 0$), and
\texttt{vdPolymerFamilies\_sum\_minus\_one\_tendsto\_zero}
($\varepsilon(t) \to 0$ as $t \to 0$). 20 new thin
direct-instantiation wrappers. The along-exhaustion and concrete
$\mathbb{Z}^d$ wrappers now live in
\texttt{AmbientLattice/SpecialCases/MayerRecurrenceHasSum.lean}
and
\texttt{Concrete/LatticeGraphCorrelation/MayerRecurrenceHasSum.lean};
the Lambda-layer wrappers live in
\texttt{IsingModel/AmbientLattice/AnalyticityLambdaMayerRecurrenceEpsilon.lean}
(narrowed out of \texttt{Analyticity.lean} in PR~\#1988).

\paragraph{\texttt{polymerFreeEnergy} tanh-bound +
ferromagnetic + \texttt{hasDerivAt} +
\texttt{eq\_log\_one\_add\_eps}, all volume layers (\S18.5,
PR~\#1599).} Lifts six abstract \texttt{polymerFreeEnergy}
lemmas to the four volume layers (\textit{Lambda}, along-ex,
$\mathbb{Z}^d$ \textit{Lambda}, $\mathbb{Z}^d$ along-ex):
\texttt{\_tanh\_le\_card\_mul} (under $0 \le \beta J$),
\texttt{\_tanh\_le\_card\_mul\_ferromagnetic},
\texttt{\_tanh\_sandwich\_ferromagnetic},
\texttt{\_tanh\_le\_card\_log\_two\_ferromagnetic} (the latter
three under $0 \le J$ and $0 < \beta$), the
$\log(1 + \varepsilon(t))$ decomposition
\texttt{\_eq\_log\_one\_add\_eps}, and \texttt{\_hasDerivAt}
giving the explicit log-derivative for $t \ge 0$. 24 new thin
direct-instantiation wrappers. The along-exhaustion and concrete
$\mathbb{Z}^d$ wrappers now live in
\texttt{AmbientLattice/SpecialCases/PolymerFreeEnergyTanhBounds.lean}
and
\texttt{Concrete/LatticeGraphCorrelation/PolymerFreeEnergyTanhBounds.lean}.
The Lambda-layer wrappers live in
\texttt{IsingModel/AmbientLattice/AnalyticityLambdaMayerPfeEdgeBounds.lean}
(narrowed out of \texttt{Analyticity.lean} in PR~\#1987).
Most ferromagnetic ℤ$^d$ /
along-ex wrappers use the abbreviated \texttt{\_ferro} suffix
to keep each declaration line under the 100-character linter
budget.

\paragraph{\texttt{polymerFreeEnergy} \texttt{at\_zero} /
\texttt{at\_one} + \texttt{analyticAt} +
\texttt{analyticOnNhd\_Ici\_zero} +
\texttt{sandwich\_of\_nonneg}, all volume layers (\S18.5,
PR~\#1598).} Lifts five abstract \texttt{polymerFreeEnergy}
lemmas to the four volume layers (\textit{Lambda}, along-ex,
$\mathbb{Z}^d$ \textit{Lambda}, $\mathbb{Z}^d$ along-ex):
$\mathrm{pFE}(G, 0) = 0$,
$\mathrm{pFE}(G, 1) = \log |\mathrm{vdCompat}(G)|$,
\texttt{analyticAt} for $t \ge 0$ (named-wrapper of
\texttt{log\_vdSum\_analyticAt}),
\texttt{AnalyticOnNhd} over $[0,\infty)$, and the sandwich
$0 \le \mathrm{pFE}(G, t) \le |E(G)| \cdot \log(1+t)$ for
$t \ge 0$. 20 new thin direct-instantiation wrappers. The
along-exhaustion and $\mathbb{Z}^d$ analytic wrappers now live in
\texttt{AmbientLattice/SpecialCases/PolymerFreeEnergyAnalyticity.lean}
and
\texttt{Concrete/LatticeGraphCorrelation/PolymerFreeEnergyAnalyticity.lean}
(moved from the legacy modules in PR~\#1861). The along-exhaustion
at-zero / at-one / sandwich wrappers and their concrete $\mathbb{Z}^d$
counterparts now live in
\texttt{AmbientLattice/SpecialCases/PolymerFreeEnergyBasic.lean}
and
\texttt{Concrete/LatticeGraphCorrelation/PolymerFreeEnergyBasic.lean}.
The Lambda-layer wrappers live in
\texttt{IsingModel/AmbientLattice/AnalyticityLambdaMayerPfeEdgeBounds.lean}
(narrowed out of \texttt{Analyticity.lean} in PR~\#1987).

\paragraph{Mayer expansion edge-cases (zero) + $n = 2$ +
\texttt{abs\_le}, all volume layers (\S18.5, PR~\#1597).}
Lifts six abstract Mayer-expansion theorems to the four volume
layers (Lambda-direct, along-exhaustion, $\mathbb{Z}^d$ Lambda-direct,
$\mathbb{Z}^d$ along-exhaustion):
\texttt{mayerExpansionTerm\_two}, \texttt{\_two\_filter},
\texttt{mayerPartialSum\_two},
\texttt{mayerPartialSum\_eq\_zero\_of\_no\_polymers},
\texttt{mayerPartialSum\_eq\_zero\_of\_edgeFinset\_empty},
and \texttt{mayerExpansionTerm\_abs\_le}. 24 new thin
direct-instantiation wrappers (6 abstracts $\times$ 4 layers).
The two `zero' wrappers establish that the partial sums vanish
on no-polymer / edgeless induced subgraphs, the `two' wrappers
provide the explicit $n=2$ /$\,N=2$ formulas, and
\texttt{abs\_le} gives the triangle-inequality bound
$|\mathrm{mayerExpansionTerm}| \le \sum |\phi^T(\omega)| \cdot
|z(t,\omega)|$. The along-exhaustion and concrete $\mathbb{Z}^d$
wrappers now live in
\texttt{AmbientLattice/SpecialCases/MayerExpansionEdgeCases.lean}
and
\texttt{Concrete/LatticeGraphCorrelation/MayerExpansionEdgeCases.lean};
the $\Lambda$-layer wrappers live in
\texttt{IsingModel/AmbientLattice/AnalyticityLambdaMayerPfeEdgeBounds.lean}
(narrowed out of \texttt{Analyticity.lean} in PR~\#1987).

\paragraph{\texttt{vdPolymerFamilies\_sum} generic-\textit{t}
bounds + decomposition, all volume layers (\S18.5, PR~\#1596).}
Adds 4-layer wraps for the four abstract generic-\textit{t}
theorems
\texttt{\_ge\_one\_of\_nonneg},
\texttt{\_le\_one\_plus\_pow\_of\_nonneg},
\texttt{\_pos\_of\_nonneg}, and \texttt{\_eq\_one\_add}
(which had no wrappers before). 16 new thin
direct-instantiation wrappers (4 abstracts $\times$ 4 layers).
The along-exhaustion and concrete $\mathbb{Z}^d$ wrappers now live in
the matching \texttt{MayerVdBounds.lean} child modules; the
$\Lambda$-layer wrappers live in
\texttt{IsingModel/AmbientLattice/AnalyticityLambdaBasicIdentities.lean}
(narrowed out of \texttt{Analyticity.lean} in PR~\#1986).
The generic-\textit{t} bounds are the non-tanh ($t \geq 0$)
counterparts of the tanh bound family from PR~\#1595; the
\texttt{\_eq\_one\_add} decomposition expresses
$\mathrm{vdSum}(G,t) = 1 + \varepsilon(t)$ where
$\varepsilon(t) = \sum_{\Gamma \neq \emptyset} \prod_{P \in
\Gamma} t^{|P|}$.

\paragraph{\texttt{vdPolymerFamilies\_sum} iff characterizations,
all volume layers (\S18.5, PR~\#1594).} The four abstract iff
characterizations (\texttt{\_eq\_one\_iff\_eps\_eq\_zero},
\texttt{\_gt\_one\_iff\_eps\_pos},
\texttt{\_tanh\_gt\_one\_iff}, \texttt{\_tanh\_eq\_one\_iff}) are
lifted to the four volume layers. The non-tanh forms had Lambda-layer
wraps already; this PR completes the along-exhaustion and
$\mathbb{Z}^d$ wraps for both, plus all four layers for the tanh
forms. 14 new thin direct-instantiation wrappers.
The along-exhaustion and concrete $\mathbb{Z}^d$ wrappers now live in
\path{IsingModel/AmbientLattice/SpecialCases/MayerVdIff.lean}
and
\path{IsingModel/Concrete/LatticeGraphCorrelation/MayerVdIff.lean};
the non-tanh $\Lambda$-layer wrappers
(\texttt{\_eq\_one\_iff\_eps\_zero}, \texttt{\_gt\_one\_iff\_eps\_pos})
live in
\path{IsingModel/AmbientLattice/AnalyticityLambdaPolymer.lean}
(narrowed out of \texttt{Analyticity.lean} in PR~\#1979);
the tanh $\Lambda$-layer wrappers
(\texttt{\_tanh\_gt\_one\_iff}, \texttt{\_tanh\_eq\_one\_iff}) live in
\path{IsingModel/AmbientLattice/AnalyticityLambdaBasicIdentities.lean}
(narrowed out of \texttt{Analyticity.lean} in PR~\#1986).

\paragraph{Basic identities at\_zero / at\_one, all volume layers
(\S18.5, PR~\#1593).} The eight abstract base-case identities
(\texttt{vdPolymerFamilies\_sum\_at\_zero} (Step~599),
\texttt{\_at\_one} (Step~639),
\texttt{mayerPartialSum\_zero}, \texttt{\_one},
\texttt{\_at\_zero} (Steps~592, 598),
\texttt{mayerExpansionTerm\_zero}, \texttt{\_one},
\texttt{\_at\_zero} (Steps~587, 598)) are lifted to the four
volume layers, giving $32$ thin direct-instantiation wrappers. The
along-exhaustion and concrete $\mathbb{Z}^d$ wrappers now live in
\path{IsingModel/AmbientLattice/SpecialCases/MayerBasicIdentities.lean}
and
\path{IsingModel/Concrete/LatticeGraphCorrelation/MayerBasicIdentities.lean};
the $\Lambda$-layer wrappers live in
\path{IsingModel/AmbientLattice/AnalyticityLambdaBasicIdentities.lean}
(narrowed out of \texttt{Analyticity.lean} in PR~\#1986).

\paragraph{Mayer identity edge-cases (no polymers / trivial /
edgeless), all volume layers (\S18.5, PR~\#1592).} The five abstract
edge-case Mayer identity statements
(\texttt{mayer\_identity\_of\_no\_polymers} (Step~618),
\texttt{\_no\_polymers\_tanh} (Step~619),
\texttt{\_of\_trivial} (Step~651),
\texttt{\_of\_edgeFinset\_empty} (Step~620),
\texttt{\_of\_edgeFinset\_empty\_tanh} (Step~622)) are lifted to
the four volume layers, giving $20$ thin direct-instantiation
wrappers. The along-exhaustion and concrete $\mathbb{Z}^d$
wrappers now live in
\texttt{AmbientLattice/SpecialCases/MayerTrivialCases.lean}
and
\texttt{Concrete/LatticeGraphCorrelation/MayerTrivialCases.lean};
the $\Lambda$-layer wrappers live in
\texttt{IsingModel/AmbientLattice/AnalyticityLambdaMayerIdentity.lean}
(narrowed out of \texttt{Analyticity.lean} in PR~\#1985).

\paragraph{\texttt{mayerPartialSum\_0} $\le$
\texttt{polymerFreeEnergy} family, all volume layers (\S18.5,
PR~\#1591).} The three abstract `mayerPartialSum 0 ≤
polymerFreeEnergy' inequalities (Steps 654-656) are lifted to the
four volume layers, giving $12$ thin direct-instantiation wrappers.
The along-exhaustion and concrete $\mathbb{Z}^d$ wrappers now live in
\texttt{AmbientLattice/SpecialCases/MayerTrivialCases.lean}
and
\texttt{Concrete/LatticeGraphCorrelation/MayerTrivialCases.lean};
the $\Lambda$-layer wrappers live in
\texttt{IsingModel/AmbientLattice/AnalyticityLambdaMayerIdentity.lean}
(narrowed out of \texttt{Analyticity.lean} in PR~\#1985).

\paragraph{Mayer identity polymer\_free\_energy variants, all
volume layers (\S18.5, PR~\#1590).} The three abstract Mayer
identity restatements
(\texttt{mayer\_identity\_at\_J\_zero\_polymer\_free\_energy},
\texttt{\_at\_beta\_zero\_polymer\_free\_energy} (Step~652),
\texttt{\_at\_either\_zero\_polymer\_free\_energy} (Step~653))
are lifted to the four volume layers, giving $12$ thin
direct-instantiation wrappers. The along-exhaustion and concrete
$\mathbb{Z}^d$ wrappers now live in
\texttt{AmbientLattice/SpecialCases/MayerEdgeCases.lean}
and
\texttt{Concrete/LatticeGraphCorrelation/MayerEdgeCases.lean};
the $\Lambda$-layer wrappers live in
\texttt{IsingModel/AmbientLattice/AnalyticityLambdaMayerIdentity.lean}
(narrowed out of \texttt{Analyticity.lean} in PR~\#1985).

\paragraph{\texttt{polymerFreeEnergy\_eq\_mayerPartialSum} at
edge-case parameter slices, all volume layers (\S18.5, PR~\#1589).}
The four abstract named-wrapper restatements
(\texttt{polymerFreeEnergy\_eq\_mayerPartialSum\_at\_zero}
(Step~611), \texttt{\_at\_betaJ\_zero},
\texttt{\_at\_beta\_zero}, \texttt{\_at\_J\_zero}
(Step~617 specialisations)) are lifted to the four volume layers,
giving $16$ thin direct-instantiation wrappers. The
along-exhaustion and concrete $\mathbb{Z}^d$ wrappers now live in
\texttt{AmbientLattice/SpecialCases/MayerEdgeCases.lean}
and
\texttt{Concrete/LatticeGraphCorrelation/MayerEdgeCases.lean};
the $\Lambda$-layer wrappers live in
\texttt{IsingModel/AmbientLattice/AnalyticityLambdaMayerIdentity.lean}
(narrowed out of \texttt{Analyticity.lean} in PR~\#1985).

\paragraph{Mayer identity at edge-case parameter slices, all
volume layers (\S18.5, PR~\#1588).} The four abstract Mayer identity
statements at trivial parameter slices
(\texttt{mayer\_identity\_at\_zero} (Step~600),
\texttt{\_at\_betaJ\_zero}, \texttt{\_at\_beta\_zero},
\texttt{\_at\_J\_zero} (Step~609 specialisations)) are lifted to
the four volume layers, giving $16$ thin direct-instantiation
wrappers. The along-exhaustion and concrete $\mathbb{Z}^d$
wrappers now live in
\texttt{AmbientLattice/SpecialCases/MayerEdgeCases.lean}
and
\texttt{Concrete/LatticeGraphCorrelation/MayerEdgeCases.lean};
the $\Lambda$-layer wrappers live in
\texttt{IsingModel/AmbientLattice/AnalyticityLambdaMayerIdentity.lean}
(narrowed out of \texttt{Analyticity.lean} in PR~\#1985).

\paragraph{\texttt{log\_vdPolymerFamilies\_sum} analyticity, all
volume layers (\S18.5, PR~\#1587).} The four abstract analyticity
statements
(\texttt{log\_vdPolymerFamilies\_sum\_analyticAt} (Step~606),
\texttt{\_analyticOnNhd\_Ici\_zero} (Step~607),
\texttt{\_tanh\_analyticAt\_\{beta,J\}} (Step~608)) are lifted to
the four volume layers, giving $16$ thin direct-instantiation
wrappers in
\texttt{IsingModel/AmbientLattice/Analyticity.lean},
\texttt{AmbientLattice/SpecialCases.lean}, and
\texttt{Concrete/LatticeGraphCorrelation.lean}. The along-exhaustion
and concrete analytic wrappers from this family now live in
\texttt{AmbientLattice/SpecialCases/VdPolymerFamiliesAnalyticity.lean}
and
\texttt{Concrete/LatticeGraphCorrelation/VdPolymerFamiliesAnalyticity.lean};
the $\Lambda$-layer wrappers live in
\texttt{IsingModel/AmbientLattice/AnalyticityLambdaVdPolymer.lean}
(moved from the legacy modules in PR~\#1863; narrowed out of
\texttt{Analyticity.lean} in PR~\#1984).

\paragraph{\texttt{vdPolymerFamilies\_sum} $\beta$ / $J$ regularity
(tanh form), all volume layers (\S18.5, PR~\#1586).} The six
abstract $\beta$/$J$ regularity statements
(\texttt{vdPolymerFamilies\_sum\_tanh\_continuous\_\{beta,J\}},
\texttt{\_differentiable\_\{beta,J\}},
\texttt{\_analyticAt\_\{beta,J\}}, Steps 556, 559, 562) are lifted
across the four volume layers. The along-exhaustion
and concrete \texttt{Continuous} / \texttt{Differentiable} wrappers from
this family now live in
\texttt{AmbientLattice/SpecialCases/MayerVdRegularity.lean}
and
\texttt{Concrete/LatticeGraphCorrelation/MayerVdRegularity.lean};
the along-exhaustion and concrete \texttt{AnalyticAt} wrappers remain in
\texttt{AmbientLattice/SpecialCases/VdPolymerFamiliesAnalyticity.lean}
and
\texttt{Concrete/LatticeGraphCorrelation/VdPolymerFamiliesAnalyticity.lean};
the $\Lambda$-layer wrappers live in
\texttt{IsingModel/AmbientLattice/AnalyticityLambdaVdPolymer.lean}
(moved from the legacy modules in PR~\#1863; narrowed out of
\texttt{Analyticity.lean} in PR~\#1984).

\paragraph{\texttt{vdPolymerFamilies\_sum} regularity in $t$, all
volume layers (\S18.5, PR~\#1585).} The four abstract regularity
statements for the polymer-family sum
(\texttt{vdPolymerFamilies\_sum\_continuous} (Step~555),
\texttt{\_differentiable} (Step~558), \texttt{\_analyticAt}
(Step~561), \texttt{\_hasDerivAt} with explicit polynomial
derivative (Step~575)) are lifted across the four volume layers.
The along-exhaustion
and concrete \texttt{Continuous} / \texttt{Differentiable} /
\texttt{HasDerivAt} wrappers from this family now live in
\texttt{AmbientLattice/SpecialCases/MayerVdRegularity.lean}
and
\texttt{Concrete/LatticeGraphCorrelation/MayerVdRegularity.lean};
the along-exhaustion and concrete \texttt{AnalyticAt} wrappers remain in
\texttt{AmbientLattice/SpecialCases/VdPolymerFamiliesAnalyticity.lean}
and
\texttt{Concrete/LatticeGraphCorrelation/VdPolymerFamiliesAnalyticity.lean};
the $\Lambda$-layer wrapper lives in
\texttt{IsingModel/AmbientLattice/AnalyticityLambdaVdPolymer.lean}
(moved from the legacy modules in PR~\#1863; narrowed out of
\texttt{Analyticity.lean} in PR~\#1984).

\paragraph{\texttt{mayerExpansionTerm} $\beta$ / $J$ regularity
(tanh form), all volume layers (\S18.5, PR~\#1584).} The six
abstract $\beta$/$J$ regularity statements (continuous,
differentiable, analyticAt, each in $\beta$ and $J$) for
\texttt{mayerExpansionTerm} composed with $\tanh(\beta J)$ are
lifted across the four volume layers. The along-exhaustion
and concrete \texttt{Continuous} / \texttt{Differentiable} wrappers from
this family now live in
\texttt{AmbientLattice/SpecialCases/MayerVdRegularity.lean} and
\texttt{Concrete/LatticeGraphCorrelation/MayerVdRegularity.lean}; the
along-exhaustion and concrete analytic wrappers remain in
\texttt{AmbientLattice/SpecialCases/MayerAnalyticity.lean} and
\texttt{Concrete/LatticeGraphCorrelation/MayerAnalyticity.lean}; the
$\Lambda$-layer wrappers remain in
\texttt{IsingModel/AmbientLattice/Analyticity.lean}
(moved from the legacy modules in PR~\#1862).

\paragraph{\texttt{mayerPartialSum} $\beta$ / $J$ regularity (tanh
form), all volume layers (\S18.5, PR~\#1583).} The eight abstract
$\beta$/$J$ regularity statements (continuous, differentiable,
analyticAt, analyticOnNhd, each in $\beta$ and $J$) for
\texttt{mayerPartialSum} composed with $\tanh(\beta J)$ are lifted
across the four volume layers. The along-exhaustion
and concrete \texttt{Continuous} / \texttt{Differentiable} wrappers from
this family now live in
\texttt{AmbientLattice/SpecialCases/MayerVdRegularity.lean} and
\texttt{Concrete/LatticeGraphCorrelation/MayerVdRegularity.lean}; the
along-exhaustion and concrete analytic wrappers remain in
\texttt{AmbientLattice/SpecialCases/MayerAnalyticity.lean} and
\texttt{Concrete/LatticeGraphCorrelation/MayerAnalyticity.lean}; the
$\Lambda$-layer wrappers remain in
\texttt{IsingModel/AmbientLattice/Analyticity.lean}
(moved from the legacy modules in PR~\#1862).

\paragraph{\texttt{mayerExpansionTerm} regularity, all volume layers
(\S18.6, PR~\#1582).} The four abstract regularity statements
(\texttt{mayerExpansionTerm\_continuous} (Step~588),
\texttt{\_differentiable} (Step~589),
\texttt{\_analyticAt} and \texttt{\_analyticOnNhd} (Step~590)) are
lifted across the four volume layers. The along-exhaustion
and concrete \texttt{Continuous} / \texttt{Differentiable} wrappers from
this family now live in
\texttt{AmbientLattice/SpecialCases/MayerVdRegularity.lean} and
\texttt{Concrete/LatticeGraphCorrelation/MayerVdRegularity.lean}; the
along-exhaustion and concrete analytic wrappers remain in
\texttt{AmbientLattice/SpecialCases/MayerAnalyticity.lean} and
\texttt{Concrete/LatticeGraphCorrelation/MayerAnalyticity.lean}; the
$\Lambda$-layer wrappers remain in
\texttt{IsingModel/AmbientLattice/Analyticity.lean}
(moved from the legacy modules in PR~\#1862).

\paragraph{\texttt{mayerPartialSum} regularity, all volume layers
(\S18.6, PR~\#1581).} The five abstract regularity statements
(\texttt{mayerPartialSum\_continuous},
\texttt{\_differentiable}, \texttt{\_analyticAt} (Step~591),
\texttt{\_continuousOn}, \texttt{\_differentiableOn} (Step~628))
are lifted across the four volume layers. The along-exhaustion
and concrete \texttt{Continuous} / \texttt{Differentiable} /
\texttt{ContinuousOn} / \texttt{DifferentiableOn} wrappers from this
family now live in
\texttt{AmbientLattice/SpecialCases/MayerVdRegularity.lean} and
\texttt{Concrete/LatticeGraphCorrelation/MayerVdRegularity.lean}; the
along-exhaustion and concrete \texttt{AnalyticAt} wrappers remain in
\texttt{AmbientLattice/SpecialCases/MayerAnalyticity.lean} and
\texttt{Concrete/LatticeGraphCorrelation/MayerAnalyticity.lean}; the
$\Lambda$-layer wrapper remains in
\texttt{IsingModel/AmbientLattice/Analyticity.lean}
(moved from the legacy modules in PR~\#1862).

\paragraph{\texttt{polymerFreeEnergy} tanh-bound family, all volume
layers (\S18.5, PR~\#1580).} The tanh-form sandwich (Step~632),
the $\le |E| \cdot \log 2$ bound for $0 \le t \le 1$ (Step~642),
its tanh-form (Step~643), and the combined double bound
(Step~645) are lifted to the four volume layers, giving $16$ thin
direct-instantiation wrappers in
\texttt{IsingModel/AmbientLattice/Analyticity.lean}, with the
along-exhaustion and concrete wrappers now in
\texttt{AmbientLattice/SpecialCases/PolymerFreeEnergyBounds.lean} and
\texttt{Concrete/LatticeGraphCorrelation/PolymerFreeEnergyBounds.lean}
(moved from the legacy modules in PR~\#1866).

\paragraph{\texttt{polymerFreeEnergy} edge-cases + comparison
family, all volume layers (\S18.5, PR~\#1579).} The four abstract
edge-case / comparison statements
(\texttt{polymerFreeEnergy\_eq\_zero\_of\_no\_polymers} (Step~621),
\texttt{polymerFreeEnergy\_eq\_zero\_of\_edgeFinset\_empty}
(Step~623),
\texttt{polymerFreeEnergy\_le\_of\_le\_of\_nonneg} (Step~649),
\texttt{polymerFreeEnergy\_le\_of\_le\_strict\_form} (Step~650))
are lifted to the four volume layers, giving $16$ thin
direct-instantiation wrappers in
\texttt{IsingModel/AmbientLattice/Analyticity.lean}, with the
along-exhaustion and concrete wrappers now in
\texttt{AmbientLattice/SpecialCases/PolymerFreeEnergyBounds.lean} and
\texttt{Concrete/LatticeGraphCorrelation/PolymerFreeEnergyBounds.lean}
(moved from the legacy modules in PR~\#1866).

\paragraph{\texttt{polymerFreeEnergy} bound family, all volume layers
(\S18.5, PR~\#1578).} The four abstract bound statements
(\texttt{polymerFreeEnergy\_nonneg\_of\_nonneg},
\texttt{polymerFreeEnergy\_le\_card\_log\_one\_plus\_of\_nonneg}
($\le |E| \cdot \log(1+t)$),
\texttt{polymerFreeEnergy\_le\_card\_mul\_of\_nonneg}
($\le |E| \cdot t$),
\texttt{polymerFreeEnergy\_monotoneOn\_Ici\_zero}) are lifted to
the four volume layers, giving $16$ thin direct-instantiation
wrappers in
\texttt{IsingModel/AmbientLattice/Analyticity.lean}, with the
along-exhaustion and concrete wrappers now in
\texttt{AmbientLattice/SpecialCases/PolymerFreeEnergyBounds.lean} and
\texttt{Concrete/LatticeGraphCorrelation/PolymerFreeEnergyBounds.lean}
(moved from the legacy modules in PR~\#1866).

\paragraph{\texttt{polymerFreeEnergy} $\beta$ / $J$ analyticity, all
volume layers (\S18.6, PR~\#1577).} The four abstract β/J real-
analyticity statements (\texttt{polymerFreeEnergy\_tanh\_analyticAt\_beta}
/ \texttt{\_J} and the corresponding \texttt{\_analyticOnNhd\_*\_Ici\_zero})
are lifted to the four volume layers ($\Lambda$ /
along-exhaustion / $\mathbb{Z}^d$ $\Lambda$ / $\mathbb{Z}^d$
along-exhaustion), giving $16$ thin direct-instantiation wrappers
in \texttt{IsingModel/AmbientLattice/Analyticity.lean},
\texttt{AmbientLattice/SpecialCases/PolymerFreeEnergyAnalyticity.lean},
and
\texttt{Concrete/LatticeGraphCorrelation/PolymerFreeEnergyAnalyticity.lean}
(moved from the legacy modules in PR~\#1861).

\paragraph{\texttt{polymerFreeEnergy} regularity, all volume layers
(\S18.5, PR~\#1576).} The four abstract regularity statements
(\texttt{polymerFreeEnergy\_continuousAt},
\texttt{polymerFreeEnergy\_differentiableAt},
\texttt{polymerFreeEnergy\_continuousOn\_Ici\_zero},
\texttt{polymerFreeEnergy\_differentiableOn\_Ici\_zero}) are lifted
to the four volume layers ($\Lambda$ / along-exhaustion /
$\mathbb{Z}^d$ $\Lambda$ / $\mathbb{Z}^d$ along-exhaustion), giving
$16$ thin direct-instantiation wrappers in
\texttt{IsingModel/AmbientLattice/Analyticity.lean}, with the
along-exhaustion and concrete wrappers now in
\texttt{AmbientLattice/SpecialCases/PolymerFreeEnergyBounds.lean} and
\texttt{Concrete/LatticeGraphCorrelation/PolymerFreeEnergyBounds.lean}
(moved from the legacy modules in PR~\#1866).

\paragraph{Strict \texttt{freeEnergy} upper bound in cluster-expansion
convergence regime, all volume layers (\S18.5, PR~\#1575).}
PR~\#1527's abstract strict bound
\texttt{freeEnergy\_lt\_log\_two\_plus\_high\_temp\_correction}
(under $0 \le \beta J,\, 0 < |\iota|,\, (1+\tanh(\beta J))^{|E|} < 2$,
giving
$f < \log 2 + (|E|/|\iota|) \log\cosh(\beta J) + \log 2/|\iota|$)
admits a ferromagnetic version
($0 \le J,\, 0 < \beta$) and is lifted to all five layers via
direct \texttt{freeEnergy\LeanLambda\_apply} / \texttt{freeEnergyAlongExhaustion}
unfoldings: 9 new theorems in
\texttt{IsingModel/AmbientLattice/Analyticity.lean},
\texttt{AmbientLattice/SpecialCases/HighTemperature.lean}, and
\texttt{Concrete/LatticeGraphCorrelation/HighTemperature.lean}. The two
$\mathbb{Z}^d$ ferromagnetic variants use the abbreviated
\texttt{\_ferro} suffix to fit the linter's 100-character budget.
PR~\#2099 further carves the four ℤ$^d$
\texttt{freeEnergy\{$\Lambda$,AlongExhaustion\}\_latticeGraph\_lt\_log\_two\_plus\_high\_temp\_correction}
wrappers (with ferromagnetic variants) out into
\texttt{Concrete/LatticeGraphCorrelation/HighTemperatureFreeEnergyCorrection.lean}.
PR~\#2100 further carves the eight ℤ$^d$ ferromagnetic
\texttt{polymerFreeEnergy\_*\_tanh\_*\_ferro} and
\texttt{vdPolymerFamilies\_sum*\_*\_\{ferromagnetic,ferro\}} wrappers
out into
\texttt{Concrete/LatticeGraphCorrelation/HighTemperatureFerromagnetic.lean}.
PR~\#2101 further carves the four ℤ$^d$ non-ferromagnetic
\texttt{vdPolymerFamilies\_sum\_\{$\Lambda$,AlongExhaustion\}\_latticeGraph\_sandwich(\_sharp)?}
wrappers out into
\texttt{Concrete/LatticeGraphCorrelation/HighTemperatureVDPolymer.lean}.

\paragraph{Ferromagnetic polymer-family sum sandwich, all volume
layers (\S18.5, PR~\#1574).} Under $0 \le J,\, 0 < \beta$ each of
the two abstract sandwich statements admits a ferromagnetic
version, derived through \texttt{mul\_nonneg h$\beta$.le hJ} and
matched at each of the five layers. The two $\mathbb{Z}^d$ sharp
wrappers use the abbreviated \texttt{\_ferro} suffix to fit the
linter's 100-character budget; the other ferromagnetic wrappers
retain \texttt{\_ferromagnetic}.

\paragraph{Polymer-family sum sandwich, all volume layers
(\S18.5, PR~\#1572).} The two abstract sandwich statements
\texttt{vdPolymerFamilies\_sum\_sandwich} ($1 \le \sum_\Gamma \le
2^{|E|}$) and \texttt{vdPolymerFamilies\_sum\_sandwich\_sharp}
($\le (1 + \tanh(\beta J))^{|E|}$), both under $0 \le \beta J$,
are now lifted to the $\Lambda$, along-exhaustion, $\mathbb{Z}^d$
$\Lambda$, and $\mathbb{Z}^d$ along-exhaustion layers --- $8$
wrappers in total --- in
\texttt{IsingModel/AmbientLattice/Analyticity.lean},
\texttt{AmbientLattice/SpecialCases/HighTemperature.lean}, and
\texttt{Concrete/LatticeGraphCorrelation/HighTemperatureVDPolymer.lean}
(carved out of \texttt{Concrete/.../HighTemperature.lean} in PR~\#2101).
Each is a thin direct-instantiation wrapper of the abstract theorem on
the corresponding induced subgraph. These high-temperature sandwich
wrappers were moved to the dedicated child modules from the legacy
monoliths in PR~\#1864.

\paragraph{Convergence radius at all volume layers (\S18.5,
PR~\#1569).} The cluster-expansion convergence-radius statements
\texttt{polymerFreeEnergy\_high\_temp\_sandwich} from \#1526 and
\texttt{polymerFreeEnergy\_hasSum\_via\_log\_of\_pow\_lt\_two} from \#1517 are
lifted to every volume layer:
\begin{itemize}
  \item \emph{Abstract tanh form (new)}:
    \texttt{polymerFreeEnergy\_tanh\_hasSum\_via\_log\_of\_pow\_lt\_two}
    --- under $0 \le \beta J$ and $(1+\tanh(\beta J))^{|E|} < 2$,
    $\mathrm{polymerFreeEnergy}\,G\,(\tanh(\beta J))$ admits the same
    explicit log-Taylor series as the abstract form, restated in the
    Ising activity $t = \tanh(\beta J)$.
  \item \emph{Lambda-layer tanh wraps}:
    \texttt{polymerFreeEnergy\_\LeanLambda\_tanh\_high\_temp\_sandwich} and
    \texttt{polymerFreeEnergy\_\LeanLambda\_tanh\_hasSum\_via\_log\_of\_pow\_lt\_two}
    in \texttt{AmbientLattice/Analyticity.lean}.
  \item \emph{Along-exhaustion wraps}:
    \texttt{polymerFreeEnergyAlongExhaustion\_\{,tanh\_\}\{high\_temp\_sandwich,hasSum\_via\_log\_of\_pow\_lt\_two\}}
    (4 theorems, per stage $n$) in
    \texttt{AmbientLattice/SpecialCases/HighTemperature.lean}.
  \item \emph{$\mathbb{Z}^d$ concrete wraps}: 8 theorems split across
    the $\Lambda$-direct family on $\Lambda \subseteq \mathbb{Z}^d$
    in
    \texttt{Concrete/LatticeGraphCorrelation/HighTemperature.lean} and
    the along-exhaustion family on
    $\mathrm{Exhaustion}\,(\mathrm{Fin}\,d \to \mathbb{Z})$ in
    \texttt{Concrete/LatticeGraphCorrelation/HighTemperatureAlongEx.lean}
    (carved out of \texttt{Concrete/.../HighTemperature.lean} in PR~\#2102).
\end{itemize}
The convergence-radius condition $(1+t)^{|E|} < 2$ is exponentially
restrictive in $|E|$; sharper polymer-density / Kotecky--Preiss
estimates that depend on the maximum vertex degree of $G$ instead of
$|E|$ remain deferred (they require chromatic-polynomial /
matrix-tree machinery not yet available in Mathlib). The moved
high-temperature child modules keep the named API available through the
public umbrellas while allowing targeted imports to avoid the legacy
monolith.

\paragraph{Decay of correlations at $h=0$, all volume layers
(\S18.7, Step~574 + PR~\#1570).} The
high-temperature exponential-decay capstone
\[
\bigl|\langle \sigma_i \sigma_j \rangle_{\beta,0}\bigr|
  \le 2^{|E|} \cdot \tanh(\beta J)^{d_G(i,j)}
\quad \text{for } 0 \le \beta J,
\]
where $d_G$ is graph distance, is proved via the
even-subgraph-boundary distance bound
\texttt{evenSubgraph\_pair\_boundary\_dist\_le} (Step~573), which
uses strong induction to construct an explicit
$G.\mathrm{Walk}\,i\,j$ from any $X$ with
$\partial X = \{i, j\}$.

PR~\#1570 lifts this capstone to every volume layer: the
$\Lambda$-direct form
\texttt{correlation\LeanLambda\_high\_temp\_h\_zero\_at\_pair\_le\_two\_pow\_edges\_mul\_tanh\_pow\_dist}
in \texttt{AmbientLattice/Defs.lean}, the per-stage
\texttt{correlationAlongExhaustion\_high\_temp\_h\_zero\_at\_pair\_le\_two\_pow\_edges\_mul\_tanh\_pow\_dist}
in
\texttt{AmbientLattice/SpecialCases/HighTemperatureBoundsDecayCapstones.lean}
(narrowed out of \texttt{HighTemperatureBounds.lean} in PR~\#2000),
plus the $\mathbb{Z}^d$ concrete forms
\texttt{correlation\LeanLambda\_latticeGraph\_high\_temp\_h\_zero\_at\_pair\_le\_two\_pow\_edges\_mul\_tanh\_pow\_dist}
and
\texttt{correlationAlongExhaustion\_latticeGraph\_high\_temp\_h\_zero\_at\_pair\_le\_two\_pow\_edges\_mul\_tanh\_pow\_dist}
in \texttt{Concrete/LatticeGraphCorrelation/HighTemperatureBounds.lean}.

PR~\#1932 further splits the concrete §18.3--§18.4 partition-function and
free-energy high-temperature expansion / closed-form / lower-bound /
upper-bound / \texttt{lower\_le\_upper} wrappers (22 theorems: the
$\mathrm{partitionFunction}\Lambda\_\mathrm{latticeGraph}\_\mathrm{high}\_\mathrm{temp}\_\mathrm{expansion}$ /
$\mathrm{partitionFunctionAlongExhaustion}\_\mathrm{latticeGraph}\_\mathrm{high}\_\mathrm{temp}\_\mathrm{expansion}$ /
$\mathrm{freeEnergy}\Lambda\_\mathrm{latticeGraph}\_\mathrm{high}\_\mathrm{temp}\_*$ /
$\mathrm{freeEnergyAlongExhaustion}\_\mathrm{latticeGraph}\_\mathrm{high}\_\mathrm{temp}\_*$ /
$\mathrm{log}\_\mathrm{partitionFunction}\Lambda\_\mathrm{latticeGraph}\_\mathrm{high}\_\mathrm{temp}\_\mathrm{expansion}\_h\_\mathrm{zero}\_\mathrm{closed}$ /
$\mathrm{log}\_\mathrm{partitionFunctionAlongExhaustion}\_\mathrm{latticeGraph}\_\mathrm{high}\_\mathrm{temp}\_\mathrm{expansion}\_h\_\mathrm{zero}\_\mathrm{closed}$ /
$\mathrm{correlation}\Lambda\_\mathrm{latticeGraph}\_\mathrm{high}\_\mathrm{temp}\_h\_\mathrm{zero}\_\mathrm{at}\_\mathrm{empty}\_A$ wrappers) out of
\texttt{Concrete/LatticeGraphCorrelation/HighTemperatureBounds.lean} into a
new child module
\texttt{Concrete/LatticeGraphCorrelation/HighTemperatureBoundsExpansion.lean}.
Sandwich and downstream wrappers remain in
\texttt{HighTemperatureBounds}; the sharper-exp wrappers were further
split out into \texttt{HighTemperatureBoundsExpSharper} in PR~\#1935;
the deviation / continuity wrappers were further split out into
\texttt{HighTemperatureBoundsDeviation} in PR~\#1936; the
ratio\_sandwich / ratio\_bound wrappers were further split out into
\texttt{HighTemperatureBoundsRatioBounds} in PR~\#1937 (with the 7
$\Lambda$-direct \texttt{triple\_ratio\_*} wrappers subsequently
narrowed into \texttt{HighTemperatureBoundsTripleRatio} in PR~\#1998
and the 12 $\Lambda$-direct
\texttt{log\_partitionFunctionΛ\_latticeGraph} /
\texttt{freeEnergyΛ\_latticeGraph} ratio wrappers narrowed into
\texttt{HighTemperatureBoundsRatioLogFe} in PR~\#1999); the
alongExhaustion correlation/sandwich basic wrappers
were further split out into
\texttt{HighTemperatureBoundsAlongExhaustionBasic} in PR~\#1938.
The legacy import path is preserved.

PR~\#1946 further splits 34 ambient §18.3--§18.4 alongExhaustion
\texttt{ratio\_sandwich} / \texttt{ratio\_bound} /
\texttt{triple\_ratio\_*} / \texttt{\_of\_nonempty} wrappers (for
$\mathrm{partitionFunctionAlongExhaustion}$,
$\mathrm{freeEnergyAlongExhaustion}$, and
$\mathrm{log}\,\mathrm{partitionFunctionAlongExhaustion}$ with
bundle / triple / ferromagnetic variants) out of
\texttt{AmbientLattice/SpecialCases/HighTemperatureBounds.lean} into
a new child module
\texttt{HighTemperatureBoundsRatioBounds.lean} (under
\texttt{AmbientLattice/SpecialCases/}). PR~\#1994 subsequently narrows
the 7 \texttt{triple\_ratio\_sandwich\_bundle} /
\texttt{triple\_ratio\_bound\_bundle} wrappers (J = 0 / $\beta$ = 0 /
ferromagnetic variants) out of that child into a separate
\texttt{HighTemperatureBoundsTripleRatio.lean}, and PR~\#1995 further
narrows the 17 \texttt{log\_partitionFunctionAlongExhaustion} and
\texttt{freeEnergyAlongExhaustion} ratio (+
\texttt{deviation\_pos} / \texttt{pow\_two\_lt}) wrappers out of the
same child into \texttt{HighTemperatureBoundsRatioLogFe.lean}, both
under the same directory. The legacy import path is preserved.

PR~\#1945 further splits 20 ambient §18.3--§18.4 alongExhaustion
\texttt{deviation\_bound\_exp} / \texttt{continuity\_bundle} /
\texttt{deviation\_sandwich} / \texttt{relative\_sandwich} /
\texttt{deviation\_pos} / \texttt{pow\_two\_lt} /
\texttt{strict\_deviation\_bundle} wrappers (for
$\mathrm{freeEnergyAlongExhaustion}$,
$\mathrm{partitionFunctionAlongExhaustion}$, and
$\mathrm{log}\,\mathrm{partitionFunctionAlongExhaustion}$ with
ferromagnetic variants) out of
\texttt{AmbientLattice/SpecialCases/HighTemperatureBounds.lean} into
a new child module
\texttt{HighTemperatureBoundsDeviation.lean} (under
\texttt{AmbientLattice/SpecialCases/}). PR~\#2018 subsequently
narrows the 10 strict-deviation wrappers
(\texttt{relative\_sandwich}, \texttt{deviation\_pos},
\texttt{pow\_two\_lt},
\texttt{log\_partitionFunctionAlongExhaustion\_*\_deviation\_pos},
\texttt{\_strict\_deviation\_bundle} with ferromagnetic variants)
out of that child into a separate
\texttt{HighTemperatureBoundsDeviationStrict.lean} under the same
directory. PR~\#2024 further narrows the 4
\texttt{freeEnergyAlongExhaustion\_*\_continuity\_*} wrappers
(\texttt{\_at\_J\_zero}, \texttt{\_at\_beta\_zero},
\texttt{\_bundle}, \texttt{\_bundle\_ferromagnetic}) out of the
same child into
\texttt{HighTemperatureBoundsDeviationContinuity.lean}. The legacy
import path is preserved.

PR~\#1944 further splits 16 ambient §18.3--§18.4 alongExhaustion
sharper-exp upper-bound / sandwich / complete-summary wrappers
(\texttt{partitionFunctionAlongExhaustion},
\texttt{freeEnergyAlongExhaustion},
\texttt{log\_partitionFunctionAlongExhaustion} \texttt{\_*\_exp} family
plus ferromagnetic variants) out of
\texttt{AmbientLattice/SpecialCases/HighTemperatureBounds.lean} into
a new child module
\texttt{HighTemperatureBoundsExpSharper.lean} (under
\texttt{AmbientLattice/SpecialCases/}). PR~\#2019 subsequently
narrows the 6 \texttt{complete\_summary\_exp} wrappers
(\texttt{freeEnergyAlongExhaustion},
\texttt{partitionFunctionAlongExhaustion},
\texttt{log\_partitionFunctionAlongExhaustion} with ferromagnetic
variants) out of that child into a separate
\texttt{HighTemperatureBoundsExpSharperComplete.lean} under the
same directory. PR~\#2023 further narrows the 5
\texttt{\_sandwich\_exp} wrappers
(\texttt{log\_partitionFunctionAlongExhaustion\_*},
\texttt{partitionFunctionAlongExhaustion\_*},
\texttt{freeEnergyAlongExhaustion\_*} plus ferromagnetic
variants for Z and f) out of that same child into
\texttt{HighTemperatureBoundsExpSharperSandwich.lean}. The legacy
import path is preserved.

PR~\#1943 further splits 20 ambient §18.3--§18.4 alongExhaustion
partition function / free energy / log-partition / correlation
expansion / closed-form / lower-bound / upper-bound / sandwich /
complete-summary wrappers (plus the
$\mathrm{one\_le\_sum\_pow\_tanh\_even\_subgraph\_alongExhaustion}$
helper) out of \texttt{AmbientLattice/SpecialCases/HighTemperatureBounds.lean}
into a new child module
\texttt{HighTemperatureBoundsExpansion.lean} (under
\texttt{AmbientLattice/SpecialCases/}). PR~\#2020 subsequently
narrows the 8 closed-form / sandwich / complete-summary wrappers
(\texttt{partitionFunctionAlongExhaustion\_*\_closed\_at\_J\_zero},
\texttt{\_closed\_at\_beta\_zero}, redundant \texttt{\_closed},
\texttt{correlationAlongExhaustion\_*\_nonneg}, \texttt{\_closed},
\texttt{partitionFunctionAlongExhaustion\_*\_sandwich},
\texttt{partitionFunctionAlongExhaustion\_*\_complete\_summary},
\texttt{freeEnergyAlongExhaustion\_*\_complete\_summary}) out of
that child into a separate
\texttt{HighTemperatureBoundsExpansionClosedForms.lean} under the
same directory. PR~\#2021 further narrows the 4
\texttt{partitionFunctionAlongExhaustion\_high\_temp\_expansion\_h\_zero}
/ \texttt{\_high\_temp\_expansion} /
\texttt{\_high\_temp\_expansion\_subset\_form} /
\texttt{one\_le\_sum\_pow\_tanh\_even\_subgraph\_alongExhaustion}
wrappers out of that same child into
\texttt{HighTemperatureBoundsExpansionVariants.lean}. PR~\#2022
further narrows the 8
\texttt{partitionFunctionAlongExhaustion} /
\texttt{freeEnergyAlongExhaustion} /
\texttt{log\_partitionFunctionAlongExhaustion}
\texttt{\_high\_temp\_expansion\_h\_zero\_closed} /
\texttt{\_high\_temp\_expansion\_h\_zero\_upper\_bound} /
\texttt{\_high\_temp\_expansion\_h\_zero\_lower\_bound} /
\texttt{\_high\_temp\_h\_zero\_upper\_bound} /
\texttt{\_high\_temp\_h\_zero\_lower\_bound} /
\texttt{\_high\_temp\_h\_zero\_lower\_le\_upper} lower/upper-bound
plus consistency wrappers out of the same child into
\texttt{HighTemperatureBoundsExpansionLowerUpper.lean}. The legacy
import path is preserved.

PR~\#1942 further splits the concrete §18.3--§18.4
$\mathrm{freeEnergyInfinite}$ high-temperature wrappers on
$\mathrm{latticeGraph}\,d$ (10 theorems with caller-supplied
\texttt{Exhaustion} BED witness and $\mathrm{cubicExhaustion}\,d$ with
the BED constant $c = d$) out of \texttt{HighTemperatureBounds.lean}
into a new child module
\texttt{HighTemperatureBoundsFreeEnergyInfinite.lean} (under
\texttt{Concrete/LatticeGraphCorrelation/}). The legacy import path is
preserved.

PR~\#1941 further splits the concrete §18.3--§18.4 alongExhaustion
$\mathrm{ratio\_sandwich\_bundle}$ / $\mathrm{ratio\_bound}$ /
$\mathrm{triple\_ratio\_*}$ wrappers (29 theorems for the
along\nobreakdash-exhaustion variants of $\mathrm{partitionFunction}$,
$\mathrm{freeEnergy}$, and
$\mathrm{log}\,\mathrm{partitionFunction}$ at $h = 0$ with bundle
/ triple\_* / \texttt{\_of\_nonempty} variants plus ferromagnetic
counterparts) out of \texttt{HighTemperatureBounds.lean} into a new
child module
\texttt{HighTemperatureBoundsAlongExhaustionRatioBounds.lean} (under
\texttt{Concrete/LatticeGraphCorrelation/}). PR~\#1996 subsequently
narrows the 7 \texttt{triple\_ratio\_sandwich\_bundle} /
\texttt{triple\_ratio\_bound\_bundle} wrappers (J = 0 / $\beta$ = 0 /
ferromagnetic variants) out of that child into a separate
\texttt{HighTemperatureBoundsAlongExhaustionTripleRatio.lean}, and
PR~\#1997 further narrows the 14
\texttt{log\_partitionFunctionAlongExhaustion\_latticeGraph} and
\texttt{freeEnergyAlongExhaustion\_latticeGraph} ratio (+
\texttt{deviation\_pos} / \texttt{pow\_two\_lt}) wrappers out of the
same child into
\texttt{HighTemperatureBoundsAlongExhaustionRatioLogFe.lean}, both
under the same directory.
PR~\#2089 further carves the two
\texttt{ratio\_sandwich\_bundle} wrappers (general and ferromagnetic)
out of the residual
\texttt{HighTemperatureBoundsAlongExhaustionRatioBounds.lean} into
\texttt{Concrete/LatticeGraphCorrelation/HighTemperatureBoundsAlongExRatioSandwichBundle.lean}.
PR~\#2090 further carves the four
\texttt{ratio\_bound} J = 0 / $\beta$ = 0 slice wrappers (general and
ferromagnetic) out of the same residual file into
\texttt{Concrete/LatticeGraphCorrelation/HighTemperatureBoundsAlongExRatioBoundSlices.lean}.
PR~\#2091 further carves the remaining two
\texttt{ratio\_bound\_bundle} wrappers (general and ferromagnetic) out
into
\texttt{Concrete/LatticeGraphCorrelation/HighTemperatureBoundsAlongExRatioBoundBundle.lean},
emptying the old parent of theorem bodies.
The legacy import path is preserved.

PR~\#1940 further splits the concrete §18.3--§18.4 alongExhaustion
\texttt{deviation\_bound\_exp} / \texttt{continuity\_bundle} /
\texttt{deviation\_sandwich} / \texttt{relative\_sandwich} /
\texttt{deviation\_pos} / \texttt{pow\_two\_lt} /
\texttt{strict\_deviation\_bundle} wrappers at $h = 0$ (18 theorems for
$\mathrm{freeEnergyAlongExhaustion}\_\mathrm{latticeGraph}$,
$\mathrm{partitionFunctionAlongExhaustion}\_\mathrm{latticeGraph}$, and
$\mathrm{log}\_\mathrm{partitionFunctionAlongExhaustion}\_\mathrm{latticeGraph}$
plus ferromagnetic variants) out of
\texttt{HighTemperatureBounds.lean} into a new child
\texttt{Concrete/LatticeGraphCorrelation/HighTemperatureBoundsAlongExhaustionDeviation.lean}.
The legacy import path is preserved. PR~\#2084 further carves the four
\texttt{deviation\_bound\_exp} and \texttt{continuity\_bundle} wrappers
(each with a ferromagnetic variant) out into
\texttt{Concrete/LatticeGraphCorrelation/HighTemperatureBoundsAlongExDeviationContinuity.lean}.
PR~\#2085 further carves the four \texttt{deviation\_sandwich} wrappers
(free energy and $\log Z$, each with a ferromagnetic variant) out into
\texttt{Concrete/LatticeGraphCorrelation/HighTemperatureBoundsAlongExDeviationSandwich.lean}.
PR~\#2086 further carves the two $Z$ \texttt{relative\_sandwich}
wrappers (general and ferromagnetic) out into
\texttt{Concrete/LatticeGraphCorrelation/HighTemperatureBoundsAlongExRelativeSandwich.lean}.
PR~\#2087 further carves the four \texttt{deviation\_pos} /
\texttt{pow\_two\_lt} strict-deviation wrappers (free energy pair, $Z$
\texttt{pow\_two\_lt}, $\log Z$ \texttt{deviation\_pos}) out into
\texttt{Concrete/LatticeGraphCorrelation/HighTemperatureBoundsAlongExDeviationPos.lean}.
PR~\#2088 further carves the remaining four wrappers
(the two \texttt{strict\_deviation\_bundle}s plus
\texttt{pow\_two\_lt\_ferromagnetic} and
\texttt{log\_*\_deviation\_pos\_ferromagnetic}) out into
\texttt{Concrete/LatticeGraphCorrelation/HighTemperatureBoundsAlongExStrictDeviation.lean},
emptying the old parent of theorem bodies.

PR~\#1939 further splits the concrete §18.3--§18.4 alongExhaustion
sharper-exp upper-bound / sandwich / complete-summary wrappers at
$h = 0$ (17 theorems for
$\mathrm{partitionFunctionAlongExhaustion}\_\mathrm{latticeGraph}\_\mathrm{high}\_\mathrm{temp}\_\mathrm{expansion}\_h\_\mathrm{zero}\_*\_\mathrm{exp}$,
$\mathrm{freeEnergyAlongExhaustion}\_\mathrm{latticeGraph}\_\mathrm{high}\_\mathrm{temp}\_h\_\mathrm{zero}\_*\_\mathrm{exp}$, and
$\mathrm{log}\_\mathrm{partitionFunctionAlongExhaustion}\_\mathrm{latticeGraph}\_\mathrm{high}\_\mathrm{temp}\_\mathrm{expansion}\_h\_\mathrm{zero}\_*\_\mathrm{exp}$
families plus ferromagnetic variants) out of
\texttt{HighTemperatureBounds.lean} into a new child
\texttt{Concrete/LatticeGraphCorrelation/HighTemperatureBoundsAlongExhaustionExpSharper.lean}.
The legacy import path is preserved.

PR~\#1938 further splits the concrete §18.3--§18.4 alongExhaustion
correlation / partition / free-energy / sandwich / pair-singleton basic
wrappers (25 theorems for
$\mathrm{correlationAlongExhaustion}\_\mathrm{latticeGraph}$,
$\mathrm{partitionFunctionAlongExhaustion}\_\mathrm{latticeGraph}$, and
$\mathrm{freeEnergyAlongExhaustion}\_\mathrm{latticeGraph}$ at $h = 0$)
out of \texttt{HighTemperatureBounds.lean} into a new child
\texttt{Concrete/LatticeGraphCorrelation/HighTemperatureBoundsAlongExhaustionBasic.lean}.
The two \texttt{\_of\_latticeAdj} along-exhaustion variants stay in
\texttt{HighTemperatureBounds} (they call the $\Lambda$-level
\texttt{\_of\_latticeAdj} wrappers which also live there). The legacy
import path is preserved.

PR~\#1933 further splits the concrete §18.3--§18.4 basic
$\mathrm{correlation}\Lambda\_\mathrm{latticeGraph}$ high-temperature
wrappers at $h = 0$ (pair nonneg, pair $\le 1$, singleton/pair trivial-slice
vanishings at $J = 0$ and $\beta = 0$, pair sandwich, singleton/pair
ferromagnetic, singleton $= 0 \wedge \le 1$, pair+singleton bundle --
11 theorems) out of \texttt{HighTemperatureBounds.lean} into a new child
\texttt{Concrete/LatticeGraphCorrelation/HighTemperatureBoundsCorrelationBasic.lean}.
The legacy import path is preserved.

PR~\#1936 further splits the concrete §18.3--§18.4 f/Z/log Z
deviation / continuity wrappers at $h = 0$ (18 theorems:
\texttt{deviation\_bound\_exp}, \texttt{continuity\_bundle},
\texttt{deviation\_sandwich}, \texttt{relative\_sandwich},
\texttt{deviation\_pos}, \texttt{pow\_two\_lt},
\texttt{strict\_deviation\_bundle} for
$\mathrm{freeEnergy}\Lambda\_\mathrm{latticeGraph}$,
$\mathrm{partitionFunction}\Lambda\_\mathrm{latticeGraph}$, and
$\mathrm{log}\_\mathrm{partitionFunction}\Lambda\_\mathrm{latticeGraph}$ plus
ferromagnetic variants) out of \texttt{HighTemperatureBounds.lean} into
a new child
\texttt{Concrete/LatticeGraphCorrelation/HighTemperatureBoundsDeviation.lean}.
The legacy import path is preserved.

PR~\#1935 further splits the concrete §18.3--§18.4 sharper-exp
upper-bound / sandwich / complete-summary wrappers at $h = 0$
(17 theorems for the
$\mathrm{partitionFunction}\Lambda\_\mathrm{latticeGraph}\_\mathrm{high}\_\mathrm{temp}\_\mathrm{expansion}\_h\_\mathrm{zero}\_*\_\mathrm{exp}$,
$\mathrm{freeEnergy}\Lambda\_\mathrm{latticeGraph}\_\mathrm{high}\_\mathrm{temp}\_h\_\mathrm{zero}\_*\_\mathrm{exp}$, and
$\mathrm{log}\_\mathrm{partitionFunction}\Lambda\_\mathrm{latticeGraph}\_\mathrm{high}\_\mathrm{temp}\_\mathrm{expansion}\_h\_\mathrm{zero}\_*\_\mathrm{exp}$
families plus ferromagnetic variants) out of
\texttt{HighTemperatureBounds.lean} into a new child
\texttt{Concrete/LatticeGraphCorrelation/HighTemperatureBoundsExpSharper.lean}.
The legacy import path is preserved.

PR~\#1934 further splits the concrete §18.7 high-temperature
pair-correlation exponential-decay capstones (16 theorems drawn from five
capstone families \texttt{tanh\_pow\_dist} / \texttt{exp\_rate\_dist} /
\texttt{exp\_highTempExpRate\_dist} / \texttt{exp\_alpha\_dist} /
\texttt{exp\_alpha\_dist\_of\_le\_highTempExpRate} in their
$\mathrm{correlation}\Lambda\_\mathrm{latticeGraph}$ /
$\mathrm{correlationAlongExhaustion}\_\mathrm{latticeGraph}$ versions and
the ferromagnetic variants that previously lived alongside them) out of
\texttt{HighTemperatureBounds.lean} into a new child
\texttt{Concrete/LatticeGraphCorrelation/HighTemperatureBoundsDecayCapstones.lean}.
Some named-rate / monotone-rate ferromagnetic variants of
\texttt{exp\_highTempExpRate\_dist} continue to live in
\texttt{Concrete/LatticeGraphCorrelation/CorrelationDecay.lean} and are
intentionally not moved. The legacy import path is preserved.

PR~\#1571 adds ferromagnetic ($0 \le J,\, 0 < \beta$) versions of
all five layers; the bridge $0 \le \beta J$ follows from
\texttt{mul\_nonneg hbeta.le hJ}. The $\mathbb{Z}^d$ along-exhaustion form
uses the abbreviated name
\texttt{correlationAlongExhaustion\_latticeGraph\_h\_zero\_at\_pair\_le\_two\_pow\_edges\_mul\_tanh\_pow\_dist\_ferro}
(dropping \texttt{\_high\_temp} and shortening
\texttt{\_ferromagnetic} to \texttt{\_ferro}) so that the
declaration line stays within the linter's 100-character budget.

PR~\#1842 repackages the same finite-volume bound in explicit
rate form.  The abstract theorem
\texttt{correlation\_high\_temp\_h\_zero\_at\_pair\_le\_two\_pow\_edges\_mul\_exp\_rate\_dist}
states
\[
\langle \sigma_i \sigma_j \rangle_{\beta,0}
  \le 2^{|E|}
      \exp\{-(-\log(\tanh(\beta J)))\,d_G(i,j)\}
\quad (0 \le \beta J).
\]
For $\tanh(\beta J)>0$ this is the equality
\(\tanh(\beta J)^n =
\exp(n\log(\tanh(\beta J)))\), proved with
\texttt{Real.exp\_log} and \texttt{Real.log\_pow}.  For the endpoint
\(\tanh(\beta J)=0\), Lean's total \texttt{Real.log 0 = 0} makes the
right-hand side equal to \(2^{|E|}\), and the result follows from
\(0^n \le 1\).  The package also adds the corresponding
\(\Lambda\)-direct, along-exhaustion, \(\mathbb{Z}^d\), and
ferromagnetic wrappers.

PR~\#1843 adds the monotone-rate companion.  If
\(\alpha \le -\log(\tanh(\beta J))\), then
\[
\langle \sigma_i \sigma_j \rangle_{\beta,0}
  \le 2^{|E|}\exp\{-\alpha\,d_G(i,j)\}.
\]
The proof uses the PR~\#1842 exact-rate form and the elementary
monotonicity step
\[
  -(-\log(\tanh(\beta J)))\,d_G(i,j)
    \le -\alpha\,d_G(i,j),
\]
which follows from \(\alpha \le -\log(\tanh(\beta J))\) because graph
distance is nonnegative.  The package again includes
\(\Lambda\)-direct, along-exhaustion, \(\mathbb{Z}^d\), and
ferromagnetic wrappers.

PR~\#1844 names this repeated rate as
\(\texttt{highTempExpRate}\,\beta\,J = -\log(\tanh(\beta J))\).  The
named API records its nonnegativity under \(0 \le \beta J\), the
ferromagnetic bridge, and named-rate versions of the exact-rate and
monotone-rate finite-volume correlation bounds.  This keeps downstream
statements from restating the logarithmic rate expression.

PR~\#1845 lifts the named-rate API through the volume layers.  It adds
\(\Lambda\)-direct, along-exhaustion, and \(\mathbb{Z}^d\) wrappers for
both the exact named-rate bound and the named monotone-rate bound
\(\alpha \le \texttt{highTempExpRate}\,\beta\,J\).  These wrappers align
the named API with the earlier explicit-rate and monotone-rate packages.

PR~\#1846 adds the matching ferromagnetic convenience wrappers for the
named exact-rate and named monotone-rate APIs.  The along-exhaustion and
concrete \(\mathbb{Z}^d\) additions are placed in lightweight
correlation-decay modules, so incremental verification does not import
the analytic cluster-expansion stack used by the broader special-case
files.

References: Glimm--Jaffe~\cite{glimm-jaffe}, \S18.4--\S18.7;
Friedli--Velenik~\cite{friedli-velenik}, \S5; for Mayer / Ursell
combinatorics see also Fern\'andez--Fr\"ohlich--Sokal
(``Random Walks, Critical Phenomena, and Triviality in Quantum
Field Theory''), Ch.~12.

\section{Peierls argument (\texttt{Peierls.lean})}

\par\noindent\textbf{Prerequisites:} \texttt{GibbsMeasure.lean},
\texttt{Lattice.lean}.

\textbf{Role:} Proves the Peierls bound and spontaneous magnetization.
This is the most ``combinatorial'' file in the project.

\textbf{What to look for:}
\begin{itemize}
  \item \texttt{phaseBoundary}: the set of ``disagreeing'' edges
    $\partial\sigma$.
  \item \texttt{hamiltonian\_boundary}: the identity
    $H = -J(|E| - 2|\partial\sigma|)$ that connects energy to contours.
  \item \texttt{peierls\_bound}: the exponential bound on contour
    probability.
  \item \texttt{prop\_5\_4\_2\_self\_contained}: the final spontaneous
    magnetization result $0 \le 1 - \langle\sigma_i\rangle_+ \le e^{-c\beta}$.
\end{itemize}

This section formalizes the Peierls argument for the existence of phase
transitions in the Ising model (Glimm--Jaffe~\cite{glimm-jaffe}, \S5.4,
pp.~80--84).

\begin{definition}[Phase boundary]
For a configuration $\sigma$ on a finite graph $G = (V, E)$, the phase
boundary is the set of edges where neighboring spins disagree:
$\partial\sigma = \{ \{i,j\} \in E : \sigma_i \ne \sigma_j \}$.
\end{definition}

\noindent\textbf{Lean:} \texttt{phaseBoundary}, \texttt{phaseBoundarySize}

\begin{theorem}[Hamiltonian--boundary identity]
For $h = 0$,
\[
  H(\sigma) = -J \bigl(|E| - 2|\partial\sigma|\bigr).
\]
\end{theorem}

\noindent\textbf{Lean:} \texttt{hamiltonian\_boundary}

\begin{definition}[Spin flip on a subset]
For $S \subseteq V$, define $\sigma^S(i) = -\sigma(i)$ if $i \in S$,
$\sigma^S(i) = \sigma(i)$ otherwise.
\end{definition}

\noindent\textbf{Lean:} \texttt{Config.flipSet}

\begin{definition}[Cut edges]
$\mathrm{cut}(S) = \{ \{i,j\} \in E : i \in S,\; j \notin S \}$.
\end{definition}

\noindent\textbf{Lean:} \texttt{cutEdges}

\begin{lemma}[Symmetric difference]
$\partial(\sigma^S) = \partial\sigma \mathbin{\triangle} \mathrm{cut}(S)$.
When $\mathrm{cut}(S) \subseteq \partial\sigma$,
$\partial(\sigma^S) = \partial\sigma \setminus \mathrm{cut}(S)$.
\end{lemma}

\noindent\textbf{Lean:} \texttt{edgeDisagrees\_flipSet\_eq},
\texttt{phaseBoundary\_flipSet\_of\_subset}

\begin{lemma}[Energy difference]
If $\mathrm{cut}(S) \subseteq \partial\sigma$ and $h = 0$, then
$H(\sigma) - H(\sigma^S) = 2J|\mathrm{cut}(S)|$.
\end{lemma}

\noindent\textbf{Lean:} \texttt{hamiltonian\_flipSet\_diff}

\begin{proposition}[Peierls bound, Prop.~5.4.1]
For any $S \subseteq V$ with $\gamma = \mathrm{cut}(S)$,
\[
  \Pr(\gamma \subseteq \partial\sigma)
  = \frac{\sum_{\sigma:\,\gamma \subseteq \partial\sigma}
          e^{-\beta H(\sigma)}}{Z}
  \le e^{-2\beta J |\gamma|}.
\]
\end{proposition}

\begin{proof}
The map $\sigma \mapsto \sigma^S$ is an involution. For $\sigma$ with
$\gamma \subseteq \partial\sigma$,
$w(\sigma) = e^{-2\beta J|\gamma|} \cdot w(\sigma^S)$.
Summing over such $\sigma$:
\[
  \sum_{\sigma:\,\gamma \subseteq \partial\sigma} w(\sigma)
  = e^{-2\beta J|\gamma|}
    \sum_{\sigma:\,\gamma \subseteq \partial\sigma} w(\sigma^S)
  \le e^{-2\beta J|\gamma|} \cdot Z.
\]
Dividing by $Z$ gives the result.
\end{proof}

\noindent\textbf{Lean:} \texttt{peierls\_sum\_bound}, \texttt{peierls\_bound}

\begin{theorem}[Contour counting bound]
For a fixed finite box of size $n$, there exist $a, b > 0$ such that
for any site $i$ and any $r$,
\[
  \#\{S \subseteq V : i \in S,\; |\mathrm{cut}(S)| = r\} \le a \cdot b^r.
\]
For the fixed box, the trivial bound $a = 2^{|V|}$, $b = 1$ suffices.
The tighter bound $N(r) \le r^d \cdot c(d)^r$ with constants independent
of $n$ (needed for the $n \to \infty$ limit) requires lattice animal
enumeration.
\end{theorem}

\noindent\textbf{Lean:} \texttt{contourCountingBound}

\begin{lemma}[Contour sum bound]
Combining the Peierls bound with the contour count:
\[
  \sum_{\substack{S : i \in S \\ |\mathrm{cut}(S)| = r}}
  \Pr(\mathrm{cut}(S) \subseteq \partial\sigma)
  \le N(r) \cdot e^{-2\beta J r}.
\]
\end{lemma}

\noindent\textbf{Lean:} \texttt{peierls\_contour\_sum\_le}

\begin{proposition}[Spontaneous magnetization, Prop.~5.4.2]
Under $+$ boundary conditions (spins fixed to $+1$ on a boundary set
$B \subseteq V$), for $h = 0$:
\[
  \langle \mathbf{1}_{\sigma_i = -1} \rangle_+
  \le \sum_{\substack{S \ni i \\ S \cap B = \emptyset}}
      e^{-2\beta J |\mathrm{cut}(S)|}.
\]
For $\beta$ sufficiently large, the RHS is exponentially small in~$\beta$,
establishing $\langle \sigma_i \rangle_+ \to 1$ as $\beta \to \infty$.
\end{proposition}

\begin{proof}
For each $+$BC configuration $\sigma$ with $\sigma_i = -1$, the set
$S = \{j : \sigma_j = -1\}$ satisfies $i \in S$, $S \cap B = \emptyset$
(boundary sites are $+1$), and $\mathrm{cut}(S) = \partial\sigma$.
So $\mathbf{1}_{\sigma_i=-1} \le \sum_{S} \mathbf{1}_{\mathrm{cut}(S) \subseteq \partial\sigma}$.
Taking expectations and applying the restricted Peierls bound (the
flip map $\sigma \mapsto \sigma^S$ preserves $+$BC when $S \cap B = \emptyset$)
gives the result.
\end{proof}

\noindent\textbf{Lean:} \texttt{spontaneous\_magnetization\_plus}

\begin{thebibliography}{9}
\bibitem{glimm-jaffe}
J.~Glimm and A.~Jaffe,
\emph{Quantum Physics: A Functional Integral Point of View},
2nd ed., Springer, 1987.

\bibitem{ruelle}
D.~Ruelle,
\emph{Characterization of Lee-Yang polynomials},
Annals of Mathematics, 171(1):589--603, 2010.

\bibitem{harcos}
G.~Harcos,
\emph{The Lee-Yang Circle Theorem},
lecture notes based on Ruelle~\cite{ruelle},
\url{https://users.renyi.hu/~gharcos/lee-yang.pdf}.

\bibitem{friedli-velenik}
S.~Friedli and Y.~Velenik,
\emph{Statistical Mechanics of Lattice Systems:
A Concrete Mathematical Introduction},
Cambridge University Press, 2017.

\bibitem{simon}
B.~Simon,
\emph{The Statistical Mechanics of Lattice Gases, Vol.~I},
Princeton University Press, 1993.
\end{thebibliography}

\section{Ambient lattice framework (\texttt{AmbientLattice.lean})}

\par\noindent\textbf{Prerequisites:} \texttt{InfiniteVolume.lean},
\texttt{FreeEnergy.lean}.

\textbf{Role:} Introduces the genuine infinite-volume framework on
an ambient type $V : \texttt{Type}$ without a $[\texttt{Fintype}~V]$
assumption.  Finite volumes are $\Lambda : \texttt{Finset}~V$;
correlations on $\Lambda$ are defined by instantiating the existing
\texttt{IsingModel} framework on the Fintype $(\uparrow\!\Lambda :
\texttt{Type})$ (via \texttt{Finset.instFintypeCoe}) and the induced
subgraph.

\subsection{Configuration and induced subgraph}

\begin{definition}[\texttt{ConfigOn}]
For $\Lambda : \texttt{Finset}~V$, $\texttt{ConfigOn}~\Lambda :=
(\uparrow\!\Lambda) \to \texttt{Spin}$, i.e., the type of
configurations on the finite volume $\Lambda$.
\end{definition}

\begin{definition}[\texttt{inducedGraph}]
For $G : \texttt{SimpleGraph}~V$ and $\Lambda : \texttt{Finset}~V$,
$\texttt{inducedGraph}~G~\Lambda := G.\texttt{induce}~(\uparrow\!\Lambda
: \texttt{Set}~V)$ is the induced subgraph of $G$ on $\Lambda$, as a
$\texttt{SimpleGraph}~(\uparrow\!\Lambda : \texttt{Type})$.
\end{definition}

\subsection{Finite-volume observables}

\begin{definition}[\texttt{partitionFunction\LeanLambda}, \texttt{correlation\LeanLambda},
  \texttt{freeEnergy\LeanLambda}]
Each is the existing \texttt{IsingModel} counterpart applied to
$\texttt{inducedGraph}~G~\Lambda$ (with $[\texttt{Fintype}~(\texttt{inducedGraph}~G~\Lambda).\texttt{edgeSet}]$).
\end{definition}

\begin{theorem}[basic properties]
$\texttt{partitionFunction\LeanLambda}_{\mathrm{pos}}$,
$\texttt{abs\_correlation\LeanLambda}_{\mathrm{le\_one}}$,
$\texttt{correlation\LeanLambda}_{\mathrm{le\_one}}$,
$\texttt{correlation\LeanLambda}_{\mathrm{nonneg}}$ (GKS-I): all follow by
direct forwarding from the corresponding existing theorems.
\end{theorem}

\subsection{Exhaustion and convergence along exhaustion}

\begin{definition}[\texttt{Exhaustion V}]
A structure capturing an increasing sequence
$\Lambda : \mathbb{N} \to \texttt{Finset}~V$ with the ``covers any
finite set eventually'' property: for any $A : \texttt{Finset}~V$
there exists $N$ such that $A \subseteq \Lambda(n)$ for $n \ge N$.
\end{definition}

\begin{definition}[\texttt{liftFinset}]
For $A \subseteq \Lambda$, the finite set
$A.\texttt{attach}.\texttt{image}~(\lambda v \mapsto v \in \Lambda)$
lifts $A$ to a $\texttt{Finset}~(\uparrow\!\Lambda)$.  This is the
finite set on which \texttt{correlation\LeanLambda} is evaluated.
\end{definition}

\begin{definition}[\texttt{correlationAlongExhaustion}]
For an \texttt{Exhaustion} $\Lambda_\bullet$ and a finite
$A : \texttt{Finset}~V$, returns a sequence
$\mathbb{N} \to \mathbb{R}$ that equals
$\texttt{correlation\LeanLambda}~G~(\Lambda(n))~p~(\texttt{liftFinset}~A)$
once $A \subseteq \Lambda(n)$, and $0$ otherwise.
\end{definition}

\begin{theorem}[\texttt{correlationAlongExhaustion\_eventually}]
For any exhaustion and any $A : \texttt{Finset}~V$, the sequence
coincides with the lifted correlation for all $n$ sufficiently large.
\end{theorem}

\begin{theorem}[\texttt{abs\_correlationAlongExhaustion\_eventually\_le\_one}]
Eventually $|\texttt{correlationAlongExhaustion}~G~\Lambda_\bullet~p~A~n|
\le 1$, inherited from the finite-volume bound.
\end{theorem}

\subsection{Ambient subgraph monotonicity on \texorpdfstring{$\Lambda$}{Lambda}}

For a fixed finite volume $\Lambda : \texttt{Finset}~V$ and
ferromagnetic parameters, the finite-volume quantities are monotone
in the ambient interaction graph.

\begin{lemma}[\texttt{inducedGraph\_mono}]
$G_1 \le G_2 \Rightarrow G_1.\texttt{induce}~\Lambda \le
G_2.\texttt{induce}~\Lambda$.
\end{lemma}

\begin{proof}
$\texttt{SimpleGraph.induce\_adj}$ reduces adjacency in the induced
subgraph to ambient adjacency; the ambient inequality then transfers
pointwise.
\end{proof}

\begin{theorem}
The following are monotone in the ambient graph direction:
\begin{itemize}
  \item \texttt{partitionFunction\LeanLambda\_monotone\_ambient\_subgraph}:
    $Z_{G_1,\Lambda} \le Z_{G_2,\Lambda}$
  \item \texttt{correlation\LeanLambda\_monotone\_ambient\_subgraph}:
    $\langle\sigma^A\rangle_{G_1,\Lambda} \le
     \langle\sigma^A\rangle_{G_2,\Lambda}$ for any
    $A : \texttt{Finset}~(\uparrow\!\Lambda)$
  \item \texttt{freeEnergy\LeanLambda\_monotone\_ambient\_subgraph}:
    $f_{G_1,\Lambda} \le f_{G_2,\Lambda}$
\end{itemize}
\end{theorem}

\begin{proof}
\texttt{inducedGraph\_mono} gives a subgraph relation on
$(\uparrow\!\Lambda)$; the existing \texttt{*\_monotone\_subgraph}
theorems (on the Fintype $(\uparrow\!\Lambda)$) then apply directly.
\end{proof}

\begin{remark}
Convergence of \texttt{correlationAlongExhaustion} (i.e., $\exists L,
\texttt{Tendsto} \ldots \texttt{(nhds }L\texttt{)}$) is the natural
next step, requiring a compatibility argument between
\texttt{correlation\LeanLambda}~on $\Lambda_n$ and $\Lambda_{n+1}$ when
$\Lambda_n \subseteq \Lambda_{n+1}$.  This is the \emph{volume}
direction of monotonicity, which requires a configuration
factorization and is left for a follow-up.
\end{remark}

\subsection{Extension of a \texorpdfstring{$\Lambda_1$}{Lambda1}-graph to \texorpdfstring{$\Lambda_2$}{Lambda2}
  (foundations for volume monotonicity)}

For $\Lambda_1 \subseteq \Lambda_2 : \texttt{Finset}\,V$,
we construct a graph on $(\uparrow\!\Lambda_2)$ whose edges are
exactly the edges of $G.\texttt{induce}\,\Lambda_1$ embedded via
the inclusion $\uparrow\!\Lambda_1 \hookrightarrow \uparrow\!\Lambda_2$.

\begin{definition}[\texttt{extendGraphFrom\LeanLambda₁}]
$(\texttt{extendGraphFrom\LeanLambda₁}\,G\,\Lambda_1\,\Lambda_2).\texttt{Adj}\,u\,v
\;\Leftrightarrow\; u.val \in \Lambda_1 \wedge v.val \in \Lambda_1 \wedge
G.\texttt{Adj}\,u.val\,v.val$, for $u, v : \uparrow\!\Lambda_2$.
\end{definition}

\begin{theorem}[\texttt{extendGraphFrom\LeanLambda₁\_le\_induce}]
$\texttt{extendGraphFrom\LeanLambda₁}\,G\,\Lambda_1\,\Lambda_2 \le
\texttt{inducedGraph}\,G\,\Lambda_2$.
\end{theorem}

\begin{proof}
Immediate from the definition: the third conjunct in the \texttt{Adj}
of the extension graph is exactly the adjacency in
\texttt{inducedGraph}.
\end{proof}

\begin{remark}[volume direction, follow-up]
The full volume-direction monotonicity
$\langle\sigma^A\rangle_{G,\Lambda_1} \le \langle\sigma^A\rangle_{G,\Lambda_2}$
requires in addition the correlation equality on the extended graph:
$$\langle\sigma^A\rangle_{\texttt{extendGraph}\,G\,\Lambda_1\,\Lambda_2,\Lambda_2}
  = \langle\sigma^A\rangle_{\texttt{inducedGraph}\,G\,\Lambda_1,\Lambda_1}.$$
This is obtained by factoring the Boltzmann weight along the
configuration decomposition
$\texttt{Config}\,(\uparrow\!\Lambda_2) \simeq
 \texttt{Config}\,(\uparrow\!\Lambda_1) \times
 \texttt{Config}\,(\uparrow\!(\Lambda_2\setminus\Lambda_1))$
(using \texttt{Equiv.piEquivPiSubtypeProd}). The argument is
staged across several follow-up PRs.
\end{remark}

\subsection{Configuration helpers for volume monotonicity}

To prepare the factorization argument, we introduce the following
helpers:

\begin{definition}[\texttt{subtypeIncl}]
For $\Lambda_1 \subseteq \Lambda_2$, the canonical injection
$\uparrow\!\Lambda_1 \to \uparrow\!\Lambda_2$ sending
$\langle v, v \in \Lambda_1\rangle$ to $\langle v, v \in \Lambda_2\rangle$.
\end{definition}

\begin{theorem}[\texttt{subtypeIncl\_injective}]
$\texttt{subtypeIncl}\,h_{12}$ is injective.
\end{theorem}

\begin{definition}[\texttt{restrictConfig}]
Given $\sigma : \uparrow\!\Lambda_2 \to \texttt{Spin}$ and $\Lambda_1 \subseteq \Lambda_2$,
the restriction $\sigma \circ \texttt{subtypeIncl}\,h_{12}
: \uparrow\!\Lambda_1 \to \texttt{Spin}$.
\end{definition}

\begin{definition}[\texttt{\LeanLambda₁subtypeEquiv}]
The equivalence $\{x : \uparrow\!\Lambda_2 \;//\; x.\text{val} \in \Lambda_1\}
\simeq \uparrow\!\Lambda_1$, which is how one bridges the subtype used
by \texttt{piEquivPiSubtypeProd} with the natural
$\uparrow\!\Lambda_1$ type.
\end{definition}

\begin{definition}[\texttt{configEquivSubtypeProd}]
The configuration space factoring
\[
  (\uparrow\!\Lambda_2 \to \texttt{Spin}) \simeq
  (\uparrow\!\Lambda_1 \to \texttt{Spin}) \times
  (\{x : \uparrow\!\Lambda_2 \;//\; x.\text{val} \notin \Lambda_1\} \to \texttt{Spin}),
\]
obtained by composing \texttt{Equiv.piEquivPiSubtypeProd} with
\texttt{\LeanLambda₁subtypeEquiv} on the first component.
\end{definition}

\begin{theorem}[\texttt{configEquivSubtypeProd\_fst}]
The first projection of \texttt{configEquivSubtypeProd} is
\texttt{restrictConfig}.
\end{theorem}

\begin{theorem}[\texttt{edgeSpin\_subtypeIncl}]
Pointwise edge-spin preservation under the subtype inclusion lift:
for $e : \texttt{Sym2}\,(\uparrow\!\Lambda_1)$,
\[
\texttt{edgeSpin}\,\sigma\,(\texttt{Sym2.map}\,(\texttt{subtypeIncl}\,h_{12})\,e)
  = \texttt{edgeSpin}\,(\texttt{restrictConfig}\,h_{12}\,\sigma)\,e.
\]
(Generic in the coefficient field~$K$.)
\end{theorem}

\begin{proof}
\texttt{Sym2.ind} reduces to the concrete-pair case, then \texttt{simp}
unfolds \texttt{edgeSpin}, \texttt{restrictConfig}, and \texttt{subtypeIncl}
to the trivial pointwise identity.
\end{proof}

\begin{theorem}[\texttt{mem\_extendGraph\_edgeSet\_of\_mem\_induce}]
If $e \in (G.\texttt{induce}\,\Lambda_1).\texttt{edgeSet}$, then
$\texttt{Sym2.map}\,(\texttt{subtypeIncl}\,h_{12})\,e \in (\texttt{extendGraphFrom\LeanLambda₁}\,G\,\Lambda_1\,\Lambda_2).\texttt{edgeSet}$.
\end{theorem}

\begin{theorem}[\texttt{exists\_induce\_edge\_of\_extendGraph}]
Every $e \in (\texttt{extendGraphFrom\LeanLambda₁}\,G\,\Lambda_1\,\Lambda_2).\texttt{edgeSet}$ is the
image of a unique $e' \in (G.\texttt{induce}\,\Lambda_1).\texttt{edgeSet}$ under
$\texttt{Sym2.map}\,(\texttt{subtypeIncl}\,h_{12})$.
\end{theorem}

These two together with \texttt{Sym2.map.injective}\,(\texttt{subtypeIncl\_injective}) give a bijection between the edge sets.

\begin{theorem}[\texttt{extendGraph\_edgeSum\_eq}]
\[
\sum_{e \in (\texttt{extendGraphFrom\LeanLambda₁}\,G\,\Lambda_1\,\Lambda_2).\texttt{edgeFinset}}
  \texttt{edgeSpin}\,\sigma\,e
= \sum_{e' \in (\texttt{inducedGraph}\,G\,\Lambda_1).\texttt{edgeFinset}}
  \texttt{edgeSpin}\,(\texttt{restrictConfig}\,h_{12}\,\sigma)\,e'.
\]
\end{theorem}

\begin{proof}
\texttt{Finset.sum\_bij} with the map
$\texttt{Sym2.map}\,(\texttt{subtypeIncl}\,h_{12})$:
membership transfer via
\texttt{mem\_extendGraph\_edgeSet\_of\_mem\_induce},
injectivity via \texttt{Sym2.map.injective} +
\texttt{subtypeIncl\_injective},
surjectivity via \texttt{exists\_induce\_edge\_of\_extendGraph},
pointwise equality via \texttt{edgeSpin\_subtypeIncl}.
\end{proof}

\begin{theorem}[\texttt{sum\_\LeanLambda₁\_subtype\_eq}]
Site-sum reindex helper: for $\Lambda_1 \subseteq \Lambda_2$,
\[
\sum_{v : \{x : \uparrow\!\Lambda_2 \mid x.\text{val} \in \Lambda_1\}}
  \texttt{Spin.sign}\,K\,(\sigma\,v.\text{val})
= \sum_{v : \uparrow\!\Lambda_1}
  \texttt{Spin.sign}\,K\,(\texttt{restrictConfig}\,h_{12}\,\sigma\,v).
\]
\end{theorem}

\begin{proof}
\texttt{Fintype.sum\_equiv} applied to \texttt{\LeanLambda₁subtypeEquiv};
pointwise equality follows by unfolding \texttt{restrictConfig}
$= \sigma \circ \texttt{subtypeIncl}$.
\end{proof}

\begin{theorem}[\texttt{siteSum\_partition}]
Specialized \texttt{Fintype.sum\_subtype\_add\_sum\_subtype}:
\[
\sum_{v : \uparrow\!\Lambda_2} \texttt{Spin.sign}\,K\,(\sigma\,v)
  = \Bigl(\sum_{v : \{x \mid x.\text{val} \in \Lambda_1\}} \texttt{Spin.sign}\,K\,(\sigma\,v.\text{val})\Bigr)
    + \sum_{v : \{x \mid \neg (x.\text{val} \in \Lambda_1)\}} \texttt{Spin.sign}\,K\,(\sigma\,v.\text{val}).
\]
\end{theorem}

\begin{theorem}[\texttt{siteSum\_split}]
Full splitting, with the $\Lambda_1$-part expressed via
\texttt{restrictConfig}:
\[
\sum_{v : \uparrow\!\Lambda_2} \texttt{Spin.sign}\,K\,(\sigma\,v)
  = \Bigl(\sum_{v : \uparrow\!\Lambda_1} \texttt{Spin.sign}\,K\,(\texttt{restrictConfig}\,h_{12}\,\sigma\,v)\Bigr)
    + \sum_{v : \{x \mid \neg (x.\text{val} \in \Lambda_1)\}} \texttt{Spin.sign}\,K\,(\sigma\,v.\text{val}).
\]
\end{theorem}

\begin{proof}
\texttt{siteSum\_partition} + \texttt{sum\_\LeanLambda₁\_subtype\_eq}.
\end{proof}

\begin{theorem}[\texttt{hamiltonian\_extendGraph\_factor}]
The Hamiltonian on \texttt{extendGraphFrom\LeanLambda₁} decomposes:
\[
H_{\texttt{extendGraphFrom\LeanLambda₁}\,G\,\Lambda_1\,\Lambda_2,\,p}(\sigma)
 = H_{\texttt{inducedGraph}\,G\,\Lambda_1,\,p}(\texttt{restrictConfig}\,h_{12}\,\sigma)
 + \Bigl(-p.h \sum_{v : \{x \mid \neg(x.\text{val} \in \Lambda_1)\}} \texttt{Spin.sign}\,\mathbb{R}\,(\sigma\,v.\text{val})\Bigr).
\]
\end{theorem}

\begin{proof}
Unfold \texttt{hamiltonian}, \texttt{interactionEnergy},
\texttt{externalFieldEnergy}. Rewrite the edge sum via
\texttt{extendGraph\_edgeSum\_eq} and the site sum via
\texttt{siteSum\_split}. \texttt{ring} closes.
\end{proof}

\begin{theorem}[\texttt{boltzmannWeight\_extendGraph\_factor}]
The Boltzmann weight on \texttt{extendGraphFrom\LeanLambda₁} factors as:
\[
w_{G_1^\text{ext},\,p}(\sigma)
  = w_{\texttt{inducedGraph}\,G\,\Lambda_1,\,p}(\texttt{restrictConfig}\,h_{12}\,\sigma)
    \cdot \exp\Bigl(\beta h \sum_{v : C} \texttt{Spin.sign}\,\mathbb{R}\,(\sigma\,v.\text{val})\Bigr).
\]
\end{theorem}

\begin{proof}
Unfold \texttt{boltzmannWeight}, apply
\texttt{hamiltonian\_extendGraph\_factor}, and use
\texttt{Real.exp\_add}. \texttt{ring} closes the exponent algebra.
\end{proof}

\begin{lemma}[\texttt{liftFinset\_eq\_image\_subtypeIncl}]
For $A \subseteq \Lambda_1 \subseteq \Lambda_2$,
$\texttt{liftFinset}\,A\,(h_A \circ h_{12})
 = (\texttt{liftFinset}\,A\,h_A).\texttt{image}\,(\texttt{subtypeIncl}\,h_{12})$.
\end{lemma}

\begin{theorem}[\texttt{spinProduct\_lift\_eq}]
For $A \subseteq \Lambda_1 \subseteq \Lambda_2$,
$\sigma : \uparrow\!\Lambda_2 \to \texttt{Spin}$:
\[
\texttt{spinProduct}\,(\texttt{liftFinset}\,A\,(h_A \circ h_{12}))\,\sigma
  = \texttt{spinProduct}\,(\texttt{liftFinset}\,A\,h_A)\,(\texttt{restrictConfig}\,h_{12}\,\sigma).
\]
\end{theorem}

\begin{proof}
Unfold \texttt{spinProduct}, rewrite the Finset via the preceding
lemma, apply \texttt{Finset.prod\_image} using
\texttt{subtypeIncl\_injective}. The remaining equality is
definitional ($\sigma \circ \texttt{subtypeIncl} = \texttt{restrictConfig}\,\sigma$).
\end{proof}

\begin{theorem}[\texttt{restrictConfig\_configEquivSubtypeProd\_symm}]
For $(\sigma_1, \sigma_2) \in (\uparrow\!\Lambda_1 \to \texttt{Spin})
\times (C \to \texttt{Spin})$,
\[
\texttt{restrictConfig}\,h_{12}\,\bigl((\texttt{configEquivSubtypeProd}\,h_{12}).\texttt{symm}\,(\sigma_1, \sigma_2)\bigr) = \sigma_1.
\]
\end{theorem}

\begin{proof}
\texttt{ext v}, \texttt{simp} with \texttt{configEquivSubtypeProd},
\texttt{\LeanLambda₁subtypeEquiv}, \texttt{piEquivPiSubtypeProd},
\texttt{subtypeIncl}, and $v.\texttt{property}$
(which selects the $\uparrow\!\Lambda_1$ branch).
\end{proof}

\begin{theorem}[\texttt{configEquivSubtypeProd\_symm\_apply\_compl}]
On the complement subtype $v \in \{x : \uparrow\!\Lambda_2 \mid x.\text{val} \notin \Lambda_1\}$,
\[
(\texttt{configEquivSubtypeProd}\,h_{12}).\texttt{symm}\,(\sigma_1, \sigma_2)\,v.\text{val}
 = \sigma_2\,v.
\]
\end{theorem}

\begin{proof}
\texttt{simp} with the same unfoldings, now $v.\texttt{property}$
selects the complement branch.
\end{proof}

\subsection{Partition function factoring}

\begin{theorem}[\texttt{partitionFunction\_extendGraph\_factor}]
\[
Z_{\texttt{extendGraphFrom\LeanLambda₁}\,G\,\Lambda_1\,\Lambda_2,\,p}
  = Z_{\texttt{inducedGraph}\,G\,\Lambda_1,\,p} \cdot F,
\]
where $F := \sum_{\sigma_2 : C \to \texttt{Spin}}
  \exp\bigl(\beta h \sum_v \texttt{Spin.sign}\,\mathbb{R}\,(\sigma_2\,v)\bigr)$
is the \texttt{complementFactor}.
\end{theorem}

\begin{proof}
\texttt{Fintype.sum\_equiv} with $(\texttt{configEquivSubtypeProd}\,h_{12}).\texttt{symm}$
to reindex into a pair sum.  \texttt{Fintype.sum\_prod\_type} splits
into nested sums.  Each summand is rewritten via
\texttt{boltzmannWeight\_extendGraph\_factor},
\texttt{restrictConfig\_configEquivSubtypeProd\_symm}, and
\texttt{configEquivSubtypeProd\_symm\_apply\_compl}.  Finally
\texttt{Finset.sum\_mul\_sum} combines into the product.
\end{proof}

\begin{theorem}[\texttt{numerator\_extendGraph\_factor}]
For $A \subseteq \Lambda_1$,
$\texttt{numerator}_{G_1^\text{ext}}(\texttt{spinProduct}\,(\texttt{liftFinset}\,A\,\ldots))
= \texttt{numerator}_{\texttt{inducedGraph}\,G\,\Lambda_1}(\texttt{spinProduct}\,(\texttt{liftFinset}\,A\,h_A)) \cdot F$.
\end{theorem}

\begin{proof}
Same pattern as partition function factoring, with
\texttt{spinProduct\_lift\_eq} additionally applied to rewrite the
spinProduct through \texttt{restrictConfig}.
\end{proof}

\begin{theorem}[\texttt{correlation\LeanLambda\_extendGraph\_eq}]
For $A \subseteq \Lambda_1$,
$\langle\sigma^{A_2}\rangle_{G_1^\text{ext}}
 = \langle\sigma^{A_1}\rangle_{\texttt{inducedGraph}\,G\,\Lambda_1}$.
\end{theorem}

\begin{proof}
Express both correlations as $\texttt{numerator} / Z$ via
\texttt{gibbsExpectation\_eq\_div}, apply the two factoring theorems,
and cancel the common complement factor $F$ using \texttt{field\_simp}.
\end{proof}

\subsection{Volume-direction monotonicity main theorem}

\begin{theorem}[\texttt{correlation\LeanLambda\_monotone\_volume}]
For ferromagnetic $p$ and $A \subseteq \Lambda_1 \subseteq \Lambda_2 : \text{Finset}\,V$,
\[
\langle\sigma^A\rangle_{G, \Lambda_1} \le \langle\sigma^A\rangle_{G, \Lambda_2}.
\]
\end{theorem}

\begin{proof}
\begin{align*}
\langle\sigma^A\rangle_{G, \Lambda_1}
&= \langle\sigma^{A_1}\rangle_{\texttt{inducedGraph}\,G\,\Lambda_1} && \text{(def)}\\
&= \langle\sigma^{A_2}\rangle_{\texttt{extendGraphFrom\LeanLambda₁}\,G\,\Lambda_1\,\Lambda_2}
   && \text{(\texttt{correlation\LeanLambda\_extendGraph\_eq})}\\
&\le \langle\sigma^{A_2}\rangle_{\texttt{inducedGraph}\,G\,\Lambda_2}
   && \text{(\texttt{correlation\_monotone\_subgraph} + \texttt{extendGraphFrom\LeanLambda₁\_le\_induce})}\\
&= \langle\sigma^A\rangle_{G, \Lambda_2}. && \text{(def)}
\end{align*}
\end{proof}

This completes the volume-direction monotonicity on the genuine
infinite-volume (\texttt{AmbientLattice}) framework.

\begin{theorem}[\texttt{correlation\LeanLambda\_shifted\_monotone\_bounded}]
For ferromagnetic $p$, an exhaustion $\Lambda$, and $N$ such that
$A \subseteq \Lambda.\text{volume}\,n$ for $n \ge N$, the sequence
\[
  n \mapsto \langle\sigma^A\rangle_{G, \Lambda.\text{volume}(n+N)}
\]
is monotone and bounded above by $1$.
\end{theorem}

\begin{proof}
Monotonicity from \texttt{correlation\LeanLambda\_monotone\_volume} applied
via $\Lambda.\text{mono}$; upper bound from \texttt{correlation\LeanLambda\_le\_one}.
\end{proof}

The full Tendsto-convergence of \texttt{correlationAlongExhaustion}
(connecting this clean shifted sequence to the dite-based definition)
is deferred to a follow-up PR.

\begin{theorem}[\texttt{correlation\LeanLambda\_shifted\_tendsto}]
The shifted correlation sequence (monotone + bounded) converges to
its supremum by \texttt{tendsto\_atTop\_ciSup}.
\end{theorem}

An alternative approach bypasses the shifted-sequence / proof-irrelevance
route and establishes global monotonicity of
\texttt{correlationAlongExhaustion} directly.

\begin{theorem}[\texttt{correlationAlongExhaustion\_monotone}]
For a ferromagnetic Ising model and any exhaustion $\Lambda_\bullet$,
the sequence $n \mapsto \texttt{correlationAlongExhaustion}\,G\,\Lambda\,p\,A\,n$
is monotone.  Specifically, for $n \le m$:
\begin{itemize}
  \item $A \subseteq \Lambda.\text{volume}\,n$: then
    $A \subseteq \Lambda.\text{volume}\,m$ follows from $\Lambda.\text{mono}$;
    monotonicity follows from \texttt{correlation\LeanLambda\_monotone\_volume}.
  \item $A \not\subseteq \Lambda.\text{volume}\,n$: LHS is 0;
    RHS is either 0 (when $A \not\subseteq \Lambda.\text{volume}\,m$)
    or $\ge 0$ by GKS-I (\texttt{correlation\LeanLambda\_nonneg}).
\end{itemize}
\end{theorem}

\begin{theorem}[\texttt{correlationAlongExhaustion\_le\_one}]
For every $n$, $\texttt{correlationAlongExhaustion}\,G\,\Lambda\,p\,A\,n \le 1$,
either because the value is 0 or via \texttt{correlation\LeanLambda\_le\_one}.
\end{theorem}

\begin{theorem}[\texttt{correlationAlongExhaustion\_tendsto\_ciSup}]
For a ferromagnetic Ising model and any exhaustion $\Lambda_n \uparrow V$
of an ambient type $V$,
\[
\text{Tendsto}\,(\text{correlationAlongExhaustion}\,G\,\Lambda\,p\,A)\,\text{atTop}\,
  \left(\text{nhds}\,\bigsqcup_n \text{correlationAlongExhaustion}\,G\,\Lambda\,p\,A\,n\right).
\]
\end{theorem}

\begin{proof}
Combine \texttt{correlationAlongExhaustion\_monotone} and
\texttt{correlationAlongExhaustion\_le\_one} and apply
\texttt{tendsto\_atTop\_ciSup}.  The proof uses only $\Lambda.\text{mono}$;
$\Lambda.\text{exhaust}$ is not required for convergence alone.  Physical
identification of the limit with the thermodynamic-limit correlation
uses $\Lambda.\text{exhaust}$ (any $A : \texttt{Finset}~V$ eventually
lies in some $\Lambda.\text{volume}~N$).
\end{proof}

\begin{theorem}[\texttt{correlationAlongExhaustion\_convergent}]
Thin wrapper: for a ferromagnetic Ising model and any exhaustion
$\Lambda_n \uparrow V$,
\[
\exists L,\ \text{Tendsto}\,(\text{correlationAlongExhaustion}\,G\,\Lambda\,p\,A)\,\text{atTop}\,(\text{nhds}\,L).
\]
Use \texttt{correlationAlongExhaustion\_tendsto\_ciSup} directly when the
identity of $L$ as a supremum is needed.
\end{theorem}

This completes the genuine infinite-volume convergence on the
\texttt{AmbientLattice} framework.

\subsection{Infinite-volume correlation function}

Reference: Glimm--Jaffe, \textit{Quantum Physics}, 2nd ed.\ (1987),
\S 4.2 Theorem 4.2.3, pp.\ 58--59 (thermodynamic-limit correlations for
the ferromagnetic Ising model).
Friedli--Velenik \S 3.2.

\begin{definition}[\texttt{correlationInfinite}]
\[
\langle \sigma^A \rangle_V^{\mathrm{FM}} \;:=\;
  \bigsqcup_{n \in \mathbb{N}} \texttt{correlationAlongExhaustion}\,G\,\Lambda\,p\,A\,n.
\]
\end{definition}

\begin{theorem}[\texttt{tendsto\_correlationAlongExhaustion\_correlationInfinite}]
$\text{Tendsto}\,\texttt{correlationAlongExhaustion}\,\text{atTop}\,
  (\text{nhds}\,\langle \sigma^A \rangle_V^{\mathrm{FM}})$.
\end{theorem}

\begin{theorem}[\texttt{correlationInfinite\_le\_one}]
$\langle \sigma^A \rangle_V^{\mathrm{FM}} \le 1$.
\end{theorem}

\begin{proof}
\texttt{correlationAlongExhaustion\_le\_one} gives 1 as an upper bound for
the supremum.
\end{proof}

\begin{theorem}[\texttt{correlationInfinite\_nonneg}]
Under ferromagnetic assumptions,
$\langle \sigma^A \rangle_V^{\mathrm{FM}} \ge 0$.
\end{theorem}

\begin{proof}
Use $\Lambda.\text{exhaust}$ to choose
$A \subseteq \Lambda.\text{volume}\,N$, then GKS-I gives
$\texttt{correlationAlongExhaustion}\,N \ge 0$. The supremum of the
monotone sequence is also nonnegative.
\end{proof}

\begin{theorem}[\texttt{tendsto\_correlation\LeanLambda\_correlationInfinite\_of\_subset}]
Given an explicit $N$ and
$h_N : \forall n \ge N,\ A \subseteq \Lambda.\text{volume}\,n$,
\[
\text{Tendsto}\,\bigl(m \mapsto \texttt{correlation\LeanLambda}~G~\Lambda(m+N)~p~
  (\texttt{liftFinset}~A~\ldots)\bigr)\,\text{atTop}\,
  (\text{nhds}\,\langle \sigma^A \rangle_V^{\mathrm{FM}}).
\]
\end{theorem}

\begin{proof}
\texttt{tendsto\_add\_atTop\_nat} shows that the shifted
\texttt{correlationAlongExhaustion} sequence has the same limit. For
$n \ge N$, unfolding the \texttt{dite} branch identifies
\texttt{correlationAlongExhaustion} with \texttt{correlation\LeanLambda}
(\texttt{Tendsto.congr}).
\end{proof}

\begin{theorem}[\texttt{tendsto\_correlation\LeanLambda\_correlationInfinite}]
Corollary using $\Lambda.\text{exhaust}$: there exist $N$ and $h_N$, so the
preceding conclusion applies.
\end{theorem}

\subsection{Exhaustion-independence}

Reference: Glimm--Jaffe \S 4.2, pp.\ 58--59 (the infinite-volume limit is
independent of the exhaustion). Friedli--Velenik \S 3.2
(Infinite-volume Gibbs states).

\begin{lemma}[\texttt{correlationAlongExhaustion\_le\_correlationInfinite\_of\_other}]
For two exhaustions $\Lambda, \Lambda'$ and every $n$,
\[
\texttt{correlationAlongExhaustion}\,G\,\Lambda'\,p\,A\,n
  \le \texttt{correlationInfinite}\,G\,\Lambda\,p\,A.
\]
\end{lemma}

\begin{proof}
If $A \subseteq \Lambda'.\text{volume}\,n$, use
$\Lambda.\text{exhaust}(\Lambda'.\text{volume}\,n)$ to choose $m$ with
$\Lambda'.\text{volume}\,n \subseteq \Lambda.\text{volume}\,m$. Chain
\texttt{correlation\LeanLambda\_monotone\_volume} (PR \#87) with \texttt{le\_ciSup}.
If $A \not\subseteq \Lambda'.\text{volume}\,n$, then the left-hand side is
0, while \texttt{correlationInfinite\_nonneg} gives a nonnegative
right-hand side.
\end{proof}

\begin{theorem}[\texttt{correlationInfinite\_indep\_exhaustion}]
$\texttt{correlationInfinite}\,G\,\Lambda\,p\,A
  = \texttt{correlationInfinite}\,G\,\Lambda'\,p\,A$.
\end{theorem}

\begin{proof}
Apply \texttt{ciSup\_le} and the preceding lemma in both directions.
\end{proof}

\subsection{Ambient-subgraph monotonicity at infinite volume}

Reference: Glimm--Jaffe \S 4.2, pp.\ 58--59 (the thermodynamic limit
preserves finite-volume monotonicity). Friedli--Velenik \S 3.6.

\begin{lemma}[\texttt{correlationAlongExhaustion\_monotone\_ambient\_subgraph}]
If $G_1 \le G_2$, then for every $n$,
$\texttt{correlationAlongExhaustion}\,G_1\,\Lambda\,p\,A\,n
  \le \texttt{correlationAlongExhaustion}\,G_2\,\Lambda\,p\,A\,n$.
\end{lemma}

\begin{proof}
If $A \subseteq \Lambda.\text{volume}\,n$, apply
\texttt{correlation\LeanLambda\_monotone\_ambient\_subgraph}; otherwise both sides
are 0.
\end{proof}

\begin{theorem}[\texttt{correlationInfinite\_monotone\_ambient\_subgraph}]
If $G_1 \le G_2$,
$\texttt{correlationInfinite}\,G_1\,\Lambda\,p\,A
  \le \texttt{correlationInfinite}\,G_2\,\Lambda\,p\,A$.
\end{theorem}

\begin{proof}
Lemma + \texttt{le\_ciSup} + \texttt{ciSup\_le}.
\end{proof}

\subsection{GKS-II at infinite volume}

Glimm--Jaffe, \textit{Quantum Physics}, 2nd ed.\ (1987),
\S 4.2 Theorem 4.2.3, pp.\ 58--59: the second Griffiths inequality
holds in the thermodynamic limit. The finite-volume version is
Friedli--Velenik
\textit{Statistical Mechanics of Lattice Systems}, Theorem 3.49
(eq.~3.55), pp.\ 127--128.

\begin{lemma}[\texttt{liftFinset\_symmDiff}]
If $A, B \subseteq \Lambda$,
$\texttt{liftFinset}\,A\,h_A \triangle \texttt{liftFinset}\,B\,h_B
  = \texttt{liftFinset}\,(A \triangle B)\,h_{A \triangle B}$.
\end{lemma}

\begin{proof}
By extensional equality (ext) and \texttt{mem\_liftFinset}.
\end{proof}

\begin{lemma}[\texttt{correlation\LeanLambda\_gks\_second}]
If $A, B \subseteq \Lambda$,
$\texttt{correlation\LeanLambda}\,G\,\Lambda\,p\,(\text{lift}\,A)
  \cdot \texttt{correlation\LeanLambda}\,G\,\Lambda\,p\,(\text{lift}\,B)
  \le \texttt{correlation\LeanLambda}\,G\,\Lambda\,p\,(\text{lift}\,(A \triangle B))$.
\end{lemma}

\begin{proof}
Apply \texttt{IsingModel.gks\_second} to
\texttt{inducedGraph}\,$G\,\Lambda$, then rewrite with
\texttt{liftFinset\_symmDiff}.
\end{proof}

\begin{theorem}[\texttt{correlationInfinite\_gks\_second}]
\[
\texttt{correlationInfinite}\,G\,\Lambda\,p\,A
  \cdot \texttt{correlationInfinite}\,G\,\Lambda\,p\,B
  \le \texttt{correlationInfinite}\,G\,\Lambda\,p\,(A \triangle B).
\]
\end{theorem}

\begin{proof}
\texttt{Tendsto.mul} shows that the left-hand sequence converges to
$\texttt{correlationInfinite}\,A \cdot \texttt{correlationInfinite}\,B$
and the right-hand sequence converges to
$\texttt{correlationInfinite}\,(A \triangle B)$. For each $n$,
use \texttt{correlation\LeanLambda\_gks\_second} when both sets are contained, or
LHS $= 0 \le$ RHS (\texttt{correlationAlongExhaustion\_nonneg})
otherwise. This gives pointwise inequalities, and
\texttt{le\_of\_tendsto\_of\_tendsto'} gives the result.
\end{proof}

\subsection{h-direction monotonicity at infinite volume}

Glimm--Jaffe, \textit{Quantum Physics}, 2nd ed.\ (1987), \S 4.2
Infinite-volume version in the $h$ direction of Proposition 4.2.4
(Exercise), p.\ 58.

\begin{lemma}[\texttt{correlation\LeanLambda\_monotone\_h}]
For $J \ge 0$, $\beta > 0$, and finite $\Lambda \subseteq V$,
$\texttt{MonotoneOn}\,(h \mapsto \texttt{correlation\LeanLambda}\,G\,\Lambda\,\langle J,h,\beta\rangle\,A)\,(\mathrm{Ici}\,0)$.
\end{lemma}

\begin{proof}
Apply \texttt{correlation\_monotone\_h} directly to
\texttt{inducedGraph}\,$G\,\Lambda$.
\end{proof}

\begin{lemma}[\texttt{correlationAlongExhaustion\_monotone\_h}]
If $0 \le h_1 \le h_2$, then for every $n$,
$\texttt{correlationAlongExhaustion}\,G\,\Lambda\,\langle J,h_1,\beta\rangle\,A\,n
  \le \texttt{correlationAlongExhaustion}\,G\,\Lambda\,\langle J,h_2,\beta\rangle\,A\,n$.
\end{lemma}

\begin{proof}
If $A \subseteq \Lambda.\text{volume}\,n$, apply the preceding lemma;
otherwise both sides are 0.
\end{proof}

\begin{theorem}[\texttt{correlationInfinite\_monotone\_h}]
$\texttt{MonotoneOn}\,(h \mapsto \texttt{correlationInfinite}\,G\,\Lambda\,\langle J,h,\beta\rangle\,A)\,(\mathrm{Ici}\,0)$.
\end{theorem}

\begin{proof}
Use pointwise $\le$ together with \texttt{ciSup\_le} and \texttt{le\_ciSup}.
\end{proof}

\subsection{β-direction monotonicity at infinite volume}

Glimm--Jaffe, \textit{Quantum Physics}, 2nd ed.\ (1987), \S 4.2
Infinite-volume version in the $\beta$ direction of Proposition 4.2.4
(Exercise), p.\ 58.

\begin{lemma}[\texttt{correlation\LeanLambda\_monotone\_beta}]
For $J \ge 0$ and $h \ge 0$,
$\texttt{MonotoneOn}\,(\beta \mapsto \texttt{correlation\LeanLambda}\,G\,\Lambda\,\langle J,h,\beta\rangle\,A)\,(\mathrm{Ioi}\,0)$.
\end{lemma}

\begin{proof}
Apply \texttt{correlation\_monotone\_beta} to
\texttt{inducedGraph}\,$G\,\Lambda$.
\end{proof}

\begin{lemma}[\texttt{correlationAlongExhaustion\_monotone\_beta}]
For $0 < \beta_1 \le \beta_2$, pointwise monotonicity holds for every $n$.
\end{lemma}

\begin{proof}
Case split and use \texttt{correlation\LeanLambda\_monotone\_beta}; otherwise both
sides are 0.
\end{proof}

\begin{theorem}[\texttt{correlationInfinite\_monotone\_beta}]
$\texttt{MonotoneOn}\,(\beta \mapsto \texttt{correlationInfinite}\,G\,\Lambda\,\langle J,h,\beta\rangle\,A)\,(\mathrm{Ioi}\,0)$.
\end{theorem}

\begin{proof}
Use pointwise $\le$ together with \texttt{ciSup\_le} and \texttt{le\_ciSup}.
\end{proof}

\subsection{J-direction monotonicity at infinite volume}

Glimm--Jaffe, \textit{Quantum Physics}, 2nd ed.\ (1987), \S 4.2
Infinite-volume version in the $J$ direction of Proposition 4.2.4
(Exercise), p.\ 58. This completes the three-parameter monotonicity symmetry
for $(J, h, \beta)$.

\begin{lemma}[\texttt{correlation\LeanLambda\_monotone\_J}]
For $h \ge 0$ and $\beta > 0$,
$\texttt{MonotoneOn}\,(J \mapsto \texttt{correlation\LeanLambda}\,G\,\Lambda\,\langle J,h,\beta\rangle\,A)\,(\mathrm{Ici}\,0)$.
\end{lemma}

\begin{proof}
Apply \texttt{correlation\_monotone\_J} to
\texttt{inducedGraph}\,$G\,\Lambda$.
\end{proof}

\begin{lemma}[\texttt{correlationAlongExhaustion\_monotone\_J}]
For $0 \le J_1 \le J_2$, pointwise monotonicity holds for every $n$.
\end{lemma}

\begin{proof}
Case split and use \texttt{correlation\LeanLambda\_monotone\_J}; otherwise both sides
are 0.
\end{proof}

\begin{theorem}[\texttt{correlationInfinite\_monotone\_J}]
$\texttt{MonotoneOn}\,(J \mapsto \texttt{correlationInfinite}\,G\,\Lambda\,\langle J,h,\beta\rangle\,A)\,(\mathrm{Ici}\,0)$.
\end{theorem}

\begin{proof}
Use pointwise $\le$ together with \texttt{ciSup\_le} and \texttt{le\_ciSup}.
\end{proof}

\subsection{Infinite-volume single-site magnetization}

Reference: Glimm--Jaffe \S 4.2 (pp.\ 57ff), \S 5.1
(p.\ 77, $m^* := \lim_{h \to 0^+} M$);
Friedli--Velenik \S 3.2 Theorem 3.16.

\begin{definition}[\texttt{magnetizationInfinite}]
\[
M^{\mathrm{FM}}_{G,\Lambda}(J, h, \beta; i)
  \;:=\; \texttt{correlationInfinite}\,G\,\Lambda\,\langle J,h,\beta\rangle\,\{i\}.
\]
\end{definition}

\begin{theorem}[\texttt{magnetizationInfinite} basic properties]
All of the following follow immediately by specializing
$\texttt{correlationInfinite\_*}$ to $A := \{i\}$:
\begin{itemize}
\item \texttt{magnetizationInfinite\_nonneg} (ferromagnetic case),
\item \texttt{magnetizationInfinite\_le\_one},
\item \texttt{magnetizationInfinite\_indep\_exhaustion},
\item \texttt{magnetizationInfinite\_monotone\_\{J,h,beta\}} (three-parameter monotonicity).
\end{itemize}
\end{theorem}

\subsection*{Z₂ spin-flip symmetry at $h = 0$}

Reference: Glimm--Jaffe \S 5.1 p.\ 77: spontaneous magnetization $m^*$ is
the $h \to 0^+$ limit with boundary-condition dependence, not the value at
$h = 0$ itself. The finite-volume input is
\texttt{IsingModel.correlation\_odd\_vanish} (\texttt{Inequalities/GHS.lean}).

\begin{theorem}[\texttt{correlationInfinite\_h\_zero}]
If $h = 0$ and $|A|$ is odd,
$\texttt{correlationInfinite}\,G\,\Lambda\,\langle J,0,\beta\rangle\,A = 0$.
\end{theorem}

\begin{proof}
Apply \texttt{correlation\_odd\_vanish} in finite volume to
\texttt{inducedGraph} (\texttt{correlation\LeanLambda\_odd\_vanish\_h\_zero}). For
each $n$, \texttt{correlationAlongExhaustion} is 0 in both \texttt{dite}
branches.
\texttt{ciSup} of zero sequence $= 0$.
\end{proof}

\begin{theorem}[\texttt{magnetizationInfinite\_zero\_at\_h\_zero}]
At $h = 0$, $M^{\mathrm{FM}} = 0$, because the singleton has odd cardinality
$|\{i\}| = 1$.
\end{theorem}

\subsection{Spontaneous correlation function (general \texorpdfstring{$A$}{A})}

Reference: Glimm--Jaffe, \textit{Quantum Physics}, 2nd ed.\ (1987),
\S 4.2 p.\ 57ff, \S 5.1 p.\ 77:
$m^*(J, \beta) := \lim_{h \to 0^+} M(J, h, \beta)$.
Friedli--Velenik \S 3.10.

Define the $h \to 0^+$ limit of the infinite-volume correlation
$\langle \sigma^A \rangle$ for a general $A : \texttt{Finset}\,V$, and record
its basic properties. The single-site case $A = \{i\}$ is the spontaneous
magnetization $m^*$ (Glimm--Jaffe \S 5.1 p.\ 77).

\begin{definition}[\texttt{spontaneousCorrelation}]
\[
\langle \sigma^A \rangle^*_{G, \Lambda}(J, \beta)
  \;:=\; \inf_{h > 0} \texttt{correlationInfinite}\,G\,\Lambda\,\langle J,h,\beta\rangle\,A.
\]
\end{definition}

\begin{theorem}[\texttt{spontaneousCorrelation} basic properties]
Under ferromagnetic assumptions, five properties follow by specializing
\texttt{correlationInfinite\_*}:
\texttt{\_nonneg},
\texttt{\_le\_one},
\texttt{\_le\_correlationInfinite},
\texttt{\_indep\_exhaustion},
\texttt{tendsto\_\ldots\_nhdsGT}.
\end{theorem}

\subsection*{Spontaneous magnetization (single-site specialization)}

Spontaneous magnetization $m^*$ is the $A = \{i\}$ specialization of
\texttt{spontaneousCorrelation}:

\begin{definition}[\texttt{spontaneousMagnetization}]
\[
\texttt{spontaneousMagnetization}\,G\,\Lambda\,J\,\beta\,i
  \;:=\; \texttt{spontaneousCorrelation}\,G\,\Lambda\,J\,\beta\,\{i\}.
\]
\end{definition}

\begin{theorem}[\texttt{spontaneousCorrelation\_singleton\_eq\_spontaneousMagnetization}]
This is equivalent by \texttt{rfl} from the definitions above.
\end{theorem}

\begin{theorem}[\texttt{spontaneousMagnetization} basic properties]
All are $A := \{i\}$ specializations of \texttt{spontaneousCorrelation\_*}
(1-line specialization):
\begin{itemize}
\item \texttt{\_nonneg}: $m^* \ge 0$,
\item \texttt{\_le\_one}: $m^* \le 1$,
\item \texttt{\_le\_magnetizationInfinite}: for $h > 0$, $m^* \le M(h)$,
\item \texttt{\_indep\_exhaustion}: independent of $\Lambda$,
\item \texttt{tendsto\_magnetizationInfinite\_spontaneousMagnetization\_nhdsGT}:
      $m^* = \lim_{h \to 0^+} M(h)$.
\end{itemize}
\end{theorem}

\subsection{Truncated 2-point correlation at infinite volume}

Reference: Glimm--Jaffe, \textit{Quantum Physics}, 2nd ed.\ (1987),
\S 4.2 p.\ 57ff (GKS-II); Friedli--Velenik \S 3.6.3.

Specialize \texttt{correlationInfinite\_gks\_second} from PR \#94 to the
two-point case.

\begin{definition}[\texttt{truncated2Infinite}]
\[
U_2(G, \Lambda; p; i, j)
  \;:=\; \langle \sigma_i \sigma_j \rangle_\infty
    - \langle \sigma_i \rangle_\infty \langle \sigma_j \rangle_\infty.
\]
\end{definition}

\begin{theorem}[\texttt{truncated2Infinite} basic properties]
\begin{itemize}
\item \texttt{\_symm}: $U_2(i, j) = U_2(j, i)$ (\texttt{Finset.pair\_comm}).
\item \texttt{\_nonneg\_of\_ne}: $i \ne j \Rightarrow U_2 \ge 0$
      (via \texttt{correlationInfinite\_gks\_second},
      $\{i,j\} = \{i\} \triangle \{j\}$).
\item \texttt{\_nonneg\_of\_eq}: $U_2(i, i) = M(i)(1 - M(i)) \ge 0$
      ($\{i, i\} = \{i\}$, $M \in [0, 1]$).
\item \texttt{\_nonneg}: for arbitrary $i, j$, $0 \le U_2(i, j)$, combining
      the two preceding cases.
\item \texttt{\_indep\_exhaustion}: independent of $\Lambda$.
\item \texttt{\_h\_zero}: at $h = 0$,
      $U_2 = \langle \sigma_i \sigma_j \rangle_\infty$ for arbitrary
      $i, j$, since singletons vanish by Z₂ via
      $\texttt{correlationInfinite\_h\_zero}$.
\end{itemize}
\end{theorem}

\subsection{Truncated 3-point correlation + GHS at infinite volume}

Reference: Glimm--Jaffe, \textit{Quantum Physics}, 2nd ed.\ (1987),
\S 4.3 Corollary 4.3.4 GHS, pp.\ 68ff.
Friedli--Velenik \S 3.6.4 (higher truncated functions).

\begin{definition}[\texttt{truncated3Infinite}]
$U_3 := \langle \sigma^{\{i,j,k\}} \rangle_\infty
  - \langle \sigma_i \rangle_\infty \langle \sigma^{\{j,k\}} \rangle_\infty
  - \langle \sigma_j \rangle_\infty \langle \sigma^{\{i,k\}} \rangle_\infty
  - \langle \sigma_k \rangle_\infty \langle \sigma^{\{i,j\}} \rangle_\infty
  + 2 \langle \sigma_i \rangle_\infty \langle \sigma_j \rangle_\infty \langle \sigma_k \rangle_\infty$.
\end{definition}

\begin{theorem}[\texttt{truncated3Infinite\_nonpos} — GHS at infinite volume]
For a ferromagnetic system and pairwise distinct $i, j, k$, $U_3 \le 0$.
\end{theorem}

\begin{proof}
Apply the finite-volume \texttt{ghs\_inequality} eventually, once
$n \ge N$ gives $\{i,j,k\} \subseteq \Lambda.\text{volume}\,n$. The five
correlation terms in $U_3$ are all assembled using
$\texttt{correlationAlongExhaustion} \to \texttt{correlationInfinite}$
(\texttt{Tendsto}), and \texttt{le\_of\_tendsto} gives the limit inequality
$\le 0$.
\end{proof}

\begin{theorem}[\texttt{truncated3Infinite\_h\_zero\_of\_distinct}]
At $h = 0$ and for pairwise distinct sites, $U_3 = 0$. All singleton terms
are odd, and $\{i, j, k\}$ is also odd when the sites are distinct, so every
term vanishes.
\end{theorem}

\begin{theorem}[\texttt{truncated3Infinite\_indep\_exhaustion}]
Independent of $\Lambda$, using \texttt{correlationInfinite\_indep\_exhaustion}
in seven places.
\end{theorem}

\subsection{Truncated 4-point correlation + \texorpdfstring{$U_4 \le 0$}{U4 <= 0} at \texorpdfstring{$h = 0$}{h = 0}}

Reference: Glimm--Jaffe, \textit{Quantum Physics}, 2nd ed.\ (1987),
\S 4.3 Corollary 4.3.3, pp.\ 68ff.

\begin{definition}[\texttt{truncated4Infinite}]
\begin{align*}
U_4 &:= \langle \sigma^{\{i,j,k,l\}} \rangle_\infty
  - \langle \sigma^{\{i,j\}} \rangle_\infty \langle \sigma^{\{k,l\}} \rangle_\infty \\
  &\quad - \langle \sigma^{\{i,k\}} \rangle_\infty \langle \sigma^{\{j,l\}} \rangle_\infty
  - \langle \sigma^{\{i,l\}} \rangle_\infty \langle \sigma^{\{j,k\}} \rangle_\infty.
\end{align*}
\end{definition}

\begin{theorem}[\texttt{truncated4Infinite\_nonpos\_h\_zero}]
At $h = 0$, for a ferromagnetic system and four pairwise distinct sites,
$U_4 \le 0$.
\end{theorem}

\begin{proof}
Apply the finite-volume \texttt{cor\_4\_3\_3} eventually, once the four sites
are contained in $\Lambda.\text{volume}\,n$. Use \texttt{Tendsto.sub/mul} for
$\texttt{truncated4AlongExhaustion} \to \texttt{truncated4Infinite}$.
\texttt{le\_of\_tendsto} gives the limit inequality $\le 0$.
\end{proof}

\begin{theorem}[\texttt{truncated4Infinite\_indep\_exhaustion}]
Independent of $\Lambda$, in seven places.
\end{theorem}

\subsection{Parameter monotonicity of spontaneous*}

Reference: Glimm--Jaffe \S 5.1 p.\ 77.
PR \#95-\#97 (correlationInfinite J/h/$\beta$) + PR \#99/\#101 (spontaneous).

\begin{theorem}[\texttt{spontaneousCorrelation\_monotone\_J}]
With $\beta > 0$ fixed,
$\texttt{MonotoneOn}\,(J \mapsto \langle \sigma^A \rangle^*)\,(\mathrm{Ici}\,0)$.
Use $\texttt{correlationInfinite\_monotone\_J}$ for each $h > 0$, then
\texttt{ciInf\_mono}.
\end{theorem}

\begin{theorem}[\texttt{spontaneousCorrelation\_monotone\_beta}]
With $J \ge 0$ fixed,
$\texttt{MonotoneOn}\,(\beta \mapsto \langle \sigma^A \rangle^*)\,(\mathrm{Ioi}\,0)$.
\end{theorem}

\begin{theorem}[\texttt{spontaneousMagnetization\_monotone\_J}]
Single-site $A = \{i\}$ specialization of
\texttt{spontaneousCorrelation\_monotone\_J}.
\end{theorem}

\begin{theorem}[\texttt{spontaneousMagnetization\_monotone\_beta}]
Single-site $A = \{i\}$ specialization of
\texttt{spontaneousCorrelation\_monotone\_beta}.
\end{theorem}

\subsection{FKG-form correlation inequality at infinite volume (nominal \S 4.4)}

Reference: Glimm--Jaffe \S 4.4 p.\ 67 (general FKG inequality, requiring
monotonicity).

Important caveat: \texttt{spinProduct} observables are not generally
monotone. For example, when $|A| = 2$, flipping both spins keeps
$\sigma^A$ at $+1$, but an intermediate configuration can make it $-1$.
Thus the general FKG form cannot be applied directly to spin products.

The following is an alias theorem that exposes the same numerical inequality
(GKS-II) under an FKG-oriented name for
(searchability + \S 4.4 coverage).

\begin{theorem}[\texttt{correlationInfinite\_fkg\_spinProduct}]
For ferromagnetic systems,
$\langle \sigma^A \rangle_\infty \langle \sigma^B \rangle_\infty
  \le \langle \sigma^{A \triangle B} \rangle_\infty$.
\end{theorem}

\begin{proof}
This is exactly \texttt{correlationInfinite\_gks\_second} (PR \#94).
GKS-II is proved through HNC / Boltzmann-weight log-supermodularity, not the
FKG lattice-condition argument.
\end{proof}

\subsection{freeEnergyAlongExhaustion (§4.6 Prop 4.6.1 scaffold)}

Reference: Glimm--Jaffe \S 4.6 Proposition 4.6.1, pp.\ 78ff
(free-energy convergence in the volume direction).

\begin{definition}[\texttt{freeEnergyAlongExhaustion}]
\[
\texttt{freeEnergyAlongExhaustion}\,G\,\Lambda\,p\,n
  \;:=\; \texttt{freeEnergy\LeanLambda}\,G\,(\Lambda.\text{volume}\,n)\,p.
\]
\end{definition}

\begin{definition}[\texttt{partitionFunctionAlongExhaustion}]
\[
\texttt{partitionFunctionAlongExhaustion}\,G\,\Lambda\,p\,n
  \;:=\; \texttt{partitionFunction\LeanLambda}\,G\,(\Lambda.\text{volume}\,n)\,p.
\]
$0 <$ positivity trivial from \texttt{partitionFunction\LeanLambda\_pos}.
\end{definition}

\begin{theorem}[\texttt{freeEnergyAlongExhaustion\_monotone\_ambient\_subgraph}]
If $G_1 \le G_2$, pointwise monotonicity holds for every $n$. Specialize
\texttt{freeEnergy\LeanLambda\_monotone\_ambient\_subgraph} to each
$\Lambda.\text{volume}\,n$.
\end{theorem}

\begin{theorem}[\texttt{partitionFunctionAlongExhaustion\_monotone\_ambient\_subgraph}]
If $G_1 \le G_2$, then for every $n$,
$Z_{G_1}(\Lambda_n) \le Z_{G_2}(\Lambda_n)$.
\end{theorem}

\begin{theorem}[\texttt{freeEnergyAlongExhaustion\_eq\_log\_div\_card}]
\[
f_n = |\Lambda.\text{volume}\,n|^{-1} \cdot \log Z_n.
\]
Unfold the definition of \texttt{IsingModel.freeEnergy}.
\end{theorem}

\begin{theorem}[\texttt{log\_partitionFunctionAlongExhaustion\_tendsto\_atTop} / \texttt{partitionFunctionAlongExhaustion\_tendsto\_atTop}]
For any exhaustion $\Lambda$ of an infinite lattice $V$ (typeclass
\texttt{[Infinite V]}) and ferromagnetic parameters
$p = \langle J, h, \beta \rangle$ with $J, h \ge 0$ and $\beta > 0$,
respectively,
\[
  \lim_{n \to \infty} \log Z_{\Lambda_n} \;=\; +\infty
\]
and
\[
  \lim_{n \to \infty} Z_{\Lambda_n} \;=\; +\infty.
\]
\end{theorem}

\begin{proof}
\texttt{Exhaustion.eventually\_volume\_nonempty} gives
$\Lambda_n \neq \emptyset$ for all sufficiently large $n$. In that range,
the existing APIs
\[
  \log 2 \le \texttt{freeEnergy\LeanLambda}\,G\,\Lambda_n\,p
  \qquad\text{(\texttt{freeEnergy\LeanLambda\_ge\_log\_two})}
\]
\[
  |\Lambda_n| \cdot \texttt{freeEnergy\LeanLambda}\,G\,\Lambda_n\,p
    = \log Z_{\Lambda_n}
  \qquad\text{(\texttt{card\_mul\_freeEnergy\LeanLambda\_eq\_log\_partitionFunction\LeanLambda\_of\_nonempty})}
\]
imply $|\Lambda_n| \cdot \log 2 \le \log Z_{\Lambda_n}$.
\texttt{Exhaustion.tendsto\_card\_atTop} (PR \#160) gives
$|\Lambda_n| \to \infty$. Since $\log 2 > 0$, the left-hand side tends to
$\infty$, and \texttt{Filter.tendsto\_atTop\_mono'} gives
$\log Z_{\Lambda_n} \to \infty$. The result
$Z_{\Lambda_n} \to \infty$ transfers by \texttt{Real.tendsto\_exp\_atTop.comp}
and \texttt{Real.exp\_log} using $Z_{\Lambda_n} > 0$.
\end{proof}

The Fekete version of free-energy-density convergence, which derives
convergence to a concrete value from super-additivity of weighted free
energy, requires translation invariance and is treated separately.

\subsubsection{Bounded edge density and uniform upper bound (PR \#123)}

\begin{definition}[\texttt{BoundedEdgeDensity}]
\[
  \exists c \in \mathbb R,\ \forall n,\
    \Lambda.\text{volume}\,n \neq \emptyset \implies
    |E(G[\Lambda_n])| \le c \cdot |\Lambda_n|.
\]
In a bounded-degree environment with $\max\deg = \Delta$, take $c = \Delta/2$.
\end{definition}

\begin{theorem}[\texttt{freeEnergyAlongExhaustion\_le\_uniform\_upper\_bound}]
Under \texttt{BoundedEdgeDensity} with constant $c$, for arbitrary $n$ with
$\Lambda_n \neq \emptyset$,
\[
  f_n \le \log 2 + |\beta| \bigl(|J| \cdot c + |h|\bigr).
\]
Substitute $|E_n|/|\Lambda_n| \le c$ into
\texttt{freeEnergyAlongExhaustion\_upper\_bound} (PR \#122).
\end{theorem}

\begin{theorem}[\texttt{BddAbove\_freeEnergyAlongExhaustion\_range}]
Under \texttt{BoundedEdgeDensity},
$\operatorname{BddAbove}(\operatorname{range}(\texttt{freeEnergyAlongExhaustion}\,G\,\Lambda\,p))$.
for empty $\Lambda_n$, $f_n = 0$ by definition. Hence
$\max\{0, \log 2 + |\beta|(|J|c+|h|)\}$ is a uniform upper bound.
\end{theorem}

\subsection{Cor 4.3.5 (inductive n-point at \texorpdfstring{$h = 0$}{h = 0}) at infinite volume}

Reference: Glimm--Jaffe, \textit{Quantum Physics}, 2nd ed.\ (1987),
\S 4.3 Corollary 4.3.5, p.\ 62.

\begin{theorem}[\texttt{correlationInfinite\_cor\_4\_3\_5\_h0}]
For a ferromagnetic system $(J \ge 0, \beta > 0)$ with $h = 0$,
$S : \texttt{Finset}\,V$, $j, k : V$, $j \notin S$, $k \notin S$, and
$j \ne k$,
\[
\langle \sigma^{\{j, k\} \cup S} \rangle_\infty
  \le \langle \sigma^S \rangle_\infty \langle \sigma^{\{j, k\}} \rangle_\infty
  + \sum_{T \in \mathcal{P}(S)}
    \langle \sigma^{\{j\} \cup T} \rangle_\infty
    \langle \sigma^{\{k\} \cup (S \setminus T)} \rangle_\infty.
\]
\end{theorem}

\begin{proof}
The along-exhaustion left-hand and right-hand sequences satisfy
\texttt{Tendsto}
(\texttt{tendsto\_finset\_sum} + \texttt{Tendsto.mul/add}).
For each $n$ such that $\{j, k\} \cup S \subseteq \Lambda.\text{volume}\,n$,
apply the finite-volume \texttt{cor\_4\_3\_5\_h0} theorem to
\texttt{inducedGraph}\,$G\,\Lambda(n)$. Use
\texttt{liftFinset\_insert} and \texttt{liftFinset\_sdiff} (PR \#107) for
compound set identifications.
\texttt{Finset.sum\_bij} reindexes $S.powerset \leftrightarrow S'.powerset$
($S' := \texttt{liftFinset}\,S$).
For other $n$, the left-hand side is $0 \le$ the right-hand side, since each
term is nonnegative. Finish with \texttt{le\_of\_tendsto\_of\_tendsto'}.
\end{proof}

\subsection{Edge-set decomposition for disjoint sum graphs (\texttt{SumGraph.lean})}

The file \texttt{IsingModel/SumGraph.lean} records the combinatorial
identities for \texttt{SimpleGraph.sum} ($G \oplus_g H$) that underlie
the super-additivity argument of Glimm--Jaffe \S 4.6 (pp.\ 70ff.):
the free-energy density on a disjoint union of lattices decomposes
additively into the two factor contributions, which in turn rests on
a disjoint splitting of the edge set.

\begin{lemma}[\texttt{SimpleGraph.sum\_eq\_map\_sup}]
For $G : \mathrm{SimpleGraph}\,V$ and $H : \mathrm{SimpleGraph}\,W$,
\[
  G \oplus_g H
    = G.\mathrm{map}(\mathrm{inl}_{V,W})
      \sqcup H.\mathrm{map}(\mathrm{inr}_{V,W}),
\]
where $\mathrm{inl}_{V,W} : V \hookrightarrow V \oplus W$ and
$\mathrm{inr}_{V,W} : W \hookrightarrow V \oplus W$ are the canonical
embeddings.
\end{lemma}

\begin{proof}
Case analysis on $a, b \in V \oplus W$: both sides evaluate
identically on \texttt{Adj}.
\end{proof}

\begin{lemma}[\texttt{SimpleGraph.edgeSet\_sum}]
\[
  (G \oplus_g H).E
    = \mathrm{inl}_*(G.E) \,\cup\, \mathrm{inr}_*(H.E),
\]
where $\mathrm{inl}_*, \mathrm{inr}_*$ denote the induced maps on
$\mathrm{Sym}^2$. The two summands are disjoint
(\texttt{disjoint\_sumInlEmb\_sumInrEmb\_edgeSet}).
\end{lemma}

\begin{proof}
Apply \texttt{edgeSet\_sup} and \texttt{edgeSet\_map} to the previous
lemma. Disjointness follows from $\mathrm{inl}\,v \ne \mathrm{inr}\,w$.
\end{proof}

\begin{proposition}[\texttt{SimpleGraph.card\_edgeFinset\_sum}]
If $V, W$ are finite and $G, H$ have decidable adjacency, then
\[
  \#(G \oplus_g H).\mathrm{edgeFinset}
    = \#G.\mathrm{edgeFinset} + \#H.\mathrm{edgeFinset}.
\]
\end{proposition}

\begin{proof}
Disjoint union of Finset images. The Finset identity
\texttt{edgeFinset\_sum} is derived from the set-level
\texttt{edgeSet\_sum} by \texttt{Finset.coe\_injective}. Cardinality
then follows from \texttt{Finset.card\_union\_of\_disjoint} combined
with \texttt{Finset.card\_map} applied to the injective
$\mathrm{Embedding.sym2Map}$.
\end{proof}

\paragraph{Context.} This file is the combinatorial Step~1 toward the
Glimm--Jaffe \S 4.6 super-additivity proof of thermodynamic-limit
convergence of the free-energy density. Subsequent steps
(\texttt{Config}-product equivalence, Hamiltonian splitting,
partition-function multiplicativity, and $\log Z$ super-additivity)
will consume the present lemmas to lift the Fekete-type convergence
of $f_{\Lambda}$ from the current \emph{discretized} subgraph
setting (\texttt{freeEnergy\_convergent\_subgraph}) to the genuine
infinite-volume exhaustion of the ambient lattice framework.

\subsection{Hamiltonian additivity on the disjoint sum graph (\texttt{SumModel.lean})}

With the combinatorial edge-set decomposition from \texttt{SumGraph.lean}
in hand (\S 0099), the Ising Hamiltonian on the disjoint sum graph
splits additively with respect to the \texttt{Sum.elim} decomposition
of configurations.

\begin{lemma}[\texttt{IsingModel.Config.sumEquiv}]
$\mathrm{Config}(\iota \oplus \iota')
  \simeq \mathrm{Config}\,\iota \times \mathrm{Config}\,\iota'$,
given by \texttt{Equiv.sumArrowEquivProdArrow}.
\end{lemma}

\begin{lemma}[\texttt{externalFieldEnergy\_sum}]
\[
  \sum_{i \in \iota \oplus \iota'}
    \mathrm{sign}(\mathrm{Sum.elim}\,\sigma_1\,\sigma_2\,i)
  = \sum_{i \in \iota} \mathrm{sign}(\sigma_1\,i)
    + \sum_{i \in \iota'} \mathrm{sign}(\sigma_2\,i),
\]
proved by \texttt{Fintype.sum\_sum\_type} together with the evaluation
of \texttt{Sum.elim} on \texttt{inl}/\texttt{inr}.
\end{lemma}

\begin{lemma}[\texttt{edgeSpin\_sumInl\_sym2Map} /
\texttt{edgeSpin\_sumInr\_sym2Map}]
For any $e \in \mathrm{Sym}^2\,\iota$,
\[
  \mathrm{edgeSpin}(\mathrm{Sum.elim}\,\sigma_1\,\sigma_2)
    (\mathrm{Sym}^2.\mathrm{map}\,\mathrm{Sum.inl}\,e)
  = \mathrm{edgeSpin}\,\sigma_1\,e,
\]
and analogously for \texttt{Sum.inr}.
\end{lemma}

\begin{proposition}[\texttt{interactionEnergy\_sum}]
For $G, H$ with finite edge sets, $J \in K$, and
$\sigma_1, \sigma_2$ configurations,
\[
  \mathrm{interactionEnergy}(G \oplus_g H)\,J\,(\mathrm{Sum.elim}\,\sigma_1\,\sigma_2)
  = \mathrm{interactionEnergy}\,G\,J\,\sigma_1
    + \mathrm{interactionEnergy}\,H\,J\,\sigma_2.
\]
\end{proposition}

\begin{proof}
Combine the edge-finset decomposition \texttt{edgeFinset\_sum} (PR \#134)
with \texttt{Finset.sum\_union} (using
\texttt{disjoint\_inl\_inr\_edgeFinset} for the disjointness side
condition), then \texttt{Finset.sum\_map} and the per-edge pullbacks
\texttt{edgeSpin\_sumInl\_sym2Map} / \texttt{edgeSpin\_sumInr\_sym2Map}.
\end{proof}

\begin{theorem}[\texttt{hamiltonian\_sum}]
For $G, H$ with finite edge sets, and parameters $p = (J, h, \beta)$,
\[
  \mathcal{H}_{G \oplus_g H}(p)(\mathrm{Sum.elim}\,\sigma_1\,\sigma_2)
  = \mathcal{H}_G(p)(\sigma_1) + \mathcal{H}_H(p)(\sigma_2).
\]
\end{theorem}

\begin{proof}
By additivity of each summand in
$\mathcal{H} = \mathrm{interactionEnergy} + \mathrm{externalFieldEnergy}$
(previous proposition and lemma).
\end{proof}

\paragraph{Step 4: partition function multiplicativity.}

\begin{proposition}[\texttt{IsingModel.partitionFunction\_sum}]
For finite lattices $G, H$ and parameters $p = (J, h, \beta)$ over $\mathbb{R}$,
\[
  Z_{G \oplus_g H}(p) = Z_G(p) \cdot Z_H(p).
\]
\end{proposition}

\begin{proof}
Unfolding $Z = \sum_\sigma e^{-\beta \mathcal{H}}$,
\begin{align*}
  Z_{G \oplus_g H}(p)
  &= \sum_{\sigma : \iota \oplus \iota' \to \mathrm{Spin}}
       e^{-\beta \mathcal{H}_{G \oplus_g H}(\sigma)}
     && \text{(\texttt{Equiv.sum\_comp}\,\texttt{Config.sumEquiv.symm})} \\
  &= \sum_{\sigma_1, \sigma_2}
       e^{-\beta (\mathcal{H}_G(\sigma_1) + \mathcal{H}_H(\sigma_2))}
     && \text{(\texttt{Fintype.sum\_prod\_type}, \texttt{hamiltonian\_sum})} \\
  &= \sum_{\sigma_1, \sigma_2}
       e^{-\beta \mathcal{H}_G(\sigma_1)}
       \cdot e^{-\beta \mathcal{H}_H(\sigma_2)}
     && \text{(\texttt{mul\_add}, \texttt{Real.exp\_add})} \\
  &= Z_G(p) \cdot Z_H(p)
     && \text{(\texttt{Finset.sum\_mul\_sum}).}
\end{align*}
\end{proof}

\begin{corollary}[\texttt{IsingModel.log\_partitionFunction\_sum}]
\[
  \log Z_{G \oplus_g H}(p) = \log Z_G(p) + \log Z_H(p).
\]
Immediate from the proposition and
\texttt{partitionFunction\_pos} / \texttt{Real.log\_mul}.
\end{corollary}

\paragraph{Step 5 preparation: super-additivity inequality.}

\begin{proposition}[\texttt{IsingModel.partitionFunction\_mul\_le\_of\_sum\_le}]
Let $G'$ be any finite simple graph on $\iota \oplus \iota'$
satisfying $G \oplus_g H \le G'$. Then for ferromagnetic parameters
$p$,
\[
  Z_G(p) \cdot Z_H(p) \le Z_{G'}(p).
\]
\end{proposition}

\begin{proof}
By Step 4, $Z_{G \oplus_g H}(p) = Z_G(p) \cdot Z_H(p)$. The
existing \texttt{partitionFunction\_monotone\_subgraph} (proved in
\texttt{FreeEnergy.lean} via GKS-I) gives
$Z_{G \oplus_g H}(p) \le Z_{G'}(p)$ under ferromagnetic parameters.
\end{proof}

\begin{corollary}[\texttt{IsingModel.log\_partitionFunction\_add\_le\_of\_sum\_le}]
Under the same assumptions,
\[
  \log Z_G(p) + \log Z_H(p) \le \log Z_{G'}(p).
\]
\end{corollary}

\begin{proof}
Positivity of partition functions and \texttt{Real.log\_mul} /
\texttt{Real.log\_le\_log}.
\end{proof}

\paragraph{Closing Step 5.} Combining this super-additivity with the
uniform upper bound (PR \#122, \#123) and Fekete's lemma produces
convergence of
$f_\Lambda = (\log Z_\Lambda)/|\Lambda|$ along arbitrary
exhaustions of the ambient lattice, completing the proof of
Glimm--Jaffe \S 4.6 Prop 4.6.1. The remaining technical piece is
the type-level identification of an induced graph on
$\Lambda_1 \cup \Lambda_2$ with a supergraph of
$\mathrm{inducedGraph}\,G\,\Lambda_1 \oplus_g \mathrm{inducedGraph}\,G\,\Lambda_2$,
which is the subject of a subsequent PR.

\subsection{Partition function invariance under graph isomorphism
(\texttt{PartitionFunctionIso.lean})}

\begin{proposition}[\texttt{IsingModel.partitionFunction\_map\_equiv}]
Let $e \colon V \simeq W$ be a type equivalence and
$G \colon \mathrm{SimpleGraph}\,V$. Then the pushforward
$G.\mathrm{map}\,e.\mathrm{toEmbedding}$ has the same partition function:
\[
  Z_{G.\mathrm{map}\,e}(p) = Z_G(p).
\]
\end{proposition}

\begin{proof}
Reindex $\sum_{\tau : W \to \mathrm{Spin}} \exp(-\beta \mathcal{H}_{G.\mathrm{map}\,e}(\tau))$ via
\texttt{Equiv.arrowCongr}\,$e$\,\texttt{Equiv.refl}, which puts it in
bijection with $\sum_{\sigma : V \to \mathrm{Spin}} \exp(-\beta \mathcal{H}_G(\sigma))$
using Hamiltonian invariance
\texttt{hamiltonian\_map\_equiv}: for $\sigma = \tau \circ e$,
$\mathcal{H}_{G.\mathrm{map}\,e}(\tau) = \mathcal{H}_G(\sigma)$. The
Hamiltonian invariance itself splits into invariance of the
interaction energy (via \texttt{edgeFinset\_map} + per-edge
spin pullback along $\mathrm{Sym}^2$-map) and invariance of the
external-field energy (\texttt{Fintype.sum\_equiv}).
\end{proof}

\paragraph{Purpose.} This is infrastructure for transporting an
\texttt{inducedGraph} on a disjoint Finset union
$\Lambda_1 \cup \Lambda_2$ back to a graph on the disjoint sum
type $\uparrow\Lambda_1 \oplus \uparrow\Lambda_2$
(via \texttt{Equiv.Finset.disjUnionEquiv}), so that the
super-additivity inequality of the previous step applies to
actual finite-volume exhaustions.

\subsection{Super-additivity on Finset disjoint union
(\texttt{AmbientLatticeSum.lean})}

\begin{theorem}[\texttt{IsingModel.log\_partitionFunction\_inducedGraph\_disjUnion\_super\_additive}]
For $G \colon \mathrm{SimpleGraph}\,V$, disjoint
$\Lambda_1, \Lambda_2 : \mathrm{Finset}\,V$, and ferromagnetic $p$,
\[
  \log Z_{\mathrm{inducedGraph}\,G\,\Lambda_1}(p)
  + \log Z_{\mathrm{inducedGraph}\,G\,\Lambda_2}(p)
  \le \log Z_{\mathrm{inducedGraph}\,G\,(\Lambda_1 \cup \Lambda_2)}(p).
\]
\end{theorem}

\begin{proof}
Let $e = \mathrm{Equiv.Finset.union}\,\Lambda_1\,\Lambda_2\,\mathrm{hd}
  \colon \uparrow\!\Lambda_1 \oplus \uparrow\!\Lambda_2
     \simeq \uparrow\!(\Lambda_1 \cup \Lambda_2)$.
\begin{align*}
  \log Z_{\mathrm{inducedGraph}\,\Lambda_1}
  + \log Z_{\mathrm{inducedGraph}\,\Lambda_2}
  &= \log Z_{\mathrm{inducedGraph}\,\Lambda_1
      \oplus_g \mathrm{inducedGraph}\,\Lambda_2}
     & \text{(\S 0101 \texttt{log\_partitionFunction\_sum})} \\
  &= \log Z_{(\mathrm{inducedGraph}\,\Lambda_1
      \oplus_g \mathrm{inducedGraph}\,\Lambda_2).\mathrm{map}\,e}
     & \text{(\S 0103 \texttt{log\_partitionFunction\_map\_equiv})} \\
  &\le \log Z_{\mathrm{inducedGraph}\,(\Lambda_1 \cup \Lambda_2)}
     & \text{(ferromagnetic subgraph monotonicity,}
\end{align*}
using the key subgraph relation \texttt{inducedGraph\_sum\_map\_le\_union}:
every edge of $(\mathrm{inducedGraph}\,\Lambda_1
\oplus_g \mathrm{inducedGraph}\,\Lambda_2)\,.\mathrm{map}\,e$ is
either an $\mathrm{inl}$-image edge (both endpoints in $\Lambda_1$)
or an $\mathrm{inr}$-image edge (both in $\Lambda_2$), and in either
case it is an edge of $\mathrm{inducedGraph}\,G\,(\Lambda_1 \cup \Lambda_2)$.
\end{proof}

\paragraph{Closing Prop 4.6.1.} Combined with the uniform upper
bound (PR \#122, \#123), Fekete's lemma now yields
convergence of $f_\Lambda = (\log Z_\Lambda)/|\Lambda|$ along
arbitrary exhaustions of the ambient lattice; this is the
Prop 4.6.1 body in the next PR.

\begin{theorem}[GJ \S 4.6 Prop 4.6.1, disjoint-tower \(+\)
\texttt{BoundedEdgeDensity} form; PR \#203]
For a graph $G : \mathrm{SimpleGraph}\, V$ with an exhaustion
$\Lambda : \mathrm{Exhaustion}\, V$, Ising parameters $p$ with
bounded edge density along $\Lambda$, and the disjoint-tower
hypotheses
\[
  |\Lambda_{m+n}| = |\Lambda_m| + |\Lambda_n|, \quad
  \log Z_{\Lambda_m} + \log Z_{\Lambda_n}
    \le \log Z_{\Lambda_{m+n}}, \quad
  |\Lambda_1| \neq 0,
\]
the exhaustion free-energy density $f_{\Lambda_n}
= (\log Z_{\Lambda_n})/|\Lambda_n|$ converges to the
infinite-volume free energy $f_\infty$:
\[
  \lim_{n\to\infty} f_{\Lambda_n}(p) = f_\infty(p).
\]
\end{theorem}

\begin{proof}
This is a strict relaxation of the base theorem
\texttt{freeEnergyAlongExhaustion\_tendsto\_of\_superadditive}
(which requires in addition the explicit hypothesis
$\mathrm{BddAbove}\,(\mathrm{range}\,(f_\Lambda))$).
The \texttt{BoundedEdgeDensity} hypothesis yields that
$\mathrm{BddAbove}$ bound via
\texttt{BddAbove\_freeEnergyAlongExhaustion\_range}: if every
$\Lambda_n$ has at most $c\,|\Lambda_n|$ edges, then each
$f_{\Lambda_n}$ is bounded by
$\log 2 + |\beta|(|J|c + |h|)$ (for nonempty stages, combined with
$0$ on empty stages via the $\max$ with $0$).
The remaining proof steps are as in the base theorem: apply
mathlib's \texttt{Subadditive.tendsto\_lim} to
$u_n := -\log Z_{\Lambda_n}$, translate via $|\Lambda_n| =
n\cdot|\Lambda_1|$ (from \texttt{hcard\_add}) to obtain
$f_{\Lambda_n} = -u_n/(n\,|\Lambda_1|)$, and identify the limit
with $f_\infty = \limsup f_{\Lambda_n}$ via
\texttt{freeEnergyInfinite\_eq\_of\_tendsto}. Reference:
Glimm--Jaffe \cite{glimm-jaffe} \S 4.6 Prop 4.6.1, p.~68.
\end{proof}

\paragraph{Translation-invariance scaffolding for GJ \S 4.6
Prop 4.6.1 (PRs \#220--\#224).}
A multi-PR programme toward dropping the disjoint-tower hypotheses
of \texttt{freeEnergyAlongExhaustion\_tendsto\_of\_disjoint\_tower}
(PR \#203). Completed steps:
\begin{enumerate}
\item \texttt{Ambient.IsTranslationInvariant G}: a
  \texttt{SimpleGraph V} is translation invariant under an
  \texttt{AddAction T V} if the edge relation is preserved
  ($G.Adj\,(t +_v u)\,(t +_v v) \leftrightarrow G.Adj\,u\,v$).
  Trivial instances for $\bot$ and $\top$.
\item \texttt{vaddFinset t A := A.image (...)} using the map
  $u \mapsto t +_v u$, with
  \texttt{vaddFinset\_card}, \texttt{vaddFinset\_zero},
  \texttt{vaddFinset\_disjoint\_of\_disjoint}.
\item \texttt{TranslationInvariantExhaustion T V}
  (\texttt{extends Exhaustion V}) with fields
  \texttt{shift} of type $\mathbb N \to T$, \texttt{shift\_zero},
  \texttt{volume\_zero}, \texttt{volume\_succ}
  ($\Lambda_{n+1} = \Lambda_n \cup (\mathrm{shift}\,n
  +_v \Lambda_1)$),
  \texttt{disjoint\_shift}. Gives automatic
  \texttt{volume\_card\_add}: $|\Lambda_{m+n}| = |\Lambda_m| +
  |\Lambda_n|$.
\item \texttt{disjointTowerHypotheses\_of\_translationInvariant}:
  given a \texttt{TranslationInvariantExhaustion} + user-supplied
  \texttt{hsuper} + \texttt{hcard\_one}, assemble a full
  \texttt{DisjointTowerHypotheses} record.
\item \texttt{freeEnergyAlongExhaustion\_tendsto\_of\_translationInvariant}:
  combined with \texttt{BoundedEdgeDensity}, yields the Fekete
  convergence to \texttt{freeEnergyInfinite}.
\end{enumerate}
\emph{Remaining}: derive \texttt{hsuper} from log-$Z$ translation
invariance. Requires Ising-Hamiltonian translation equivariance
(separate multi-PR step).
Reference: Glimm--Jaffe \cite{glimm-jaffe} \S 4.6 Prop 4.6.1,
p.~68.

\begin{theorem}[Finset-based susceptibility at $J = 0$ closed form;
PR \#231]
For the repo-level Finset-based susceptibility
$\chi := \sum_j U_2(i,j)$, where $U_2(i,j) = \langle
\sigma^{\{i,j\}}\rangle - \langle\sigma^{\{i\}}\rangle
\langle\sigma^{\{j\}}\rangle$,
for any ambient graph $G$ with finite $\iota$ and any
$h, \beta \in \mathbb{R}$,
\[
  \chi^{\langle 0, h, \beta\rangle}_G(i)
  = \tanh(\beta h) \cdot \big(1 - \tanh(\beta h)\big).
\]
(Caveat: this differs from the physics response-function
$dM/dh = \beta(1-t^2)$ at the diagonal, because the Finset
collapse $\{i,i\} = \{i\}$ gives
$\langle\sigma^{\{i\}}\rangle - \langle\sigma^{\{i\}}\rangle^2
= t - t^2$ rather than the physics
$\langle\sigma_i^2\rangle - \langle\sigma_i\rangle^2 = 1 - t^2$.)
\end{theorem}

\begin{proof}
$\chi(i) := \sum_j U_2(i,j)$. For $j \neq i$,
\texttt{truncated2\_J\_zero\_of\_ne} (PR \#207) gives summand zero.
For $j = i$, the Finset $\{i,i\}$ collapses to $\{i\}$, so
$U_2(i,i) = \langle \sigma^{\{i\}}\rangle -
\langle \sigma^{\{i\}}\rangle^2 = t - t^2$ with $t = \tanh(\beta h)$
(\texttt{correlation\_J\_zero} + card$\,\{i\}=1$). Factor to
$t(1-t)$. Reference: Glimm--Jaffe \cite{glimm-jaffe} \S 4.1
(non-interacting slice); \S 5.1 pp.~76--77 (susceptibility).
\end{proof}

\begin{theorem}[Finite-volume magnetization at $J = 0$
$= \tanh(\beta h)$; PR \#219]
For any graph $G$ with finite $\iota$, any $h, \beta \in \mathbb{R}$,
and any site $i$, the finite-volume magnetization satisfies
$M^{\langle 0, h, \beta\rangle}_G(i) = \tanh(\beta h)$.
\end{theorem}

\begin{proof}
$\mathrm{magnetization}(p; i) := \mathrm{correlation}(p; \{i\})$.
Apply the finite-volume closed form
$\mathrm{correlation}(\langle 0, h, \beta\rangle; A)
= \tanh(\beta h)^{|A|}$ (\texttt{correlation\_J\_zero}) with
$A = \{i\}$: $|A| = 1$, so the power reduces to
$\tanh(\beta h)$.
Reference: Glimm--Jaffe \cite{glimm-jaffe} \S 4.1
(non-interacting $J = 0$ slice); \S 5.1 pp.~76--77
(magnetization).
\end{proof}

\begin{theorem}[$\infty$-vol magnetization at $J = 0$
ferromagnetic = $\tanh(\beta h)$; PR \#218]
For a ferromagnetic parameter $\langle 0, h, \beta \rangle$ and
any graph $G$ with exhaustion $\Lambda$, the infinite-volume
magnetization at a site $i$ is
\[
  M^\infty(G, \Lambda; \langle 0, h, \beta\rangle; i)
  = \tanh(\beta h).
\]
\end{theorem}

\begin{proof}
$\mathrm{magnetizationInfinite}(p; i) :=
\mathrm{correlationInfinite}(p; \{i\})$. Apply the closed form
$\mathrm{correlationInfinite}(\langle 0, h, \beta\rangle; A) =
\tanh(\beta h)^{|A|}$ (PR \#210) with $A = \{i\}$, so
$|A| = 1$ and the power reduces to $\tanh(\beta h)$.
Reference: Glimm--Jaffe \cite{glimm-jaffe} \S 4.1
(non-interacting $J = 0$ slice; $\beta$ is unrestricted here, so
this is distinct from the infinite-temperature $\beta = 0$
slice); \S 5.1 pp.~76--77 (magnetization).
\end{proof}

\begin{theorem}[Lebowitz 4-point $J = 0$ closed form
$-2\,t^4$ (finite and $\infty$-volume); PR \#215]
For any ambient graph $G$, any $h, \beta \in \mathbb{R}$, and
pairwise distinct sites $i, j, k, l$,
\[
  U_4^{\text{Lebowitz}}(G, \langle 0,h,\beta\rangle; i,j,k,l)
  = -2\,\tanh(\beta h)^4.
\]
The same formula holds for the infinite-volume Lebowitz 4-point
under a ferromagnetic hypothesis.
\end{theorem}

\begin{proof}
At $J = 0$,
$\langle\sigma^A\rangle_{\langle 0,h,\beta\rangle}
= \tanh(\beta h)^{|A|}$
by \texttt{correlation\_J\_zero} (finite volume) or
\texttt{correlationInfinite\_J\_zero} (ferromagnetic
infinite volume, PR \#210). With pairwise distinct
$i, j, k, l$ the seven Finset cardinals are
$4$ (once) and $2$ (six times), so
\[
  U_4 = t^4 - 3(t^2 \cdot t^2) = t^4 - 3\,t^4 = -2\,t^4,
  \quad t := \tanh(\beta h).
\]
Note $-2\,t^4 \le 0$ unconditionally, consistent with
$U_4 \le 0$ (Glimm--Jaffe Cor. 4.3.3).
Reference: Glimm--Jaffe \cite{glimm-jaffe} \S 5.1 pp.~72--74
(cluster context); \S 4.3 Cor. 4.3.3 / Lebowitz.
\end{proof}

\begin{theorem}[Lebowitz 4-point $\beta = 0$ vanishing (finite and
$\infty$-volume); PR \#214]
For any ambient graph $G$, any $J, h \in \mathbb{R}$, any sites
$i, j, k, l$:
$U_4^{\text{Lebowitz}}(G, \langle J,h,0\rangle; i,j,k,l) = 0$
in finite volume, and
$U_4^{\text{Lebowitz}}(\infty; i,j,k,l) = 0$ at infinite volume
along any exhaustion.
\end{theorem}

\begin{proof}
At $\beta = 0$ every Boltzmann weight is $1$, so
$\langle \sigma^A \rangle = 0$ for any nonempty Finset
$A \subseteq \iota$
(\texttt{correlation\_beta\_zero\_vanish\_of\_nonempty\_A}). The
Lebowitz 4-point
$U_4 = \langle \sigma^{\{i,j,k,l\}}\rangle
- \langle\sigma^{\{i,j\}}\rangle\langle\sigma^{\{k,l\}}\rangle
- \langle\sigma^{\{i,k\}}\rangle\langle\sigma^{\{j,l\}}\rangle
- \langle\sigma^{\{i,l\}}\rangle\langle\sigma^{\{j,k\}}\rangle$
has all seven Finset correlations over nonempty subsets (each
contains at least one supplied site), so every term is $0$ and
$U_4 = 0$.

$\infty$-volume version is identical with
\texttt{correlationInfinite\_beta\_zero\_vanish\_of\_nonempty\_A}
in place of the finite-volume lemma.

Note: at $J = 0$ the Lebowitz 4-point is $-2\,t^4$ where
$t = \tanh(\beta h)$ for pairwise-distinct sites, which is
non-zero when $\beta h \neq 0$, so the $J = 0$ slice does
\emph{not} vanish.
Reference: Glimm--Jaffe \cite{glimm-jaffe} \S 5.1 pp.~72--74
(cluster context); \S 4.3 Cor. 4.3.3 / Lebowitz.
\end{proof}

\begin{theorem}[$\infty$-vol truncated-3 trivial slices; PR \#212]
For a ferromagnetic $\langle 0, h, \beta\rangle$ and pairwise
distinct $i, j, k$, the $\infty$-vol Ursell 3-point function
$\langle \sigma_i; \sigma_j; \sigma_k \rangle_\infty^{\langle
0,h,\beta\rangle}$ vanishes. For $\langle J, h, 0\rangle$ and any
sites, the $\infty$-vol Ursell 3-point function also vanishes.
\end{theorem}

\begin{proof}
$J = 0$: Substitute
$\langle \sigma^A \rangle_\infty = \tanh(\beta h)^{|A|}$ (the PR \#210
closed form) into each of the seven Finset correlations, use
$|\{i,j,k\}| = 3, |\{i\}|=|\{j\}|=|\{k\}|=1,
|\{i,j\}|=|\{j,k\}|=|\{i,k\}|=2$, and reduce the Ursell
combination to $t^3 - 3\cdot t^3 + 2\cdot t^3 = 0$ with
$t := \tanh(\beta h)$.

$\beta = 0$: Every Finset correlation in the Ursell expression is
over a nonempty Finset, so the PR \#211 vanishing
$\langle \sigma^A \rangle_\infty^{\langle J,h,0\rangle} = 0$
forces each term to zero.

Reference: Glimm--Jaffe \cite{glimm-jaffe} \S 5.1 pp.~72--74
(cluster context); \S 4.3 (Ursell / GHS).
\end{proof}

\begin{theorem}[$\infty$-vol $\beta = 0$ correlation vanishing and
truncated-2 zero; PR \#211]
For any ambient graph $G$, any exhaustion $\Lambda :
\mathrm{Exhaustion}\,V$, any $J, h \in \mathbb{R}$, and any
nonempty $A \subseteq V$,
\[
  \langle \sigma^A \rangle_\infty^{G,\Lambda,\langle J,h,0\rangle}
  = 0.
\]
In particular, for any sites $i, j : V$ (distinct or not),
$\langle \sigma_i; \sigma_j \rangle_\infty^{G,\Lambda,\langle
J,h,0\rangle} = 0$.
No ferromagnetic hypothesis is required.
\end{theorem}

\begin{proof}
At $\beta = 0$ every Boltzmann weight equals $1$, so
\texttt{correlation\_beta\_zero\_vanish\_of\_nonempty\_A} gives
finite-volume vanishing on any nonempty Finset. Pulling this to
the along-exhaustion sequence: when $A \subseteq \Lambda.volume\,n$
the lifted subset is still nonempty (via \texttt{liftFinset\_card})
and the finite-volume lemma applies; when
$A \not\subseteq \Lambda.volume\,n$ the
$\mathrm{correlationAlongExhaustion}$ is $0$ by definition. Hence
the sequence is pointwise zero, so
$\mathrm{correlationInfinite} = \bigsqcup_n 0 = 0$.
The truncated-2 vanishing follows by substitution at
$\{i,j\}, \{i\}, \{j\}$ — all three are nonempty, and $i = j$ is
handled by the Finset collapse $\{i,j\} = \{i\}$.
Reference: Glimm--Jaffe \cite{glimm-jaffe} \S 5.1 pp.~72--74
(cluster context); \S 4.1 (infinite-temperature slice).
\end{proof}

\begin{theorem}[$\infty$-vol $J = 0$ correlation closed form and
truncated-2 vanishing; PR \#210]
For a ferromagnetic parameter $\langle 0, h, \beta \rangle$
(i.e. $h \ge 0$, $0 < \beta$; the strict-$\beta$ condition comes
from \texttt{Ferromagnetic.hβ}), any graph $G :
\mathrm{SimpleGraph}\, V$ and exhaustion $\Lambda :
\mathrm{Exhaustion}\, V$, and any finite $A \subseteq V$,
\[
  \langle \sigma^A \rangle_\infty^{G,\Lambda,\langle 0,h,\beta\rangle}
  = \tanh(\beta h)^{|A|}.
\]
For distinct $i \neq j$ the infinite-volume truncated 2-point
function vanishes:
$\langle \sigma_i; \sigma_j \rangle_\infty^{G,\Lambda,\langle 0,h,\beta\rangle}
= 0$.
\end{theorem}

\begin{proof}
By the finite-volume closed form $\langle \sigma^A
\rangle_\Lambda^{\langle 0,h,\beta\rangle} = \tanh(\beta h)^{|A|}$
(see \texttt{correlation\_J\_zero}), when $A \subseteq \Lambda.volume\,n$
we have
$\text{correlationAlongExhaustion}\,G\,\Lambda\,\langle 0,h,\beta\rangle\,A\,n
= \tanh(\beta h)^{|A|}$
(the Finset lift preserves cardinality). By $\Lambda.\mathrm{exhaust}$,
this holds for all sufficiently large $n$, so the sequence is
eventually constant at $\tanh(\beta h)^{|A|}$. On the other hand,
by \texttt{correlationAlongExhaustion\_tendsto\_ciSup}
(ferromagnetic, monotone + bounded), the sequence also tends to
$\mathrm{correlationInfinite}\,G\,\Lambda\,\langle 0,h,\beta\rangle\,A$.
Uniqueness of limits gives the claimed closed form.
The truncated-2 vanishing is immediate by substitution at
$\{i,j\}, \{i\}, \{j\}$ and the cardinalities $2, 1, 1$.
Reference: Glimm--Jaffe \cite{glimm-jaffe} \S 5.1 pp.~72--74
(cluster property context); \S 4.1 (infinite-temperature slice).
\end{proof}

\begin{theorem}[Trivial-slice vanishing of the truncated 3-point
(Ursell) function; PR \#209]
For any ambient graph $G$, any $h, \beta \in \mathbb{R}$, and
pairwise distinct sites $i, j, k$,
$\langle \sigma_i; \sigma_j; \sigma_k
\rangle_{G,\langle 0,h,\beta\rangle} = 0$.
For any $J, h \in \mathbb{R}$ and any sites $i, j, k$ — distinctness
is not required, since repeated indices collapse at the Finset
level inside the Lean definition of \texttt{truncated3}, paralleling
\texttt{truncated2\_beta\_zero} —
$\langle \sigma_i; \sigma_j; \sigma_k
\rangle_{G,\langle J,h,0\rangle} = 0$.
\end{theorem}

\begin{proof}
$J = 0$ case. By \texttt{correlation\_J\_zero},
$\langle \sigma^A \rangle_{\langle 0,h,\beta\rangle} =
\tanh(\beta h)^{|A|}$. With $t := \tanh(\beta h)$ and pairwise
distinct $i, j, k$ one has $|\{i,j,k\}| = 3$ and each pair has
card $2$. The Ursell combination becomes
\[
  t^3 - 3\,(t \cdot t^2) + 2\,t^3
  = t^3 - 3\,t^3 + 2\,t^3 = 0.
\]

$\beta = 0$ case. By
\texttt{correlation\_beta\_zero\_vanish\_of\_nonempty\_A}, every
nonempty-subset correlation in the Ursell combination vanishes,
so the linear combination is $0$.

Reference: Glimm--Jaffe \cite{glimm-jaffe} \S 5.1 pp.~72--74
(cluster property context); \S 4.3 (Ursell functions / GHS).
\end{proof}

\begin{theorem}[Infinite-temperature ($\beta=0$) vanishing of the
truncated 2-point function; PR \#208]
For any ambient graph $G$, any $J, h \in \mathbb{R}$, and any sites
$i, j : \iota$ (distinct or not),
\[
  \langle \sigma_i; \sigma_j \rangle_{G,\langle J,h,0\rangle} = 0.
\]
\end{theorem}

\begin{proof}
At $\beta = 0$, the Boltzmann weight $\exp(-\beta H) \equiv 1$, so
$\langle \sigma^A \rangle_{G,\langle J,h,0\rangle}$ is the uniform
spin average. By
\texttt{correlation\_beta\_zero\_vanish\_of\_nonempty\_A}, this
vanishes on any nonempty $A$. Hence each of
$\langle\sigma^{\{i,j\}}\rangle$,
$\langle\sigma^{\{i\}}\rangle$, $\langle\sigma^{\{j\}}\rangle$
is $0$, so the truncated value is $0 - 0 \cdot 0 = 0$.
Note that unlike the $J = 0$ slice, here no $i \neq j$ hypothesis
is required: if $i = j$ the Finset first term
$\langle\sigma^{\{i,j\}}\rangle$ collapses to
$\langle\sigma^{\{i\}}\rangle$ (not the physics product
$\langle\sigma_i\sigma_i\rangle = 1$) and again vanishes at
$\beta = 0$.
Reference: Glimm--Jaffe \cite{glimm-jaffe} \S 5.1,
pp.~72--74 (cluster property context); §4.1 infinite-temperature
slice of the correlation function is the direct source of
\texttt{correlation\_beta\_zero\_vanish\_of\_nonempty\_A}.
\end{proof}

\begin{theorem}[Non-interacting ($J=0$) factorisation of the
truncated 2-point function; PR \#207]
For any ambient graph $G : \mathrm{SimpleGraph}\, \iota$ with
$\iota$ finite, any $h, \beta \in \mathbb{R}$, and any distinct
sites $i \neq j$ in $\iota$,
\[
  \langle \sigma_i; \sigma_j \rangle_{G,\langle 0,h,\beta\rangle}
  = \langle \sigma_i \sigma_j \rangle - \langle \sigma_i\rangle
  \langle \sigma_j \rangle = 0.
\]
No distance / separation hypothesis is needed: the
factorisation is identically true for any two distinct sites.
\end{theorem}

\begin{proof}
By \texttt{correlation\_J\_zero}, for any subset
$A \subseteq \iota$,
$\langle \sigma^A \rangle_{G,\langle 0,h,\beta\rangle} =
\tanh(\beta h)^{|A|}$.
With $i \neq j$ one has $|\{i,j\}| = 2$ and $|\{i\}| = |\{j\}| = 1$,
so
\[
  \langle \sigma^{\{i,j\}} \rangle = \tanh(\beta h)^{2}, \qquad
  \langle \sigma^{\{i\}} \rangle\langle \sigma^{\{j\}}\rangle
  = \tanh(\beta h)^{1} \cdot \tanh(\beta h)^{1}
  = \tanh(\beta h)^{2},
\]
and the difference vanishes.
Reference: Glimm--Jaffe \cite{glimm-jaffe} \S 5.1, pp.~72--74
(cluster property context). This is the trivial non-interacting
slice of the cluster property discussion; note that $\beta$ is
arbitrary, so this is distinct from the high-temperature
($\beta$ small) regime. Non-zero-$J$ decay at large separation
in pure phases remains follow-up work.
\end{proof}

\begin{theorem}[GJ \S 4.6 Prop 4.6.1, log-linear-card builder and
$\beta = 0$ instance; PR \#206]
(i) If $\log Z_\Lambda = |\Lambda| \cdot c$ for all stages, then the
super-additivity field of
\texttt{DisjointTowerHypotheses} is automatic under
\texttt{hcard\_add} — this is the generic builder
\texttt{DisjointTowerHypotheses.of\_log\_linear\_card}.

(ii) At $\beta = 0$ one has $Z_\Lambda = 2^{|\Lambda|}$, hence
$\log Z_\Lambda = |\Lambda| \cdot \log 2$, fitting the
log-linear pattern. So
\texttt{DisjointTowerHypotheses.of\_beta\_zero} and the Fekete
corollary
\texttt{freeEnergyAlongExhaustion\_beta\_zero\_tendsto\_of\_hcard\_add}
follow immediately, parallel to the $J = 0$ case.
\end{theorem}

\begin{proof}
(i) Super-additivity reduces to
$(|\Lambda_m| + |\Lambda_n|) \cdot c \le |\Lambda_{m+n}| \cdot c$;
under \texttt{hcard\_add} the two sides are \emph{equal}, so the
proof does not rely on any sign assumption on $c$ — it is a pure
equality reduction, not a monotonicity argument.
(ii) \texttt{IsingModel.partitionFunction\_beta\_zero} +
\texttt{card\_config\_eq\_two\_pow} + \texttt{Real.log\_pow} give
$\log Z_\Lambda = |\Lambda| \cdot \log 2$. The rest is
\texttt{of\_log\_linear\_card} applied with $c = \log 2$.
Reference: Glimm--Jaffe \cite{glimm-jaffe} \S 4.6 Prop 4.6.1, p.~68.
\end{proof}

\begin{theorem}[GJ \S 4.6 Prop 4.6.1, concrete $J = 0$ instance; PR \#205]
At $J = 0$, with bounded edge density along $\Lambda$, a
cardinality-additive exhaustion
($|\Lambda_{m+n}| = |\Lambda_m| + |\Lambda_n|$) with
$|\Lambda_1| \neq 0$, the exhaustion free-energy density
converges to the infinite-volume free energy:
\[
  f_{\Lambda_n}(0, h, \beta) \longrightarrow f_\infty(0, h, \beta).
\]
\end{theorem}

\begin{proof}
At $J = 0$, \texttt{IsingModel.partitionFunction\_J\_zero} gives
$Z_\Lambda = (2\cosh(\beta h))^{|\Lambda|}$, so
$\log Z_\Lambda = |\Lambda| \cdot \log(2\cosh(\beta h))$. The
\texttt{DisjointTowerHypotheses} record's super-additivity field
becomes
\[
  (|\Lambda_m| + |\Lambda_n|) \log(2\cosh(\beta h))
  \le |\Lambda_{m+n}| \log(2\cosh(\beta h)),
\]
which holds with equality under \texttt{hcard\_add}. Combined with
\texttt{BoundedEdgeDensity}, the bundled Fekete theorem
\texttt{freeEnergyAlongExhaustion\_tendsto\_of\_disjointTowerHypotheses}
closes the proof.  Reference: Glimm--Jaffe \cite{glimm-jaffe} \S 4.6
Prop 4.6.1, p.~68. This is the first concrete
\texttt{DisjointTowerHypotheses} instance; non-$J\!=\!0$ instances
require translation invariance and remain follow-up.
\end{proof}

\begin{theorem}[GJ \S 4.6 Prop 4.6.1,
\texttt{DisjointTowerHypotheses} bundle; PR \#204]
The three structural hypotheses $|\Lambda_{m+n}| = |\Lambda_m|
+ |\Lambda_n|$, $\log Z_{\Lambda_m} + \log Z_{\Lambda_n} \le
\log Z_{\Lambda_{m+n}}$, and $|\Lambda_1| \neq 0$ are packaged
as a single \texttt{DisjointTowerHypotheses} record. Given
\texttt{BoundedEdgeDensity}$\,G\,\Lambda$ and such a record,
$f_{\Lambda_n}(p) \to f_\infty(p)$.
\end{theorem}

\begin{proof}
Definitional unpack of the
\texttt{DisjointTowerHypotheses} record into the three
hypotheses of
\texttt{freeEnergyAlongExhaustion\_tendsto\_of\_disjoint\_tower}.
Reference: Glimm--Jaffe \cite{glimm-jaffe} \S 4.6 Prop 4.6.1,
p.~68. This is an API-site convenience wrapper; no new
mathematical content.
\end{proof}

\begin{theorem}[Direct $\limsup$ corollary of GJ \S 5.4 Prop 5.4.2
along an exhaustion; PR \#213]
Under the same per-stage hypotheses as PR \#202, the
$\limsup_n (1 - \langle \sigma_{i_n} \rangle_{+}^{G_n, B_n})$ along
\texttt{atTop} is bounded above by $\exp(-c\beta)$.
\end{theorem}

\begin{proof}
From the per-stage bound of PR \#202
($\texttt{prop\_5\_4\_2\_along\_exhaustion}$): at every $n$,
$0 \le 1 - \langle \sigma_{i_n} \rangle_{+}^{G_n, B_n} \le
\exp(-c\beta)$. Apply mathlib's
\texttt{Filter.limsup\_le\_of\_le}: the eventually-uniform upper
bound $\exp(-c\beta)$ holds pointwise, and the
\texttt{IsCoboundedUnder (\cdot \le \cdot)} side condition is
satisfied by the pointwise lower bound $0$
(\texttt{isCoboundedUnder\_le\_of\_eventually\_le} with $x = 0$).
Reference: Glimm--Jaffe \cite{glimm-jaffe} \S 5.4 Prop 5.4.2,
p.~83. This is a direct $\limsup$ corollary; strengthening to a
genuine infinite-volume $+$-BC expectation \emph{inequality}
remains follow-up work.
\end{proof}

\begin{theorem}[Corollary/wrapper of GJ \S 5.4 Prop 5.4.2 along an
exhaustion (per-stage); PR \#202]
For a graph $G : \mathrm{SimpleGraph}\, V$ with an exhaustion
$\Lambda : \mathrm{Ambient.Exhaustion}\, V$, suppose that for every
stage $n$ the induced graph $G_n := \mathrm{Ambient.inducedGraph}\,G\,(\Lambda_n)$
is preconnected, $B_n \subseteq \Lambda_n$ is a non-empty set of
$+$-boundary-condition sites,
$i_n \in \Lambda_n$, and the common exponential bound
\[
  2 \cdot 2^{\#\Lambda_n} \exp(-2\beta J) \le \exp(-c\beta)
\]
holds uniformly in $n$. Then, for $J, \beta > 0$ and $h = 0$, the
Peierls bound
\[
  0 \;\le\; 1 - \langle \sigma_{i_n} \rangle_{+}^{G_n, B_n}
  \;\le\; \exp(-c\beta)
\]
holds at every $n$.
\end{theorem}

\begin{proof}
Direct application of \texttt{prop\_5\_4\_2\_self\_contained}
(Glimm--Jaffe \cite{glimm-jaffe} \S 5.4 p.~83) at each stage $n$;
the hypotheses at stage $n$ are the finite-volume hypotheses of
that theorem specialised to $G_n$, $B_n$, $i_n$. This is a
scaffolding step: subsequent PRs on the same branch will pick a
canonical $(B_n, i_n)$ from the ambient geometry and take the
$\limsup$ / $\liminf$ to recover the genuine infinite-volume
bound $0 \le 1 - \langle\sigma_i\rangle_+^\infty \le \exp(-c\beta)$
(Glimm--Jaffe \cite{glimm-jaffe} \S 5.4 p.~83).
\end{proof}

\begin{theorem}[Groundwork for GJ \S 4.6 Prop 4.6.1 — concrete
translation-invariant lattice; PR \#244]
For every $d : \mathbb{N}$, the Ising lattice graph
$\texttt{latticeGraph}\ d : \operatorname{SimpleGraph}(\mathrm{Fin}\ d \to \mathbb{Z})$
(edges between points at $\ell^1$-distance $1$) is invariant under
the canonical self-action of $\mathrm{Fin}\ d \to \mathbb{Z}$ on
itself:
\[
  \texttt{Ambient.IsTranslationInvariant}\ (\mathrm{Fin}\ d \to \mathbb{Z})\ (\texttt{latticeGraph}\ d).
\]
\end{theorem}
\begin{proof}
On the self-action, $t +_v u = t + u$ pointwise, hence
$(t +_v u) - (t +_v v) = u - v$ in every coordinate. The defining
ℓ¹ condition $\sum_i |x_i - y_i| = 1$ is therefore preserved in
both directions.
\end{proof}

\paragraph{Remark.} This is the first non-trivial
\texttt{IsTranslationInvariant} instance — earlier examples
$(\bot, \top)$ are structurally trivial (no edges / all pairs).
The lattice graph $\texttt{latticeGraph}\ d$ is the natural ambient
graph behind the \S 4.6 Prop 4.6.1 translation-invariant exhaustion
programme. The natural two-sided cubic exhaustion
$\Lambda_n = \{x : \mathrm{Fin}\ d \to \mathbb{Z} \mid \forall i,\ -n \le x_i \le n\}$
does not fit the current
\texttt{TranslationInvariantExhaustion} structure, which demands
that each stage differ from the previous by a single translated base
block (\texttt{volume\_succ}) with \texttt{shift} an additive
homomorphism $\mathbb{N} \to T$. Assembling a compatible concrete
exhaustion requires either relaxing the single-block constraint or
altering the \texttt{Exhaustion} interface, and is deferred to a
subsequent PR.
Reference: Glimm--Jaffe \cite{glimm-jaffe} \S 4.6 Prop 4.6.1, p.~68.

\begin{theorem}[GJ \S 4.6 concrete: cubic \texttt{Ambient.Exhaustion} on
$\mathbb{Z}^d$; PR \#245]
For every $d : \mathbb{N}$, the \emph{cubic box}
\[
  \mathrm{cubicBox}\ d\ n
  := \mathrm{Fintype.piFinset}\ (\lambda \_\ \Rightarrow\
    \mathrm{Finset.Icc}\ (-(n:\mathbb{Z}))\ n)
\]
determines an \texttt{Ambient.Exhaustion} of the integer lattice
$\mathrm{Fin}\ d \to \mathbb{Z}$ via
$\mathrm{cubicExhaustion}\ d .\mathrm{volume} := \mathrm{cubicBox}\ d$,
together with the monotonicity lemma $m \le n \Rightarrow
\mathrm{cubicBox}\ d\ m \subseteq \mathrm{cubicBox}\ d\ n$ and the
exhaust lemma: every finite $A \subset \mathbb{Z}^d$ is contained in
some $\mathrm{cubicBox}\ d\ N$ where $N$ is the maximum absolute
coordinate across points of $A$.
\end{theorem}
\begin{proof}[Exhaust lemma]
Let $S := \{|a_i|^{\mathrm{nat}} : a \in A,\ i : \mathrm{Fin}\ d\}$
as a \texttt{Finset $\mathbb{N}$}. In the non-degenerate case
$A \ne \emptyset$ and $d \ne 0$, $S$ is non-empty and we take
$N := \mathrm{Finset.max'}\ S$. For any $n \ge N$ and $a \in A$,
$i : \mathrm{Fin}\ d$, we have $|a_i|^{\mathrm{nat}} \le N \le n$,
so $|a_i| \le n$ in $\mathbb{Z}$, hence $-n \le a_i \le n$ and
$a \in \mathrm{cubicBox}\ d\ n$. In the degenerate case $S = \emptyset$
(equivalent to $A = \emptyset \vee d = 0$), take $N := 0$; the
membership goal $A \subseteq \mathrm{cubicBox}\ d\ n$ reduces,
under the assumption $a \in A$, to producing an element of $S$ (from
any $i : \mathrm{Fin}\ d$ supplied by $\mathrm{mem\_cubicBox}$'s
intro), contradicting $S = \emptyset$.
\end{proof}

\paragraph{Remark.} This is the first concrete
\texttt{Ambient.Exhaustion} on the physical integer lattice — prior
instances exist only on abstract $V$. It makes
$\mathrm{correlationInfinite}\ (\mathrm{latticeGraph}\ d)\
(\mathrm{cubicExhaustion}\ d)\ p\ A$ an explicit object. A matching
\texttt{TranslationInvariantExhaustion} is not provided (structural
obstruction noted above); bounded-edge-density for
$\mathrm{latticeGraph}\ d$ along this exhaustion is deferred to a
subsequent PR.
Reference: Glimm--Jaffe \cite{glimm-jaffe} \S 4.6, p.~67.

\begin{theorem}[Bounded edge density for $\mathrm{latticeGraph}\,d$
on cubic boxes; PR \#246]
For every $d : \mathbb{N}$,
\[
  \mathrm{Ambient.BoundedEdgeDensity}\ (\mathrm{IsingModel.latticeGraph}\,d)\
  (\mathrm{Ambient.cubicExhaustion}\,d)
\]
holds with constant $c = d$: $|E(\mathrm{latticeGraph}\,d[\Lambda_n])| \le
d \cdot |\Lambda_n|$ where $\Lambda_n = \mathrm{cubicBox}\,d\,n$.
\end{theorem}
\begin{proof}
\emph{Vertex degree bound.} For every vertex $v : \mathrm{Fin}\,d \to \mathbb{Z}$,
the neighbours in $\mathrm{latticeGraph}\,d$ lie in the $2d$-element set
$\{\mathrm{Function.update}\,v\,i\,(v_i \pm 1) : i : \mathrm{Fin}\,d\}$
(formalised as \texttt{latticeNeighborEnum}). The ℓ¹-distance-$1$ condition
$\sum_j |v_j - w_j| = 1$ forces exactly one coordinate $i$ to differ from $v$,
with $|v_i - w_i| = 1$. Restricting to $w \in \Lambda_n$ only lowers the
degree, so $\deg_{\mathrm{induced}}(v) \le 2d$ for every $v \in \Lambda_n$.

\emph{Handshake.} By the degree-sum formula
$\sum_{v \in \Lambda_n} \deg(v) = 2|E|$, we get
$2|E| \le 2d |\Lambda_n|$, hence $|E| \le d |\Lambda_n|$ (divide by $2$).
\end{proof}

\paragraph{Remark.}
Combined with $\mathrm{cubicExhaustion}\,d$ (PR~\#245), this closes the
concrete-ambient-lattice infrastructure for GJ \S 4.6 modulo the still
missing super-additivity ingredient (a compatible
$\mathrm{TranslationInvariantExhaustion}$). The per-stage
$\mathrm{freeEnergyAlongExhaustion}$ bound
$\le \log 2 + |\beta|(|J|\cdot d + |h|)$ now has an explicit constant.
Reference: Glimm--Jaffe \cite{glimm-jaffe} \S 4.6, p.~67.

\begin{theorem}[Concrete $\mathbb{Z}^d$ Ising free-energy sandwich; PR \#247]
For every $d : \mathbb{N}$, ferromagnetic parameters $0 \le J$,
$0 \le h$, $0 < \beta$, and any stage $n$ with $(\mathrm{cubicExhaustion}\,d)
.\mathrm{volume}\,n$ non-empty,
\[
  \log 2
    \le \mathrm{freeEnergyAlongExhaustion}\ (\mathrm{latticeGraph}\,d)\
        (\mathrm{cubicExhaustion}\,d)\ \langle J, h, \beta\rangle\ n
    \le \log 2 + |\beta| \cdot (|J| \cdot d + |h|).
\]
\end{theorem}
\begin{proof}
The lower bound is the ferromagnetic, nonempty-stage specialisation
of $\mathrm{freeEnergyAlongExhaustion\_ge\_log\_two}$ in
$\texttt{AmbientLattice.lean}$. The upper bound applies
$\mathrm{freeEnergyAlongExhaustion\_le\_uniform\_upper\_bound}$ with
the concrete edge-density constant $c = d$ furnished by
$\mathrm{inducedLatticeGraph\_card\_edgeFinset\_le}$ (PR~\#246).
\end{proof}

\paragraph{Remark.}
Capstone for the concrete $\mathbb{Z}^d$ infrastructure developed
across PRs~\#244 ($\mathrm{IsTranslationInvariant}$),
\#245 ($\mathrm{cubicExhaustion}\,d$), and \#246
($\mathrm{BoundedEdgeDensity}$ with $c = d$). The Fekete convergence
$f_n \to f_\infty$ still requires a compatible
$\mathrm{TranslationInvariantExhaustion}$ on $\mathbb{Z}^d$, deferred
pending a structure refactor.

\begin{theorem}[Shifted exhaustion; PR \#249]
Let $T$ be an additive group acting on $V$ (with $\mathrm{DecidableEq}\,V$).
For every $\mathrm{Ambient.Exhaustion}\,V\ \Lambda$ and $t : T$,
\[
  \Lambda.\mathrm{shift}\,t : \mathrm{Ambient.Exhaustion}\,V,
  \quad (\Lambda.\mathrm{shift}\,t).\mathrm{volume}\,n := t +_v \Lambda.\mathrm{volume}\,n,
\]
inherits monotonicity and the exhaust property; moreover
$|(\Lambda.\mathrm{shift}\,t).\mathrm{volume}\,n| = |\Lambda.\mathrm{volume}\,n|$
and $\Lambda.\mathrm{shift}\,0 = \Lambda$ pointwise on volumes.
\end{theorem}
\begin{proof}
Monotonicity: the translate $t +_v \cdot$ preserves inclusion.
Exhaust: given a finite $A$, pull back to $(-t) +_v A$, apply
$\Lambda$'s exhaust, and push forward.
\end{proof}

\paragraph{Remark.}
Infrastructure for translation-invariance lifts of
$\mathrm{correlationInfinite}$ / $\mathrm{freeEnergyInfinite}$.
Composing $\Lambda.\mathrm{shift}\,t$ with exhaustion-independence
corollaries is the planned mechanism to establish
$\mathrm{correlationInfinite}\,G\,\Lambda\,p\,A =
\mathrm{correlationInfinite}\,G\,\Lambda\,p\,(t +_v A)$ when $G$
is $\mathrm{IsTranslationInvariant}$.

\section{Concrete \texorpdfstring{$\mathbb Z^d$}{Z d} Fekete toolkit for \S 4.6 Prop 4.6.1}

This section summarises the concrete $\mathbb Z^d$ Fekete toolkit
(PRs \#638--\#667) that applies the generic-Finset Fekete framework
of \texttt{freeEnergy\_of\_finset\_sequence\_tendsto\_of\_superadditive}
(PR \#638) to four concrete Finset sequences on
$\mathrm{latticeGraph}\,d$:

\begin{itemize}
\item \texttt{linearBox}\,$n : \mathrm{Finset}(\mathrm{Fin}\,1 \to \mathbb Z)$
  $= [0, n)$ — 1D half-line (\texttt{Concrete/LinearBrick.lean}, PRs \#639--\#640).
\item \texttt{stripeBrick2D}\,$w\,n : \mathrm{Finset}(\mathrm{Fin}\,2 \to \mathbb Z)$
  $= [0, n) \times [0, w)$ — 2D rectangular stripe
  (\texttt{Concrete/StripeBrick2D.lean}, PR \#641).
\item \texttt{slabBrick}\,$(\mathrm{widths} : \mathrm{Fin}\,d \to \mathbb N)\,n :
  \mathrm{Finset}(\mathrm{Fin}\,(d+1) \to \mathbb Z)$
  $= [0, n) \times \prod_j [0, \mathrm{widths}\,j)$ — general $(d+1)$-dim
  single-sided slab (\texttt{Concrete/SlabBrick.lean}, PRs \#642--\#643).
\item \texttt{centeredSlab}\,$\mathrm{widths}\,n : \mathrm{Finset}(\mathrm{Fin}\,(d+1) \to \mathbb Z)$
  $= [-n, n) \times \prod_j [0, \mathrm{widths}\,j)$ — two-sided
  centered slab (\texttt{Concrete/CenteredSlab.lean}, PR \#644).
\end{itemize}

Each family supplies the four Fekete hypotheses
(\texttt{hcard\_add}, \texttt{hsuper}, \texttt{hbdd},
\texttt{hcard\_one}) via combinatorial lemmas plus disjoint-union
super-additivity at the $\Lambda$-layer, translation invariance of
$\mathrm{partitionFunction}_\Lambda$, and the subsingleton-congruence
bridge \texttt{partitionFunction\LeanLambda\_congr\_finset}. The uniform upper
bound $\log 2 + |\beta|\cdot((d+1)\cdot|J|+|h|)$ is transported from
\texttt{freeEnergy\_upper\_bound} + the lattice edge-count bound
\texttt{inducedLatticeGraph\_card\_edgeFinset\_le}.

\subsection{Named infinite-volume limits (\S 4.6)}

Each family has a named limit defined via \texttt{Classical.choose}
on the Fekete-convergence existential:
\[
  \mathrm{freeEnergyInfinite\_}\ast\,(\mathrm{hw})\,p\,(\mathrm{hf})
  := \mathrm{Classical.choose}(\mathrm{freeEnergy\_}\ast\,\mathrm{tendsto}\,\ldots),
\]
with the corresponding \texttt{freeEnergy\_*\_tendsto\_freeEnergyInfinite}
convergence lemma.
(PRs \#646 for slab/centered, \#648 for linearBox/stripeBrick2D.)

\subsection{Named-limit properties}

For all four families, under ferromagnetic $0 \le J, 0 \le h, 0 < \beta$:

\begin{itemize}
\item \textbf{Sandwich:} $\log 2 \le \mathrm{freeEnergyInfinite\_}\ast
  \le \log 2 + |\beta|\cdot((d+1)\cdot|J|+|h|)$
  (PR \#650). Sharpened lower bound
  $\log(2\cosh(\beta h)) \le \mathrm{freeEnergyInfinite\_}\ast$
  (PR \#659), combined into a cosh-form sandwich (PR \#660).
\item \textbf{Positivity:} $0 < \mathrm{freeEnergyInfinite\_}\ast$
  (PR \#666), hence nonvanishing (PR \#667).
\item \textbf{Trivial slice:} at $J = 0$,
  $\mathrm{freeEnergyInfinite\_}\ast\,\langle 0, h, \beta\rangle =
  \log(2\cosh(\beta h))$ (PRs \#647, \#648) via per-stage
  $\mathrm{freeEnergy\_J\_zero}$ + \texttt{tendsto\_nhds\_unique}.
\item \textbf{Monotonicity:} $\mathrm{freeEnergyInfinite\_}\ast$ is
  monotone on $\mathrm{Ioi}\,0$ in $\beta$ (PR \#661), and on
  $\mathrm{Ici}\,0$ in $h$ and $J$ (PR \#662). Proof via
  \texttt{le\_of\_tendsto\_of\_tendsto'} on per-stage
  \texttt{freeEnergy\_monotone\_\{$\beta$,h,J\}}.
\item \textbf{Translation invariance of the Fekete limit:} any
  coordinate shift of the sequence converges to the same named limit
  (PR \#652), via \texttt{freeEnergy\LeanLambda\_vaddFinset\_eq} (PR \#255).
\item \textbf{Arbitrary-$h$ convergence via $|h|$-symmetry:} for any
  real $h$, the sequence at $\langle J, h, \beta\rangle$ converges to
  $\mathrm{freeEnergyInfinite\_}\ast\,\langle J, |h|, \beta\rangle$
  (PR \#663), bridging the ferromagnetic infrastructure to arbitrary
  real $h$ via \texttt{freeEnergy\_eq\_abs\_h}.
\end{itemize}

\subsection{Cross-family equivalences}

All four named limits are equal at matching widths (PRs \#651,
\#654, \#655, \#656):
\[
  \mathrm{freeEnergyInfinite\_linearBox} =
  \mathrm{freeEnergyInfinite\_slabBrick}\,(\mathrm{Fin.elim0}), \quad
\]
\[
  \mathrm{freeEnergyInfinite\_stripeBrick2D}\,w =
  \mathrm{freeEnergyInfinite\_slabBrick}\,(\lambda\_,w),
\]
\[
  \mathrm{freeEnergyInfinite\_centeredSlab}\,\mathrm{widths} =
  \mathrm{freeEnergyInfinite\_slabBrick}\,\mathrm{widths}.
\]
The last equivalence uses the Finset identity
$\mathrm{centeredSlab}\,\mathrm{widths}\,n =
\mathrm{shift}_{-n}\,(\mathrm{slabBrick}\,\mathrm{widths}\,(2n))$
plus subsequence convergence at $n \mapsto 2n$.

%% -------------------------------------------------------------------
\section{Random current representation (GJ §5.1 / FV §3.7, §9.7)}
%% -------------------------------------------------------------------

\subsection{Simon--Lieb prep: source connectivity + edge peeling (PRs \#894--\#898)}

Reference: Glimm--Jaffe \cite{glimm-jaffe} \S5.1 (pp.~127--130);
Friedli--Velenik \cite{friedli-velenik} Prop.~9.31 (pp.~399--400).

\paragraph{Sources have even cardinality (PR \#894, step 98).}

\begin{theorem}[\texttt{Current.sources\_card\_even}]
  For any current $n$ and $0 \le \beta J$, $|\partial n|$ is even.
\end{theorem}
\begin{proof}
  Cast to $\mathbb{Z}/2\mathbb{Z}$: $|\partial n| \equiv \sum_{v \in \Lambda} \mathrm{parity}(v) \equiv 0$
  (by the global handshake \texttt{sum\_parity\_eq\_zero}, PR \#893).
  Even follows via \texttt{ZMod.natCast\_eq\_zero\_iff\_even}.
\end{proof}

\paragraph{Source connectivity (PR \#896, step 100).}

\begin{theorem}[\texttt{Current.sources\_reachable\_of\_sources\_eq\_pair}]
  If $\partial n = \{i, j\}$ with $i \ne j$, then $i$ and $j$ are connected
  in the simple graph $n.\mathrm{toSimpleGraph}$.
\end{theorem}
\begin{proof}
  Let $R = \{v : \mathrm{Reachable}(i, v)\}$.
  By the local handshake within $R$ (\texttt{sum\_parity\_closed\_eq\_zero}),
  $|\partial n \cap R|$ is even.
  Since $i \in \partial n \cap R$ (reflexivity), $|\partial n \cap R| \ge 2$.
  Since $|\partial n| = 2$, we get $\partial n \subseteq R$, so $j \in R$.
\end{proof}

\paragraph{Weight peeling identity (PR \#897, step 101).}

\begin{theorem}[\texttt{Current.weight\_pred\_edge}]
  For $n$ with $n_{e_0} \ge 1$:
  \[
    w(n) = \frac{\beta J}{n_{e_0}} \cdot w(n - \mathbf{1}_{e_0}).
  \]
\end{theorem}
\begin{proof}
  Split the weight product at $e_0$ using \texttt{Finset.mul\_prod\_erase}.
  The key identity $\frac{(\beta J)^k}{k!} = \frac{\beta J}{k} \cdot \frac{(\beta J)^{k-1}}{(k-1)!}$
  at $k = n_{e_0}$ is proved by \texttt{pow\_succ} + \texttt{Nat.factorial\_succ} + \texttt{field\_simp}.
\end{proof}

\paragraph{Bridge lemma (PR \#898, step 102).}

\begin{theorem}[\texttt{Current.sources\_sub\_edge\_symmDiff}]
  For $n$ with $n_{e_0} \ge 1$:
  \[
    \partial(n - \mathbf{1}_{e_0}) = \partial n \,\triangle\, \{u, v\}
  \]
  where $e_0 = \{u, v\}$.
\end{theorem}
\begin{proof}
  Since $n_{e_0} \ge 1$, we have $\mathbf{1}_{e_0} \le n$ pointwise.
  Apply \texttt{sub\_sources\_eq\_symmDiff} (PR \#870) with
  $m = \mathbf{1}_{e_0}$, then \texttt{fromEdgeFinset\_singleton\_sources} (PR \#813).
\end{proof}

\paragraph{Edge-peeling bound for \texttt{weightSum} (PR \#898, step 102, full proof).}

\begin{theorem}[\texttt{Current.weightSum\_pair\_le\_edge\_sum}; GJ \S5.1 / FV Prop.~9.31]
  For $i \ne j$ in $\Lambda$ and $0 \le \beta J$:
  \[
    \rho^{\Lambda}_{\{i,j\}} \;\le\; \beta J \sum_{\substack{e \in E(\Lambda) \\ i \in e}}
      \rho^{\Lambda}_{\{i,j\} \triangle e},
  \]
  where $\rho^{\Lambda}_A = \sum_{\partial n = A} w(n)$.
\end{theorem}
\begin{proof}
Apply \texttt{Real.tsum\_le\_of\_sum\_le} to reduce to finite-Finset bounds.
For each finite set $u$ of currents and each $n$ with $\partial n = \{i,j\}$,
pick an active edge $e_0 \ni i$ via \texttt{supportAt\_nonempty\_of\_mem\_sources},
bound $w(n) \le \beta J \cdot w(n - \mathbf{1}_{e_0})$ via \texttt{weight\_le\_mul\_pred\_edge},
absorb into the edge sum via \texttt{Finset.single\_le\_sum},
interchange sums via \texttt{Finset.sum\_comm}, then bound each per-edge finite
sum against $\sum' m$ with $\partial m = \{i,j\} \triangle e$ via
\texttt{Summable.sum\_le\_tsum} using the injection $n \mapsto n - \mathbf{1}_e$
(injectivity by \texttt{fun x \_ y \_ h => funext ...}) and the bridge lemma
\texttt{sources\_sub\_edge\_symmDiff} (establishes $\partial(n - \mathbf{1}_e) = \{i,j\} \triangle e$
when $n_e \ge 1$).
Summability of source-filtered weights is proved via \texttt{summable\_of\_sum\_le}
with the bound from \texttt{sum\_weight\_boundedFinset\_le} (uses \texttt{Finset.sum\_map}
to handle the image sum without \texttt{DecidableEq (Current G \LeanLambda)}, then
\texttt{Fintype.prod\_sum} for the product-sum exchange).
\end{proof}
\noindent\textit{Status}: Full proof. No sorry.

\paragraph{Switching Lemma (PR \#899, step 103).}

\begin{theorem}[\texttt{Current.subFinset\_with\_source\_card\_switching}; GJ \S5.1 Thm.~5.1.2 / FV Thm.~9.35]
  If $\text{symmDiff}(\partial n, A) = B$, then $|\mathcal{S}(n,A)| = |\mathcal{S}(n,B)|$,
  where $\mathcal{S}(n,A) = \{m \leq n : \partial m = A\}$.
\end{theorem}
\begin{proof}
The map $\varphi: m \mapsto n - m$ is an involution on $\mathcal{S}(n,A) \to \mathcal{S}(n,B)$.
Well-definedness: if $m \leq n$ and $\partial m = A$, then $n-m \leq n$ (\texttt{sub\_le\_self})
and $\partial(n-m) = \partial n \triangle A = B$ (\texttt{sub\_hasSources\_iff} +
$\text{symmDiff}(\text{symmDiff}(A,B), A) = B$ by \texttt{symmDiff\_symmDiff\_cancel\_left}).
The inverse direction uses $\text{symmDiff}(\partial n, B) = A$ (same logic).
Involution: $n-(n-m)=m$ by \texttt{sub\_sub\_self\_of\_le}.
Implemented via \texttt{Finset.card\_nbij'}.
\end{proof}

\begin{theorem}[\texttt{Current.sum\_subFinset\_with\_source\_weight\_mul\_weight\_switching}]
  Under the same condition, $\sum_{m \in \mathcal{S}(n,A)} w(m)\,w(n-m) = \sum_{m \in \mathcal{S}(n,B)} w(m)\,w(n-m)$.
\end{theorem}
\begin{proof}
Under $\varphi: m \mapsto n-m$, the summand transforms as
$w(n-m)\cdot w(n-(n-m)) = w(n-m)\cdot w(m) = w(m)\cdot w(n-m)$.
Implemented via \texttt{Finset.sum\_nbij'}.
\end{proof}
\noindent\textit{Status}: Both lemmas fully proved. No sorry.

%-------------------------------------------------------
\subsection*{PR \#900 --- Simon-Lieb Inequality (GJ~\S5.1 / FV Prop.~9.31, p.~428)}
%-------------------------------------------------------

\begin{theorem}[\texttt{Current.weightSum\_eq\_iSup}]
  For $0 \le \beta J$,
  $\mathrm{weightSum}(G,\Lambda,A,\beta,J) = \sup_N \mathrm{CurrentBounded.weightSum}(G,\Lambda,N,A,\beta,J)$.
\end{theorem}
\begin{proof}
  Uniqueness of limits: both sides are the limit of $N\mapsto\mathrm{CurrentBounded.weightSum}(N)$.
  The first limit uses \texttt{tendsto\_weightSum\_atTop\_currentWeightSum} (requires \texttt{summable\_weight\_if\_sources}, PR~\#898); the second uses \texttt{tendsto\_weightSum\_atTop\_iSup\_of\_nonneg}.
\end{proof}

\begin{theorem}[\texttt{boltzmannWeight\_inducedGraph\_eq\_prod\_exp\_of\_h\_zero}]
  For $h=0$: $e^{-\beta H(\sigma)} = \prod_{e \in \mathrm{edgeSet}} e^{\beta J \cdot \mathrm{spinEdgeProduct}(\sigma,e)}$.
\end{theorem}
\begin{proof}
  Unfold at $h=0$: $H = -J\sum_e \mathrm{edgeSpin}(\sigma,e)$.
  Apply \texttt{Real.exp\_sum} then \texttt{Finset.prod\_subtype} (edgeFinset $\to$ edgeSet),
  and replace $\mathrm{edgeSpin}$ by $\mathrm{spinEdgeProduct}$ via the bridge lemma.
\end{proof}

\begin{theorem}[\texttt{correlation\_inducedGraph\_eq\_weightSum\_ratio}]
  For $h=0$, $0\le\beta J$:
  $\langle\sigma^A\rangle = \mathrm{weightSum}(A)\,/\,\mathrm{weightSum}(\emptyset)$.
\end{theorem}
\begin{proof}
  Numerator $= 2^{|\Lambda|}\cdot\mathrm{weightSum}(A)$, denominator $= 2^{|\Lambda|}\cdot\mathrm{weightSum}(\emptyset)$;
  cancel $2^{|\Lambda|}>0$ via \texttt{field\_simp}.
\end{proof}

\begin{theorem}[\texttt{correlation\_inducedGraph\_simon\_lieb}; GJ~\S5.1 / FV Prop.~9.31]
  For $h=0$, $0\le\beta J$, $i\neq j\in\Lambda$:
  \[
    \langle\sigma_i\sigma_j\rangle \;\le\;
    \beta J \sum_{e\ni i} \langle\sigma^{\{i,j\}\triangle e}\rangle.
  \]
\end{theorem}
\begin{proof}
  By \texttt{correlation\_inducedGraph\_eq\_weightSum\_ratio}, both sides are
  $\mathrm{weightSum}(\cdot)/\mathrm{weightSum}(\emptyset)$.
  Since $\mathrm{weightSum}(\emptyset)>0$ (zero current has weight~$1$),
  divide both sides by it via \texttt{div\_le\_div\_iff\_of\_pos\_right}.
  The resulting inequality is exactly \texttt{weightSum\_pair\_le\_edge\_sum} (PR~\#898).
\end{proof}
\noindent\textit{Status}: All lemmas fully proved. No sorry. File: \texttt{Inequalities/SimonLieb.lean}.

%=======================================================
\subsection*{PR \#901 --- High-Temperature Susceptibility Bound (GJ~\S5.1 / FV~\S3.7.3)}
%-------------------------------------------------------

\begin{lemma}[\texttt{symmDiff\_pair\_pair\_of\_ne}; private]
  Let $\alpha$ be a type with decidable equality, $i,j,u\in\alpha$, $i\neq j$, $i\neq u$, $u\neq j$.
  Then $\{i,j\}\triangle\{i,u\} = \{u,j\}$ as \texttt{Finset}s.
\end{lemma}
\begin{proof}
  Rewrite via \texttt{symmDiff\_def}.
  Show $\{i,j\}\setminus\{i,u\}=\{j\}$ and $\{i,u\}\setminus\{i,j\}=\{u\}$
  by \texttt{ext} + \texttt{simp} (\texttt{mem\_sdiff}, \texttt{not\_or})
  + \texttt{rintro}/\texttt{rcases} case analysis.
  Union via \texttt{simp [or\_comm]}.
\end{proof}

\begin{lemma}[\texttt{correlation\_inducedGraph\_nonneg}; private]
  For $0\le\beta J$ and any $A\subseteq\Lambda$,
  $\langle\sigma^A\rangle_\Lambda \ge 0$.
\end{lemma}
\begin{proof}
  By \texttt{correlation\_inducedGraph\_eq\_weightSum\_ratio}:
  $\langle\sigma^A\rangle = \mathrm{weightSum}(A)/\mathrm{weightSum}(\emptyset)$.
  Apply \texttt{div\_nonneg} with \texttt{weightSum\_nonneg} (numerator $\ge 0$)
  and \texttt{weightSum\_empty\_pos} (denominator $> 0$).
\end{proof}

\begin{theorem}[\texttt{correlation\_sum\_le\_of\_high\_temp}; GJ~\S5.1 / FV~\S3.7.3]
  Let $G$ be a graph on $\Lambda$ with $h=0$, $0\le\beta J$, and
  $D\in\mathbb{N}$ bounding the incident-edge count of every vertex.
  If $\beta J D < 1$ then for every $i\in\Lambda$,
  \[
    \sum_{j\neq i}\langle\sigma_i\sigma_j\rangle_\Lambda
    \;\le\;
    \frac{\beta J D}{1-\beta J D}.
  \]
\end{theorem}
\begin{proof}
  Define $T_k = \sum_{j\neq k}\langle\sigma_k\sigma_j\rangle$
  and $M = \max_k T_k$ (exists by \texttt{Finset.exists\_max\_image}).

  \textbf{Step A.}  Apply \texttt{correlation\_inducedGraph\_simon\_lieb} to each $k$:
  \[
    T_k \;\le\; \beta J\sum_{e\ni k}\sum_{j\neq k}
    \langle\sigma^{\{k,j\}\triangle e}\rangle.
  \]

  \textbf{Step B.}  For a fixed edge $e = \{k,u\}$ incident to $k$, expand the inner sum
  over $j\neq k$ by splitting at $j=u$:
  \begin{itemize}
    \item $j=u$: $\{k,u\}\triangle\{k,u\}=\emptyset$, $\langle\sigma^\emptyset\rangle=1$.
    \item $j\neq u$, $j\neq k$: by \texttt{symmDiff\_pair\_pair\_of\_ne},
      $\{k,j\}\triangle\{k,u\}=\{u,j\}$, contributing $\langle\sigma_u\sigma_j\rangle$.
  \end{itemize}
  Hence $\sum_{j\neq k}\langle\sigma^{\{k,j\}\triangle\{k,u\}}\rangle
  \le 1 + T_u \le 1+M$
  (via \texttt{Finset.sum\_erase\_add} + \texttt{Finset.sum\_le\_sum\_of\_subset\_of\_nonneg}).

  \textbf{Step C.}  Summing over edges incident to $k$ (at most $D$ edges):
  $T_k \le \beta J D (1+M)$.

  \textbf{Step D.}  Take $k=k_0 = \operatorname{argmax}T$:
  $M \le \beta J D(1+M)$, so $M(1-\beta J D)\le\beta J D$, giving
  $M\le\beta J D/(1-\beta J D)$ by \texttt{le\_div\_iff\textsubscript{0}}.
  The conclusion $T_i\le M$ follows from the argmax bound.
\end{proof}

\noindent\textit{Status}: All lemmas fully proved. No sorry.
File: \texttt{Inequalities/SimonLieb.lean} (branch \texttt{feat/random-current-simon-lieb-sum-bound}, PR~\#901).
References: Glimm--Jaffe \S5.1~(pp.~73--74); Friedli--Velenik \S3.7.3.

%=======================================================
\subsection*{PR \#902 --- Infinite-Volume Susceptibility Bound (GJ~\S5.1 / FV~\S3.7.3)}
%-------------------------------------------------------

\begin{lemma}[\texttt{edgeFilter\_card\_eq\_degree}; private]
  For a simple graph $G$ on a finite vertex type $V$ and any vertex $v$,
  the number of edges of $G$ incident to $v$ (counted as a Finset filter over
  the subtype $G.\mathrm{edgeSet}$) equals $G.\deg(v)$.
\end{lemma}
\begin{proof}
  Bijection with $G.\mathrm{incidenceFinset}(v)$ via $e \mapsto e.\mathrm{val}$,
  using \texttt{SimpleGraph.edge\_mem\_incidenceSet\_iff}
  and \texttt{card\_incidenceFinset\_eq\_degree}.
\end{proof}

\begin{theorem}[\texttt{susceptibility\LeanLambda\_le\_of\_high\_temp}; GJ~\S5.1 / FV~\S3.7.3]
  For $h=0$, $0\le\beta J$, $D$ bounding the incident-edge count, $\beta J D<1$:
  $\chi_\Lambda(i)
  := \sum_{j\in\Lambda} U_2(i,j)
  = \sum_{j\in\Lambda}\langle\sigma_i\sigma_j\rangle_\Lambda
  \le \beta J D/(1-\beta J D)$.

  \emph{Caveat (Finset collapse)}: in this repo \texttt{susceptibility} sums
  \texttt{truncated2} over \emph{all} $j$, including $j=i$;
  the diagonal $\{i,i\}$ collapses to $\{i\}$ as a Finset, so
  $\langle\sigma_i\sigma_i\rangle_{\rm Finset} = \langle\sigma_i\rangle = 0$ (not $1$).
  Hence the repo bound equals $\sum_{j\ne i}\langle\sigma_i\sigma_j\rangle$,
  and the physical susceptibility bound is $1+\text{(this)} = 1/(1-\beta J D)$.
\end{theorem}
\begin{proof}
  By \texttt{truncated2\_h\_zero}: $U_2(i,j)=\langle\sigma_i\sigma_j\rangle$ at $h=0$.
  By \texttt{magnetization\_zero\_at\_h\_zero} (and Finset collapse):
  diagonal term $\langle\sigma_i\sigma_i\rangle_{\rm Finset}=0$.
  By \texttt{Finset.sum\_filter\_add\_sum\_filter\_not}: split sum.
  Off-diagonal sum $\le \beta J D/(1-\beta J D)$ by
  \texttt{correlation\_sum\_le\_of\_high\_temp} (PR~\#901).
\end{proof}

\begin{theorem}[\texttt{susceptibilityAlongExhaustion\_le\_of\_high\_temp}]
  Per-stage bound: for each $n$,
  $\chi_{\Lambda_n}(i) \le \beta J D/(1-\beta J D)$.
\end{theorem}
\begin{proof}
  If $i\in\Lambda_n$: apply above; else $0 \le \text{bound}$.
\end{proof}

\begin{theorem}[\texttt{susceptibilityInfinite\_le\_of\_high\_temp}]
  $\chi_\infty(i) = \sup_n\chi_{\Lambda_n}(i) \le \beta J D/(1-\beta J D)$.
\end{theorem}
\begin{proof}
  \texttt{ciSup\_le} from the per-stage bound.
\end{proof}

\begin{theorem}[\texttt{susceptibilityInfinite\_latticeGraph\_le\_of\_high\_temp};
  Z\textasciicircum d concrete]
  For the $d$-dimensional lattice graph, $0\le\beta J$, $\beta J\cdot 2d<1$:
  $\chi_\infty(i) \le \beta J\cdot 2d/(1-\beta J\cdot 2d)$.
\end{theorem}
\begin{proof}
  Apply the abstract theorem with $D=2d$;
  degree bound from \texttt{edgeFilter\_card\_eq\_degree}
  $+$ \texttt{inducedLatticeGraph\_degree\_le}.
\end{proof}

\noindent\textit{Status}: All lemmas fully proved. No sorry.
File: \texttt{Inequalities/HighTemp.lean} (branch \texttt{feat/random-current-susceptibility-high-temp}, PR~\#902).
References: Glimm--Jaffe \S5.1~(pp.~73--74); Friedli--Velenik \S3.7.3.

\subsection*{PR \#903 --- High-Temperature Cluster Property (GJ~\S5.1 / FV~\S3.7.3)}

Steps~107--108.  Under \(\mathrm{Ferromagnetic}\,\langle J,0,\beta\rangle\) (i.e.\
$0\le J$, $0<\beta$) and $\beta J D < 1$:

\begin{theorem}[\texttt{truncated2Infinite\_summable\_of\_high\_temp}]
The map $j\mapsto U_2(i,j;\langle J,0,\beta\rangle)$ is summable over $V$.
\end{theorem}

\paragraph{Proof sketch.}
Use \texttt{summable\_of\_sum\_le} with constant $\beta J D/(1-\beta J D)$.
For any finite $s\subseteq V$, the partial sum
$\sum_{j\in s} U_2(i,j) = \sum_{j\in s}\mathrm{correlationInfinite}(\{i,j\})$ (at $h=0$,
by \texttt{truncated2Infinite\_h\_zero}).
The sequence $n\mapsto \sum_{j\in s}\mathrm{correlationAlongExhaustion}(\{i,j\},n)$
converges to this partial sum (\texttt{tendsto\_finset\_sum} +
\texttt{tendsto\_correlationAlongExhaustion\_correlationInfinite}).
For $n$ large enough so that $s\cup\{i\}\subseteq\Lambda_n$, the private helper
\texttt{sum\_correlationAlongExhaustion\_le\_susceptibilityAlongExhaustion} gives
\[
  \textstyle\sum_{j\in s}\mathrm{correlationAlongExhaustion}(\{i,j\},n)
  \;\le\;
  \mathrm{susceptibilityAlongExhaustion}(i,n)
  \;\le\; \frac{\beta J D}{1-\beta J D}.
\]
By \texttt{le\_of\_tendsto}, the partial sum $\sum_{j\in s} U_2(i,j)$ is bounded.

\paragraph{Helper lemma.}
\texttt{sum\_correlationAlongExhaustion\_le\_susceptibilityAlongExhaustion}:
for $i\in\Lambda_n$ and $s\subseteq\Lambda_n$,
$\sum_{j\in s}\mathrm{correlationAlongExhaustion}(\{i,j\},n)
\le \mathrm{susceptibilityAlongExhaustion}(i,n)$.
Proof: at $h=0$, each $\mathrm{correlationAlongExhaustion}(\{i,j\},n)
= \mathrm{truncated2}(\Lambda_n)\,\langle i,h_i\rangle\,\langle j,h_j\rangle$
(via \texttt{correlationAlongExhaustion\_of\_subset}, \texttt{correlation\LeanLambda\_apply},
\texttt{liftFinset} pair form, \texttt{truncated2\_h\_zero}).
After \texttt{Finset.sum\_attach} + \texttt{Finset.sum\_congr}, bound
$\sum_{j\in s.\mathtt{attach}} \le \sum_{j':\uparrow\Lambda_n}$
via \texttt{Finset.sum\_image} injectivity + \texttt{Finset.sum\_le\_sum\_of\_subset\_of\_nonneg}.

\begin{theorem}[\texttt{clusterProperty\_of\_high\_temp}]
Under the same hypotheses, $\mathrm{clusterProperty}\,G\,\Lambda\,\langle J,0,\beta\rangle$.
\end{theorem}

\begin{proof}
Immediate from \texttt{truncated2Infinite\_summable\_of\_high\_temp} +
\texttt{clusterProperty\_of\_summable}.
\end{proof}

\begin{corollary}[\texttt{clusterProperty\_latticeGraph\_of\_high\_temp}]
For the $\mathbb{Z}^d$ lattice graph with cubic exhaustion,
$\mathrm{Ferromagnetic}\,\langle J,0,\beta\rangle$, and $\beta J\cdot 2d < 1$:
the high-temperature cluster property holds.
Proof: apply with $D=2d$, using \texttt{inducedLatticeGraph\_degree\_le}.
\end{corollary}

\paragraph{Physical significance.}
At high temperature ($\beta J \cdot 2d < 1$ on $\mathbb{Z}^d$),
$U_2(i,j)\to 0$ as $|i-j|\to\infty$: the pure-phase cluster property holds.
The exponential-decay rate is formalized in PR~\#916 (Step~110).

File: \texttt{Inequalities/HighTemp.lean} (branch \texttt{feat/random-current-cluster-high-temp}, PR~\#903).
References: Glimm--Jaffe \S5.1~(pp.~72--74); Friedli--Velenik \S3.7.3.

%--------------------------------------------------------------------
\subsection*{PR \#915 / \#916 --- Infinite-Volume Simon-Lieb + Exponential Decay
(GJ~\S5.1~pp.~76--79, pp.~74--75; FV~Prop.~9.31~p.~428)}

\subsubsection*{Step~109: Infinite-volume Simon-Lieb (PR~\#915)}

\begin{theorem}[\texttt{correlationInfinite\_simon\_lieb}]
For $h=0$, $0\le\beta J$, $i\ne j$, and $\neg G.\mathit{Adj}\,i\,j$:
\[
  \langle\sigma_i\sigma_j\rangle_\infty
  \;\le\;
  \beta J\cdot\sum_{k\sim i}\langle\sigma_k\sigma_j\rangle_\infty.
\]
\end{theorem}

\begin{proof}
\texttt{ciSup\_le} reduces to per-stage bounds.
In the subset case, \texttt{correlation\_inducedGraph\_simon\_lieb}
(finite-vol peeling, PR~\#900) is applied; the edge-filter sum maps
injectively into $G.\mathrm{neighborFinset}\,i$ via \texttt{Sym2.Mem.other};
each term is bounded by \texttt{correlationAlongExhaustion\_le\_correlationInfinite};
sum monotonicity closes the goal.
\end{proof}

Concrete $\mathbb{Z}^d$ wrapper: \texttt{correlationInfinite\_simon\_lieb\_latticeGraph}.
File: \texttt{Inequalities/HighTemp.lean} (PR~\#915).

\subsubsection*{Step~110: High-Temperature Exponential Decay (PR~\#916)}

\begin{theorem}[\texttt{hasExponentialDecay\_of\_high\_temp}]
For $0\le\beta J$ and $\beta J D < 1$ ($D = 2d$ on $\mathbb{Z}^d$):
\[
  \mathrm{HasExponentialDecay}\;d\;(\mathrm{cubicExhaustion}\;d)\;
  \langle J,0,\beta\rangle\;(-\log(\beta J D)),
\]
with witness constant $C = 1/(1-\beta J D)$.
Explicitly: $|\langle\sigma_i\sigma_j\rangle_\infty|
  \le C\cdot(\beta J D)^{d_\ell(i,j)}$ for all $i\ne j$.
\end{theorem}

\begin{proof}
Let $P(n)$ be: for all $i\ne j$ with $d_\ell(i,j)\ge n+1$,
$\langle\sigma_i\sigma_j\rangle_\infty\le(\beta J D)^n\cdot\frac{\beta J D}{1-\beta J D}$.

\textit{Base $n=0$}: $d_\ell(i,j)\ge 1$, i.e.\ $i\ne j$.
For each stage $n'$, if $\{i,j\}\subseteq\Lambda_{n'}$:
$\mathrm{correlation}\{i',j'\}
=\mathrm{truncated2}\,i'\,j'$
(by \texttt{truncated2\_h\_zero})
$\le\sum_{k:\uparrow\Lambda_{n'}}\mathrm{truncated2}\,i'\,k$
(by \texttt{Finset.single\_le\_sum} + nonnegativity via
\texttt{correlation\_inducedGraph\_eq\_weightSum\_ratio} + \texttt{Current.weightSum\_nonneg})
$=\chi_{\Lambda_{n'}}(i)
\le\beta J D/(1-\beta J D)$ (Step~106).
Not-in-stage gives $0\le\beta J D/(1-\beta J D)$.
Then \texttt{ciSup\_le} closes the base case.

\textit{Step $n\to n+1$}: $d_\ell(i,j)\ge n+2\ge 2\Rightarrow\neg\mathrm{Adj}\,i\,j$.
Step~109 (Simon-Lieb): $\langle\sigma_i\sigma_j\rangle_\infty\le\beta J\cdot\sum_{k\sim i}\langle\sigma_k\sigma_j\rangle_\infty$.
For each $k\sim i$: $d_\ell(k,j)\ge d_\ell(i,j)-1\ge n+1$ (triangle).
By $P(n)$: $\langle\sigma_k\sigma_j\rangle_\infty\le(\beta J D)^n\cdot\frac{\beta J D}{1-\beta J D}$.
Summing over $|\mathrm{neighborFinset}\,i|\le D=2d$:
$\beta J\cdot D\cdot(\beta J D)^n\cdot\frac{\beta J D}{1-\beta J D}
=(\beta J D)^{n+1}\cdot\frac{\beta J D}{1-\beta J D}$.

Apply $P(N-1)$ where $N=d_\ell(i,j)$:
$\langle\sigma_i\sigma_j\rangle_\infty\le C\cdot(\beta J D)^N$.
For $\beta J D>0$: $(\beta J D)^N=\exp(\log(\beta J D)\cdot N)=\exp(-\alpha N)$
(by \texttt{Real.exp\_log} + \texttt{Real.log\_pow}), $\alpha=-\log(\beta J D)$.
For $\beta J D=0$ ($J=0$ or $\beta=0$): $0^N\le 1=\exp(0)$.
Conclude $\langle\sigma_i\sigma_j\rangle_\infty\le C\cdot e^{-\alpha N}$.
\end{proof}

\noindent\textit{Status}: Fully proved. No sorry. Linter warnings: 0.
File: \texttt{Concrete/LatticeGraphCorrelation/LatticeMassHighTemperature.lean} (moved from \texttt{Concrete/LatticeGraphCorrelation/Inequalities.lean} in PR~\#1926; original PR~\#916).
References: Glimm--Jaffe \S5.1~pp.~74--75; Friedli--Velenik Prop.~9.31~p.~428.

%--------------------------------------------------------------------
\subsection*{PR \#917 --- Positive Lattice Mass at High Temperature
(GJ~\S17.5~pp.~304--306)}

Step~111.  For $0 < \beta J$ and $\beta J D < 1$ ($D = 2d$):

\begin{theorem}[\texttt{latticeMass\_pos\_of\_high\_temp}]
$0 < \mathrm{latticeMass}\;d\;(\mathrm{cubicExhaustion}\;d)\;\langle J,0,\beta\rangle$.
\end{theorem}

\begin{proof}
\textit{Case $d = 0$}: $\mathrm{Fin}\,0\to\mathbb{Z}$ is a singleton; the predicate
$\mathrm{HasExponentialDecay}\;0\;\Lambda\;p\;1$ holds vacuously (no pairs $i \ne j$),
so rate $1 \in \{\alpha\}$ and $\mathrm{latticeMass} \ge 1 > 0$.

\textit{Case $d \ge 1$}: $\beta J (2d) > 0$, hence
$\alpha_0 := -\log(\beta J D) > 0$.
By Step~110, $\mathrm{HasExponentialDecay}\;d\;\Lambda\;p\;\alpha_0$,
so $\alpha_0 \in \{\alpha : \mathrm{NNReal} \mid \mathrm{HasExponentialDecay}\;\cdots\}$
and $\mathrm{latticeMass} \ge (\alpha_0 : \mathrm{ENNReal}) > 0$.
\end{proof}

\noindent\textit{Status}: Fully proved. No sorry.
File: \texttt{Concrete/LatticeGraphCorrelation/LatticeMassHighTemperature.lean} (moved from \texttt{Concrete/LatticeGraphCorrelation/Inequalities.lean} in PR~\#1926; original PR~\#917).
Reference: Glimm--Jaffe \S17.5~pp.~304--306.

\subsection*{PR \#919 --- Lattice Mass Antitonicity in $\beta$ and $J$
(GJ~\S17.5~pp.~304--306)}

Step~112.  For $J \ge 0$ and $0 < \beta_1 \le \beta_2$:

\begin{theorem}[\texttt{latticeMass\_antitone\_beta}]
$\mathrm{latticeMass}\;d\;\Lambda\;\langle J,0,\beta_2\rangle \le
\mathrm{latticeMass}\;d\;\Lambda\;\langle J,0,\beta_1\rangle$.
\end{theorem}

\begin{theorem}[\texttt{latticeMass\_antitone\_J}]
For $0 \le J_1 \le J_2$ and $\beta > 0$:
$\mathrm{latticeMass}\;d\;\Lambda\;\langle J_2,0,\beta\rangle \le
\mathrm{latticeMass}\;d\;\Lambda\;\langle J_1,0,\beta\rangle$.
\end{theorem}

\begin{proof}
We show the $\beta$-case; the $J$-case is identical with $J$-monotonicity in place of $\beta$-monotonicity.

Suppose $\mathrm{HasExponentialDecay}(d,\Lambda,\langle J,0,\beta_2\rangle,\alpha)$
with constant $C \ge 0$.  We construct a witness for $\beta_1$.
At $h = 0$, the truncated two-point function reduces to the correlation:
$U_2^\infty(\beta_i; i,j) = \langle\sigma_i\sigma_j\rangle_\infty^{(\beta_i)}$
(by $Z_2$ symmetry, $\langle\sigma_i\rangle_\infty = 0$, so \texttt{truncated2Infinite\_h\_zero}).
By GKS-I (\texttt{correlationInfinite\_nonneg\_of\_hβJ}), both values are $\ge 0$.
By GKS-II $\beta$-monotonicity (GJ Prop.~4.2.4,
\texttt{correlationInfinite\_monotone\_beta}):
$\langle\sigma_i\sigma_j\rangle_\infty^{(\beta_1)} \le
\langle\sigma_i\sigma_j\rangle_\infty^{(\beta_2)}$.
Hence $|U_2^\infty(\beta_1;i,j)| = U_2^\infty(\beta_1;i,j) \le
U_2^\infty(\beta_2;i,j) = |U_2^\infty(\beta_2;i,j)| \le C\,e^{-\alpha\,\mathrm{dist}(i,j)}$.
So $\mathrm{HasExponentialDecay}(\beta_1,\alpha)$ holds with the same $C$.

Since the valid-rate set for $\beta_2$ is a subset of that for $\beta_1$,
taking $\mathrm{sSup}$ gives the inequality.
\end{proof}

\noindent\textit{Status}: Fully proved. No sorry.
File: \texttt{Concrete/LatticeGraphCorrelation/LatticeMassHighTemperature.lean} (moved from \texttt{Concrete/LatticeGraphCorrelation/Inequalities.lean} in PR~\#1926; original PR~\#919).
References: Glimm--Jaffe \S17.5~pp.~304--306; \S4.2~Props.~4.2.3--4.2.4~pp.~57--60.

\subsection*{Step 113 — \texttt{twoPointFunction\_ge\_tanh\_betaJ\_of\_adj} (GJ \S17.5)}

\begin{theorem}[\texttt{twoPointFunction\_ge\_tanh\_betaJ\_of\_adj}]
Let $d \in \mathbb{N}$, $J \ge 0$, $\beta > 0$, $h = 0$.
Let $r \in \mathbb{Z}^d$ be adjacent to $0$ in $\mathrm{latticeGraph}(d)$.  Then
\[
  \tanh(\beta J) \;\le\; \langle\sigma_0\sigma_r\rangle_\infty^{(J,0,\beta)}.
\]
\end{theorem}

\begin{proof}
Let $G_{\mathrm{single}}$ be the graph on $\mathbb{Z}^d$ with exactly one edge $\{0,r\}$.
Since $0$ and $r$ are adjacent in $\mathrm{latticeGraph}(d)$, we have $G_{\mathrm{single}} \le \mathrm{latticeGraph}(d)$.

We compute $\langle\sigma_0\sigma_r\rangle_\infty^{G_{\mathrm{single}}}$ directly:
for any finite $\Lambda \ni 0,r$ the Hamiltonian on $G_{\mathrm{single}}$ depends only on $(\sigma_0, \sigma_r)$,
so the partition function and numerator factorize with a common factor $2^{|\Lambda|-2}$.
The ratio is $\tanh(\beta J)$, independently of $\Lambda$.
By the eventually-constant convergence lemma for $\mathrm{correlationAlongExhaustion}$
and \texttt{correlationAlongExhaustion\_tendsto\_ciSup},
$\mathrm{correlationInfinite}(G_{\mathrm{single}}) = \tanh(\beta J)$.

Finally, GKS-II subgraph monotonicity (\texttt{correlationInfinite\_monotone\_ambient\_subgraph})
gives the result.
\end{proof}

\noindent\textit{Status}: Fully proved. No sorry.
File: \texttt{Concrete/LatticeGraphCorrelation/LatticeMassHighTemperature.lean} (moved from \texttt{Concrete/LatticeGraphCorrelation/Inequalities.lean} in PR~\#1926; original PR~\#920).
References: Glimm--Jaffe \S17.5~pp.~304--306 (2nd~ed.); \S4.2~(GKS-II subgraph monotonicity).

\subsection*{Step 114 — \texttt{twoPointFunction\_ge\_tanh\_betaJ\_pow\_dist} (GJ \S17.5)}

\begin{theorem}[\texttt{twoPointFunction\_ge\_tanh\_betaJ\_pow\_dist}]
Let $d \in \mathbb{N}$, $J \ge 0$, $\beta > 0$, $h = 0$, and $r \in \mathbb{Z}^d$ with $r \neq 0$.
Write $n = \mathrm{latticeDistance}(d,0,r)$.  Then
\[
  \tanh(\beta J)^n \;\le\; \langle\sigma_0\sigma_r\rangle_\infty^{(J,0,\beta)}.
\]
\end{theorem}

\begin{proof}
Strong induction on $n = \mathrm{latticeDistance}(d,0,r) \ge 1$.

\textbf{Auxiliary lemma} (\texttt{exists\_latticeDistance\_succ\_adj}).
If $\mathrm{latticeDistance}(d,0,r) = n+1$, there exists $v \in \mathbb{Z}^d$ with
$(v,r) \in \mathrm{latticeGraph}(d)$ and $\mathrm{latticeDistance}(d,0,v) = n$.
Proof: find a coordinate $i_0$ with $|r_{i_0}| \ge 1$ (since the $\ell^1$ sum is $\ge 1$);
set $v = r[i_0 \mapsto r_{i_0} \mp 1]$.

\textbf{Base case} ($n = 0$, giving $r$ adjacent to $0$). Apply Step~113 directly.

\textbf{Inductive step} ($n \ge 1$).
Obtain $v$ from the auxiliary lemma.  Since $\mathrm{latticeDistance}(d,0,v) = n \ge 1$, $v \neq 0$,
and the inductive hypothesis gives $\tanh(\beta J)^n \le \langle\sigma_0\sigma_v\rangle_\infty$.

Translation invariance (\texttt{correlationInfinite\_latticeGraph\_cubicExhaustion\_vaddFinset}):
translating by $-v$ gives
\[
  \langle\sigma_v\sigma_r\rangle_\infty
  = \langle\sigma_0\sigma_{r-v}\rangle_\infty
  = \mathrm{twoPointFunction}(d,\langle J,0,\beta\rangle, r-v).
\]
Since $(v,r)$ is an edge and the lattice graph is translation-invariant,
$(0, r-v)$ is also an edge, so Step~113 gives
$\tanh(\beta J) \le \langle\sigma_v\sigma_r\rangle_\infty$.

GKS-II product inequality (\texttt{correlationInfinite\_latticeGraph\_cubicExhaustion\_gks\_second}):
\[
  \langle\sigma_0\sigma_v\rangle_\infty \cdot \langle\sigma_v\sigma_r\rangle_\infty
  \;\le\; \langle\sigma_{\{0,v\}\triangle\{v,r\}}\rangle_\infty
  = \langle\sigma_0\sigma_r\rangle_\infty,
\]
where $\{0,v\}\triangle\{v,r\} = \{0,r\}$ because $0 \neq v$, $0 \neq r$, $v \neq r$.

Combining: $\tanh(\beta J)^{n+1} \le \langle\sigma_0\sigma_v\rangle_\infty \cdot \langle\sigma_v\sigma_r\rangle_\infty \le \langle\sigma_0\sigma_r\rangle_\infty$.
\end{proof}

\noindent\textit{Status}: Fully proved. No sorry.
File: \texttt{Concrete/LatticeGraphCorrelation/LatticeMassHighTemperature.lean} (moved from \texttt{Concrete/LatticeGraphCorrelation/Inequalities.lean} in PR~\#1926; original PR~\#921).
References: Glimm--Jaffe \S17.5~pp.~304--306 (2nd~ed.); \S4.2~(GKS-II).

\subsection*{Step 115 — \texttt{latticeMass\_le\_neg\_log\_tanh\_betaJ} (GJ \S17.5)}

\begin{theorem}[\texttt{latticeMass\_le\_neg\_log\_tanh\_betaJ}]
Let $d \ge 1$, $J > 0$, $\beta > 0$, $h = 0$.  Then
\[
  \mathrm{latticeMass}(d,\mathrm{cubicExhaustion}(d),\langle J,0,\beta\rangle)
  \;\le\; -\log\tanh(\beta J).
\]
Combined with Step~111, this gives the two-sided bound
$-\log(\beta J D) \le \mathrm{latticeMass} \le -\log\tanh(\beta J)$
(where $D = 2d$) in the high-temperature regime $\beta J D < 1$.
\end{theorem}

\begin{proof}
Set $c := -\log\tanh(\beta J) > 0$ (positive since $0 < \tanh(\beta J) < 1$).

Fix $\alpha : \mathbb{R}_{\ge 0}$ with $\mathrm{HasExponentialDecay}(d,\Lambda,p,\alpha)$ witnessed by $C \ge 0$.
We show $\alpha \le c$.

Suppose for contradiction $\alpha > c$, i.e., $\varepsilon := \log\tanh(\beta J) + \alpha > 0$.

\textbf{Archimedean step.}
By the Archimedean property, choose $n_0 \in \mathbb{N}$ with $C < \varepsilon \cdot n_0$.

\textbf{Axis point.}
Let $r_{n_0} := (n_0, 0, \ldots, 0) \in \mathbb{Z}^d$.
The auxiliary lemma \texttt{latticeDistance\_coord\_eq} gives
$\mathrm{latticeDistance}(d,0,r_{n_0}) = n_0$.
Since $n_0 \ge 1$, $r_{n_0} \neq 0$.

\textbf{Lower bound (Step~114).}
\[
  \tanh(\beta J)^{n_0} \;\le\; \mathrm{twoPointFunction}(d,\langle J,0,\beta\rangle,r_{n_0}).
\]

\textbf{Upper bound (HasExponentialDecay).}
At $h = 0$, $\mathrm{twoPointFunction} = |\mathrm{truncated2Infinite}|$, so
\[
  \mathrm{twoPointFunction}(d,p,r_{n_0}) \;\le\; C\,e^{-\alpha n_0}.
\]

\textbf{Rearrangement.}
Combining:
\[
  e^{\varepsilon n_0}
  = \tanh(\beta J)^{n_0} \cdot e^{\alpha n_0}
  \;\le\; C\,e^{-\alpha n_0} \cdot e^{\alpha n_0} = C.
\]

\textbf{Contradiction.}
$e^{\varepsilon n_0} \ge 1 + \varepsilon n_0 > 1 + C > C$.
This contradicts $e^{\varepsilon n_0} \le C$.

\textbf{Conclusion.}
Every $\alpha$ in the supremand satisfies $\alpha \le c$, hence
$\mathrm{latticeMass} = \sup\{\alpha \mid \mathrm{HasExponentialDecay}(\alpha)\} \le c$.
\end{proof}

\noindent\textit{Status}: Fully proved. No sorry.
File: \texttt{Concrete/LatticeGraphCorrelation/LatticeMassHighTemperature.lean} (moved from \texttt{Concrete/LatticeGraphCorrelation/Inequalities.lean} in PR~\#1926; original PR~\#923).
References: Glimm--Jaffe \S17.5~pp.~304--306 (2nd~ed.).

\subsection*{Step 116 — \texttt{truncated4Infinite\_ge\_neg\_pair\_correlations} (GJ \S17.3 (17.3.1))}

\begin{theorem}[\texttt{truncated4Infinite\_ge\_neg\_pair\_correlations}]
Let $G$ be a simple graph, $\Lambda$ an exhaustion, $J > 0$, $\beta > 0$, $h = 0$,
and $i,j,k,l \in V$ pairwise distinct.  Then
\[
  -\bigl(\langle\sigma_i\sigma_k\rangle_\infty \cdot \langle\sigma_j\sigma_l\rangle_\infty
  + \langle\sigma_i\sigma_l\rangle_\infty \cdot \langle\sigma_j\sigma_k\rangle_\infty\bigr)
  \;\leq\; U_4^\infty(i,j,k,l) \;\leq\; 0.
\]
The right inequality is \texttt{truncated4Infinite\_nonpos\_h\_zero} (Cor~4.3.3 / Lebowitz).
The left inequality is Step~116.
\end{theorem}

\begin{proof}[Proof of left inequality]
Unfold $U_4^\infty(i,j,k,l) = \langle\sigma_i\sigma_j\sigma_k\sigma_l\rangle - \langle\sigma_i\sigma_j\rangle\langle\sigma_k\sigma_l\rangle - \langle\sigma_i\sigma_k\rangle\langle\sigma_j\sigma_l\rangle - \langle\sigma_i\sigma_l\rangle\langle\sigma_j\sigma_k\rangle$.

GKS-II (\texttt{correlationInfinite\_gks\_second}): since $\{i,j\}$ and $\{k,l\}$ are disjoint
(all four sites distinct), $\{i,j\} \triangle \{k,l\} = \{i,j,k,l\}$, so
$\langle\sigma_i\sigma_j\rangle \cdot \langle\sigma_k\sigma_l\rangle
\leq \langle\sigma_i\sigma_j\sigma_k\sigma_l\rangle$.

Substituting: $U_4^\infty(i,j,k,l) \geq -\langle\sigma_i\sigma_k\rangle\langle\sigma_j\sigma_l\rangle - \langle\sigma_i\sigma_l\rangle\langle\sigma_j\sigma_k\rangle$.
\end{proof}

\noindent\textit{Status}: Fully proved. No sorry.
Files: \texttt{AmbientLattice.lean} and
\path{IsingModel/Concrete/LatticeGraphCorrelation/InfiniteVolumeCorrelationInequalities.lean}
(PR~\#925; moved to a narrow child module in PR~\#1916).
References: Glimm--Jaffe \S17.3~p.~308~(17.3.1)~(2nd~ed.); \S4.2~Thm~4.2.3~(GKS-II).

\subsection*{Step 117a — \texttt{hasDerivAt\_correlation\_beta} (GJ \S17.5 pp.~310--311)}

\begin{theorem}[\texttt{hasDerivAt\_correlation\_beta}]
For a finite simple graph $G$, parameters $\langle J, 0, \beta\rangle$, and $A \subseteq V$:
\[
  \frac{d}{d\beta}\langle\sigma^A\rangle_\beta
  = J \sum_{e=\langle u,v\rangle \in E}
    \bigl[\langle\sigma^{A\triangle\{u,v\}}\rangle_\beta
          - \langle\sigma^A\rangle_\beta \cdot \langle\sigma_u\sigma_v\rangle_\beta\bigr].
\]
\end{theorem}

\begin{proof}
\textbf{Step 1 (Quotient rule).}
Since the Hamiltonian $H(\sigma) = -J\sum_e\sigma_{e+}\sigma_{e-}$ is independent of $\beta$,
\[
  \frac{d}{d\beta}\mathrm{bw}(\sigma) = -H(\sigma)\cdot\mathrm{bw}(\sigma), \qquad
  \frac{d}{d\beta}Z = \sum_\sigma(-H)\cdot\mathrm{bw}.
\]
By the quotient rule:
$\dfrac{d}{d\beta}\langle F\rangle = \langle F\cdot(-H)\rangle - \langle F\rangle\langle{-H}\rangle$.

\textbf{Step 2 (Expand $-H$).}
At $h=0$: $-H(\sigma) = J\sum_{e\in E}\mathrm{edgeSpin}_e(\sigma)$.

\textbf{Step 3 (Spin product identity).}
By \texttt{spinProduct\_mul}: $\sigma^A \cdot \sigma^{\{u,v\}} = \sigma^{A\triangle\{u,v\}}$.
Hence $\langle\sigma^A\cdot\mathrm{edgeSpin}_e\rangle = \langle\sigma^{A\triangle\{e_1,e_2\}}\rangle$.

\textbf{Conclusion.}
$\dfrac{d}{d\beta}\langle\sigma^A\rangle = J\sum_e\bigl[\langle\sigma^{A\triangle\{e_1,e_2\}}\rangle - \langle\sigma^A\rangle\langle\sigma_{e_1e_2}\rangle\bigr]$.
\end{proof}

\noindent\textit{Status}: Fully proved. No sorry.
File: \texttt{BetaDerivative.lean} (PR~\#926).
References: Glimm--Jaffe \S17.5~pp.~310--311~(2nd~ed.)~(implicit in pseudo-mass proof).

\subsection*{Step 117b — \texttt{correlation\_beta\_deriv\_le\_lebowitz} (GJ \S17.5 p.312)}

\begin{theorem}[\texttt{summand\_le\_lebowitz\_of\_disjoint} + \texttt{correlation\_beta\_deriv\_le\_lebowitz}]
For pairwise distinct $r,s,u,v \in V$ and ferromagnetic $\langle J,0,\beta\rangle$:
\[
  \langle\sigma_r\sigma_s\sigma_u\sigma_v\rangle
    - \langle\sigma_{rs}\rangle\langle\sigma_{uv}\rangle
  \;\leq\;
  \langle\sigma_{ru}\rangle\langle\sigma_{sv}\rangle
  + \langle\sigma_{rv}\rangle\langle\sigma_{su}\rangle.
\]
Hence for the $\beta$-derivative and $r \neq s$:
\[
  \frac{d}{d\beta}\langle\sigma_r\sigma_s\rangle_{\beta}
  \;\leq\;
  J\sum_{e \in E(G)}
    \bigl[\langle\sigma_r\sigma_{e_1}\rangle\langle\sigma_s\sigma_{e_2}\rangle
        + \langle\sigma_r\sigma_{e_2}\rangle\langle\sigma_s\sigma_{e_1}\rangle\bigr]
  + J\cdot|E(G)|.
\]
The additive $J\cdot|E(G)|$ accounts for degenerate edges (incident to $r$ or $s$);
a tighter bound is $J\cdot(\deg r + \deg s)$.
\end{theorem}

\begin{proof}[Proof of \texttt{summand\_le\_lebowitz\_of\_disjoint}]
Since $r,s,u,v$ are pairwise distinct, $\{r,s\}\cap\{u,v\}=\emptyset$, so
$\{r,s\}\triangle\{u,v\}=\{r,s,u,v\}$.
Apply Cor.~4.3.3 (\texttt{cor\_4\_3\_3}): $U_4(r,s,u,v)\le 0$.
Unfolding $U_4 = \langle\sigma_{rsuv}\rangle - \langle\sigma_{rs}\rangle\langle\sigma_{uv}\rangle
- \langle\sigma_{ru}\rangle\langle\sigma_{sv}\rangle - \langle\sigma_{rv}\rangle\langle\sigma_{su}\rangle\le 0$
gives the result.

The full derivative bound follows from Step~117a (\texttt{hasDerivAt\_correlation\_beta})
and a case split per edge: non-degenerate edges use the above; degenerate ones
are bounded via $\langle\sigma_{r,s}\triangle\{u,v\}\rangle - \langle\sigma_{rs}\rangle\langle\sigma_{uv}\rangle\le 1$
(\texttt{summand\_le\_one}).
\end{proof}

\noindent\textit{Status}: Fully proved. No sorry.
File: \texttt{BetaDerivative.lean} (PR~\#926).
References: Glimm--Jaffe \S17.5~p.~312~(2nd~ed.); Cor.~4.3.3~(Lebowitz).

\subsection*{Step 117c — Pseudo-mass profile \texttt{pseudoMassG} (GJ \S17.5 pp.~311--312)}

\begin{definition}[\texttt{pseudoMassG}]
For $\alpha \in \mathbb{N}$, $\alpha \ge 1$ (a natural-number specialisation of GJ's general $\alpha > d/2$),
$r > 0$, and $t \ge 0$, define the pseudo-mass profile
\[
  g(t, r, \alpha) \;=\; \frac{2\,e^{-tr}}{1 + (tr)^{\alpha}}.
\]
The pseudo-mass $m^-(x,y,\beta,A)$ is then the unique $t \ge 0$ satisfying
$g(t, \mathrm{dist}(x,y), \alpha) = \langle\sigma_x\sigma_y\rangle_{\beta,A}$ whenever the
correlation lies in $(0,2)$.
\end{definition}

\begin{theorem}[\texttt{pseudoMassG\_zero}, \texttt{pseudoMassG\_pos}, \texttt{pseudoMassG\_le\_two}]
$g(0,r,\alpha) = 2$; for $t \ge 0$, $r > 0$: $0 < g(t,r,\alpha) \le 2$.
\end{theorem}

\begin{theorem}[\texttt{pseudoMassG\_strictAntiOn}]
For $r > 0$ and $\alpha \ge 1$, $g(\,\cdot\,,r,\alpha)$ is strictly decreasing on $[0,\infty)$.
\end{theorem}
\begin{proof}
For $s < t$ in $[0,\infty)$: the numerator $2e^{-tr} \le 2e^{-sr}$ (exp is decreasing)
and the denominator $1+(tr)^\alpha > 1+(sr)^\alpha$ (power strictly increasing).
Hence $N(t)/D(t) \le N(s)/D(t) < N(s)/D(s)$.
Applied via \texttt{div\_lt\_div₀'} in Lean.
\end{proof}

\begin{theorem}[\texttt{pseudoMassG\_tendsto\_zero}]
$g(t,r,\alpha) \to 0$ as $t \to \infty$ for $r > 0$.
\end{theorem}
\begin{proof}
Squeeze: $0 < g(t,r,\alpha) \le 2e^{-tr} \to 0$ since $tr \to +\infty$.
\end{proof}

\begin{theorem}[\texttt{pseudoMassG\_exists\_of\_mem\_Ioo}]
For $c \in (0,2)$, $r > 0$, $\alpha \ge 1$, there exists $t \ge 0$ with $g(t,r,\alpha) = c$.
\end{theorem}
\begin{proof}
$g$ is continuous on $[0,\infty)$, $g(0)=2 > c$, and $g(T) < c$ for large $T$
(by the previous theorem, take $T$ with $g(T) < c/2$).
The Intermediate Value Theorem on $[0,T]$ yields the required $t$.
Applied via \texttt{intermediate\_value\_Icc'} in Lean.
\end{proof}

\begin{theorem}[\texttt{pseudoMassG\_unique}]
The solution is unique: if $g(t_1,r,\alpha) = g(t_2,r,\alpha) = c$ with $t_1, t_2 \ge 0$,
then $t_1 = t_2$.
\end{theorem}
\begin{proof}
Immediate from strict antitonicity (\texttt{StrictAntiOn.injOn}).
\end{proof}

\noindent\textit{Status}: Fully proved. No sorry.
File: \texttt{PseudoMass.lean} (PR~\#926).
References: Glimm--Jaffe \S17.5~pp.~311--312~(2nd~ed.), equation~(17.5.3).

\subsection*{Step 117d — \texttt{pseudoMassG\_hasDerivAt} + \texttt{pseudoMass} definition (GJ \S17.5)}

\begin{theorem}[\texttt{pseudoMassG\_hasDerivAt}]
For $\alpha \in \mathbb{N}$, $r > 0$, $t \ge 0$:
\[
  \frac{d}{dt}\,g(t,r,\alpha) \;=\;
  \frac{-2r\,e^{-(tr)}\,(1+(tr)^\alpha) - 2\,e^{-(tr)}\,\alpha(tr)^{\alpha-1}\,r}
       {(1+(tr)^\alpha)^2}.
\]
The hypothesis $\alpha \ge 1$ is \emph{not} required; the formula holds for all $\alpha \in \mathbb{N}$.
\end{theorem}
\begin{proof}
Quotient rule: numerator $f(t) = 2e^{-(tr)}$ has derivative $f'=-2r\,e^{-(tr)}$;
denominator $h(t) = 1+(tr)^\alpha$ has derivative $h'=\alpha(tr)^{\alpha-1}\cdot r$.
Applied via \texttt{HasDerivAt.div} in Lean.
\end{proof}

\begin{theorem}[\texttt{pseudoMassG\_deriv\_neg}]
For $\alpha \in \mathbb{N}$, $r > 0$, $t > 0$: $g'(t,r,\alpha) < 0$.
\end{theorem}
\begin{proof}
Numerator $= -2r\,e^{-(tr)}\,(1+(tr)^\alpha) - 2e^{-(tr)}\,\alpha(tr)^{\alpha-1}\,r < 0$
since $r,e^{-(tr)}>0$ and $1+(tr)^\alpha+\alpha(tr)^{\alpha-1}\ge 1>0$.
Denominator $(1+(tr)^\alpha)^2>0$.
\end{proof}

\begin{definition}[\texttt{pseudoMass}]
For $\alpha \ge 1$, $r>0$, $c\in(0,2)$, the pseudo-mass $m^-(c,r)$ is defined as the
$\mathsf{Classical.choose}$ of the existence theorem \texttt{pseudoMassG\_exists\_of\_mem\_Ioo}:
the unique $t\ge 0$ with $g(t,r,\alpha)=c$.
Key lemmas: \texttt{pseudoMass\_spec} (defining equation), \texttt{pseudoMass\_nonneg}
($0 \le m^-$), \texttt{pseudoMass\_eq\_iff} (uniqueness characterisation).
\end{definition}

\noindent\textit{Status}: Fully proved. No sorry.
File: \texttt{PseudoMass.lean} (PR~\#926).
References: Glimm--Jaffe \S17.5~pp.~311--312~(2nd~ed.).

\subsection*{Step 117e — \texttt{pseudoMass\_deriv\_formula} (GJ \S17.5 p.312, implicit differentiation)}

\begin{theorem}[\texttt{pseudoMass\_deriv\_formula}]
Let $\alpha \in \mathbb{N}$, $r > 0$.
Suppose $h, c : \mathbb{R} \to \mathbb{R}$ satisfy:
\begin{itemize}
  \item $\mathrm{HasDerivAt}\;h\;h'\;\beta$ and $\mathrm{HasDerivAt}\;c\;c'\;\beta$,
  \item $\forall\,\beta',\; g(h(\beta'), r, \alpha) = c(\beta')$ (defining equation),
  \item $h(\beta) > 0$.
\end{itemize}
Then
\[
  h' \;=\; \frac{c'(\beta)}{g'(h(\beta), r, \alpha)}.
\]
\end{theorem}
\begin{proof}
Chain rule: \texttt{HasDerivAt.comp} gives $\mathrm{HasDerivAt}\,(g\circ h)\,(g'(h\beta)\cdot h')\,\beta$.
Since $g \circ h = c$, uniqueness of derivatives gives $g'(h\beta)\cdot h' = c'(\beta)$.
Then $h' = c'/g'$ since $g'(h\beta) < 0 \ne 0$ (by \texttt{pseudoMassG\_deriv\_neg}).
Applied via \texttt{field\_simp} + \texttt{linarith}.
\end{proof}

\begin{remark}
The full GJ §17.5 Lipschitz estimate $m^{-\,2\alpha}\,\frac{dm^-}{d\sigma} \le \mathrm{const}$
additionally requires bounding the Lebowitz sum using the lattice convolution inequality
$\sum_{z \in \mathbb{Z}^d} \frac{1}{|x-z|^\alpha|y-z|^\alpha} \le C\,|x-y|^{d-2\alpha}$
(for $\alpha > d/2$).  This convolution estimate (Step~117f+) is left for future work.
\end{remark}

\noindent\textit{Status}: Fully proved. No sorry.
File: \texttt{PseudoMass.lean} (PR~\#926).
References: Glimm--Jaffe \S17.5~p.~312~(2nd~ed.).

%% ============================================================================
%% Steps 117f--h: Lipschitz estimate + HLS + Theorem 17.5.1 (in progress)
%% ============================================================================

\subsection*{Step 117f -- \texttt{pseudoMassG\_deriv\_abs\_ge}: Derivative lower bound}

\begin{theorem}[\texttt{pseudoMassG\_deriv\_abs\_ge}]
For $\alpha \in \mathbb{N}$, $r > 0$, $t \ge 0$:
\[
  r \cdot g(t,r,\alpha) \le |g'(t,r,\alpha)|
\]
where $g(t,r,\alpha) = \frac{2e^{-tr}}{1+(tr)^\alpha}$.
\end{theorem}

\begin{proof}
Algebraic computation: $|g'| - r \cdot g = \frac{2e^{-tr} \cdot \alpha \cdot (tr)^{\alpha-1} \cdot r}{(1+(tr)^\alpha)^2} \ge 0$
for all $t \ge 0, r > 0, \alpha \in \mathbb{N}$.
\end{proof}

\noindent\textit{Status}: Fully proved. No sorry.
File: \texttt{PseudoMass.lean} (PR~\#927).
References: Glimm--Jaffe \S17.5~p.~312~(2nd~ed.).

\subsection*{Step 117g -- \texttt{pseudoMass\_pos} and \texttt{pseudoMass\_strictAnti}}

\begin{theorem}[\texttt{pseudoMass\_pos}]
For $\alpha \ge 1$, $r > 0$, $c \in (0,2)$:
\[
  0 < \mathrm{pseudoMass}(\alpha, r, c)
\]
\end{theorem}

\begin{theorem}[\texttt{pseudoMass\_strictAnti}]
For $\alpha \ge 1$, $r > 0$, $c_1, c_2 \in (0,2)$ with $c_1 < c_2$:
\[
  \mathrm{pseudoMass}(\alpha, r, c_2) < \mathrm{pseudoMass}(\alpha, r, c_1)
\]
\end{theorem}

\begin{proof}
\texttt{pseudoMass\_pos}: Since $g(0) = 2 > c$ and $g$ is strictly decreasing (Step~117f--g),
by IVT there exists unique $t > 0$ with $g(t) = c$, hence $\mathrm{pseudoMass} > 0$.

\texttt{pseudoMass\_strictAnti}: If $c_1 < c_2$, then by strict antitonicity $m_1^- > m_2^-$,
hence the result.
\end{proof}

\noindent\textit{Status}: Fully proved. No sorry. Partial Lemma 17.5.2 (lower bound).
File: \texttt{PseudoMass.lean} (PR~\#928).
References: Glimm--Jaffe \S17.5~p.~311~(2nd~ed.).

\subsection*{Step 117h -- Discrete HLS inequality + Theorem 17.5.1 (in progress)}

\subsubsection*{Discrete Hardy--Littlewood--Sobolev inequality (Axiom)}

Since the discrete HLS inequality is not available in Mathlib, we axiomatize it:

\begin{axiom}[\texttt{discrete\_hls\_inequality}]
For $\alpha \in \mathbb{N}$ with $\alpha \ge 1$ (specializing GJ's $\alpha > d/2$) and $2\alpha > d$:
\[
  \sum_{z \in \mathbb{Z}^d} \frac{1}{|x-z|^\alpha |y-z|^\alpha} \le C_{\alpha,d} \cdot |x-y|^{d-2\alpha}
\]
for all $x, y \in \mathbb{Z}^d$ with $x \ne y$.
\end{axiom}

\textit{Reference}: Glimm--Jaffe 2nd~ed. (1987), \S17.5 (pp.~345--347).

\subsubsection*{Lemma 17.5.2 and Theorem 17.6.1 (sketch)}

Combining Steps 117a--g with the HLS axiom:

\begin{theorem}[Lemma 17.5.2 (partial)]
Lower bound: $0 < m^-(\beta)$ for all $c \in (0, 2)$ (Step 117g).

Upper bound: $m(\beta) \le C \cdot m^-(\beta)$ for some constant $C > 0$ (via HLS axiom).
\end{theorem}

\begin{theorem}[Theorem 17.6.1: Continuity at critical point]
The lattice mass function $m(\beta)$ is continuous at the critical point $\beta = \beta_c$.

\textbf{Proof sketch}: Lemma 17.5.2 bounds + implicit differentiation (Step 117e)
+ derivative bounds (Step 117f) + discrete HLS $\Rightarrow$ Lipschitz in $\beta$ $\Rightarrow$ continuity at $\beta_c$.
\end{theorem}

\noindent\textit{Status}: Sketch formalization (PR~\#929). Full Lipschitz proof deferred.
File: \texttt{PseudoMass.lean}.
References: Glimm--Jaffe 2nd~ed. (1987), \S17.5 pp.~345--347, \S17.6 pp.~348--351.

%% ---------------------------------------------------------------
\subsection*{Step 118 -- Theorem 17.6.1: $\partial/\partial h$ of correlations (Done)}

\subsubsection*{Mathematical structure}

The external field $h$ enters the Ising Hamiltonian as:
\[
  H(\sigma; J, h, \beta) = -J\sum_{\{i,j\} \in E} \sigma_i\sigma_j - h \sum_{i \in \Lambda} \sigma_i
  = \text{interactionEnergy} + \text{externalFieldEnergy}(h, \sigma)
\]

where $\text{externalFieldEnergy}(h, \sigma) = -h \cdot M(\sigma)$ and $M(\sigma) = \sum_{i \in \Lambda} \text{sign}(\sigma_i)$ is the \emph{total magnetization}.

\subsubsection*{Derivative formulas}

Since $\partial_h H = -M(\sigma)$, the chain rule gives:
\[
  \frac{\partial}{\partial h} e^{-\beta H(\sigma)} = \beta M(\sigma) \cdot e^{-\beta H(\sigma)}
\]

Applying the quotient rule to $\langle F \rangle_h = Z(h)^{-1} \sum_\sigma F(\sigma) e^{-\beta H(\sigma; h)}$:
\[
  \frac{\partial}{\partial h} \langle F \rangle_h
  = \beta \left( \langle F \cdot M \rangle_h - \langle F \rangle_h \cdot \langle M \rangle_h \right)
\]

In particular, for spin products:
\[
  \frac{\partial}{\partial h} \langle \sigma^A \rangle_h
  = \beta \left( \langle \sigma^A \cdot M \rangle_h - \langle \sigma^A \rangle_h \cdot \langle M \rangle_h \right)
\]

\subsubsection*{Lean 4 formalization}

\begin{theorem}[\texttt{hasDerivAt\_correlation\_field}]
For any finite graph $G$, coupling $J$, field $h$, inverse temperature $\beta$, and index set $A$:
\[
  \mathtt{HasDerivAt}\ (\lambda h' \Rightarrow \langle \sigma^A \rangle_{J,h',\beta})
  \ (\beta \cdot (\langle \sigma^A M \rangle - \langle \sigma^A \rangle \langle M \rangle))\ h
\]
\end{theorem}

\noindent\textbf{Proof}: Compose with \texttt{hasDerivAt\_gibbsExpectation\_field} via quotient rule.
Intermediate steps: \texttt{externalFieldEnergy\_eq}, \texttt{hasDerivAt\_hamiltonian\_field},
\texttt{hasDerivAt\_boltzmannWeight\_field} (chain rule via \texttt{HasDerivAt.comp}),
\texttt{hasDerivAt\_partitionFunction\_field}, \texttt{hasDerivAt\_weightedBoltzmannSum\_field}.
Ring closure uses \texttt{simp\_rw} + \texttt{Finset.mul\_sum} for $\beta$-factoring.

\noindent\textit{Status}: Complete formalization (PR~\#931).
File: \texttt{FieldDerivative.lean}.
Reference: Glimm--Jaffe 2nd~ed. (1987), \S17.6 pp.~348--351.

\subsection*{Step 123 -- Antitonicity of pseudo-mass in $\beta$ (Done)}

\subsubsection*{Mathematical content}

Combining Steps 117g and 122:
\begin{itemize}
  \item $\beta \mapsto c(\beta) = \langle \sigma^A \rangle_\beta$ is monotone increasing (Step 122).
  \item $c \mapsto \mathrm{pseudoMass}(c)$ is strictly decreasing on $(0,2)$ (Step 117g).
  \item Therefore $\beta \mapsto \mathrm{pseudoMass}(c(\beta))$ is antitone (decreasing) on $\{0 < \beta\}$.
\end{itemize}

This captures the physical picture: higher inverse temperature $\beta$ leads to
stronger correlations but smaller mass (mass decreases toward zero at $\beta_c$).

\subsubsection*{Lean 4 formalization}

\begin{theorem}[\texttt{pseudoMass\_comp\_corr\_antitoneOn\_beta}]
For $J \geq 0$ and $c(\beta) \in (0,2)$ for all $\beta > 0$:
\[
  \mathrm{AntitoneOn}\;(\lambda\beta \Rightarrow
    \mathrm{pseudoMass}(c(\beta)))\;(\mathrm{Ioi}\;0)
\]
\end{theorem}

\noindent\textit{Status}: Complete (PR~\#937, Step 123).
File: \texttt{PseudoMass.lean}.
Reference: Derived from Steps 117g + 122; implicit in Glimm--Jaffe 2nd~ed., \S17.5, pp.~311--312.

\subsection*{Step 127 -- Lebowitz--exponential product bound (Done)}

\noindent\textbf{Theorems} (\texttt{summable\_truncated2Infinite\_prod\_of\_hasExponentialDecay},
\texttt{tsum\_truncated2Infinite\_prod\_le}):
Under \texttt{HasExponentialDecay} with constant $C \geq 0$ and rate $\alpha > 0$:
\begin{enumerate}
\item The product sum $\sum_{z \in \mathbb{Z}^d} U_2^\infty(x,z) \cdot U_2^\infty(y,z)$
  is summable over $\mathbb{Z}^d$.
\item Explicit bound:
  $\sum_z U_2^\infty(x,z) U_2^\infty(y,z) \leq (C+1)^2 \cdot 2C(\alpha/2,d) \cdot e^{-\alpha d(x,y)/4}$,
  where $C(\alpha/2,d) = \sum_z e^{-(\alpha/2)d(0,z)}$.
\end{enumerate}

\noindent\textbf{Proof sketch}:
Uniform bound: the private lemma \texttt{truncated2Infinite\_le\_hDecay\_uniform}
gives $U_2^\infty(i,z) \leq (C+1) e^{-(\alpha/2)d(i,z)}$ for all $z$ (including $i=z$),
using $U_2^\infty(i,i) \leq 1 \leq C+1$ (GKS-II) for the diagonal and
$C e^{-\alpha d} \leq (C+1) e^{-(\alpha/2)d}$ for $d \geq 0$ off-diagonal.
The product is bounded by $(C+1)^2 e^{-(\alpha/2)d(x,z)} e^{-(\alpha/2)d(y,z)}$.
Summability follows from \texttt{summable\_exp\_neg\_dist} with rate $\alpha/2$.
The tsum bound applies \texttt{lattice\_exp\_sum\_conv\_le} (Step~126) with rate $\alpha/2$.

\noindent\textit{Status}: Complete (PR~\#943, Step 127, Issue~\#940 PR~N+2).
File: \texttt{Concrete/LatticeGraphCorrelation/LatticeMassPseudoMassTransfer.lean} (moved from \texttt{Concrete/LatticeGraphCorrelation/Inequalities.lean} in PR~\#1927; original PR~\#943).
Reference: GJ 1st~ed., \S17.5, applying Lemma~17.5.2 (pp.~311--312).

\subsection*{Step 128 -- Discrete polynomial HLS summability (Done)}

\noindent\textbf{Statement.}
For $\gamma > d$, the function $z \mapsto (1 + d(0,z))^{-\gamma}$ is summable over $\mathbb{Z}^d$.

\noindent\textbf{Proof outline.}
\begin{enumerate}
\item \emph{Step A -- $\mathbb{Z}$ summability} (\texttt{summable\_pow\_neg\_Z}):
  For $\beta > 1$, $\sum_{n:\mathbb{Z}} (1+|n|)^{-\beta} < \infty$.
  Proof: split into $\mathbb{N}$ part and negative part via
  \texttt{Summable.of\_nat\_of\_neg\_add\_one}, compare each with
  \texttt{Real.summable\_one\_div\_nat\_add\_rpow}.

\item \emph{Step B -- Product summability} (\texttt{summable\_prod\_pow\_neg\_lattice}):
  For $\beta > 1$,
  $\sum_{z:\mathbb{Z}^d} \prod_i (1+|z_i|)^{-\beta} < \infty$.
  By induction on $d$ using \texttt{summable\_prod\_of\_nonneg}
  (Fubini criterion for nonneg product sums).

\item \emph{Step C -- AM-GM comparison} (\texttt{one\_add\_dist\_rpow\_neg\_le}, private):
  For $d \geq 1$ and $\gamma > 0$:
  $(1 + d(0,z))^{-\gamma} \leq \prod_i (1+|z_i|)^{-\gamma/d}$.
  Proof: AM-GM (\texttt{geom\_mean\_le\_arith\_mean\_weighted}) gives
  $\prod (1+|z_i|)^{1/d} \leq 1 + \sum|z_i|/d \leq 1 + d(0,z)$.
  Raise to $-\gamma$ using \texttt{Real.rpow\_le\_rpow\_iff\_of\_neg}.

\item \emph{Step D -- Main} (\texttt{summable\_pow\_neg\_latticeDistance}):
  Apply comparison with $\beta = \gamma/d > 1$ (since $\gamma > d$).
\end{enumerate}

\noindent\textit{Status}: Complete (PR~\#944, Step 128, Issue~\#940 PR~N+3).
File: \texttt{IsingModel/PolyDecay.lean}.
Reference: GJ 1st~ed., \S17.5, pp.~310--312.

\subsection*{Step 135 -- Coupling derivative formula for bond subsets (Done)}

\noindent\textbf{Reference.} Glimm--Jaffe, \emph{Quantum Physics}, 2nd ed., \S17.8 pp.~316--318.

\noindent\textbf{Setup.}
For a finite graph $G=(V,E)$, edge subset $E_0\subseteq E$, Ising parameters $(J,h,\beta)$ with $\beta,J>0$,
and $s\in[0,1]$, the \emph{scaled Boltzmann weight} is:
\[w_s(\sigma) := w_G(\sigma)\cdot\exp\!\bigl(-\beta(1-s)J\sum_{e\in E_0}\sigma_e\bigr),\]
where $\sigma_e = \sigma_k\sigma_l$ for $e=\{k,l\}$.
At $s=0$: weight with $E_0$-bonds absent. At $s=1$: full model.

\noindent\textbf{Main results.}
\begin{itemize}
  \item \texttt{hasDerivAt\_scaledBoltzmannWeight}:
    $\dfrac{d}{ds}w_s(\sigma) = \beta J\Bigl(\sum_{e\in E_0}\sigma_e\Bigr)w_s(\sigma)$.
  \item \texttt{hasDerivAt\_scaledPartitionFunction}:
    $\dfrac{d}{ds}Z_s = \beta J\,Z_s\bigl\langle\sum_{E_0}\sigma_e\bigr\rangle_s$.
  \item \texttt{hasDerivAt\_scaledCorrelation} (main):
    \[\frac{d}{ds}\langle\sigma^A\rangle_s = \beta J\sum_{e=(u,v)\in E_0}
      \Bigl(\langle\sigma^{A\triangle\{u,v\}}\rangle_s - \langle\sigma^A\rangle_s\cdot\langle\sigma^{\{u,v\}}\rangle_s\Bigr).\]
\end{itemize}

\noindent\textbf{Proof sketch.}
Apply the chain rule to $\exp(-\beta(1-s)J X_\sigma)$ where $X_\sigma=\sum_{E_0}\sigma_e$.
Then use the quotient rule for the Gibbs expectation $Z_s^{-1}\sum_\sigma F(\sigma)w_s(\sigma)$,
and the identity $\sigma^A\cdot\sigma^e=\sigma^{A\triangle e}$ (\texttt{spinProduct\_mul}).

\noindent\textit{Status}: Complete (PR~\#952 Step 135, Issue~\#951 Part~1).

\subsection*{Step 136 -- Ball-boundary Simon-Lieb inequality (Done)}

\noindent\textbf{Reference.} Glimm--Jaffe, \emph{Quantum Physics}, 2nd ed., \S17.8 pp.~316--318.

\noindent\textbf{New infrastructure.}
\begin{itemize}
  \item \texttt{scaledCorrelation\_nonneg}: GKS-I for the scaled model ($\langle\sigma^A\rangle_s\ge 0$
    for ferromagnetic $p$ and $s\ge 0$), proved via the product form of $w_s$ and
    \texttt{hasNonnegCorrelations\_edge\_site\_product} with per-edge couplings.
  \item \texttt{scaledCorrelation\_gks\_second}: GKS-II for the scaled model
    ($\langle\sigma^A\rangle_s\langle\sigma^B\rangle_s\le\langle\sigma^{A\triangle B}\rangle_s$),
    proved by the duplicate-variable trick adapted to the scaled modified weight.
  \item \texttt{scaledCorrelation\_monotoneOn}: $\langle\sigma^A\rangle_s$ is non-decreasing in $s$,
    from $d/ds\ge 0$ using GKS-II for the scaled model.
\end{itemize}

\begin{theorem}[Ball-boundary Simon-Lieb, weak form]
For ferromagnetic $p=(J,0,\beta)$, $E_0\subseteq E(G)$, distinct $r,s\in V$ not appearing
as endpoints of any $E_0$-bond, and $\langle\sigma_r\sigma_s\rangle_0=0$:
\[\langle\sigma_r\sigma_s\rangle\le\beta J\sum_{(k,l)\in E_0}
  \bigl(\langle\sigma_r\sigma_k\rangle\langle\sigma_s\sigma_l\rangle
        +\langle\sigma_r\sigma_l\rangle\langle\sigma_s\sigma_k\rangle
        +\langle\sigma_r\sigma_s\rangle\langle\sigma_k\sigma_l\rangle\bigr).\]
The extra $\langle\sigma_r\sigma_s\rangle\langle\sigma_k\sigma_l\rangle$ term vanishes in GJ's
formula (eq.~17.8.4) where Lebowitz for the scaled model removes it.
\end{theorem}

\noindent\textbf{Proof sketch.}
FTC: $\langle\cdot\rangle_1-\langle\cdot\rangle_0=\int_0^1(d/ds)\langle\cdot\rangle_s\,ds$.
With $\langle\sigma_r\sigma_s\rangle_0=0$, use the Step~135 derivative formula, GKS-I (drop negative
term), 4-pt monotonicity (GKS-II for scaled model), and full-model Lebowitz
(\texttt{summand\_le\_lebowitz\_of\_disjoint}).
The MVT via \texttt{norm\_image\_sub\_le\_of\_norm\_deriv\_le\_segment\_01'} concludes.

\noindent\textit{Status}: Complete (PR~\#953 Step 136, Issue~\#951 Part~2).

\subsection*{Step 137 -- GJ \S17.8 Thm 17.8.1: polynomial decay $\Rightarrow$ exponential decay (Partial)}

\noindent\textbf{Reference}: Glimm--Jaffe \S17.8 pp.\ 316--318 (GJ 1987).

\noindent\textbf{Statement.}
\begin{itemize}
\item \texttt{HasPolynomialDecay d $\Lambda$ p}:
  $\lim_{|x|\to\infty}\langle\sigma_0\sigma_x\rangle_\infty\cdot|x|^{d-1}=0$.
\item \texttt{correlationInfinite\_polynomial\_implies\_exponential}:
  polynomial decay $\Rightarrow$ $\exists\,m>0$, \texttt{HasExponentialDecay} (sorry'd, deferred).
\end{itemize}

\noindent\textbf{New supporting material in PR \#954.}
\begin{itemize}
\item \texttt{scaledCorrelation\_odd\_vanish}: at $h=0$, odd-cardinality scaled correlations
  vanish (global spin-flip symmetry for the scaled model).
\item \texttt{cor\_4\_3\_3\_scaled} (axiom): Lebowitz 4-site inequality for the scaled model
  (proved via $\phi^4$ approximation, same structure as \texttt{lebowitz\_four}).
\item \texttt{ball\_boundary\_simon\_lieb\_tight}: tight ball-boundary bound
  (no extra $\langle\sigma_r\sigma_s\rangle\cdot\langle\sigma_k\sigma_l\rangle$ term),
  using \texttt{cor\_4\_3\_3\_scaled}.
\item \texttt{scaledCorrelation\_at\_zero\_of\_sep} (axiom): when $E_0$ separates $r$ from $s$
  and $h=0$, $\langle\sigma_r\sigma_s\rangle_{s=0}=0$.
\item \texttt{latticeBall d r}, \texttt{latticeBallBoundaryEdges d r}: lattice ball and
  boundary bond set definitions.
\end{itemize}

\noindent\textbf{Proof sketch.}
Choose $R$ large so that $\alpha^{-1}:=\beta J\sum_{k\in\partial B_R}\langle\sigma_0\sigma_k\rangle_\infty<1$
(possible by polynomial decay). Apply \texttt{scaledCorrelation\_at\_zero\_of\_sep} (origin inside $B_R$,
target outside) and \texttt{ball\_boundary\_simon\_lieb\_tight} to get
$\langle\sigma_0\sigma_x\rangle_\infty\le \alpha^{-1}\cdot M_{|x|-R-1}$ where
$M_n=\sup_{|y|\ge n}\langle\sigma_0\sigma_y\rangle_\infty$.
Induction: $M_n\le\alpha^{-n/(R+1)}\to 0$ exponentially. Full induction argument sorry'd.

\noindent\textit{Status}: Partial (PR~\#954 Step 137, Issue~\#951 Part~3; main theorem sorry'd).

\subsection*{Step 130 -- Constant HLS convolution bound (Done)}

\noindent\textbf{Statement.}
For $2\alpha > d$, any $x, y \in \mathbb{Z}^d$:
\[
\sum_{z \in \mathbb{Z}^d} (1 + d(x,z))^{-\alpha} \cdot (1 + d(y,z))^{-\alpha}
\leq \sum_{z \in \mathbb{Z}^d} (1 + d(0,z))^{-2\alpha}.
\]

\noindent\textbf{Proof outline.}
\begin{enumerate}
\item \texttt{latticeDistance\_translate\_eq}: $d(x,z) = d(0,z-x)$.
\item \texttt{tsum\_pow\_neg\_translate}: $\sum_z (1+d(x,z))^{-\gamma} = \sum_z (1+d(0,z))^{-\gamma}$
  by change of variable $z' = z - x$ via \texttt{Equiv.addRight}.
\item \texttt{tsum\_pow\_neg\_conv\_le\_const}: apply AM-GM $ab \leq (a^2+b^2)/2$
  pointwise, sum, and apply translation invariance to both $x$- and $y$-translated sums.
\item \texttt{discrete\_hls\_convolution\_constant}: packages the explicit
  witness $C=\sum_z(1+d(0,z))^{-2\alpha}$ with strict positivity (from the
  $z=0$ term) and the uniform convolution inequality for all $x,y$.
\item \texttt{lemma\_17\_5\_2\_hls\_convolution\_constant}: exposes the same
  bound under the GJ Lemma~17.5.2 namespace for the later
  Lipschitz/HLS upper-bound composition.
\end{enumerate}

\noindent\textit{Status}: Complete (PR~\#946, Step 130, with current
PR~\#2665 \#1645 follow-up packaging).
Files: \texttt{IsingModel/PolyDecay.lean} and
\texttt{Concrete/LatticeGraphCorrelation/Lemma\_17\_5\_2.lean}.
Reference: GJ 1st~ed., \S17.5, pp.~310--312.

\subsection*{Step 134 -- Lipschitz continuity of $(m^-)^{2\alpha+1}$ on interval (Done)}

\noindent\textbf{Statement} (\texttt{pseudoMass\_pow\_succ\_lipschitz}).
Under the hypotheses of Steps~131--133, uniform on $[\beta_1,\beta_2]$:
\[|(h(\beta_2))^{2\alpha+1} - (h(\beta_1))^{2\alpha+1}|
  \leq \lceil 2\alpha+1\rceil\cdot K/r \cdot (\beta_2-\beta_1).\]

\noindent\textbf{Proof outline.}
Apply \texttt{norm\_image\_sub\_le\_of\_norm\_deriv\_le\_segment'} (MVT) to
$f(\beta') = (h(\beta'))^{2\alpha+1}$.
At each $\beta'\in[\beta_1,\beta_2]$: \texttt{HasDerivAt.fun\_pow} gives the derivative;
\texttt{pseudoMass\_power\_deriv\_le} (Step~131b) gives the bound $\|(f')(\beta')\|\leq C$.

\noindent\textit{Status}: Complete (PR~\#950, Step 134).
File: \texttt{IsingModel/PseudoMass.lean}.
Reference: GJ 1st~ed., \S17.5, Theorem~17.5.1, pp.~311--312.

\subsection*{Step 133 -- Lipschitz derivative of $(m^-)^{2\alpha+1}$ (Done)}

\noindent\textbf{Statement} (\texttt{pseudoMass\_pow\_succ\_deriv\_bound}).
Given the same hypotheses as Step~131b, there exists $d$ with
$\mathrm{HasDerivAt}\;(\beta'\mapsto(h(\beta'))^{2\alpha+1})\;d\;\beta$
and $|d|\leq\lceil 2\alpha+1\rceil\cdot K/r$.

\noindent\textbf{Proof outline.}
Chain rule (\texttt{HasDerivAt.fun\_pow}): $d=(2\alpha+1)(h(\beta))^{2\alpha}h'$.
From Step~131b: $(h(\beta))^{2\alpha}|h'|\leq K/r$.
Hence $|d|=(2\alpha+1)(h(\beta))^{2\alpha}|h'|\leq(2\alpha+1)K/r$.

\noindent\textit{Status}: Complete (PR~\#949, Step 133).
File: \texttt{IsingModel/PseudoMass.lean}.
Reference: GJ 1st~ed., \S17.5, used in the proof of Theorem~17.5.1, p.~312.

\subsection*{Step 132 -- Pseudo-mass upper bound and correlation decay (Done)}

\noindent\textbf{Statements.}
\begin{itemize}
  \item \texttt{pseudoMassG\_le\_two\_div\_one\_add\_pow} (Step 132a):
    $\mathrm{pseudoMassG}(\alpha,r,t) \leq 2/(1+(tr)^\alpha)$.
  \item \texttt{pseudoMassG\_le\_two\_div\_one\_add\_pow\_of\_preimage\_le} (Step 132b):
    If $\mathrm{pseudoMassG}(\alpha,1,m_1)=c$ and $0\leq m_0\leq m_1$, then $c\leq 2/(1+m_0^\alpha)$.
\end{itemize}

\noindent\textbf{Proof outline.}
Step 132a: $e^{-(tr)}\leq 1$ for $tr\geq 0$, so $2e^{-tr}/(1+(tr)^\alpha)\leq 2/(1+(tr)^\alpha)$.
Step 132b: \texttt{pseudoMassG\_strictAntiOn} gives $c=\mathrm{pseudoMassG}(m_1)\leq\mathrm{pseudoMassG}(m_0)$;
Step 132a gives $\mathrm{pseudoMassG}(m_0)\leq 2/(1+m_0^\alpha)$.

\noindent\textbf{GJ context.} Step 132b is the abstract form of the bound used in GJ p.~312:
$c_{x,z}=\langle\phi(x)\phi(z)\rangle/A\leq 2/(1+(m^-_{\rm global}\cdot d(x,z))^\alpha)$,
since $m^-_{\rm global}\leq m^-_{x,z}$ (inf) and $\mathrm{pseudoMassG}(m^-_{x,z}\cdot d)=c_{x,z}$.
Combined with the Lebowitz bound and HLS (Step~130), this will yield the \texttt{hc\_der}
hypothesis for \texttt{pseudoMass\_power\_deriv\_le} (Step~131b).

\noindent\textit{Status}: Complete (PR~\#948, Step 132).
File: \texttt{IsingModel/PseudoMass.lean}.
Reference: GJ 1st~ed., \S17.5, Theorem~17.5.1 proof, pp.~311--312.

\subsection*{Step 131 -- Abstract Lipschitz bound for pseudo-mass (Done)}

\noindent\textbf{Statements.}
\begin{itemize}
  \item \texttt{pseudoMass\_deriv\_abs\_le} (Step 131a):
    $|h'| \leq |c'| / (r \cdot c(\beta))$.
  \item \texttt{pseudoMass\_power\_deriv\_le} (Step 131b):
    $(h(\beta))^{2\alpha} \cdot |h'| \leq K/r$,
    assuming $|c'| \leq K \cdot c(\beta) / (h(\beta))^{2\alpha}$.
\end{itemize}

\noindent\textbf{Proof outline.}
From \texttt{pseudoMass\_deriv\_formula} (Step~117e): $h' = c'/g'$ where
$g' = d/dt\,\mathrm{pseudoMassG}(\alpha, r, \cdot)|_{t=h(\beta)}$.
From \texttt{pseudoMassG\_deriv\_abs\_ge} (Step~117f): $r \cdot c(\beta) \leq |g'|$.
Since $g' < 0$ (\texttt{pseudoMassG\_deriv\_neg}), $|g'| > 0$, giving
$|h'| = |c'|/|g'| \leq |c'|/(r \cdot c(\beta))$.
The power bound follows by multiplying by $(h(\beta))^{2\alpha}$ and using the HLS hypothesis.
The HLS hypothesis $|c'| \leq K \cdot c(\beta)/(h(\beta))^{2\alpha}$ is motivated by
Step~130 (\texttt{tsum\_pow\_neg\_conv\_le\_const}) via Lebowitz's inequality.
The concrete high-temperature finite-stage input is now exposed as
\texttt{inducedLatticeGraph\_beta\_deriv\_abs\_le\_high\_temp}, with the
GJ-named wrapper
\texttt{lemma\_17\_5\_2\_beta\_deriv\_abs\_le\_high\_temp}.

\noindent\textit{Status}: Complete (PR~\#947, Step 131).
File: \texttt{IsingModel/PseudoMass.lean}.
Reference: GJ 1st~ed., \S17.5, Theorem~17.5.1 proof, p.~312.

\subsection*{Step 126 -- Discrete exponential convolution bound (Done)}

\noindent\textbf{Theorems} (\texttt{summable\_exp\_neg\_int\_natAbs},
\texttt{summable\_exp\_neg\_latticeDistance},
\texttt{tsum\_exp\_neg\_dist\_eq},
\texttt{summable\_exp\_neg\_dist},
\texttt{lattice\_exp\_sum\_conv\_le}):
For $m > 0$ and $x, y \in \mathbb{Z}^d$:
\begin{enumerate}
\item $\sum_{n \in \mathbb{Z}} e^{-m|n|} < \infty$.
\item $\sum_{z \in \mathbb{Z}^d} e^{-m \|z\|_1} < \infty$.
\item $\sum_z e^{-m d(x,z)} = \sum_z e^{-m d(0,z)} =: C(m,d)$ (translation invariance).
\item $\sum_z e^{-m d(x,z)} < \infty$ for any basepoint $x$.
\item Convolution bound: $\sum_z e^{-m d(x,z)} e^{-m d(y,z)} \leq 2 C(m,d) \cdot e^{-m d(x,y)/2}$.
\end{enumerate}

\noindent\textbf{Proof sketch}:
Item (1): split $\mathbb{Z}$ into positive and negative parts, each bounded by the geometric series $\sum e^{-mn}$.
Item (2): induction on $d$ using $\mathbb{Z}^{d+1} \cong \mathbb{Z} \times \mathbb{Z}^d$ and product summability.
Item (3): translation invariance of lattice distance.
Items (4): immediate from (2) via translation.
Item (5): key pointwise inequality $e^{-m a} e^{-m b} \leq e^{-m d/2} e^{-m \min(a,b)}$,
which follows from $d(x,y) \leq a + b = \max(a,b) + \min(a,b)$, hence $\max(a,b) \geq d(x,y)/2$.
Sum using $e^{-m \min(a,b)} \leq e^{-m a} + e^{-m b}$ to bound $\sum e^{-m \min}$ by $2C$.

\noindent\textit{Status}: Complete (PR~\#940, Step 126).
File: \texttt{LatticeExpSum.lean}.
Reference: Glimm--Jaffe 2nd~ed., \S17.5 Lemma~17.5.2 (pp.~311--312).

\subsection*{Step 125 -- GHS consequence at infinite volume: $U_2^\infty(i,j)$ antitone in $h$ (Done)}

\noindent\textbf{Theorem} (\texttt{truncated2Infinite\_antitoneOn\_h\_of\_ne}):
For distinct sites $i \neq j$, ferromagnetic coupling $J \geq 0$, and $\beta > 0$,
the infinite-volume truncated 2-point function
$h \mapsto U_2^\infty(i,j)(h) = \langle\sigma_i\sigma_j\rangle_\infty(h) - \langle\sigma_i\rangle_\infty(h)\langle\sigma_j\rangle_\infty(h)$
is antitone on $[0,\infty)$.

\noindent\textbf{Proof sketch}:
For each stage $n$ with $\{i,j\} \subseteq \Lambda_n$, Step~124 gives
$U_2^{(n)}(h_2) \leq U_2^{(n)}(h_1)$ for $0 \leq h_1 \leq h_2$.
By $\texttt{tendsto\_truncated2AlongExhaustion\_truncated2Infinite}$,
$U_2^{(n)} \to U_2^\infty$ as $n \to \infty$.
Apply \texttt{le\_of\_tendsto\_of\_tendsto} to pass the pointwise inequality to the limit.

\noindent\textit{Status}: Complete (PR~\#939, Step 125).
File: \texttt{AmbientLattice.lean}.
Reference: Glimm--Jaffe 2nd~ed., \S4.3 Cor.~4.3.4; Friedli--Velenik \S3.6.3.

\subsection*{Step 124 -- GHS consequence: $\langle\sigma_i;\sigma_j\rangle_T$ antitone in $h$ (Done)}

\noindent\textbf{Theorem} (\texttt{truncated2\_antitoneOn\_h\_of\_ne}):
For distinct sites $i \neq j$, ferromagnetic coupling $J \geq 0$, and $\beta > 0$,
the truncated 2-point function
$h \mapsto \langle\sigma_i;\sigma_j\rangle_T(h) = \langle\sigma_i\sigma_j\rangle_h - \langle\sigma_i\rangle_h\langle\sigma_j\rangle_h$
is antitone on $[0,\infty)$.

\noindent\textbf{Proof sketch}: Compute $\tfrac{d}{dh}\langle\sigma_i;\sigma_j\rangle_T = \beta\sum_{k\in\Lambda} s_k$
where $s_k = \langle\sigma_i;\sigma_j;\sigma_k\rangle_T \leq 0$ for $k \notin \{i,j\}$ (GHS inequality),
and $s_i = -2\langle\sigma_i\rangle\langle\sigma_i;\sigma_j\rangle_T \leq 0$,
$s_j = -2\langle\sigma_j\rangle\langle\sigma_i;\sigma_j\rangle_T \leq 0$ (GKS-I + GKS-II).
Hence $\tfrac{d}{dh}\langle\sigma_i;\sigma_j\rangle_T \leq 0$, giving antitonicity via MVT.

\noindent\textit{Status}: Complete (PR~\#938, Step 124).
File: \texttt{FieldDerivative.lean}.
Reference: Glimm--Jaffe 2nd~ed., \S4.3 Cor.~4.3.4; Friedli--Velenik \S3.6.3.

\subsection*{Step 122 -- $\beta$-monotonicity: $\partial_\beta \langle \sigma^A \rangle \geq 0$ (Done)}

\subsubsection*{Mathematical content}

Using the derivative formula from Step 117a and GKS-II:
\[
  \frac{d}{d\beta} \langle \sigma^A \rangle_\beta
  = J \sum_{\{u,v\} \in E} \bigl(
    \langle \sigma^{A \triangle \{u,v\}} \rangle - \langle \sigma^A \rangle \cdot \langle \sigma^{\{u,v\}} \rangle
  \bigr) \geq 0
\]
By the mean value theorem, $\beta \mapsto \langle \sigma^A \rangle_\beta$ is monotone on $\beta \geq 0$.

\subsubsection*{Lean 4 formalization}

\begin{theorem}[\texttt{correlation\_beta\_deriv\_nonneg}]
For ferromagnetic $\mathtt{hf} = (J \geq 0, h = 0, \beta > 0)$: $0 \leq J\sum_e(\ldots)$.
\end{theorem}

\begin{theorem}[\texttt{correlation\_monotoneOn\_beta}]
For $J \geq 0$: $\mathtt{MonotoneOn}\;(\lambda\beta \Rightarrow \langle\sigma^A\rangle_\beta)\;(\mathrm{Ici}\;0)$.
\end{theorem}

\noindent\textit{Status}: Complete (PR~\#936, Step 122).
File: \texttt{BetaDerivative.lean}.
Reference: Friedli--Velenik \S3.7, Lemma~3.31 part 2, p.~107 (generalized to $\sigma^A$);
Glimm--Jaffe 2nd~ed. (1987), \S17.5 pp.~345--347.

\subsection*{Step 121 -- GKS-II monotonicity: $\partial_h \langle \sigma^A \rangle \geq 0$ (Done)}

\subsubsection*{Mathematical content}

Using the derivative formula from Step 118 and GKS-II:
\[
  \frac{\partial}{\partial h} \langle \sigma^A \rangle_h
  = \beta \sum_{i \in V} \bigl(
    \langle \sigma^{A \triangle \{i\}} \rangle - \langle \sigma^A \rangle \cdot \langle \sigma_i \rangle
  \bigr) \geq 0
\]
where GKS-II gives each term $\langle \sigma^{A \triangle \{i\}} \rangle \geq \langle \sigma^A \rangle \cdot \langle \sigma_i \rangle$.

\subsubsection*{Lean 4 formalization}

\begin{theorem}[\texttt{correlation\_field\_deriv\_nonneg}]
For ferromagnetic parameters $\mathtt{hf}$ and $\beta \geq 0$:
\[
  0 \leq \beta \cdot \bigl(
    \langle \sigma^A \cdot M \rangle - \langle \sigma^A \rangle \cdot \langle M \rangle
  \bigr)
\]
\end{theorem}

\noindent\textbf{Proof sketch}:
Helper lemmas rewrite $\langle \sigma^A \cdot M \rangle = \sum_i \langle \sigma^{A\triangle\{i\}}\rangle$
and $\langle M \rangle = \sum_i \langle \sigma_i \rangle$.
Then $\langle \sigma^{A\triangle\{i\}}\rangle - \langle \sigma^A \rangle \cdot \langle \sigma_i \rangle \geq 0$
follows from \texttt{gks\_second}, and the sum is nonneg by \texttt{Finset.sum\_nonneg}.

\noindent\textit{Status}: Complete formalization (PR~\#935, Step 121).
File: \texttt{FieldDerivative.lean}.
Reference: Friedli--Velenik, \emph{Statistical Mechanics of Lattice Systems} (2017),
\S4.2, Proposition~4.2.4, p.~58;
Glimm--Jaffe 2nd~ed. (1987), \S17.6 pp.~348--351 (derivative formula).

\end{document}
